% Framework source amendment v3.068: scoped geometry/operator additions, physical occurrence unchanged.
\documentclass[12pt,a4paper]{article}
\usepackage{graphicx}
\usepackage{slashed}
\usepackage{tikz}
\usepackage{colortbl}
\usepackage{calc}


% Macros-v2_203: v3.077 source-current bindings for unconditionality source integration. No mathematical macro changes.
%% Macros-v2_201: downstream full-carrier geometry update, five current root bindings, no mathematical macro change.
%% Macros-v2_200: Paper III-F v1.280 current-source binding for framework v3.074; mathematical macro definitions unchanged.
%% Macros-v2_198: current project v3.072 III-D and bibliography binding; mathematical macros unchanged.
%% Current project v3.070 source-reference bindings; previous macro bodies unchanged.
% Current II-D reference root for PapersMinimal v3.063; mathematical macros otherwise unchanged.
% Macros-v2_188: re-based current paper identities and A--G family; no change to geometry.
% Macros-v2_184: active Paper III-E v1.200/III-F v1.267 binding for framework v3.031; mathematical macros unchanged.
% Macros-v2_183: updated shared cross-reference bindings for Paper1B/III-E/III-F v3.030; mathematical macros unchanged.
% Macros-v2_182: active Paper II-F/V-A cross-references for framework v3.029; definitions unchanged.
% Paper-series v2.381: new I-B torsion/modular results, bibliography and current source roots.
% Paper-series v2.380: new exceptional bibliography and synchronized paper references.
% Paper-series v2.379 current 31-paper macro/aux registry.
% Paper-series v2.378: current 31-paper AUX/bibliography registry updated for Paper I-B v1.2 literature additions.
% Framework v2.371: current-source reference registry refreshed; old theorem labels retained.
% Framework v2.370: current 30-paper source/AUX root rebinding; native mathematical macros unchanged.
% Framework v2.368: synchronized merged 30-root registry; mathematical notation unchanged.
% Framework v2.358: active 30-paper source registry synchronized; existing mathematical macros unchanged.
% Framework v2.357: selective new source-owned certificates; active reference roots updated for five revised papers.
% Framework v2.356: active 30-root registry synchronized to the 30 papers in this release.
% Framework v2.275: exact positive dense-twin word cubature, finite repeatability and eighth-moment classification, nested sampling frames, joint action and POVM-jet exports; preserves v2.274.
% Framework v2.274: filename and root-source synchronization.
% Framework v2.273: restored missing Problem 3 theorems from v2.268; exact exceptional frames, twin geometry, convolution repeatability and projector/occurrence firewalls; merged on v2.272 without overwriting independent branches.
% Framework v2.272: merged v2.270 covariant native-readout mathematics with v2.271 Problem 4 finite-action Hopf hypercharge/holonomy corrections; no universal bulk-gauge uniqueness claim.
% Independent branches: v2.270 readout/guard exports and v2.270 K16 action-consumed charge closure; all prior scoped assumptions retained.
% Framework v2.269: synchronized 30 active AUX roots; mathematics and bibliography preserved from Macros-v2_161.
% Framework v2.266: Paper II-C converse-derived outcome resources; II-E, III-C, IV-A/B, V-A and VI-B downstream mathematical closures; 30-root registry synchronized.
% Framework v2.265: II-C coordinate-stage constitution and exact first-symbol host inheritance; prior mathematical domains retained.
% Framework v2.264: paper II-B reverse/audit and IV-A observable-minimal exchange; synchronized external references to 30 active roots.
% Framework v2.263: Paper I-C adds bridge-controlled root-carry parity and source-relative old-cap cancellation; other root mathematical bodies retained; active 30-root references synchronized.
% Framework v2.262: integrates Paper II-F radius, Casimir and readout results; active 30-root references synchronized.
% Framework v2.261: merged Paper II-B GP2-level/geometric repairs; preserves v2.260 II-D repairs; active references synchronized.
% Framework v2.260: Paper II-D signed RS, finite-symbol typing, and conditional curvature/Jordan repairs; references synchronized.
% Framework v2.259: conditional native SU(2) atlas and explicit later discharge; v2.258 mathematical bodies retained elsewhere.
% Framework v2.258: preserves parallel II-A/dynamics update and adds complete readout/frequency closures.
% Framework v2.257: exact minimal complete position readout and arbitrary-preparation frequency law.
% Framework v2.256: exact native inertia and sharp bounded ordered-source closure; preserves concurrent v2.255 work.
% Framework v2.255: preserves parallel detector integration and adds exact finite readout closures.
% Framework v2.254: Pasterski bibliography and synchronized 30-paper external-reference registry.
% Framework v2.253: typed raw coordinate defect in I-A and synchronized VI-A import; active references synchronized.
% Framework v2.252: synchronized source registry; optional III-G shell-frame calculation.
% Framework v2.251: forward common-action execution, joint source jets and scoped coupled control.
% Framework v2.251: exact full-heat variations, fermion source jets and explicit physical completion frontier.
% Framework v2.251: certified interacting source jets, pole stability and scoped completion proofs.
% Framework v2.251: synchronized active source registry; mathematical notation unchanged.
% Framework v2.251: synchronized active source registry; mathematical notation unchanged.
% Framework v2.251: synchronized active source registry; mathematical notation unchanged.
% Framework v2.251: synchronized source registry; no mathematical notation changes.
% Framework v2.244: synchronized 30-paper AUX registry and rooted phase-source action-order tests.
% Framework v2.242: signed rooted phase lift, complete transverse return and exact family cubic.
% Framework v2.241: proved rooted conormal geometry, conditional cap realization and scoped action tests.
% Framework v2.239: all 30 active source roots synchronized; no mathematical macro changes.
% Framework v2.238: Paper IV A/B/C split and 30-root reference registry.
% Framework v2.237: Paper VII-D technical split and 28-root reference registry.
% Framework v2.236: Paper VI-C technical-volume reference added; all current root bindings advanced.
% Framework v2.235: 26-root AUX and bibliography bindings synchronized; mathematical macros unchanged.
% BEGIN INLINE SOURCE: Paper00-TeXBuildCompatibility-v1_0.tex
% Rendering-only repairs for inherited long appendices and bitmap text fonts.
% Integer counters through 26 retain the original labels; later values use AA,
% AB, ... . This file changes no action coefficient, state or theorem premise.
\RequirePackage{alphalph}
\makeatletter
\let\BHPrimitiveAlph\@Alph
\newalphalph{\BHExtendedAlph}{\BHPrimitiveAlph}{26}
\let\@Alph\BHExtendedAlph
\AtBeginDocument{%
  \@ifpackageloaded{microtype}{\microtypesetup{expansion=false}}{}%
}
\makeatother

% END INLINE SOURCE: Paper00-TeXBuildCompatibility-v1_0.tex

% Framework v2.164: current reference/bibliography profile; mathematical definitions retained.
% Framework v2.162: finite metric, functional Hilbert and scoped Ward integration.
% Framework v2.161: inherits Macros-v2_57.tex exactly except the current external-reference binding.
% Macros v2.57 / framework v2.156: current reference registry only; mathematical definitions unchanged.
% Framework v2.143 / Macros v2.55: current external-reference bindings; mathematical notation unchanged.
% Framework v2.110 / Macros v2.54: current external-reference registry and bibliography; mathematical notation inherited.
% Framework v2.67 / Macros-v2_53: adds Born-action-cell and joint-EH readout symbols; absolute unit realization is metrological; MATCH_{J/scal}^{EH} remains the sole G24 normalization bridge.
% Framework v2.35 / Macros-v2_52: grouped paper numbering, active AUX registry, and Born-v6_20 bibliography binding synchronized; notation unchanged.
% Framework v2.26: active-source AUX and bibliography bindings synchronized; notation unchanged.
% Framework v2.09: reviewed early-reciprocity integration; scoped physical imports.
% Framework v2.08: Born--Euler mathematics and conditional complete-action potential criterion.
% Framework v2.07: active AUX registry includes Paper I-B v1.2; shared bibliography Born-v6_18.
% Macros-v2_46 / framework v2.04: mathematical atlas and downstream realization correction.
% Macros-v2_46 / framework v2.03: scoped structure integration and current dependency bindings.
% Macros-v2_46 / framework v2.01: shared source bindings; prior revision comments are historical.
% v2.42 / framework v1.99: Synchronized RG/AS owner and bibliography bindings.
% v2.41 / framework v1.99: active manuscript AUX bindings and geometric-closure bibliography synchronized.
% Release-v1.80 centralized source synchronization for scalar-circle/comparison-band integration.
% Release v1.79: Hurwitz hierarchy and kinetic-normalized action response integration; physical mass matching remains conditional.
% Release v1.78: approved M8 closure-batch integration; physical and analytic domains retained.
% Current v1.78 registry: all active TeX roots are synchronized below.
% Earlier release comments are historical scope records, not current status.
% Release-v1.77 synchronization: native same-owner Noether closure, generator--flow/history typing, analytic-certification firewalls, scalaron/K16, LF/QCD, quantum-Ward, and RAR boundary-owner transport are registered; shared mathematical notation is unchanged.
% Release-v1.76 synchronization: registry points to the v1.76 active paper roots.
% Release-v1.76 synchronization: registers the local A+Phi+g Noether/constraint closure and dynamics-frontier propagation; shared mathematical notation is unchanged.
% Release-v1.75 synchronization: registry points to the v1.75 active paper roots.
% Release-v1.75 synchronization: integrates the joint-Dirac Einstein bulk real form and the local odd-scalar heat-coefficient action on explicit analytic first-jet/interior domains; v1.74 Yang--Mills closure is retained. Fermionic reality, pure-metric curvature-squared/exact-nonlocal terms, noncanonical boundaries, physical vacuum selection, and actual solution-history tokens remain open.
% Release-v1.74 synchronization: registry points to the v1.74 active paper roots and retains the v1.73 theorem labels.
% Release-v1.74 synchronization: canonical Macros-v2_34.tex; generated local nonlinear Yang--Mills action, constrained connection-field phase space, source-free Gauss momentum map, and Legendre closure are integrated on their proved domains. The remaining non-gauge ACT_cl sectors and actual solution-history selection remain open.
% Release-v1.71 synchronization: closes the former JOINT_x/JOINT_5
% connection-occurrence gap by canonical structure-group extension of the
% nonflat Born connection, faithful associated representation, and the local
% cotangent--Clifford operator constructor.  Actual fluctuation members,
% actions, vacua, solutions, calibrations, and analytic QFT domains remain
% separately typed.  The active-AUX registry is advanced once.
% Release-v1.71 canonical macro and active-AUX registry.
% ======================================================================
% Macros-v2_39.tex -- canonical publication macro file
% Release v1.71 retains the v1.70 null-ray/QCD boundaries and additionally
% removes the duplicate structural connection actuality premise.  It does not
% claim global bundle topology, select a gauge/Higgs member, or promote the
% finite construction to an interacting continuum theory.
% Release v1.68 constructs the finite light-front total-fraction operators and
% parton number/momentum operator measures, proves their simplex and Hausdorff-limit
% theorems, and separates positive cut measures from scheme-renormalized
% light-ray distributions.  It also separates finite real Schur poles, stable
% norm-resolvent/Riesz pole-residue limits, and continued resonance
% pole/residue limits.  The Paper-VII-C scattering package is claim-indexed:
% local QFT and spectral claims import no S-matrix or infrared premise, while
% asymptotic amplitudes, LSZ, S-matrix poles, completeness, and cross sections
% import only their named massive or massless/long-range clauses.
% Release v1.67 factorizes BND48_H into a same-action Hessian/Schur
% descendant and still-open typed boundary-trace, critical-lift, and
% refinement-limit conditions.  It proves the finite-collar Dirichlet-to-
% Neumann law sqrt(y)tanh(L sqrt(y)), its nonuniform critical limit, and the
% finite-return ultra-infrared turnover.  The Paper-I-A Hardy constant mode is
% an exact geometric near-hit, not yet the physical metric boundary mode.
% The current local f(R) Hessian remains massive and therefore still refutes
% a RAR prediction on the minimal action branch.
% Release v1.66 derives the transverse depth-six normalizer bridge from the
% already generated C6 Reynolds image and Paper-I-A binary-dihedral normalizer.
% The bridge is the unique least 2D3-stable scalar-containing realization of
% the faithful quotient 2D3/C6=C2, has trace defect six and hence K=48, but is
% transverse to the nested active carrier.  Exact K=44 and K=52 contraction,
% semigroup, and energy-midpoint identities are registered as structural
% readouts rather than projector events.  The active E48 carrier gap remains
% unchanged.  A same-action boundary-symbol/Calderon intertwiner is isolated
% as the necessary and still unproved owner bridge to the static DtN sort;
% representation equivalence alone neither derives RAR nor creates a K=48
% action junction, pole, occupation, or cosmological transition.
% Release v1.65 derives the native E_48-centered carrier mass gap from the
% GP1--GP3-generated active spectrum, distinguishes its carrier-free interval
% from a particle located at E_48, and withdraws the unsupported promotion to
% a positive continuum Yang--Mills mass gap.  The latter remains a valid
% conditional transfer theorem on an enlarged continuum branch, but it is not
% an active framework prediction or an unfinished native-gap proof.
% Release v1.64 synchronizes Derivation Batches 14--22.  It registers the
% GP1--GP3 heat-profile orbit, same-Hessian electroweak response and pole
% ownership, K16/K48 action-response corrections, perturbation state and
% non-Gaussianity typing, finite-vacuum/gap/QCD limit repairs, factorized
% position-area ownership, the actual static metric response, and the ordered
% Wilson/string-tension limits.  Strict continuum convergence, positive
% physical gap, surviving higher cumulants, and positive tension remain open.
% Release v1.61 closes TIMEWIDTH16 as an independent direct-production datum.
% JUNC_{2|7} continues to own the finite intrinsic microscopic profile/width,
% while the actual same-action background/production history determines its
% physical-time pullback, support, and duration.  The current Paper-I-A--VII-C AUX
% registry is advanced once; shared mathematical notation is unchanged.
% Release v1.60 derives the final-octet primitive-remnant depletion, canonical
% Hopf-current representative, exact Hopf descent to dOmega/(4pi), and
% closed-FLRW a^{-3} neutral-current shape.  It retypes K=48 as the onset
% of the final one-unit geometric transfer and records cDM/RAR as distinct
% physical projections of that common neutral parent while preserving
% occurrence, Hilbert-owner, stress-matching, and static-kernel premises.
% Release v1.59 transports the v1.58 exceptional-shell/K16 closure through
% Papers III-F and VI-B--VII-C.  It replaces Paper VI-B's obsolete mixed-Hessian
% output typing by the third-variation pair vertex and Bogoliubov occupation
% support, derives the post-seed Higgs-pair feedback vertex and its exact
% alpha--delta kernel factors, propagates the unique K16/K48 action-junction
% split, identifies the QCD delta count with the upstream 2O shell atom only
% as a common-scalar numerical compatibility, and narrows THERM16 to dynamical
% approach/beta selection because the finite Gibbs/KMS target is already
% unique on a supplied simple observable algebra once its Hamiltonian and
% inverse temperature are supplied; direct-sum central weights remain typed.
% Release v1.58 registers the exceptional nearest-shell closure, the derived
% transverse K=16 action event, the unique interior action junction, the
% type-correct ordered pair vertex, the thin-layer squeeze law, and the
% radiation-like/matter-like kinematic split.  REG_0, transfer magnitude,
% physical layer width, and thermalization remain explicitly conditional.
% Release v1.57 registers the active-depth-six obstruction, the native
% K=48 carrier desert, the typed 7|9 owner reduction, and the dynamical
% reheating handoff across Papers III-B--VI-B.
% Release v1.56 registers the least-support n_2 inflation dressing of the
% persistent n_7 Laurent response, the analytic global-convexity and positive
% endpoint-measure proofs, the derived inflation core and K=16 coefficient
% freeze, and the no-particle dark-sector split into static RAR plus the
% conditional homogeneous effective-curvature response at reflected K=48.
% There is no n_6 action, scalaron-DM mass, or second Einstein--Hilbert term.
% Release v1.55 registers the protected-octet resolved/residual reflection
% q_sp+Z_W=8, its j<->8-j and K<->64-K scale covariance, and the derived
% rank-one-exit quadratic grade C_1^2/U=C_2 in Papers III-C, III-E, and VIII.
% The reflection is a typed trace/scale theorem, not a time reversal or a
% physical event-occurrence theorem; reheating transfer remains conditional.
% Release v1.54 registers the exact positive-part Laurent indices, the
% active-domain/terminal-endpoint action split, the no-second-EH firewall, and
% the q_DM=6/K_DM=48 scalaron readout typing across Papers III-D/E/F, VI,
% VIII, and IX.
% Release v1.53 registers the common nonnegative Laurent obstacle projection
% for n_1^+ and n_7^+, discharges the three former exit/freeze assumptions,
% and advances the Paper-III-D/E/F, VIII, and IX consumers.  Shared notation
% and bibliography ownership are unchanged.
% Release v1.51 registers the PQTI derivation batch and its updated
% Paper-III-A/C/D/E/F, VIII, and IX consumers.  Shared notation and
% bibliography ownership are unchanged.
% Release v1.50 advances the Paper-III-B AUX target to v1.40 after the
% rigorous audit repair; shared notation and bibliography ownership are unchanged.
% Release v1.49 advances the Paper-II-E/F AUX target to v3.14 after the
% minimal audit repair; shared notation and bibliography ownership are unchanged.
% Release v1.46 synchronizes Paper II-A's minimal frozen constructor-origin
% theorem and its five dependent AUX targets.  Shared notation is unchanged.
% Release v1.45 updates the Paper-II-B AUX target to v2.23 after the
% GP3-rigidity and appendix-placement repair; shared notation is unchanged.
% Release v1.44 advances the centralized Paper-I-A--VII-C AUX registry after the
% Paper-I-A v5.16 protected-completion scope repair.  Every current
% paper imports this one file; no paper-level opt-in or \externaldocument list
% is permitted.  The bibliography remains centralized here through Born-v6_13.
% Shared bibliography, theorem infrastructure, colors, notation alignment,
% and current Paper-I-A--VII-C external-reference registration.
% This file centralizes common imports, bibliography sorting, colors, theorem
% infrastructure, and the notation shared by the complete series.
% v2.9 (release v1.43) made the Paper-I-A--VII-C AUX registry automatic and
% centralized the bibliography and all external-reference ownership.
% v2.8 release-registry refresh (v1.33) registers the rooted generated-parent
% cascade and the POL/REG/free-law/11+4/Born-group closure revisions.  No
% shared mathematical notation or Paper-II-A fundamental principle is changed.
% v2.8 release-registry refresh (v1.32) registers the universal generated-
% continuation revisions in Papers II-A, III-A, III-C, IV, and VIII.  No
% shared mathematical notation or Paper-II-A fundamental principle is changed.
% v2.8 release-registry refresh (v1.30) registers the CRP4/projective-state
% source revisions and their direct registry-only consumers.  No shared
% mathematical notation or Paper-II-A axiom is changed.
% v2.7 release-registry refresh (v1.27) updates only the current external
% Paper-I-A--VII-C basenames after the placement/assumption synchronization; no
% shared mathematical notation or Paper-II-A axiom is changed.
% v2.7 synchronizes the geometric defect-repair revision: Paper-II-A v3.42,
% Paper-III-C v1.26, Paper-III-D v1.27, Paper-III-E v1.32, and Paper-III-F
% v1.34.  It registers the common trace-defect/spectral-filtration theorem,
% its owner-typed response pullback, and the reduced GRED and TERM interfaces.
% No shared mathematical notation is changed.
% v2.6 synchronizes Paper-III-E v1.31's cone-type repair, minimal-response
% classification, and explicit classical action-role boundary, together with
% Paper-III-F v1.32's downstream response-atlas synchronization.  No shared
% mathematical notation is changed.
% v2.5 synchronizes the Paper-III-D v1.26 faithful-trace-duality proof,
% exact ordered-owner joint locus, and typed alternative-presentation
% bookkeeping repair, together with its Paper-III-E and closure-frontier
% references.  No shared mathematical notation is changed.
% v2.4 synchronizes the Paper-III-C v1.25 orbit/process/owner repair, its
% downstream Papers III-B--F, V, VI, VIII, and IX revisions, and Born-v6_13.
% v2.3 synchronizes the Paper-II-A v3.41 witness/support/principal inheritance
% layer, the Paper-III-A v1.24 joint-realizability/readout firewall, all dependent
% host-rate revisions, and the current Paper-I-A--VII-B source registry.
% v2.2 synchronizes the generated active-realization, determinant-depth,
% dimension-weighted hierarchy, and typed cosmological-event interfaces.
% v2.1 synchronizes the exceptional-support interface revisions: Paper I-A
% v5.9, Paper II-A v3.30, Paper II-E/F v3.7, and Paper IV v4.16.
% v2.0 introduces the stage-safe G0 occurrence notation, synchronizes the
% active Paper-I-A--IV source registry (current rigorous Paper-III closure:
% III-A v1.17, III-B v1.22, III-C--E v1.17, III-F v1.14, IV v4.14), and makes
% external-reference self-skip depend on logical paper role rather than a
% versioned job name.
% ======================================================================

% ======================================================================
% macros.tex — Shared macros for the Born Geometry series
% ======================================================================

% ======================================================================
% 1. PACKAGES, BIBLIOGRAPHY, COLORS, AND UNIT DEFINITIONS
% ======================================================================
% Core packages shared across the complete Paper-I-A--VII-C series.  Loading them here makes
% the macro file a genuine series preamble while remaining compatible with existing
% paper preambles that load the same packages before \input.
\usepackage{amsmath}
\usepackage{amssymb}
\usepackage{amsthm}
\usepackage{mathrsfs}
\usepackage{array}
\usepackage{booktabs}
\usepackage{xcolor}
\makeatletter
\@ifpackageloaded{geometry}{}{\usepackage{geometry}\geometry{margin=1in}}
\makeatother
\usepackage{mdframed}

% Shared support packages required by the macros below and by most papers.
\usepackage{ragged2e}
\usepackage{longtable}
\usepackage{tabularx}
\usepackage{enumitem}
\usepackage[most]{tcolorbox}

% Series-wide numeric bibliography convention using the portable BibTeX path.
% All Papers I-A--VII-C use the ordinary \cite interface.
\usepackage[numbers,sort&compress]{natbib}
\providecommand{\printbibliography}{%
  \bibliographystyle{plainnat}%
  \bibliography{Born-v6_60}%
}

% Shared colors.  \providecolor avoids clobbering a paper-specific override.
\providecolor{postulatecolor}{RGB}{230,242,255}
\providecolor{compare}{RGB}{230,245,230}
\providecolor{definitioncolor}{RGB}{230,242,255}

\providecolor{openriskcolor}{RGB}{255,235,235}
\providecolor{bridgecolor}{RGB}{240,240,255}
\providecolor{gold}{RGB}{255,215,0}

% ======================================================================
% 2. CUSTOM ENVIRONMENTS
% ======================================================================
% Common table column types used in all papers.  Guarded so that a paper may
% define its own versions before loading this file.
\makeatletter
\@ifundefined{NC@find@L}{%
  \newcolumntype{L}[1]{>{\RaggedRight\arraybackslash}p{#1}}%
}{}
\@ifundefined{NC@find@C}{%
  \newcolumntype{C}[1]{>{\Centering\arraybackslash}p{#1}}%
}{}
\@ifundefined{NC@find@Y}{%
  \newcolumntype{Y}{>{\RaggedRight\arraybackslash}X}%
}{}
\makeatother

% ======================================================================
% 3. THEOREM STYLES AND ENVIRONMENTS
% ======================================================================
% Custom underlined theorem style (useful for remarks or notes)
\newtheoremstyle{underlined}
  {3pt}   % Space above
  {3pt}   % Space below
  {}      % Body font (empty means normal/roman text, typical for remarks)
  {}      % Indent amount
  {}      % Theorem head font (can be \itshape, \bfseries, etc.)
  {.}     % Punctuation after theorem head
  { }     % Space after theorem head
  {\underline{\thmname{#1}\thmnumber{ #2}}\thmnote{ (#3)}} % Theorem head formatting

% Standard plain theorem style environments.  Each environment is guarded so
% the shared macro file can be loaded after a paper-specific theorem setup.
\theoremstyle{plain}
\makeatletter
\@ifundefined{theorem}{\newtheorem{theorem}{Theorem}[section]}{}
\@ifundefined{lemma}{\newtheorem{lemma}{Lemma}[section]}{}
\@ifundefined{proposition}{\newtheorem{proposition}{Proposition}[section]}{}
\@ifundefined{corollary}{\newtheorem{corollary}{Corollary}[section]}{}
\makeatother

% Definition style environments (typically non-italicized body)
\theoremstyle{definition}
\makeatletter
\@ifundefined{definition}{\newtheorem{definition}{Definition}[section]}{}
\@ifundefined{postulate}{\newtheorem{postulate}{Postulate}[section]}{}
\@ifundefined{closure}{\newtheorem{closure}{Closure}[section]}{}
\makeatother

% Applying the custom underlined style
\theoremstyle{underlined}
\makeatletter
\@ifundefined{remark}{\newtheorem{remark}{Remark}[section]}{}
\@ifundefined{convention}{\newtheorem{convention}{Convention}[section]}{}
\@ifundefined{principle}{\newtheorem{principle}{Principle}[section]}{}
\makeatother

% Framed proof environment.
% Every \begin{proof} ... \end{proof} is automatically framed/shaded.
% This avoids mdframed/trivlist nesting, which can drop the top border after
% theorem-like environments.
\makeatletter
\renewenvironment{proof}[1][\proofname]{%
  \par\addvspace{6pt}%
  \noindent
  \begin{tcolorbox}[%
    enhanced,
    breakable,
    colback=black!4,
    colframe=black!35,
    boxrule=0.35pt,
    arc=1pt,
    left=6pt,
    right=6pt,
    top=4pt,
    bottom=4pt,
    boxsep=0pt,
    before skip=0pt,
    after skip=6pt,
    parbox=false
  ]%
  \pushQED{\qed}%
  \normalfont\itshape #1\@addpunct{.}\enspace
  \normalfont\ignorespaces
}{%
  \unskip\nobreak\hfill\popQED
  \end{tcolorbox}%
  \par\addvspace{2pt}%
}
\makeatother

\mdfdefinestyle{postulatebox}{
  backgroundcolor=yellow!10,
  linecolor=black,
  linewidth=0.4pt,
  skipabove=\baselineskip,
  skipbelow=\baselineskip,
  nobreak=true
}

% ======================================================================
% PROOF-READABILITY HELPERS
% ======================================================================

\newcommand{\ProofPlan}[1]{%
  \par\smallskip
  \noindent\textbf{Proof plan. }#1%
  \par\smallskip
}

\newlist{proofsteps}{enumerate}{2}
\setlist[proofsteps,1]{
  label=\textbf{(\alph*)},
  ref=\alph*,
  leftmargin=2.2em,
  itemsep=0.45\baselineskip,
  topsep=0.35\baselineskip
}
\setlist[proofsteps,2]{
  label=\textbf{(\roman*)},
  ref=\roman*,
  leftmargin=2.0em,
  itemsep=0.25\baselineskip,
  topsep=0.2\baselineskip
}


% ======================================================================
% 4. NUMBER SYSTEMS AND ALPHABETS
% ======================================================================
\providecommand{\RR}{\mathbb{R}} % Real numbers
\providecommand{\PP}{\mathbb{P}} % Projective space / Probability
\providecommand{\CC}{\mathbb{C}} % Complex numbers
\providecommand{\ZZ}{\mathbb{Z}} % Integers
\providecommand{\NN}{\mathbb{N}} % Natural numbers
\providecommand{\HH}{\mathbb{H}} % Quaternions / Hyperbolic space
\providecommand{\LFRayGeom}{\mathsf{LF\!\mbox{-}RAY\!\mbox{-}GEOM}}
\providecommand{\LFRep}{\mathsf{LF\!\mbox{-}REP}}
\providecommand{\LFState}{\mathsf{LF\!\mbox{-}STATE}}
\providecommand{\LFLim}{\mathsf{LF\!\mbox{-}LIM}}
\providecommand{\LFScat}{\mathsf{LF\!\mbox{-}SCAT}}
% Number systems / projective spaces
\providecommand{\CP}{\mathbb{CP}}
\providecommand{\RP}{\mathbb{RP}}

% ======================================================================
% 5. GROUPS AND LIE THEORY
% ======================================================================
\providecommand{\SO}{\ensuremath{\mathrm{SO}}}     % Special Orthogonal group
\providecommand{\SU}{\ensuremath{\mathrm{SU}}}     % Special Unitary group
\providecommand{\Spin}{\ensuremath{\mathrm{Spin}}} % Spin group
\providecommand{\CS}{\ensuremath{\mathrm{CS}}}     % Chern-Simons (capitalized)
\providecommand{\cs}{\ensuremath{\mathrm{cs}}}     % Chern-Simons (lowercase)

% ======================================================================
% 6. COMMON MANIFOLDS AND TOPOLOGY
% ======================================================================
\providecommand{\Sone}{S^1}                 % Standard circle

% ======================================================================
% 7. MATH OPERATORS (Declared via \DeclareMathOperator)
% ======================================================================
\DeclareMathOperator{\End}{End}   % Endomorphisms
\DeclareMathOperator{\Aut}{Aut}   % Automorphisms
\DeclareMathOperator{\spec}{spec} % Lowercase spectrum
\DeclareMathOperator{\tr}{tr}     % Lowercase trace (if you already have \Tr)
\DeclareMathOperator{\supp}{supp} % Support
\DeclareMathOperator{\TrA}{Tr_{A}} % Partial trace over subsystem A
\DeclareMathOperator{\TrB}{Tr_{B}} % Partial trace over subsystem B
% Lie/algebra operators
\DeclareMathOperator{\Ad}{Ad}
\DeclareMathOperator{\Sym}{Sym}
\DeclareMathOperator{\su}{su}
% ======================================================================
% 8. STANDARD OPERATOR-LIKE SYMBOLS (Safe definitions via \providecommand)
% ======================================================================
\providecommand{\vol}{\operatorname{vol}}   % Volume form
\providecommand{\Vol}{\operatorname{Vol}}   % Total volume
\providecommand{\Ric}{\operatorname{Ric}}   % Ricci tensor
\providecommand{\Rm}{\operatorname{Rm}}     % Riemann tensor
\providecommand{\Tr}{\operatorname{Tr}}     % Uppercase trace
\providecommand{\Spec}{\operatorname{Spec}} % Uppercase spectrum
\providecommand{\diag}{\operatorname{diag}} % Diagonal matrix
\providecommand{\rank}{\operatorname{rank}} % Rank
\providecommand{\scal}{\mathcal R} % Scalar curvature
\providecommand{\Cl}{\operatorname{Cl}}     % Clifford algebra
\providecommand{\diam}{\operatorname{diam}} % Diameter

% ======================================================================
% 9. FUNCTIONS, TRIGONOMETRY, AND CALCULUS
% ======================================================================

% Codifferential naming conventions
% (Avoid \ddag: commonly defined by physics.sty and symbol packages)
\providecommand{\codiff}{\mathrm{d}^{\dagger}} % Codifferential operator

% Quantum mechanics bra-ket notation (needed for basis tables)
\providecommand{\ket}[1]{\lvert #1 \rangle}
\providecommand{\bra}[1]{\langle #1 \rvert}
\providecommand{\braket}[2]{\langle #1 \mid #2 \rangle}

% ======================================================================
% 10. PHYSICAL SCALES AND CONSTANTS (Paper 1c / Born Geometry)
% ======================================================================
\providecommand{\Rp}{\ensuremath{R_{\mathrm P}}}     % matched common EH curvature scale (not defined from observed G)
\providecommand{\hBorn}{\ensuremath{h_{\mathrm B}}} % physical representative 2\pi\hbar of one primitive Born scalar action cell
\providecommand{\UEHjoint}{\ensuremath{U_{J}}} % joint-Dirac EH curvature-coordinate scale 8 kappa_J
\providecommand{\Gread}{\ensuremath{G_{\rm scal}}} % scalar-action coordinate c_op^3/(h_B U); physical matching explicit
\providecommand{\GreadJoint}{\ensuremath{G_{J,\mathrm{raw}}}} % raw ordinary-trace coordinate = G_scal/96
\providecommand{\Rds}{R_{\mathrm{dS}}}    % de Sitter Curvature
\providecommand{\Mpl}{M_{\mathrm{Pl}}}    % Planck mass
\providecommand{\lp}{\ell_{\mathrm{Pl}}}  % Planck length
\providecommand{\Pmax}{P_{\mathrm{max}}}  % Maximum momentum
\providecommand{\Mplbar}{\bar{M}_{\mathrm{Pl}}} % Reduced Planck mass
\providecommand{\eff}{\mathrm{eff}}       % Subscript for effective quantities
\providecommand{\Cosmo}{\ensuremath{\Lambda}}                % Cosmological constant (use \Cosmo in source)
\providecommand{\LCDM}{\ensuremath{\Cosmo\mathrm{CDM}}}  % \LambdaCDM (typesets with \Cosmo)
\providecommand{\Lcut}{\ensuremath{\Lambda_{\mathrm{cut}}}}  % Cutoff scale (brace when used as script)
\providecommand{\Rcs}{\ensuremath{R_{\mathrm{cs}}}}      % Configuration-space radius
\providecommand{\CosmoEff}{\ensuremath{\Cosmo_{\eff}}}   % Effective cosmological constant

% Disambiguation: \Rmom is the momentum-space (S^3) radius. 
% Reserve plain R for the 4D curvature scalar in f(R).
\providecommand{\Rmom}{\ensuremath{R_{\mathrm{mom}}}} 
\providecommand{\Punit}{P_{0}} % Momentum unit set by \Rmom (= \hbar/\Rmom)
\providecommand{\Rvec}{\ensuremath{R_{\mathrm{vec}}}} % invariant vector-space radius in the reduced round sector
% Specific numeric values and expressions
\providecommand{\Vlogcap}{L_{\mathrm{Hopf}}^{(1)}}
\providecommand{\Hopf}{\operatorname{Hopf}}
% Scalar phase (S^1) sector + universal cover notation (Paper I-A)
\providecommand{\tper}{T_\star}      % Time period/scale
\providecommand{\Eper}{E_\star}      % Energy period/scale
\providecommand{\ttilde}{\tilde t}   % Universal cover time
\providecommand{\Etilde}{\tilde E}   % Universal cover energy

% ======================================================================
% 10A. SERIES-WIDE LOGICAL / NOTATIONAL DISAMBIGUATIONS
% ======================================================================
% Use these symbols to avoid collisions among: (i) the Paper--III
% no-insertion/geometric-closure class, (ii) Paper--I-A trimer diagonals,
% (iii) Paper--V-A/V-C family-space diagonals, (iv) scalar-index running,
% and (v) protected shell-clock labels.
\providecommand{\Rprim}{\ensuremath{\mathsf R_{\mathrm{prim}}^{\Sigma}}} % primitive finite structural readouts
\providecommand{\Rloop}{\ensuremath{\mathsf R_{\circ}^{\Sigma}}} % closed-loop readout role
\providecommand{\Rbdry}{\ensuremath{\mathsf R_{\partial}^{\Sigma}}} % boundary-conditioned readout role
\providecommand{\Rgen}{\ensuremath{\mathsf R_{\mathrm{fin}}^{\Sigma}}} % finite generated closure of primitive roles
\providecommand{\PhysOcc}{\operatorname{PhysOcc}}
\providecommand{\Actual}{\operatorname{Actual}}
\providecommand{\GeomLic}{\operatorname{GeomLic}}
\providecommand{\InvDist}{\operatorname{InvDist}}
\providecommand{\Dcl}{\ensuremath{\mathsf{D}_1}}      % unmarked positive Dirac-even closed spectral-readout form class
\providecommand{\Dheat}{\ensuremath{\mathsf D_{\mathrm{heat}}}} % one-parameter heat-semigroup subclass
\providecommand{\DtriI}{\ensuremath{D^{\mathrm{tri}}_1}} % Paper--I-A cyclic-trimer diagonal
\providecommand{\DtriII}{\ensuremath{D^{\mathrm{tri}}_2}}
\providecommand{\DfamI}{\ensuremath{D^{\mathrm{fam}}_1}} % Family-space diagonal selector
\providecommand{\DfamII}{\ensuremath{D^{\mathrm{fam}}_2}}
\providecommand{\alpharun}{\ensuremath{\alpha_{\mathrm{run}}}} % scalar spectral-index running
\providecommand{\Ksh}{\ensuremath{K_{\mathrm{sh}}}} % protected shell-clock coordinate
\providecommand{\ddef}{\ensuremath{\delta_{\mathrm{def}}}} % universal defect scalar after Paper--III realization

% ======================================================================
% 10B. SERIES-WIDE LOGICAL STATUS TAGS
% ======================================================================
% Upright small-caps wrapper for audit/status labels.  It prevents LaTeX from
% trying to use unavailable small-caps italic font shapes inside theorem bodies
% such as TU/lmr/m/scit.
\providecommand{\AuditLabel}[1]{}
\renewcommand{\AuditLabel}[1]{{\normalfont\scshape #1}}

% These tags are bookkeeping labels for consistency audits. They do not
% create new mathematical hypotheses; they record how a datum is being used
% in the proof chain.
\providecommand{\StatusPrimitive}{}
\renewcommand{\StatusPrimitive}{\AuditLabel{Primitive}}
\providecommand{\StatusDerived}{}
\renewcommand{\StatusDerived}{\AuditLabel{Derived}}
\providecommand{\StatusImported}{}
\renewcommand{\StatusImported}{\AuditLabel{Imported}}
\providecommand{\StatusConditional}{}
\renewcommand{\StatusConditional}{\AuditLabel{Conditional}}
\providecommand{\StatusBridge}{}
\renewcommand{\StatusBridge}{\AuditLabel{Bridge-derived}}
\providecommand{\StatusComparison}{}
\renewcommand{\StatusComparison}{\AuditLabel{Comparison-only}}

% ======================================================================
% 10C. SERIES-WIDE CLAIM-WORDING AUDIT LABELS
% ======================================================================
% These are textual audit labels used to keep modal language consistent
% across the series.  They introduce no new mathematical assumptions.
\providecommand{\ClaimForced}{}
\renewcommand{\ClaimForced}{\AuditLabel{forced}}
\providecommand{\ClaimUnique}{}
\renewcommand{\ClaimUnique}{\AuditLabel{unique}}
\providecommand{\ClaimCanonical}{}
\renewcommand{\ClaimCanonical}{\AuditLabel{canonical}}
\providecommand{\ClaimMinimal}{}
\renewcommand{\ClaimMinimal}{\AuditLabel{minimal}}
\providecommand{\ClaimClosed}{}
\renewcommand{\ClaimClosed}{\AuditLabel{closed}}
\providecommand{\ClaimPrediction}{}
\renewcommand{\ClaimPrediction}{\AuditLabel{prediction}}

% ======================================================================
% 10D. SERIES-WIDE THEOREM-CHAIN AUDIT LABELS
% ======================================================================
% These labels mark the direction of load-bearing theorem chains. They are
% audit annotations only; they introduce no new mathematical assumptions.
\providecommand{\ChainInput}{}
\renewcommand{\ChainInput}{\AuditLabel{Chain input}}
\providecommand{\ChainOutput}{}
\renewcommand{\ChainOutput}{\AuditLabel{Chain output}}
\providecommand{\ChainGuard}{}
\renewcommand{\ChainGuard}{\AuditLabel{Chain guard}}
\providecommand{\ChainNoBackflow}{}
\renewcommand{\ChainNoBackflow}{\AuditLabel{No backflow}}

% ======================================================================
% 10E. SERIES-WIDE LOCAL THEOREM-AUDIT LABELS
% ======================================================================
% These labels mark a statement-level dependency audit around individual
% theorem/definition nodes. They are bookkeeping labels only.
\providecommand{\LocalAudit}{}
\renewcommand{\LocalAudit}{\AuditLabel{Local audit}}
\providecommand{\UsesOnly}{}
\renewcommand{\UsesOnly}{\AuditLabel{Uses only}}
\providecommand{\Exports}{}
\renewcommand{\Exports}{\AuditLabel{Exports}}
\providecommand{\DoesNotUse}{}
\renewcommand{\DoesNotUse}{\AuditLabel{Does not use}}

% ======================================================================
% 10F. SERIES-WIDE NUMERICAL-ROW AUDIT LABELS
% ======================================================================
% These labels classify individual rows in numerical, benchmark, status,
% and comparison tables.  They are audit annotations only and do not change
% any theorem, value, formula, or comparison number.
\providecommand{\RowAudit}{}
\renewcommand{\RowAudit}{\AuditLabel{Numerical-row audit}}
\providecommand{\RowTheorem}{}
\renewcommand{\RowTheorem}{\AuditLabel{Theorem-fixed row}}
\providecommand{\RowBridge}{}
\renewcommand{\RowBridge}{\AuditLabel{Bridge-fixed row}}
\providecommand{\RowRepresentative}{}
\renewcommand{\RowRepresentative}{\AuditLabel{Representative row}}
\providecommand{\RowComparison}{}
\renewcommand{\RowComparison}{\AuditLabel{Comparison row}}
\providecommand{\RowCalibrated}{}
\renewcommand{\RowCalibrated}{\AuditLabel{Calibrated row}}
\providecommand{\RowEstimate}{}
\renewcommand{\RowEstimate}{\AuditLabel{Estimate row}}
\providecommand{\RowOpen}{}
\renewcommand{\RowOpen}{\AuditLabel{Open row}}
\providecommand{\RowNoSelector}{}
\renewcommand{\RowNoSelector}{\AuditLabel{No selector}}

% ======================================================================
% 10G. TABLE HELPERS
% ======================================================================
% Table helpers used by Paper V-C.  The unused readout-label preset has been
% removed; only active table helpers remain.
\providecommand{\gpsection}[1]{%
  \addlinespace[0.55em]
  \multicolumn{3}{@{}l}{\textsc{#1}}\\[-0.25em]
  \midrule
}
\providecommand{\mathcell}[1]{%
  \(\begin{gathered}#1\end{gathered}\)%
}

\newcommand{\DerBy}[1]{%
  \begin{tabular}[t]{@{}l@{}}#1\end{tabular}%
}

\newcommand{\TableNotes}[1]{%
  \par\medskip
  \begin{minipage}{0.94\linewidth}
  \footnotesize\emph{Notes.} #1
  \end{minipage}\par
}
% ======================================================================
% 10H. PAPER VI-A / QCD NOTATION
% ======================================================================
% Shared symbols used by Paper VI-A and downstream QCD/Yang--Mills handoffs.
\providecommand{\Pqcd}{\mathfrak P_{\rm VIII}^{\rm QCD}}
\providecommand{\Iqcd}{\mathfrak I_{\rm QCD}^{(7)}}
\providecommand{\muQCD}{\mu_{\rm QCD}^{\rm int}}
\providecommand{\LambdaQCD}{\Lambda_{\rm QCD}}
\providecommand{\LambdaBH}{\Lambda_{\rm BH}^{\rm QCD}}
\providecommand{\HQCD}{\mathcal H_{\rm QCD}^{\rm BH}}
\providecommand{\HQCDfin}{H_{\rm QCD}^{\rm fin}}
\providecommand{\Hraw}{\mathcal H_{\rm raw}}
\providecommand{\HphysQCD}{\mathcal H_{\rm phys}^{\rm QCD}}
\providecommand{\Vcol}{V_C}
\providecommand{\Ps}{P_{\rm sing}}
\providecommand{\PiOct}{\Pi_8}
\providecommand{\MSbar}{\overline{\rm MS}}
\providecommand{\YM}{\mathrm{YM}}
\providecommand{\Ktri}{K_{\triangle}}
\providecommand{\KYM}{K_{\YM}^{\rm BH}}
\providecommand{\aBH}{a_{\rm BH}}
\providecommand{\epsBH}{\varepsilon_{\rm BH}}
\newcommand{\RowConstructive}{\textnormal{\textsc{Constructive}}}
\newcommand{\RowObligation}{\textnormal{\textsc{Bridge obligation}}}
\newcommand{\RowTarget}{\textnormal{\textsc{Target row}}}
\newcommand{\RowValidation}{\textnormal{\textsc{Validation row}}}
\providecommand{\slashed}[1]{\not{#1}}

% ======================================================================
% 10I. PAPER VI-B / CONTINUOUS-QFT NOTATION
% ======================================================================
% Shared symbols used by Paper VI-B and any downstream continuous-QFT or
% Yang--Mills handoff paper.  All definitions are guarded so individual
% papers may override them locally before loading this file if needed.
\providecommand{\A}{\mathfrak A}
\providecommand{\Hphys}{\mathcal H_{\rm phys}}
\providecommand{\QCD}{\mathrm{QCD}}
\providecommand{\BH}{\mathrm{BH}}
\providecommand{\loc}{\mathrm{loc}}

\providecommand{\RBH}{\mathcal R_{\rm BH}}
\providecommand{\DG}{\mathfrak D_G}
\providecommand{\PW}{\mathcal{P\!W}} % Peter--Weyl spectral-symbol subspace
\providecommand{\Toe}{\operatorname{Toe}} % Toeplitz compression/representative
\providecommand{\Hol}{\operatorname{Hol}} % Holonomy

% ======================================================================
% 11. LATE-LOADED PACKAGES (Must be at the end of the preamble)
% ======================================================================
% Load late packages only if a paper has not already loaded them, avoiding
% option clashes in paper-specific preambles.
\makeatletter
\@ifpackageloaded{hyperref}{}{\usepackage{hyperref}}
\makeatother
% --- Series-wide disambiguations (Pass 3) ---
\providecommand{\diracdelta}{\delta_{\rm D}} % Dirac delta distribution (delta reserved for hierarchy)
\providecommand{\Hg}{H_{\mathrm g}} % Hodge harmonic projector (keep H for Hubble)
\providecommand{\Gr}{G_{\mathrm r}} % Green operator (keep G for Newton)
\providecommand{\FD}{\mathrm{D}} % functional/variational derivative symbol (delta reserved)
\providecommand{\Id}{\mathbf{1}} % identity/Kronecker tensor (avoid \delta_{ij})


% BEGIN INLINE SOURCE: Paper00-ActiveExternalReferences-v2_130.tex
% Framework v2.191: current reference binding; existing definitions retained.
% Framework v2.190: current22-root registry with native clock-evolution proofs.
% Framework v2.189: exact current22-root external registry, including retained v2.188 hardening owners.
% Framework v2.187: current registry synchronization only; source release v2.186 preserved.
% Framework v2.186: authoritative current 22-paper registry.
% Framework v2.182: all current 22 logical paper AUX bindings.
% Framework v2.170: current I-B v1.43 external AUX binding.
% Framework v2.167: exact current 22-root registry after CP--BL2 integration.
% Framework v2.166: exact current 22-root compact-release AUX registry.
% Framework v2.165: authoritative current root bindings.
% Framework v2.164: exact current 22 manuscript AUX bindings.
% Framework v2.162: finite metric, functional Hilbert and scoped Ward integration.
% Framework v2.161: current 22 manuscript AUX bindings; source only.
% Framework v2.156 / registry v2.113: synchronized current 22 manuscript roots; source-only registry.
% Framework v2.143 / registry v2.111: active manuscript AUX basenames; no generated AUX or PDF included.
% ======================================================================
% Central external-reference imports, numbers only, no PDF jumping
% ======================================================================

% xr-hyper is loaded centrally so that papers may retain their existing
% \usepackage{xr-hyper} lines without redefining \externaldocument after this
% registry is installed.  If plain xr was loaded earlier, retain it.
\makeatletter
\@ifpackageloaded{xr-hyper}{}{%
  \@ifpackageloaded{xr}{}{\usepackage{xr-hyper}}%
}
\makeatother

% Make external-document registration idempotent.  The registry identifies a
% logical paper by its Paper<number><optional letter> prefix.  Prefix-qualified
% imports remain independent.  Current paper roots contain no local
% \externaldocument calls; this guard also makes accidental duplicates harmless.
\ExplSyntaxOn
\prop_new:N \g_bh_external_refs_prop
\cs_new_eq:NN \BHOriginalExternalDocument \externaldocument

\cs_new_protected:Npn \bh_external_document:nnnn #1#2#3#4
 {
  \tl_set:Nn \l_tmpa_tl {#3}
  \regex_replace_once:nnN
   { \A (?: xref/ )? (Paper[0-9]+[A-G]?) .* \Z }
   { \1 }
   \l_tmpa_tl
  \tl_set:Nx \l_tmpb_tl { #1 | \tl_use:N \l_tmpa_tl }
  \prop_if_in:NVTF \g_bh_external_refs_prop \l_tmpb_tl
   {
    \PackageInfo{BornHopf}
     {Skipping~duplicate~external~refs~for~#3}
   }
   {
    \prop_gput:NVn \g_bh_external_refs_prop \l_tmpb_tl {#3}
    \IfNoValueTF{#4}
     { \BHOriginalExternalDocument[#1][#2]{#3} }
     { \BHOriginalExternalDocument[#1][#2]{#3}[#4] }
   }
 }

\RenewDocumentCommand{\externaldocument}{O{} O{} m o}
 { \bh_external_document:nnnn {#1}{#2}{#3}{#4} }

\NewDocumentCommand{\BHExternalRefs}{m}
 {
  % Do not import the aux file of the logical paper currently being compiled,
  % even when the job and registry use different version suffixes.
  \tl_set:Nx \l_tmpa_tl { \c_sys_jobname_str }
  \regex_replace_once:nnN
   { \A (Paper[0-9]+[A-G]?) .* \Z }
   { \1 }
   \l_tmpa_tl
  \tl_set:Nn \l_tmpb_tl { #1 }
  \regex_replace_once:nnN
   { \A (Paper[0-9]+[A-G]?) .* \Z }
   { \1 }
   \l_tmpb_tl
  \str_if_eq:eeTF { \tl_use:N \l_tmpa_tl } { \tl_use:N \l_tmpb_tl }
   {
    \PackageInfo{BornHopf}
     {Skipping~external~refs~for~#1~because~this~is~the~current~document}
   }
   {
    \IfFileExists{xref/#1.aux}
     { \externaldocument[][nocite]{xref/#1}[#1.pdf] }
     {
      \IfFileExists{#1.aux}
       { \externaldocument[][nocite]{#1} }
       {
        \PackageInfo{BornHopf}
         {External~aux~file~#1.aux~not~found;~no~labels~imported}
       }
     }
   }
 }
\ExplSyntaxOff

% External references are automatic.  This file is the sole owner of both the
% bibliography and the current Paper-I-A--VII-D AUX registry (30 roots including I-C).  Missing auxiliaries
% are informational on the first pass and are populated by the source-only
% two-pass series build.  Basenames deliberately omit temporary download/copy
% suffixes such as ``(1)'' or ``-1''.
  % Paper I-A: round S3 compendium.
  \BHExternalRefs{Paper1A-Geometry-v5_181}

  % Paper I-B: complete exceptional geometry and finite spectral structures.
  \BHExternalRefs{Paper1B-ExceptionalGeometry-v1_27}
  % Paper I-D: technical companion of I-A through I-C.
  \BHExternalRefs{Paper1D-TechnicalMaterial-v1_69}
  % Paper I-C: reciprocal M8 compendium.
  \BHExternalRefs{Paper1C-M8Geometry-v1_146}

\ifdefined\BHMathOnly
  % Papers I-A and I-B import only the mathematical atlas labels.
\else
  % Paper II: five physical-mathematical main papers plus one technical proof volume.
  \BHExternalRefs{Paper2A-GeometricConstitution-v3_227}
  \BHExternalRefs{Paper2B-MomentumCarrier-v2_147}
  \BHExternalRefs{Paper2C-ReciprocalLocalization-v2_177}
  \BHExternalRefs{Paper2D-SymmetrySpectra-v2_145}
  % The former II-F export labels remain in II-E; current II-F is the technical volume.
  \BHExternalRefs{Paper2E-KinematicExports-v3_140}
  % Paper II-F: complete technical proof and calculation owner for II-A--II-E.
  \BHExternalRefs{Paper2F-TechnicalMaterial-v1_71}

  % Paper III: Dynamics and Intrinsic Geometric Constitution (six components).
  \BHExternalRefs{Paper3A-GeneratedDynamics-v1_181}
  \BHExternalRefs{Paper3B-SpectralAction-v1_192}
  \BHExternalRefs{Paper3C-ProtectedHierarchy-v1_196}
  \BHExternalRefs{Paper3D-RunningResponse-v1_168}
  \BHExternalRefs{Paper3E-GravityAction-v1_204}
  \BHExternalRefs{Paper3F-CompleteDynamics-v1_281}
  % Paper III-G: complete preserved technical proof owner.
  \BHExternalRefs{Paper3G-TechnicalMaterial-v1_58}

  % Paper IV-A: generated finite quantum mechanics.
  \BHExternalRefs{Paper4A-FiniteQuantumMechanics-v4_217}
  % Paper IV-B: complete finite quantum processes and physical response.
  \BHExternalRefs{Paper4B-QuantumResponse-v1_55}
  % Paper IV-C: extended proofs and finite technical calculations.
  \BHExternalRefs{Paper4C-TechnicalMaterial-v1_50}

  % Paper V-A: The Standard Model.
  \BHExternalRefs{Paper5A-FiniteMatterGeometry-v3_219}

  % Paper V-B: Finite Spectrum and Flavor Closure.
  \BHExternalRefs{Paper5B-TerminalSpectrum-v5_202}

  % Paper V-C: Finite Interaction Boundaries of the Standard Model.
  \BHExternalRefs{Paper5C-InteractionResponse-v2_185}
  % Paper V-D: preserved technical proof and calculation owner.
  \BHExternalRefs{Paper5D-TechnicalMaterial-v1_58}

  % Paper VI-A: Gravity and Closed Cosmology.
  \BHExternalRefs{Paper6A-GravityCosmology-v2_268}

  % Paper VI-B: Geometric Dark Sector and Cosmological Readouts.
  \BHExternalRefs{Paper6B-CosmologyReadouts-v1_215}

  % Paper VI-C: complete technical proof and comparison owner.
  \BHExternalRefs{Paper6C-TechnicalMaterial-v1_53}

  % Paper VII-A: Finite QCD Kernel and Nonperturbative Closure.
\BHExternalRefs{Paper7A-QCDKernel-v1_185}

  % Paper VII-B: Hadronic, Thermal, and Nuclear QCD Readouts.
  \BHExternalRefs{Paper7B-QCDReadouts-v1_169}

  % Paper VII-C: Continuous Quantum Field Theory.
  \BHExternalRefs{Paper7C-ContinuousQFT-v2_208}

  % Paper VII-D: complete technical material for VII-A through VII-C.
  \BHExternalRefs{Paper7D-TechnicalMaterial-v1_60}

\fi

% END INLINE SOURCE: Paper00-ActiveExternalReferences-v2_130.tex


\hypersetup{hidelinks}


\usetikzlibrary{positioning, calc} 

\title{\textbf{Born--Hopf Geometry II-C: Reciprocal Operators and Localization}\\[0.5em]
\large Generated position, protected Clifford coordinates, and effects}

\author{Peter Kugelmann\thanks{\texttt{pkugelmann@proton.me}}}
\date{Version 2.177 --- 20 November 2026 \\ Framework v3.077}

\setlength{\emergencystretch}{3em}

\begin{document}

\maketitle

\begin{abstract}
The constituted momentum role $\mathcal H_p\cong\mathcal H_1$ with
$m_p=1$ admits a generated reciprocal round position role. Its
primitive scalar-resolution seed is
$\mathcal P_x(1)=E_0^x\oplus E_1^x$, of dimension five, not a second
fifteen-dimensional momentum carrier. We construct finite covariant
coordinate operators, classify the intrinsic and ambient completion
relations, and derive exact operator-algebra, commutator, leakage,
and reciprocal-harmonic transport statements. The intrinsic tangential
class $\mathsf C_D$ is distinct from the radial cone class
$\mathsf C_{\rm cone}$; optional ambient completions do not become
physical coordinate or kinetic owners merely by existing.

The generated kinematic algebra has complex dimension $43$. The
intrinsic coordinate closure is
$\mathcal A_D\cong M_{11}(\mathbb C)\oplus M_4(\mathbb C)$;
a full $M_{15}(\mathbb C)$ envelope requires an additional generated
cross-sector bridge. Exact intrinsic localization moments are attainable
at the sharp classical radius $\sqrt3\,\Rmom$, including a seven-outcome
finite realization. The native readout span hierarchy has dimensions
$65$, $107$ and $137$ on its respective marked-observable domains.
Further sharp mathematical results distinguish the seven-outcome
position test, the $22$-effect protected endpoint comparison, and the
$136$-dimensional real endpoint quotient; these are distinct tasks,
not interchangeable outcome counts. Conditional Petersen, F40 and
marked finite comparison families are treated on their stated domains.

These are mathematical finite-operator and effect results. No abstract
POVM, frame, ambient completion or reconstruction certificate by itself
selects a physical preparation, detector, instrument, event-frequency
law or time history. The corresponding occurrence and calibration
requirements remain explicitly separate.

 A separate finite reciprocal-position scalar comparison determines the
complete compressed-coordinate matrix algebra and sharp one-setting
readout limitations, without claiming the corresponding apparatus,
coherent evolution, or scalar sharp-frontier occurrence.
\end{abstract}

\tableofcontents
\listoftables
\newpage

\section{Reciprocal-role inputs and observable scope}
\label{paper2c:sec:organized-scope}

\subsection{Introduction}
\label{p2c:sec:introduction}

\paragraph{Full technical derivations.}
The active auxiliary classifications, exact calculations, and negative
results supporting this paper are collected without abridgment in
Paper~II-F, Sections~\ref{p2ftech:sec:operators} and
\ref{p2ftech:sec:localization}.
That companion is a proof owner, not an additional constructor,
physical occurrence, state, or calibration assumption.

Paper~II-C uses the derived compact parent from Paper~II-A and the finite
equivariant module package organized in Paper~II-B.  The carrier alone does
not determine an arbitrary ambient coordinate completion.  Paper~I-A supplies the intrinsic round and Hopf geometry; Paper~I-B
contains the finite round-symbol class and exceptional comparison results,
including the intrinsic/cone fingerprints and their ancestry classification.  The native problem of this paper is to place those mathematical
exports in the occurring finite-coordinate target, classify wider ambient
completions, and derive the resulting compact kinematic algebra.  In
particular, physical co-realization of the \emph{constituted} intrinsic
first-symbol comparison does not imply that the bare II-B carrier alone
chooses an arbitrary Hermitian coordinate completion; the precise
coordinate-stage boundary is recorded in
Remark~\ref{p2c:rem:v265-coordinate-constitution-boundary}.

\subsubsection{Logical scope and ownership}
The parent/host and finite-seed theorems are imported, not reproved.  The
  constitutional chain is supplied by Papers~II-A and~II-B.  Paper~II-A supplies the stage-safe G0 geometric-actualization
certificate, the round parent, generated-role/co-realization rules, and the
constitutional GP triptych (finite/bounded, exact/faithful,
irreducible/least); Paper~II-B supplies the constructor-relative momentum
host theorem and its exact generated type-support,
\[
        \mathcal G^{\rm car}_{\Sigma_p}\cong\mathcal H_1,
        \qquad m_p=1.
\]
Full-endpoint multiplicity-one is conditional on the independent
whole-structured certificate of
Theorem~\ref{p2b:thm:momentum-carrier-rigidity}.
The native problem of II-C is therefore operator-theoretic: classify
finite-coordinate actions and ambient completions on that seed, prove the
canonical retraction for a fixed ambient action, derive the protected-block
formulas and exact compact commutators, and record the generic auxiliary
finite-\(m\) statements needed by later recovery.  Heat/zeta/complete spectral
analysis and Hopf-sector topology are owned by II-D; dynamics and internal
symmetry lie later in the series.  No later Paper-II theorem is used as a
premise.

\subparagraph{Branch status.}
All momentum-role seed statements on \(\mathcal H_1\) are structural
\([\mathrm{HOST}]\) outputs imported from II-B, not level-discharge
corollaries.  Their role-level physicality is inherited from the same
Paper~II-A primitive-transfer occurrence; no second momentum witness is
introduced.  Generic auxiliary results on a mathematical \(\mathcal H_m\)
remain \([\mathrm{FIN}{-}m]\), and a specialization to a fixed auxiliary
level may be marked \([\mathrm{SPEC}{-}m_0]\).  The Paper-I-A Toeplitz
family \(\mathsf C_{\rm PI}(\rho_{\rm prot})\), its protected normalization
\(\varrho_{\rm sp}=\rho_{\rm prot}/2\), an optional scale match
\(\rho_{\rm prot}=\Rmom\), and the choice of an ambient completion remain
logically independent of \(\mathcal H_p\cong\mathcal H_1\) and \(m_p=1\).
The scale-matched Paper-I-A point is denoted \(\mathsf C_{\rm cone}\).  The
separate intrinsic tangential class \(\mathsf C_D\) and the two-branch
primitive ancestry classification are imported from Paper~I-A.  The
coordinate-operation stage of this paper explicitly allocates the intrinsic
fixed-radius first-symbol comparison to a finite coordinate target sort
\emph{before its finite coefficients are evaluated}.  Paper~II-A's GP2 host
faithfulness and generated-role co-realization then retain its exact output
without a second occurrence or a numerical selector.  This stage allocation
is not a claim that the initial II-A momentum seed or II-B carrier by itself
already selects a finite coordinate operation; the radial-cone word is not
part of the constituted fixed-radius target grammar.

\begin{remark}[Generated six-type support versus full-endpoint carrier]
\label{p2c:rem:v283-seed-vs-full-carrier}
The one-copy $\mathcal H_1$ is the canonical representation of the
GP2-generated six-type support; $m_p=1$ is its word/level label and not
an endpoint multiplicity theorem.  Statements formulated on
$\mathcal H_1$ are exact for that canonical finite model; statements
about an arbitrary complete physical endpoint require the GP3
whole-structured comparison of
Theorem~\ref{p2b:thm:momentum-carrier-rigidity} or a separately certified
compatible extension.  The GP2 co-realization of its six generated
constituents needs no additional momentum occurrence root.
\end{remark}

\begin{remark}[Paper-II-C host and auxiliary-finite contract]
\label{p2c:rem:v2-finite-level-contract}
The structural momentum-role carrier is fixed before II-C begins:
\[
        \mathcal H_p\cong\mathcal H_1,
        \qquad m_p=1.
\]
Every theorem stated for this carrier is therefore a \([\mathrm{HOST}]\)
kinematic theorem, possibly relative to a specified
\(\mathsf C\in\mathfrak M_{\rm coord}\), including the representative
\(\mathsf C_{\rm PI}\), when stated explicitly.  By contrast, symbolic
finite-\(m\) statements are auxiliary \([\mathrm{FIN}{-}m]\) results and do
not assert a second physical momentum tower.  The reciprocal position side
instead uses the scalar windows
\(\mathcal P_x(m)=\bigoplus_{k\le m}E_k^x\).  Its primitive coordinate
frontier has the derived resolution \(m_x^{\rm res}=1\).  A later value
\(m_x[n]>1\) is available only after a new finite position-test frontier is
generated; a physical-time trajectory additionally requires a dynamics
theorem relating time to those cumulative frontiers.  Large resolution and
local recovery still require their own band, curvature, and comparison
estimates.
\end{remark}

\subparagraph{Reciprocity at the cutoff.}
For the scalar-window/enriched-carrier comparison
\(\mathcal P_x(m)\otimes\mathcal H_{1,p}\), the unequal level is a retained
mark, so the selected presentation is not internally Born invariant.
Paper~II-A, Proposition~\ref{p2a:prop:typed-born-branch-transport}, supplies
the physical permission and typing for transport to the separately certified
reciprocal presentation.  Paper~I-C,
Theorem~\ref{p1b:thm:unequal-exchange}, then supplies the concrete
mixed-family projector/operator intertwiner: it maps to the presentation with
the two labels exchanged.  The least invariant comparison is the sum of these
two carriers.  Compressing exchange to just one carrier is an involution only
on its exchange-closed part.  This supplies a precise comparison without
changing \(m_p=1\), creating a second physical momentum tower, or identifying
a scalar window with an enriched carrier.


\begin{remark}[Imported G0 occurrence and readout boundary]
\label{p2c:rem:g0-occurrence-readout-boundary}
All physical interfaces below are role-specific specializations of
Paper~II-A's G0 certificate.  Each asserted occurrence must exhibit a
contingent token of a geometrically licensed, claim-relative invariant
distinction with result-independent provenance; a new role enters only through
a separately declared downstream dependency frontier.  The generated algebra, effects, reciprocal evaluation spectrum, scalar
resolution filtration, and canonical first-symbol position--transvection
comparison are structural theorem outputs.  Calling a generated test an
apparatus observable, selecting or preparing a state, interpreting an
evaluation as an observed probability, realizing a nonconstant position
history, or performing a particular first-symbol comparison requires the
corresponding G0 occurrence/readout token.  A stronger finite-band or
local--global reciprocal identification requires its own explicit bridge,
domain, and error topology, and a dimensionful identification requires a
separate calibration.  None of these token/bridge/calibration certificates is
needed to constitute the role-level first-symbol pairing itself. For a supplied stronger finite comparison, the effect-frame and common-readout tests of Section~\ref{p2c:sec:v2210-lit-effect-frames} determine precisely which source differences require separation; they do not turn an unperformed apparatus into a generated token.  Every such
certificate must be generated from the constituted geometry and fixed
independently of its outcome; its tag alone supplies no witness and it cannot
enlarge or reselect the frozen momentum seed.  Analytic
recovery estimates remain separate theorem-domain hypotheses and are not
derived by G0.
Accordingly, any later phrase ``physical momentum carrier'' denotes the
co-realized structural carrier inherited from the Paper~II-A/II-B momentum
provenance chain.  It does not by itself assert a particular measurement,
prepared token, numerical momentum value, or calibration.
\end{remark}

\begin{remark}[Dependency record for the finite coordinate package]
\label{p2c:rem:dependency-record}
Four inputs remain distinct.  First, II-A and II-B supply the round momentum
host, its \(\Spin(4)\) action, and the derived carrier
\(\mathcal H_p\cong\mathcal H_1\).  Second, II-C introduces the auxiliary
multiplicity-free ambient universe
\(\mathcal H_\infty\) and the mathematical integer-indexed finite family
\(\mathcal H_m\) for completion and recovery questions; these spaces are not
GP2-generated momentum-role carriers.  Third, the finite completion
family \(\mathfrak M_{\rm coord}\) is classified below, with
\(\mathsf C_{\rm PI}(\rho_{\rm prot})\) supplying the Paper-I-A protected-scale
family, \(\mathsf C_{\rm cone}\) its scale-matched ambient-symbol point, and
\(\mathsf C_D\) the intrinsic tangential point imported from Paper~I-A.  Fourth,
extending a finite point of that family to an ambient
equivariant coordinate action remains nonunique.
Proposition~\ref{p2:prop:ambient-completion-existence-nonuniqueness}
proves existence and nonuniqueness of such completions.  None of these later
ambient choices alters the already frozen momentum-seed carrier.
\end{remark}

\begin{remark}[Active-scope status of the II-C interfaces]
\label{p2c:rem:selector-status-classification}
The II-C registry separates four logically different classes.
\begin{enumerate}[label=(\arabic*), leftmargin=2.2em]
\item \textbf{closed on every active Papers-I-A--VI-B use:}
\(\mathcal H_1\), \(\mathcal A_{\rm kin}\), the reciprocal role
\(S^3_B\), the canonical first-symbol and operator-germ pairings, the scalar
filtration \(\mathcal P_x(m)\), the least-resolution functor, the
baseline-relative zero-tail ambient section, generated all-band reciprocal
covariance, intrinsic same-parent compression/recovery, the scalar and
unsharp mathematical POVMs, and the exact optimal classical moment radius;
\item \textbf{unused stronger or post-IX extensions:}
a single level-independent full-observable physical ambient action, a
universal all-resolution quantization into the frozen momentum carrier, and
recovery/comparison between independently constructed parent and finite
operators.  These are not active residual roots through Paper~VI-B;
\item \textbf{one actual preparation/measurement occurrence package:}
an apparatus or channel, a prepared state, a performed test, a physical
ownership map for its labels/effects, an outcome/readout, and any actually
performed first-symbol comparison;
\item \textbf{role-specific physical, calibration, or analytic data:}
a physical kinetic/Hamiltonian owner, an actual nonconstant frontier history,
numerical scale identifications, interacting-continuum dynamics, and any
claim-specific approximation topology.
\end{enumerate}
The raw compatible completion cone and its ambient lifts remain useful host
mathematics, but the fixed-radius branch already has the unique physically
generated finite symbol descendant \(\mathsf C_D\).  Mathematical
availability in classes (2)--(4) neither reselects the momentum seed nor
actualizes a physical token.  In particular, the stronger ambient-action,
stronger pairing, and independent-owner recovery warnings are removed from
the active Papers-I-A--VI-B residual count rather than replacing their solved
weaker targets.
\end{remark}

\subparagraph{Optional tower infrastructure.}
The word and lower-ideal towers retained in Paper~II-B remain available as
mathematical/later-role infrastructure, but they no longer supply the
momentum-seed selector.  The former \(\mathcal E_9\) alternative-selector
analysis is retained only as the diagnostic in
Paper~II-F, Section~\ref{p2c:app:e9-diagnostic}.

\label{p2:part:dictionary-operator-algebra}

This part fixes the representation dictionary and operator algebra used later
in the series.  The entries are classified by logical status: some are native Paper-II
outputs, some are Paper-I-A theorem-bank consequences used as kinematic infrastructure, some are
worked representatives, and some are results reserved for later papers.  No
dynamical law, renormalization flow, or physical threshold identification is
introduced in this part as an additional premise.


\section{Reciprocal dictionary and generated position role}
\label{paper2c:sec:organized-dictionary}

\subsection{Constitutive Mappings (Core Dictionary)}


\subsubsection*{Status convention}

The dictionary records constitutive mappings from intrinsic kinematic
structures to physical representatives; it introduces no dynamical law.  It
uses the series statuses \([\mathrm{HOST}]\), \([\mathrm{PARENT}]\),
\([\mathrm{FIN}{-}m]\), \([\mathrm{SPEC}{-}m_0]\), and
\([\mathrm{DEFERRED}]\).  Within those types it further marks a construction
as derived/canonical, completion-relative, or a worked representative.  A
completion-relative result is exact after an ambient equivariant action has
been declared; a worked representative carries no uniqueness claim beyond
its stated convention.

\subparagraph{Later-paper boundary.}
\label{p2:sec:p4-closure-map}
II-C supplies finite representation spaces, bounded coordinate observables,
sector labels, and exact kinematic commutators.  It does not assume or import
later Born rules, Fock constructions, Hamiltonians, superselection principles,
dispersion relations, thresholds, vacuum energies, or thermal interpretations.
Any such reading requires the later paper that constructs it and is not part
of the coordinate mathematics proved here.
\subsubsection{Kinematic physics dictionary with explicit logical status} 
\label{p2:sec:physics-dictionary}
Tables~\ref{p2:tab:physics-dictionary-kinematic} and
\ref{p2:tab:physics-dictionary2} record the nonredundant kinematic and
fibration/spin mappings.  Spectral formulas are collected separately in
the II-D-owned spectral table; II-C does not duplicate a complete-spectrum or
threshold atlas.  A
physical representative named in a dictionary row is not a standalone
Paper-II dynamical interpretation; any later structural or dynamical
realization must be cited separately.  Protected spinor-coordinate invariants are stated below
in terms of \(\varrho_{\rm sp}=(\alpha^2+\beta^2)^{1/2}\).  The symbol
\(\gamma=\varrho_{\rm sp}\) is used only after fixing the relative chiral
phase, and each numerical normalization is identified as a separate
convention.


The constitutive mappings are extended with the rigorously derived identifications from the geometry of the vector-momentum constitution.


\begin{table}[!htbp]
\centering
\footnotesize
\begingroup
\setlength{\tabcolsep}{2.5pt}
\renewcommand{\arraystretch}{1.18}

\begin{tabular}{|
  L{3.0cm}|
  L{7.0cm}|
  L{\dimexpr\linewidth-3.0cm-7.0cm-6\tabcolsep-4\arrayrulewidth\relax}|
}
\hline
\textbf{Structure in the vector-momentum constitution}
& \textbf{Geometric Property}
& \textbf{Physical Representative} \\
\hline

\multicolumn{3}{|c|}{\textbf{\textit{Kinematic Identifications}}} \\
\hline

\(\SU(2)\cong S^3(1)\cong S^3(\Rmom)\) by \(q\mapsto \Rmom q\)
& Compact spherical momentum cover with transported product
  \(q\odot r=\Rmom^{-1}qr\), identity
  \(e_{\Rmom}=\Rmom\mathbf1\), and separate generator scale
  \(\Punit=\hbar/\Rmom\)
& Momentum space \\
\hline

Embedding \(S^3 \subset \RR^4\)
& Classical scalar coordinate functions \(x_a\); after choosing an ambient
  equivariant completion \(\mathsf X_a\), its unique finite return is
  \[
    \hat X_a=\Pi_m\mathsf X_a\Pi_m
  \]
& Completion-relative coordinate representatives.  Scalar multiplication
  coordinates commute, while finite compressed coordinates need not; see
  Theorems~\ref{p2:thm:derived-ambient-coordinate-action} and
  \ref{p2:thm:finite-toeplitz-coordinate-uniqueness}. \\
\hline

Dimensionless chiral generators \(L^{\pm}_i\)
& Left/right \(\mathfrak{su}(2)\) generators.  Either chiral copy generates
  group-valued transfer composition; their difference gives the
  homogeneous-space transvection operator, which is not a subgroup generator
& Homogeneous-space displacement (symmetry-momentum) operator
  \[
    \hat{P}_i
    =
    \frac{\hbar}{\Rmom}(L^+_i-L^-_i)
  \] \\
\hline

\(\SO(4)\) Casimir \(\mathcal C\)
& Selected scalar homogeneous-space Laplacian/Casimir; the full
  \(\Spin(4)\) enveloping algebra also has independent chiral Casimirs
& Scaled Casimir/Laplacian representative after choosing \(M\ne0\)
  \[
    \hat H
    =
    \frac{\hbar^2}{2M\Rmom^2}\mathcal C
  \] \\
\hline

Scalar Peter--Weyl basis \(D^j_{qq'}\); full lifted basis
\(V_{j_L}\otimes V_{j_R}\)
& Complete scalar coefficient system, with non-diagonal sectors interpreted as
  lifted vector/spinor/bundle module modes
& Casimir and component-weight modes / finite module modes \\
\hline

\end{tabular}

\caption[Kinematic physics dictionary]{Kinematic physics dictionary.}
\label{p2:tab:physics-dictionary-kinematic}
\TableNotes{%
This table records the logically classified kinematic mappings from vector-momentum geometry to physical representatives.
}
\endgroup
\end{table}





\subparagraph{Supplementary dictionary.}
The coordinate type is \((\tfrac12,\tfrac12)\); the ambient universe is
multiplicity-free.  The detailed fibration/contact/spin dictionary is
deferred to Paper~II-F, Section~\ref{p2c:app:deferred-host-infrastructure} and supplies
no seed selector.

\subsubsection{Canonical operator representatives}

\begin{theorem}[Canonical transvection bridge from the regular action]
\label{p2:thm:canonical-transvection-bridge}
Let \(G=SU(2)\) carry its round bi-invariant metric, and let
\(G_L\times G_R\) act by
\[
        (a,b)\cdot q=aqb^{-1}.
\]
At \(e\), the stabilizer algebra and its canonical complement are
\[
 \mathfrak k=\{(X,X):X\in\mathfrak g\},
 \qquad
 \mathfrak p=\{(X,-X):X\in\mathfrak g\}.
\]
They are respectively the \(+1\) and \(-1\) eigenspaces of the involution
induced by group inversion.  The evaluation map is
\[
        \operatorname{ev}_e(X,Y)=X-Y,
\]
and its restriction \(\operatorname{ev}_e|_{\mathfrak p}:\mathfrak p\to
T_eG\) is an isomorphism.  Its inverse
\[
        X\longmapsto \left(\frac X2,-\frac X2\right)
\]
is the unique diagonal-\(G\)-equivariant, inversion-odd right inverse of
evaluation.  Thus the difference of the left and right generators is not a
convention imported from the flat limit: it is the canonical unframed
representative of a tangent displacement.  The momentum unit
\(\hbar/\Rmom\) supplies its physical normalization.
\end{theorem}

\begin{proof}
The stabilizer condition follows from \(aeb^{-1}=e\).  Conjugating the action
by inversion exchanges the two factors, so its differential has eigenspaces
\(\mathfrak k\) and \(\mathfrak p\).  Differentiating
\(\exp(tX)e\exp(-tY)\) gives \(X-Y\); hence evaluation sends
\((X,-X)\) to \(2X\).  Every inversion-odd right inverse must take values in
\(\mathfrak p\), where evaluation is already an isomorphism, proving
uniqueness.  Diagonal equivariance follows from the adjoint action.  No
left/right chirality is selected.
\end{proof}

\begin{definition}[Transfer-composition generators versus transvections]
\label{p2:def:composition-versus-transvection}
Two operator triples must not be conflated.
\begin{enumerate}[label=(\roman*), leftmargin=2em]
\item Group-valued transfer composition on \(SU(2)\) is generated by one
regular chiral action.  In the Pauli dictionary below, the plus
chiral example is the geometric right regular action:
\[
 \hat T_i^{\,+}:=\frac{\hbar}{\Rmom}L_i^+,
 \qquad
 [\hat T_i^{\,+},\hat T_j^{\,+}]
 =\frac{i\hbar}{\Rmom}\epsilon_{ijk}\hat T_k^{\,+},
\]
The other regular action is equally available after its convention is fixed.
\item The diagonal homogeneous-space splitting of
\(\mathfrak{so}(4)\) instead gives
\[
 \hat P_i^{\,\mathrm{tr}}
 :=\frac{\hbar}{\Rmom}(L_i^+-L_i^-).
\]
The span of these transvections is not a Lie subalgebra:
\[
 [\hat P_i^{\,\mathrm{tr}},\hat P_j^{\,\mathrm{tr}}]
 =\frac{i\hbar}{\Rmom^2}\epsilon_{ijk}\hat J_k.
\]
\end{enumerate}
The first triple generates the intrinsic transfer-composition group.  The
second generates geodesic displacement in the symmetric-space description and
is the symmetry-momentum operator used below.  Theorem~
\ref{p2:thm:canonical-transvection-bridge} derives the second as the unique
unframed tangent-displacement representative.  It does not identify it with a
single chiral composition generator; the two operational roles remain
distinct.
\end{definition}

\begin{definition}[Derived homogeneous-space transvection operators]
\label{p2:def:compact-momentum-generators}
The operators \(\hat P_i\) used in the remaining spectral
formulas are the compact transvections
\(\hat P_i:=\hat P_i^{\,\mathrm{tr}}\) in the symmetric-pair splitting
\[
\mathfrak{so}(4)=\mathfrak{so}(3)\oplus\mathfrak p,
\qquad
\mathfrak p=\mathrm{span}\{\hat J_{i4}\}_{i=1}^3,
\qquad
\hat P_i=\frac{1}{\Rmom}\hat J_{i4}.
\]
They generate geodesic displacements on the compact momentum sphere, not the
intrinsic group-composition action of
Definition~\ref{p2:def:composition-versus-transvection} and not a global
abelian translation group.  In a local tangent chart, the unitary
\[
U(\mathbf a)=\exp\left(-\frac{i}{\hbar}\mathbf a\cdot\hat{\mathbf P}\right)
\]
implements infinitesimal geodesic displacement, while globally the compact
bracket remains nonzero:
\[
[\hat P_i,\hat P_j]
=\frac{i\hbar}{\Rmom^2}\epsilon_{ijk}\hat J_k.
\]
The relation has an algebraic flat contraction as
\(\Rmom^{-2}\to0\).  At fixed \(m\), however,
\(\hat P_i=(\hbar/\Rmom)(L_i^+-L_i^-)\to0\).  A nontrivial flat
translation representation would require a separately specified family of
representations and states with an appropriate scaling limit; no such family
is constructed or assumed here.
\end{definition}


\paragraph{Derivation of Specific Mappings}

\begin{proposition}[Transvection operators within the diagonal \texorpdfstring{$SU(2)$}{SU(2)} symmetric-pair complement]
\label{p2:prop:momentum-unique}
Let \(L_i^\pm\) denote the dimensionless chiral \(\mathfrak{su}(2)\) generators.  The physical diagonal rotation generators are
\[
\hat J_i:=\hbar(L_i^+ + L_i^-),
\]
and the complementary dimensionless vector triple is
\[
K_i:=L_i^+-L_i^-.
\]
Under the diagonal rotation algebra generated by \(\hat J_i\), both triples
\[
\mathrm{span}\{\hat J_i\}
\qquad\text{and}\qquad
\mathrm{span}\{K_i\}
\]
transform as vectors.  Within the fixed symmetric-pair splitting
\(\mathfrak{so}(4)=\mathfrak{so}(3)_{\rm diag}\oplus\mathfrak p\), the
complement \(\mathfrak p\) is a single irreducible diagonal vector
representation.  More precisely, let an admissible symmetry-momentum triple
be the image of the standard vector basis under a diagonal-
\(SU(2)\)-equivariant linear map into \(\mathfrak p\).  Then, after the
rotation generators are fixed to be \(\hat J_i\), every such triple has the
form
\[
\hat P_i = \kappa_P\,K_i,
\]
and evaluation normalization together with the momentum unit fixes
\[
\boxed{\hat P_i = \frac{\hbar}{\Rmom}\,(L_i^+ - L_i^-)}
\]
up to the conventional sign \(\hat P_i\mapsto -\hat P_i\).
\end{proposition}

\begin{proof}
On \(S^3\simeq SU(2)\), the left- and right-invariant vector fields generate
\[
\mathfrak{su}(2)_L\oplus\mathfrak{su}(2)_R \cong \mathfrak{so}(4).
\]
With dimensionless \(L_i^\pm\) satisfying
\[
[L_i^\pm,L_j^\pm]=i\epsilon_{ijk}L_k^\pm,
\qquad
[L_i^+,L_j^-]=0,
\]
one obtains
\[
[\hat J_i,\hat J_j]=i\hbar\,\epsilon_{ijk}\hat J_k,
\qquad
[\hat J_i,K_j]=i\hbar\,\epsilon_{ijk}K_k.
\]
Thus \(\mathrm{span}\{K_i\}\) is a vector module under the diagonal rotation algebra. Since
\[
\mathfrak{so}(4)=\mathrm{span}\{L_i^+ + L_i^-\}\oplus \mathrm{span}\{L_i^+-L_i^-\},
\]
and both summands are irreducible vector copies for the diagonal \(SU(2)\), the
complement satisfies
\[
\dim\operatorname{Hom}_{SU(2)_{\rm diag}}(V_1,\mathfrak p)=1.
\]
Schur's lemma therefore makes the complementary equivariant vector triple
unique up to scale within this fixed symmetric-pair decomposition.  Theorem~
\ref{p2:thm:canonical-transvection-bridge} fixes the relative left/right
coefficient by evaluation.  The overall normalization is fixed by the momentum unit
\[
\Punit=\hbar/\Rmom,
\]
and the remaining sign is conventional.
\end{proof}

\subparagraph{Supplementary cone/Hopf geometry.}
The cone-dilation, scaling-weight, Hopf/Reeb, Chern, and contact-form
calculations are supplementary to the finite-coordinate proof and are
deferred to Paper~II-F, Section~\ref{p2c:app:supplementary-geometric-interfaces}.

\subsubsection{Derived momentum-carrier representatives}
For the structural momentum-role carrier \(\mathcal H_p\cong\mathcal H_1\), the
momentum operators admit explicit \(15\times15\) Hermitian representatives.
The basis and matrices, including the diagonal form of \(\widehat P_3\), are
recorded in Paper~II-F, Section~\ref{p2:app:explicit-operators}; the main text uses only
their basis-independent algebraic consequences.

\subparagraph{Optional kinetic representative and scale firewall.}
An optional scaled Casimir/Laplacian representative is recorded in
Paper~II-F, Section~\ref{p2c:app:deferred-host-infrastructure}.  It neither selects the
seed nor supplies a Hamiltonian occurrence.  In particular,
\[
 \boxed{
 \mathcal H_p,\ m_p,\ S^3_B,\ \mathcal P_x,\ \Cl(0,3)
 \not\Rightarrow
 \Rmom^{\rm numerical},\
 \rho_{\rm prot}=\Rmom,\
 M,\
 H_{\rm kin}^{\rm physical}.}
\]
Numerical scale identification and a physical kinetic/Hamiltonian meaning
must be constituted by their downstream calibration or dynamical owner.


\section{Finite coordinate operators and ambient completions}
\label{paper2c:sec:organized-operators}

\subsection{Finite-symbol operators and the kinematic algebra}
\label{p2:sec:derived-algebra}
This section separates the operation-first momentum algebra from
completion-relative coordinate data.  It introduces neither dynamics nor
semiclassical assumptions.

\begin{table}[htbp]
\centering
\small
\begingroup
\renewcommand{\arraystretch}{1.3}
\begin{tabular}{|
  L{2.75cm}|
  L{3.75cm}|
  L{\dimexpr\linewidth-2.75cm-3.75cm-6\tabcolsep-4\arrayrulewidth\relax}|
}
\hline
\textbf{Geometric Input} & \textbf{Operator Form} & \textbf{Rigorous Consequence} \\
\hline
\multicolumn{3}{|c|}{\textbf{\textit{Commutation Relations}}} \\
\hline
Embedding
\newline coordinates & $\hat{x}_i$ (multiplication) & $[\hat{x}_i, \hat{x}_j] = 0$ \\
\hline
Chiral generators & \(\hat{P}_i = \frac{\hbar}{\Rmom}(L^+_i - L^-_i)\) & \([\hat{P}_i, \hat{P}_j] = \frac{i\hbar}{\Rmom^2}\epsilon_{ijk}\hat{J}_k\) \\
\hline
$SO(4)$
\newline fundamental rep & $[\hat{P}_i, \hat{x}_j]$ & $= -\frac{i\hbar}{\Rmom} \delta_{ij} \hat{x}_4$ \\
\hline
SO(3) subalgebra & $\hat{J}_i$ generators & $[\hat{J}_i, \hat{J}_j] = i\hbar\epsilon_{ijk}\hat{J}_k$ \\
\hline
Vector
\newline transformation & $[\hat{J}_i, \hat{P}_j]$,
\newline $[\hat{J}_i, \hat{x}_j]$ & $= i\hbar\epsilon_{ijk}\hat{P}_k$,
\newline $= i\hbar\epsilon_{ijk}\hat{x}_k$ \\
\hline
\end{tabular}
\caption[Core coordinate and momentum identities]{Core coordinate and momentum identities.}
\label{p2:tab:derived-algebra-complete}
\TableNotes{%
This table combines the II-B-owned rotation--transvection algebra with the
II-C-owned coordinate covariance and transvection--coordinate identities.
Commutativity of scalar multiplication coordinates is definitional.
Spectral and Hopf consequences are owned by Paper~II-D and are not duplicated
here.
}
\endgroup
\end{table}


\subsubsection{Momentum algebra}
\label{p2:sec:momentum-algebra}

The momentum operators defined in
Definition~\ref{p2:def:compact-momentum-generators} are
\begin{equation}
\hat{P}_i = \frac{\hbar}{\Rmom}(L^+_i - L^-_i),
\end{equation}
where \(L^{\pm}_i\) are the dimensionless chiral generators satisfying \([L^{\pm}_i,L^{\pm}_j]=i\epsilon_{ijk}L^{\pm}_k\) and \([L^+_i,L^-_j]=0\).

The following rotation--transvection relation is the II-B compact-algebra
export.  It is rederived here only to keep the normalization used by the
coordinate commutators locally explicit.

\begin{proposition}[Momentum Commutation Relations]
\label{p2:prop:momentum-commutation-relations}
The momentum operators satisfy
\begin{equation}
[\hat{P}_i, \hat{P}_j] = \frac{i\hbar}{\Rmom^2}\epsilon_{ijk}\hat{J}_k,
\end{equation}
where \(\hat{J}_k=\hbar(L^+_k+L^-_k)\) is the physical angular momentum generating the diagonal \(\SO(3)\subset\SO(4)\).
\end{proposition}

\begin{proof}
Direct computation using the chiral algebra:
\begin{align}
[\hat{P}_i, \hat{P}_j] &= \frac{\hbar^2}{\Rmom^2}[(L^+_i - L^-_i), (L^+_j - L^-_j)] \\
&= \frac{\hbar^2}{\Rmom^2}\bigl([L^+_i, L^+_j] + [L^-_i, L^-_j]\bigr) \\
&= \frac{\hbar^2}{\Rmom^2}\bigl(i\epsilon_{ijk}L^+_k + i\epsilon_{ijk}L^-_k\bigr) \\
&= \frac{i\hbar^2}{\Rmom^2}\epsilon_{ijk}(L^+_k + L^-_k) \\
&= \frac{i\hbar}{\Rmom^2}\epsilon_{ijk}\hat{J}_k.
\end{align}
\end{proof}

\begin{definition}[Generated level-one kinematic operation algebra]
\label{p2c:def:generated-kinematic-test-algebra}
On the derived carrier \(\mathcal H_p\cong\mathcal H_1\), put
\[
 \mathsf J_i:=\widehat J_i/\hbar,
 \qquad
 \mathsf P_i:=\Rmom\widehat P_i/\hbar,
\]
and define the generated unital \(*\)-algebra
\[
 \mathcal A_{\rm kin}
 :=\operatorname{alg}_{\CC}^{*}
 \bigl(\mathbf1,\mathsf J_1,\mathsf J_2,\mathsf J_3,
              \mathsf P_1,\mathsf P_2,\mathsf P_3\bigr)
 \subseteq\operatorname{End}(\mathcal H_1).
\]
This is the GP2 operation-first closure of the constituted
rotation--transvection operations.  Together with its positive unital
functionals and effects it supplies the canonical finite state/effect
calculus.  Calling any element an actually implemented apparatus test requires
the appropriate G0 occurrence certificate; the algebra is not defined to
contain every mathematically available endomorphism of the carrier.
\end{definition}

\begin{theorem}[Exact level-one kinematic algebra]
\label{p2c:thm:exact-level-one-kinematic-algebra}
Let
\[
 \Lambda_1
 :=\left\{(0,0),(\tfrac12,0),(0,\tfrac12),(\tfrac12,\tfrac12),
             (1,0),(0,1)\right\}
\]
be the multiplicity-free sector support of \(\mathcal H_1\).  Then
\[
 \boxed{
 \mathcal A_{\rm kin}
 =\bigoplus_{\lambda\in\Lambda_1}\operatorname{End}(V_\lambda)
 \cong
 \CC\oplus M_2(\CC)\oplus M_2(\CC)\oplus M_4(\CC)
     \oplus M_3(\CC)\oplus M_3(\CC).}
\]
Consequently
\[
 \dim_{\CC}\mathcal A_{\rm kin}=43
 <225=\dim_{\CC}\operatorname{End}(\mathcal H_1).
\]
\end{theorem}

\begin{proof}
The source definitions give
\[
 L_i^\pm=\tfrac12(\mathsf J_i\pm\mathsf P_i),
\]
so \(\mathcal A_{\rm kin}\) is the algebra generated by the two chiral
\(\mathfrak{su}(2)\) actions.  They preserve every irreducible
\((j_L,j_R)\) sector, giving the inclusion into the displayed block algebra.

Put
\[
 C_L:=\sum_i(L_i^+)^2,
 \qquad
 C_R:=\sum_i(L_i^-)^2,
 \qquad
 D:=C_L+10C_R.
\]
On the sectors in the displayed order, \(D\) has eigenvalues
\[
 0,\quad \tfrac34,\quad \tfrac{15}{2},\quad
 \tfrac{33}{4},\quad 2,\quad20,
\]
which are pairwise distinct.  Hence every central sector projector is the
generated polynomial
\[
 P_\lambda
 =\prod_{\substack{\mu\in\Lambda_1\\\mu\ne\lambda}}
   \frac{D-d_\mu\mathbf1}{d_\lambda-d_\mu}
 \in\mathcal A_{\rm kin}.
\]
On \(P_\lambda\mathcal H_1=V_\lambda\), the
\(\mathfrak{su}(2)_L\oplus\mathfrak{su}(2)_R\) representation is complex
irreducible.  Burnside's theorem therefore gives
\[
 P_\lambda\mathcal A_{\rm kin}P_\lambda
 =\operatorname{End}(V_\lambda).
\]
Because the central projectors themselves are generated, the blocks can be
selected and summed independently.  This proves equality, and
\[
 1^2+2^2+2^2+4^2+3^2+3^2=43
\]
gives the dimension.
\end{proof}

\begin{corollary}[Generated level-one tensor operator handoff]
\label{p2c:cor:generated-level-one-tt-operators}
The operation-generated $\mathcal A_{\rm kin}$ is the sector-preserving
$\mathcal A_1$ of Paper~I-A. Therefore
Corollary~\ref{p1:cor:level-one-shear-tt-operators} exports one $(1,1)$
shear operator type and one each of $(2,0),(0,2)$ at unchanged $m_p=1$.
These are operator directions of the generated algebra, not added object
sectors or tests from its full $225$-dimensional presentation envelope.
\end{corollary}
\begin{proof}
The preceding exact-algebra theorem identifies the two block algebras.
Apply the conjugation decomposition of the cited corollary.
\end{proof}
Paper~III-F's Corollary~\ref{p3f:cor:level-one-quadrupole-incidence}
supplies the actual Dirac action incidence of the two $M_3$ quadrupoles;
an apparatus implementation or source occupation keeps its own occurrence
requirements.


\begin{corollary}[Contained quantum-kinematic structural existence]
\label{p2c:cor:contained-quantum-kinematic-existence}
With the full positive-normalized state-space convention of Paper~II-A,
\[
 \boxed{\begin{gathered}
 \mathsf{QEX}_{\rm kin}:\\
 (\mathcal H_1,\mathcal A_{\rm kin},S(\mathcal A_{\rm kin}))\\
 \text{is a canonically generated noncommutative}\\
 \text{finite state--test system}.
 \end{gathered}}
\]
In particular it contains incompatible effects and a non-Boolean projection
lattice.  This conclusion uses neither coordinate normalization nor a full
\(M_{15}(\CC)\) operator envelope.  Its empirical event, preparation, and
probability readings are G0 certificates rather than additional algebraic
postulates.
\end{corollary}

\begin{proof}
On \(V_{\frac12,0}\), one has
\[
 \mathsf P_i=\mathsf J_i=\sigma_i/2,
 \qquad
 [\mathsf P_1,\mathsf P_2]=i\sigma_3/2\ne0.
\]
The central projector \(P_{\frac12,0}\) belongs to
\(\mathcal A_{\rm kin}\) by
Theorem~\ref{p2c:thm:exact-level-one-kinematic-algebra}.  Therefore the
block-supported projections
\[
 E_z:=\tfrac12P_{\frac12,0}(\mathbf1+2\mathsf P_3)P_{\frac12,0},
 \qquad
 E_x:=\tfrac12P_{\frac12,0}(\mathbf1+2\mathsf P_1)P_{\frac12,0}
\]
are effects in \(\mathcal A_{\rm kin}\) and satisfy \([E_z,E_x]\ne0\).
Paper~II-A, Theorem~\ref{p2:thm:canonical-finite-state-test-evaluation},
supplies their separating state evaluations.  Assigning those evaluations
empirical quantum-event probability meaning and selecting an actually prepared
state remain later G0 readout/preparation certificates; the noncommutative
finite state--test system itself is already derived.
\end{proof}

\begin{remark}[Geometric realization and raw-field convention]
Use the Pauli embedding \(g=y_4I+i y_i\sigma_i\), and fix
\(\tau_i=-i\sigma_i\), so that
\([\tau_i,\tau_j]=2\epsilon_{ijk}\tau_k\).  Define
\[
 X_i^R f(g):=\left.\frac{d}{dt}\right|_{t=0}f(e^{-t\tau_i}g),
 \qquad
 X_i^L f(g):=\left.\frac{d}{dt}\right|_{t=0}f(ge^{t\tau_i}).
\]
The minus sign in the first pullback is part of the convention.  These fields
satisfy
\[
 [X_i^R,X_j^R]=2\epsilon_{ijk}X_k^R,\qquad
 [X_i^L,X_j^L]=2\epsilon_{ijk}X_k^L,\qquad [X_i^R,X_j^L]=0.
\]
Let \(\mathcal V_i^L\) and \(\mathcal V_i^R\) differentiate
\(f(e^{it\sigma_i}g)\) and \(f(ge^{it\sigma_i})\), respectively.
Then \(X_i^R=\mathcal V_i^L\), \(X_i^L=-\mathcal V_i^R\), and the
realization compatible with the embedding-coordinate commutators is
\[
 \boxed{L_i^+=\frac{i}{2}X_i^L,\qquad
        L_i^-=\frac{i}{2}X_i^R.}
\]
Both \(L^\pm\) obey \([L_i^\pm,L_j^\pm]=i\epsilon_{ijk}L_k^\pm\).
The geometric left circle \(g\mapsto e^{it\sigma_i}g\) has Hermitian
generator \(-2L_i^-\), while the geometric right circle
\(g\mapsto ge^{it\sigma_i}\) has generator \(2L_i^+\).
Thus the plus/minus labels name the fixed algebraic chiral factors.
An atlas whose pair labels instead name geometric multiplication factors
uses \((j_{\rm geom,L},j_{\rm geom,R})=(j_R,j_L)\).
The compact transvection remains
\(\widehat P_i=(\hbar/\Rmom)(L_i^+-L_i^-)\);
the coordinate brackets and abstract compact algebra are unchanged.
\end{remark}

This is the constituted compact isotropic momentum--rotation algebra
compatible with bounded momentum spectra, with the transvection representative
derived in Theorem~\ref{p2:thm:canonical-transvection-bridge} and
Proposition~\ref{p2:prop:momentum-unique}.  The statement is the compact
constitution of the \((\hat J,\hat P)\) Lie-algebra classification, not a
classification of all possible coordinate--momentum Heisenberg deformations.

\subsubsection{Finite coordinate operators and compact algebra}
\label{p2:sec:position-operators}

For ambient-extension and recovery purposes, introduce the auxiliary
multiplicity-free representation universe and its integer-indexed mathematical
finite family.  These objects are not the GP2 physical generated closure of
the momentum seed; the physical seed remains
\(\mathcal H_p\cong\mathcal H_1\).  II-C now defines the Hilbert direct-sum
completion
\[
\mathcal H_\infty
=
\widehat{\bigoplus}_{j_L,j_R\in\frac12\mathbb Z_{\ge0}}
V_{j_L}\otimes V_{j_R}
\cong
\overline{\bigcup_{m\in\mathbb Z_{\ge0}}\mathcal H_m},
\]
and its finite initial segments
\[
\mathcal H_m
=
\bigoplus_{j_L+j_R\le m}V_{j_L}\otimes V_{j_R}.
\]
This \(\mathcal H_\infty\) is local ambient-completion notation; it is not a
claim that Paper~II-B selected an infinite realized carrier.  Likewise the
family \(\mathcal H_m\) below is auxiliary mathematical infrastructure, not
the constitutional source of the physical \(\mathcal H_1\) seed.
Whenever \(m=1\) is used for the momentum role below, it refers to the
already derived \([\mathrm{HOST}]\) carrier \(\mathcal H_p\cong\mathcal H_1\);
statements for symbolic \(m\) remain auxiliary \([\mathrm{FIN}{-}m]\) results.
The scalar coordinate symbol determines multiplication on the diagonal
Peter--Weyl subcarrier.  It does not determine the off-diagonal continuation
to all of \(\mathcal H_\infty\).  This section first classifies that freedom
and then proves that, once an ambient action is fixed, its return to
\(\mathcal H_m\) is unique.

\begin{proposition}[Projected coordinate graph at \(m=1\)]
\label{p2:prop:h1-coordinate-coupling-graph}
The retained-edge statement is the projected coordinate-coupling result
already proved in Paper~I-A~\cite{Kugelmann2026a}; its short proof is reproduced
here so that II-C's subsequent forest classification and completion analysis
are locally checkable.

Let \(V_{\rm coord}=(\tfrac12,\tfrac12)\).  Projection of
\(V_{\rm coord}\otimes\mathcal H_1\) back to \(\mathcal H_1\) leaves exactly
\[
(0,0)\leftrightarrow(\tfrac12,\tfrac12),\qquad
(\tfrac12,\tfrac12)\leftrightarrow(1,0),\ (0,1),
\]
and
\[
(\tfrac12,0)\leftrightarrow(0,\tfrac12).
\]
Thus the projected graph has the invariant components
\[
\mathcal H_{\rm geo}
=
(0,0)\oplus(\tfrac12,\tfrac12)\oplus(1,0)\oplus(0,1)
\]
and
\[
\mathcal H_{\rm sp}
=
(\tfrac12,0)\oplus(0,\tfrac12).
\]
\end{proposition}

\begin{proof}
The Clebsch--Gordan rule is
\[
(j_L,j_R)\otimes(\tfrac12,\tfrac12)
=
\bigoplus_{\epsilon_L,\epsilon_R=\pm1}
(j_L+\tfrac{\epsilon_L}{2},j_R+\tfrac{\epsilon_R}{2}),
\]
with negative labels omitted.  Applying it to the six summands of
\(\mathcal H_1\) and retaining only total label at most one gives precisely
the displayed edges.  The omitted targets are
\((1,1)\), \((\tfrac32,\tfrac12)\),
\((\tfrac12,\tfrac32)\), \((1,\tfrac12)\), and
\((\tfrac12,1)\).
\end{proof}

\begin{corollary}[Complete finite coordinate classification at \(m=1\)]
\label{p2:cor:h1-coordinate-classification}
Fix normalized intertwiners on the four unoriented edges of
Proposition~\ref{p2:prop:h1-coordinate-coupling-graph}.  The projected graph
is a forest with two connected components.  Consequently sector rephasings
remove every finite edge phase, and there is no cycle sign or phase holonomy
on \(\mathcal H_1\).

Matching the diagonal scalar multiplication action fixes the reduced
coefficient on
\[
        (0,0)\leftrightarrow(\tfrac12,\tfrac12).
\]
Modulo sector rephasing, every compatible Hermitian equivariant coordinate
vector on \(\mathcal H_1\) is therefore classified by exactly three
nonnegative reduced magnitudes,
\[
        r_{1,0},\qquad r_{0,1},\qquad \varrho_{\rm sp},
\]
on the respective edges from \((\tfrac12,\tfrac12)\) to \((1,0)\), from
\((\tfrac12,\tfrac12)\) to \((0,1)\), and between the two protected chiral
sectors.  This is a classification of finite coordinate actions, not a
uniqueness theorem for their ambient completions.
Equivalently,
\[
 \boxed{
 \mathfrak M_{\rm coord}(\mathcal H_1)/{\sim}
 \cong\mathbb R_{\ge0}^{3},
 \qquad
 [\widehat X]\longleftrightarrow
 (r_{1,0},r_{0,1},\varrho_{\rm sp}) .}
\]
This cone is the complete residual finite realization family after the
unmarked sector rephasings have been quotiented.  Its three reduced
magnitudes are not presentation labels: their zero/nonzero pattern can change
the generated coupling graph and operator closure.  Therefore a theorem that
depends on a particular raw finite coordinate action must state or derive its
point in this cone; GP3/no-new-choice does not select one.

Invariant cycle products remain genuine on larger ambient graphs, and a
nonzero product carries a phase holonomy.  For example,
\[
(0,\tfrac12)\to(\tfrac12,0)\to(1,\tfrac12)
\to(\tfrac12,1)\to(0,\tfrac12)
\]
is a four-cycle whose terminal total weight is \(3/2\).  It is absent from
\(\mathcal H_1\) and is present by \(\mathcal H_2\), the next member of the
integer-indexed family.
\end{corollary}

\begin{proof}
Hermiticity identifies the coefficients on the two orientations of each
edge, leaving one complex reduced coefficient per unoriented edge.  On a
forest, choose a root in each connected component and rephase sectors
successively away from the root; this makes every edge coefficient
nonnegative.  Equivalently, the vertex--edge incidence map has rank equal to
the number of edges, so it has no cycle-phase quotient.  Scalar matching
fixes one of the four resulting magnitudes, leaving the three displayed
parameters.  The final cycle follows by applying the same
Clebsch--Gordan rule to its four displayed sectors.
\end{proof}

\begin{theorem}[Imported intrinsic tangential round-symbol class]
\label{p2c:thm:intrinsic-tangential-coordinate-class}
Paper~I-B, Theorem~\ref{p1:thm:intrinsic-tangential-coordinate-class}, proves
the following geometric export.  Use the retained smooth realization
\(J_R:\mathcal H_1\to\mathcal L_R\), identified with the generated carrier by
Theorem~\ref{p2b:thm:canonical-retained-smooth-realization}.  On differential
forms plus spinors over the fixed-radius sphere put
\[
 \mathscr D_{\rm LC}:=i(d-d^\dagger)\oplus D_{S^3}
\]
and define the finite first-symbol quadruple
\[
 X_a^D:=\frac{\Rmom}{\sqrt3}\,
 J_R^*\,i[\mathscr D_{\rm LC},M_{x_a}]\,J_R,
 \qquad a=1,\ldots,4.
\]
An overall Clifford-sign convention changes only removable sector phases.
The class \(\mathsf C_D:=[X^D]\) is independent of every zero-order term
added to \(\mathscr D_{\rm LC}\), is unique on this intrinsic tangential
symbol branch, and has the sectorwise quadratic fingerprint
\[
 Q_D:=\sum_{a=1}^4(X_a^D)^2
 =\left(
 \Rmom^2,\ \frac{11\Rmom^2}{12},\
 \frac{4\Rmom^2}{9},\ \frac{4\Rmom^2}{9},\
 \frac{3\Rmom^2}{4},\ \frac{3\Rmom^2}{4}
 \right)
\]
in the ordered sectors
\[
 (0,0),\quad(\tfrac12,\tfrac12),\quad(1,0),\quad(0,1),
 \quad(\tfrac12,0),\quad(0,\tfrac12).
\]
Fix the raw exterior-edge atlas convention invariantly by
\[
 Q|_{(1,0)}=2r_{1,0}^2I,\qquad
 Q|_{(0,1)}=2r_{0,1}^2I,\qquad
 Q|_{S_+\oplus S_-}=4\varrho_{\rm sp}^2I.
\]
Then
\[
 \boxed{
 (r_{1,0},r_{0,1},\varrho_{\rm sp})_D
 =\left(\frac{\sqrt2\Rmom}{3},
        \frac{\sqrt2\Rmom}{3},
        \frac{\sqrt3\Rmom}{4}\right).}
\]
In particular all four allowed forest edges are nonzero and
\(\varrho_{\rm sp}>0\).
\end{theorem}

\begin{proof}
Paper~I-B proves the form-edge coefficients, the two exact spin-gap equations,
the fingerprint, and phase-qualified uniqueness.  The II-B theorem identifies
the complete generated carrier with the same retained smooth unitary class,
so no new geometric premise or recalculation is required here.
\end{proof}

\begin{corollary}[Cone--tangential inequivalence]
\label{p2c:cor:cone-tangential-coordinate-inequivalence}
Paper~I-A's ambient metric-cone class \(\mathsf C_{\rm cone}\) has fingerprint
\[
 Q_{\rm cone}=\left(
 \Rmom^2,\ \Rmom^2,\ \frac{\Rmom^2}{2},\
 \frac{\Rmom^2}{2},\ \Rmom^2,\ \Rmom^2\right).
\]
Hence
\[
 \boxed{\mathsf C_D\ne\mathsf C_{\rm cone}}
\]
in \(\mathfrak M_{\rm coord}(\mathcal H_1)/\!\sim\).  The discrepancy is
exactly the radial Clifford component retained by the cone symbol and removed
by fixed-radius pullback; it is not a normalization phase.
\end{corollary}

\begin{proof}
This is Paper~I-A,
Corollary~\ref{p1:cor:cone-tangential-coordinate-inequivalence}, imported into
the generated-carrier notation.  Paper~I-B also proves that the distinction is
the radial/tangential ancestry split and not a normalization convention.
\end{proof}

\begin{remark}[Constitution boundary of the finite coordinate operation]
\label{p2c:rem:v265-coordinate-constitution-boundary}
This is the \emph{coordinate-stage} allocation already used by
Corollary~\ref{p2c:cor:intrinsic-coordinate-class-corealization}, not a
retrospective addition to Paper~II-A's initial momentum-seed grammar.
The actual round momentum role is the fixed-radius, oriented metric host;
Paper~II-B has already co-realized its finite smooth module
\(J_R:\mathcal H_1\to\mathcal L_R\).  At this later stage, the coordinate
operation target is constituted as the exact finite matrix-element image of
the \emph{intrinsic first metric symbol} applied to the four generated
embedding-coordinate covectors, with the existing scalar-to-vector comparison
fixing the sole common normalization.  Before evaluating those images, the
typed operations are: fixed-radius pullback to \(T^*S^3\), metric Clifford
symbol in the smooth forms/spinors operation sort, the already retained
\(J_R\)-matrix-element pairing, and finite Riesz representation in the
declared Hermitian equivariant coordinate-operation sort.  The output is
retained for subsequent coordinate-role comparisons.  These are the inputs,
operations, retention, and target of the frozen \(\Sigma_{\rm II}^{\rm coord}\)
word, not a claim that arbitrary ambient pointwise multiplication, arbitrary
compression, or a Hamiltonian is generated.

The allocation is explicit constitutional content at this chronological
stage; its values, uniqueness, and absence of a free scale are mathematical
consequences, not additional constitutive parameters.  Host preservation
applies to this \emph{declared} word and its finite matrix elements.  It does
not require downstream GP2 to regenerate the content allocated here, and it
does not imply that every natural operator on the ambient host is a physical
constructor.  The radial-cone symbol uses a separate normal covector
\(dr\), absent from the actual fixed-radius symbol input, and therefore
is not a second output of this same constructor.  A particular apparatus,
prepared state, test performance, or numerical readout is not constituted by
this structural allocation.
\end{remark}

\begin{proposition}[Unique exact finite inheritance of the metric first symbol]
\label{p2c:prop:v265-matrix-element-host-inheritance}
Let \(\mathscr K_R=L^2\Omega^\bullet(S^3;\mathbb C)\oplus
L^2(S^3,\mathbb S)\) be the smooth-parent Hilbert completion and
\(J_R:\mathcal H_1\hookrightarrow\mathscr K_R\) the retained isometry,
whose finite range is \(\mathcal L_R\).  Let
\(C_a=i[\mathscr D_{\rm LC},M_{x_a}]\) be the Hermitian fixed-radius
metric first-symbol operators.  Put \(s_R=\Rmom/\sqrt3\).  For all
\(u,v\in\mathcal H_1\), the coordinate-stage pairing is
\[
 B_a(u,v)=s_R\langle J_Ru,C_aJ_Rv\rangle_{\mathscr K_R}.
\]
There exists exactly one Hermitian finite operator \(Y_a\) representing this
pairing for every pair \((u,v)\), namely
\[
 \boxed{Y_a=s_RJ_R^*C_aJ_R=X_a^D.}
\]
Consequently the physically co-realized coordinate \emph{class} on the
nonobstructed coordinate stage is \(\mathsf C_D\); its reduced magnitudes
are those of Theorem~\ref{p2c:thm:intrinsic-tangential-coordinate-class}.
In particular no raw equivariant family member with independently rescaled
protected block \(X_a^{(\lambda)}=X_a^{D,{\rm ext}}\oplus
\lambda X_a^{D,{\rm spin}}\), \(\lambda\geq0\), can represent the
same retained pairing unless \(\lambda=1\).  This result is relative to
the declared metric-symbol word and does not claim selection among all raw
equivariant coordinate completions by the carrier alone.
\end{proposition}

\begin{proof}
On the fixed-radius round parent, the metric Clifford relation gives
\(\{C_a,C_b\}=2\langle d_{S^3}x_a,d_{S^3}x_b\rangle I\) and
\(\sum_a C_a^2=3I\).  The scalar-to-vector edge matching of
Theorem~\ref{p2c:thm:intrinsic-tangential-coordinate-class} fixes
\(s_R=\Rmom/\sqrt3\) independently of the protected block.  The finite
Riesz representation theorem applied to each \(B_a\) gives the displayed
\(Y_a\), and nondegeneracy of the Hilbert pairing makes it unique:
if \(\langle u,(Y_a-Z_a)v\rangle=0\) for all \(u,v\), then \(Y_a=Z_a\).
The inherited \(\Spin(4)\) action transports the entire tuple, and changes
of unmarked II-B sector phases conjugate it without changing its class.
Paper~I-B's exact symbol computation therefore determines all three reduced
magnitudes, including \(\varrho_{\rm sp}=\sqrt3\Rmom/4\).

For the protected-rescaling family the invariant compressed square is
\[
 \left.\sum_a(X_a^{(\lambda)})^2\right|_{S_+\oplus S_-}
 =\frac{3\lambda^2\Rmom^2}{4}I_4.
\]
The exact inherited pairing has the same expression with \(\lambda=1\).
Since \(\lambda\geq0\), equality of the operators (hence of this necessary
invariant) requires \(\lambda=1\).  The family is a valid counterexample to
a \emph{seed-alone} selection inference but not to the constituted
metric-symbol word.  Conversely, this word does not preserve arbitrary
Clifford products under compression: writing \(P_R=J_RJ_R^*\),
\[
 \sum_a(X_a^D)^2
 =\Rmom^2 I_{\mathcal H_1}
 -\frac{\Rmom^2}{3}\sum_aJ_R^*C_a(I-P_R)C_aJ_R.
\]
The summand is positive because it is
\(\sum_a((I-P_R)C_aJ_R)^*((I-P_R)C_aJ_R)\).
This exact positive leakage identity distinguishes finite matrix-element
faithfulness from the false assertion that the finite subspace is invariant
under every \(C_a\).  It has the sectorwise values already recorded in the
intrinsic fingerprint.  Thus neither a full ambient operator representation
nor a new numerical scale is inserted by finite inheritance.
\end{proof}

\begin{corollary}[Structural co-realization of the intrinsic finite class]
\label{p2c:cor:intrinsic-coordinate-class-corealization}
Use the coordinate-stage allocation and retained matrix-element target
of Remark~\ref{p2c:rem:v265-coordinate-constitution-boundary}.
The target is the Hermitian equivariant coordinate-operator class on the
already co-realized carrier, not the entire raw completion cone.  The complete
occurring root tuple
\[
 \mathfrak R_D=
 (S^3_A(\Rmom),g_{\Rmom},\mathcal H_1,J_R,
   \sigma_1(\mathscr D_{\rm LC}),
   e_0\longleftrightarrow E_1\ \text{under }M_{x_a})
\]
has the required output \(F_D(\mathfrak R_D)=\mathsf C_D\).  In the explicitly frozen
II-C coordinate-stage grammar, set
\[
 w_D:=
 \left(
 \mathfrak R_D
 \xrightarrow{\ i_{\Rmom}^*,\,c_T,\,J_R^*(\cdot)J_R\ }
 \mathsf C_D
 \right).
\]
The fixed-radius pullback, tangential Clifford-symbol, and retained
finite matrix-element operations are constituted in this chronological
coordinate stage with their output and reuse sorts prior to evaluating the
coefficients.  Proposition~\ref{p2c:prop:v265-matrix-element-host-inheritance}
proves that the required output is uniquely \(\mathsf C_D\).  Thus
\[
 \operatorname{Orig}_{\Sigma_{\rm II}^{\rm coord}}
 (w_D;\mathfrak R_D,\mathfrak M_{\rm coord}(\mathcal H_1)/{\sim})
\]
holds.  Paper~II-A's rooted generated-image theorem therefore gives
\[
 \PhysOcc(\mathscr O_{\rm tr})
 \Longrightarrow
 \operatorname{CoReal}\!\left(
 S^3_A(\Rmom),\mathcal H_1,\mathsf C_D
 \mid\mathscr O_{\rm tr}\right).
\]
This is finite structural co-realization only.  It selects no member of the
ambient-lift fibre, no POVM with first moments \(X_a^D\), no prepared state,
and no physical kinetic/Hamiltonian owner.
\end{corollary}

\begin{proof}
Every input port of \(w_D\) is filled by the displayed complete root tuple;
its symbol and compression outputs are retained and reusable in their native
sorts.  Candidate independence and naturality follow from the oriented round
metric, the II-B realization, and the principal symbol.  These facts verify
the frozen word but do not license an undeclared ambient constructor; the
coordinate-stage constitution of the complete word and target is the
allocative premise, as recorded in
Remark~\ref{p2c:rem:v265-coordinate-constitution-boundary}.  The coefficient
\(\Rmom/\sqrt3\) introduces no scale: it is uniquely forced by conservative
transport of the already occurring normalized scalar comparison
\(M_{x_a}e_0=(\Rmom/2)\phi_a\) into the declared finite target.  Theorem~
\ref{p2c:thm:intrinsic-tangential-coordinate-class} proves that the exact
essential image is the singleton class \(\mathsf C_D\), with every allowed
finite edge nonzero and no coefficient outside \(\mathcal H_1\).  The
displayed II-B finite realization is GP1-admissible and GP2-exact, so the FRC
obstruction alternative is false.  Theorem~
\ref{p2a:thm:rooted-generated-image-corealization} then retains and
co-realizes the required output \(\mathsf C_D\).  Zero joint GP2 defect is a
property of that realization, not a premise licensing the constructor.
Ambient coefficients, apparatus data, state preparation, and dynamical
ownership are absent from the constructor and therefore cannot be inherited
by this proof.
\end{proof}

\subsection{Intrinsic coordinate comparison and fixed-radius ancestry}

\begin{corollary}[Complete coordinate-role comparison word]
\label{p2c:cor:intrinsic-joint-comparison-role-export}
On the occurring intrinsic coordinate branch, retain the complete Paper-I-A tuple
\[
 D_0=(u=1,H_{\rm Hur}=2T,n_0=\mathbf i,
             h_0=(1+\mathbf i+\mathbf j+\mathbf k)/2)
\]
and its simultaneous $SU(2)$-conjugation orbit, with defining spin transport on
the positive Hopf line. At the retained axial mark $q=h$, the licensed word
\[
 (X_D,u,q)\longmapsto A_D(u,q)
 \longmapsto\mathcal T_u(q;\,\cdot\,)
 \longmapsto(\mathcal O,N_{\mathcal O},m_{\mathcal O},
                         \nabla_{N_{\mathcal O}})
\]
has the geometric targets and transport laws of
Theorems~\ref{p1:thm:intrinsic-even-hopf-hurwitz-comparison},
\ref{p1:thm:hopf-hurwitz-quadratic-comparison},
\ref{p1:thm:joint-orbit-quadratic-descent} and
\ref{p1:thm:joint-comparison-unitary-connection}.
In particular, $\|m_{\mathcal O}\|^2=1/2$, and the tensor-line connection is
the unique Hermitian connection preserving that nonzero quadratic.
This connection is flat, with holonomy $-1$ around the generator of
$\pi_1(\mathcal O)$. Its relative determinant presentation is the natural isomorphism in
Proposition~\ref{p1:prop:axial-relative-determinant-line}.
Each target is co-realized with the full specified input in its declared
geometric sort.
\end{corollary}
\begin{proof}
The intrinsic coordinate-class export supplies
$\varrho_D^2=3\Rmom^2/16>0$, which cancels in the normalized even word.
Write $A_D(u,h)=r+v$ in a compatible quaternion model. Paper I-A gives
\[
 \|m\|^2=4r^2\bigl(|v|^2-(n\cdot v)^2\bigr).
\]
For $D_0$, $r=1/2$ and the transverse squared norm is $1/2$, so
$\|m\|^2=1/2$ throughout the simultaneous-conjugation orbit. The cited
descent and connection theorems therefore apply on their exact domain.
Products, invariant contractions, restriction to the marked lines and the
displayed homogeneous-orbit construction are the explicit natural operations
proved in Paper I-A. Composing their source and target certificates with the
frozen constructor-origin and rooted co-realization rules proves the export.
Nonparallelness alone is not used as a substitute for the nonzero scalar
factor in the displayed norm. Section~\ref{p2c:app:marked-comparison-domain}
records the exact nonzero stratum and the sign-line connection argument.
\end{proof}

This export does not assert a physical charged realization on a new
base, or retention of only one projected output after its companions
have been forgotten.  Such an application needs its own target action,
base map and complete retained diagram.  The reciprocal position role
is still generated by evaluation of coordinate roles.  Evaluation
constructs the role and its action; it does not constrain the later
position--momentum configuration space to a graph \(x=x(p)\).


\begin{corollary}[Physical fixed-radius ancestry exhaustion]
\label{p2c:cor:physical-fixed-radius-ancestry-exhaustion}
Relative to the current II-A/II-B momentum-role constructor grammar, the
essential image in the finite round-symbol target is
\[
 \boxed{\operatorname{EssIm}_{\rm phys}^{\rm coord}
 =\{\mathsf C_D\}.}
\]
The intrinsic class is therefore the unique physically generated finite
symbol descendant of the occurring momentum root, up to the already
quotiented sector phases.  The ambient class \(\mathsf C_{\rm cone}\) remains
a Paper-I-A host comparison unless a cone/radial parent is separately
constituted and actualized.  No generic raw atlas point or ambient-lift fibre
member is selected by this conclusion.
\end{corollary}

\begin{remark}[Fixed-radius nondegeneracy is not source actuality]
\label{p2c:rem:v270-fixed-radius-source-firewall}
The singleton physical coordinate image $\{\mathsf C_D\}$ is nondegenerate;
in particular the raw completion face with vanishing protected coordinate
scale is not a physical fixed-radius descendant.  This does not select a
nonzero fermion source member.  Coordinate nondegeneracy and an odd finite
$\mathsf{SrcVal}$ are distinct target sorts, and Paper~II-A's source-incidence
certificate is required for any cross-sort promotion.
\end{remark}


\begin{proof}
The current occurring root contains
\[
 S^3_A(\Rmom),\quad g_{\Rmom},\quad\mathcal H_1,\quad J_R,\quad
 \sigma_1(\mathscr D_{\rm LC}),
\]
and the normalized scalar comparison.  Hence the typed word
\[
 \mathfrak R_D
 \xrightarrow{\ i_{\Rmom}^*,\,c_T,\,J_R^*(\cdot)J_R\ }
 \mathsf C_D
\]
is licensed, retained, and co-realized by
Corollary~\ref{p2c:cor:intrinsic-coordinate-class-corealization}.
Paper~I-A,
Theorem~\ref{p1:thm:primitive-round-symbol-ancestry-exhaustion}, proves that
the only other output of the primitive mathematical grammar is the ambient
parallel cone word.  That word requires the radial covector \(dr\) and the
ambient flat-cone symbol target, neither of which occurs in the frozen
II-A/II-B root or is produced by its finite-coordinate constructor.  It is
therefore host mathematics at this stage, not a second physical output.
Paper~I-A's sharpness remark shows that every further equivariant branch
requires an additional radial weight or other enlarged ancestry datum.  Thus
the current physical grammar has no omitted or excess finite symbol class.
\end{proof}


\subsection{Sampled coordinate representations and coupling-graph closures}
\begin{corollary}[Inherited realization of the full sampled carrier]
\label{p2c:cor:v204-sampled-realization}
On the root already realized in
Corollary~\ref{p2c:cor:intrinsic-coordinate-class-corealization}, choose the
same retained Hurwitz support and geometric frame representatives and put
$T_R=\mathsf S_RJ_R:\mathcal H_1\to\mathcal V_R$.
The sampled intrinsic coordinates are exactly
\[
 \widetilde X_a=T_RX_a^DT_R^*,\qquad
 \widetilde{\mathcal A}=T_R\mathcal A T_R^*
\]
for every already generated algebra $\mathcal A$ on this carrier.
All its relations, gradings, projectors, commutants and generated geometric
unitaries are represented faithfully. This is a presentation of the same
realized finite object and retains the co-realization of $\mathsf C_D$;
there is no new root occurrence or additional carrier to constitute.
\end{corollary}
\begin{proof}
The cited co-realization corollary supplies the source root and
Corollary~\ref{p2c:cor:physical-fixed-radius-ancestry-exhaustion} its
intrinsic ancestry. Paper~I-B, Theorem~\ref{p1:thm:v203-full-carrier-sampling}
and Corollary~\ref{p1:cor:v203-full-symbol-transport}, construct its concrete
sampling map and coordinate matrices. The full sampling and image-algebra
transport are collected in Paper~1B,
Theorems~\ref{p1b:thm:full-sampling} and~\ref{p1b:thm:image-transport}
\cite{Kugelmann2026Paper1B}. Apply them to $T_R=\mathsf S_RJ_R$.
The inverse is $T_R^*$ on the sampled image; its identity in the ambient
sample space is $T_RT_R^*$, and all constraints and commutants are taken
on that image. Thus the transported diagram has the existing root,
sector phases and support-frame orbit.
\end{proof}
The scalar comparison algebra on $\mathcal P_x(2)$ is transported by its
own rank-fourteen sampling isometry, with the exact compressed symbols of
Theorem~\ref{p1:thm:v203-scalar-observable-compiler} and
Corollary~\ref{p1:cor:v203-quadratic-top-block-correction}. The full
$\mathcal H_1$ and scalar comparison band remain distinct. Subsequent
action, state and process maps use this inherited representation of their
existing operators, not a second admission of the intrinsic coordinate root.

\begin{remark}[Normalization conversion for intrinsic coordinates]
Paper~I-B, Remark~\ref{p1:rem:intrinsic-edge-normalization-conversion}, owns the
following convention dictionary.  The raw wedge map satisfies
\(TT^*=2I_{\Lambda^2}\).  If instead one calls
the coisometric Clebsch--Gordan map \(T_0=2^{-1/2}T\) ``normalized,'' the same
intrinsic edge is written with coefficient \(2\Rmom/3\) rather than
\(\sqrt2\Rmom/3\).  The invariant quadratic fingerprint and the atlas
definitions in the theorem remove this conventional ambiguity.
\end{remark}

\begin{theorem}[Exact sector-coupling graph closure]
\label{p2c:thm:sector-coupling-graph-closure}
Let
\[
 \mathcal H=\bigoplus_{\lambda\in\Lambda}V_\lambda,
 \qquad
 \mathcal A_0=\bigoplus_{\lambda\in\Lambda}\operatorname{End}(V_\lambda),
\]
and let \(X\subseteq\operatorname{End}(\mathcal H)\) be a generated operator
family.  With \(P_\lambda\) the block projectors, define the undirected
nonzero-coupling graph \(G_X\) by
\[
 \{\lambda,\mu\}\in E(G_X)
 \quad\Longleftrightarrow\quad
 \exists x\in X:\
 P_\lambda xP_\mu\ne0\ \text{ or }\ P_\mu xP_\lambda\ne0.
\]
For \(C\in\pi_0(G_X)\), put
\(\mathcal H_C:=\bigoplus_{\lambda\in C}V_\lambda\).  Then
\[
 \boxed{
 \operatorname{alg}_{\CC}^{*}(\mathcal A_0,X)
 =\bigoplus_{C\in\pi_0(G_X)}\operatorname{End}(\mathcal H_C).}
\]
In particular,
\[
 \boxed{
 \operatorname{alg}_{\CC}^{*}(\mathcal A_0,X)
 =\operatorname{End}(\mathcal H)
 \quad\Longleftrightarrow\quad
 G_X\text{ is connected}.}
\]
\end{theorem}

\begin{proof}
No generator has a block between two different connected components, so every
generator and every element of the generated \(*\)-algebra preserves each
\(\mathcal H_C\).  This gives the inclusion into the displayed componentwise
full algebra.

Conversely, let \(\{\lambda,\mu\}\) be an edge and, after taking an adjoint if
needed, choose a nonzero block
\[
 T=P_\lambda xP_\mu
 \in\operatorname{Hom}(V_\mu,V_\lambda).
\]
Because the full endpoint matrix algebras already lie in \(\mathcal A_0\),
the span
\[
 \operatorname{End}(V_\lambda)\,T\,
 \operatorname{End}(V_\mu)
\]
is all of \(\operatorname{Hom}(V_\mu,V_\lambda)\): matrix units at the two
endpoints isolate one nonzero entry of \(T\) and move it to any ordered pair
of basis vectors.  Adjoint closure supplies the reverse rectangular block.
Multiplying these rectangular blocks along a path in \(G_X\) generates every
\(\operatorname{Hom}(V_\nu,V_\lambda)\) between vertices in the same
component.  Together with the diagonal blocks, these are all matrix units of
\(\operatorname{End}(\mathcal H_C)\).  Hence no componentwise endomorphism is
omitted, proving equality and the connectedness criterion.
\end{proof}

\begin{corollary}[Present coordinate closure and the exact full-operator-envelope condition]
\label{p2c:cor:present-coordinate-envelope-condition}
For the level-one coordinate family
\(\widehat X=(\widehat X_1,\ldots,\widehat X_4)\), let
\(G_{\widehat X}\) retain precisely the nonzero blocks of the allowed graph in
Proposition~\ref{p2:prop:h1-coordinate-coupling-graph}.  Then
\[
 \operatorname{alg}_{\CC}^{*}
 (\mathcal A_{\rm kin},\widehat X_1,\ldots,\widehat X_4)
 =\bigoplus_{C\in\pi_0(G_{\widehat X})}
   \operatorname{End}(\mathcal H_C).
\]
Every present coordinate completion preserves
\[
 \mathcal H_{\rm geo}
 =(0,0)\oplus(\tfrac12,\tfrac12)\oplus(1,0)\oplus(0,1),
 \qquad \dim\mathcal H_{\rm geo}=11,
\]
and
\[
 \mathcal H_{\rm sp}
 =(\tfrac12,0)\oplus(0,\tfrac12),
 \qquad \dim\mathcal H_{\rm sp}=4.
\]
Therefore, as an unconditional algebraic inclusion,
\[
 \boxed{
 \operatorname{alg}_{\CC}^{*}
 (\mathcal A_{\rm kin},\widehat X_1,\ldots,\widehat X_4)
 \subseteq M_{11}(\CC)\oplus M_4(\CC).}
\]
If all four allowed forest edges have nonzero reduced coefficients, equality
holds.  For any later generated operator family \(Y\), the full
\(M_{15}(\CC)\) operator envelope is generated exactly when the total nonzero-coupling
graph \(G_{\widehat X\cup Y}\) is connected.
\end{corollary}

\begin{proof}
Theorem~\ref{p2c:thm:exact-level-one-kinematic-algebra} supplies the full local
sector blocks, and Theorem~\ref{p2c:thm:sector-coupling-graph-closure} applies
to the actual nonzero coordinate blocks.  The allowed edges of
Proposition~\ref{p2:prop:h1-coordinate-coupling-graph} never connect
\(\mathcal H_{\rm geo}\) to \(\mathcal H_{\rm sp}\), proving the inclusion.
When every allowed edge is nonzero, each of those two subgraphs is connected,
so the graph theorem gives \(M_{11}\oplus M_4\).  Applying the same theorem
after adjoining \(Y\) gives the final if-and-only-if statement.  Thus a later
generated operator may derive the full operator envelope only by deriving a
nonzero bridge.  Its physical typing is decided by the later grading and
quantum-role theorems; an independently supplied carrier-resolution interface
would be a distinct event extension, not evidence that the coordinate family
generated such a bridge.
\end{proof}


\begin{theorem}[Multiplicity-stable native test representations and coordinate-component locking]
\label{p2c:thm:v294-native-multiplicity-algebra}
Let $T_p$ be the six-type support of
Theorem~\ref{p2:thm:momentum-host-H1}, and let a finite unitary
$\Spin(4)$ realization have
\[
 \mathcal H=\bigoplus_{\lambda\in T_p}V_\lambda\otimes M_\lambda,
 \qquad m_\lambda:=\dim M_\lambda\geq1.
\]
Suppose its represented rotation and transvection operations are precisely
$L_i^\pm\otimes I_{M_\lambda}$ on each isotypic summand, without
additional multiplicity-sensitive generator ports. Then their unital
$*$-algebra is \emph{exactly}
\[
 \boxed{\mathcal A_{\rm kin}(\mathcal H)
 =\bigoplus_{\lambda\in T_p}
 (\End(V_\lambda)\otimes I_{M_\lambda})
 \cong\mathcal A_{\rm kin},\qquad
 \dim_{\mathbb C}\mathcal A_{\rm kin}(\mathcal H)=43.}
\]
In particular, all native rotation--transvection polynomial tests are
multiplicity blind on \emph{every} such realization, not only on the
fifteen-dimensional one-copy presentation. The state/effect quotient of
Theorem~\ref{p2b:thm:v293-operational-multiplicity-quotient} therefore
requires no extra test-blindness hypothesis \emph{for this specified
kinematic algebra}. This does not characterize the complete generated test
family of an arbitrary whole structured candidate.

More generally, suppose an actual coordinate-role implementation is a
faithful unital $*$-representation of the intrinsic coordinate algebra
$\mathcal A_X=M_{11}(\mathbb C)\oplus M_4(\mathbb C)$ of
Corollary~\ref{p2c:cor:present-coordinate-envelope-condition} and
Theorem~\ref{p2c:thm:operational-symmetry-filtration}, with the
$\Spin(4)$ action given by the represented native inclusion. Then there
are positive integers $k_g,k_s$ and a unitary identification
\[
 \boxed{\mathcal H\cong
 (\mathcal H_{\rm geo}\otimes\mathbb C^{k_g})
 \oplus(\mathcal H_{\rm sp}\otimes\mathbb C^{k_s}),\qquad
 \pi(\mathcal A_X)=
 (M_{11}\otimes I_{k_g})\oplus(M_4\otimes I_{k_s}).}
\]
Thus the four geometric types have the \emph{same} multiplicity $k_g$
and the two spinorial types have the \emph{same} multiplicity $k_s$.
The coordinate-state quotient is explicitly
\[
 Q_X(\rho)=\operatorname{Tr}_{\mathbb C^{k_g}}(P_g\rho P_g)
 \oplus\operatorname{Tr}_{\mathbb C^{k_s}}(P_s\rho P_s),
\]
which preserves the expectation of \emph{every} represented
$M_{11}\oplus M_4$ test; the blockwise maximally mixed section is a
canonical right inverse. No equation $k_g=k_s$ follows from the present coordinate algebra.
If a further \emph{physically test-typed} full $M_{15}(\mathbb C)$
extension has a faithful representation with the same embedded
$\Spin(4)$ action, the six multiplicities must all coincide; this
condition alone still does not force that common multiplicity to be one.
All these conclusions are about represented exact algebras, not a
newly generated GP3 whole-candidate comparison.
\end{theorem}
\begin{proof}
The chiral generators are $L_i^\pm=\tfrac12(\mathsf J_i\pm\mathsf P_i)$.
As in Theorem~\ref{p2c:thm:exact-level-one-kinematic-algebra}, the
polynomial $D=C_L+10C_R$ has six distinct eigenvalues. Its Lagrange
polynomials yield the central isotypic projections $P_\lambda$, now
acting as $I_{V_\lambda}\otimes I_{M_\lambda}$, irrespective of
$m_\lambda$. On each $V_\lambda$ the chiral generators act
irreducibly; Burnside's theorem makes their algebra $\End(V_\lambda)$.
Every generator is the identity on $M_\lambda$, so no product or adjoint
can create a multiplicity-sensitive matrix. Conversely the generated
central projections isolate all six full $\End(V_\lambda)$ blocks.
This gives the stated exact equality and dimension; the operational
quotient follows from Theorem~\ref{p2b:thm:v293-operational-multiplicity-quotient}.

For the coordinate statement, use matrix units of each simple block
$M_d$. In any unital $*$-representation, the images of the $d$ diagonal
matrix units are equivalent mutually orthogonal projections and have
a common finite rank $k$; the off-diagonal matrix units identify their
ranges. Hence the represented block is unitarily $M_d\otimes I_k$.
Faithfulness forces both central blocks to be nonzero and thus
$k_g,k_s\ge1$. The embedded $\Spin(4)$ action then decomposes each
geometric irreducible with multiplicity $k_g$ and each spinorial one
with multiplicity $k_s$. The state/effect assertion is the blockwise partial-trace identity;
tensoring each block density with $I_{k_C}/k_C$ supplies the right
inverse. The same matrix-unit argument applied to
$M_{15}$ forces one common multiplicity. Nothing in that argument
requires $k=1$ or constructs a GP3 comparison.
\end{proof}

\begin{proposition}[Finite-basis certificate for generated process descent]
\label{p2c:prop:v294-finite-process-descent}
Let $\mathcal A\subseteq M_N(\mathbb C)$ be any of the above
multiplicity-blind represented finite $C^*$-algebras, and let
$\{F_1,\ldots,F_d\}$ be a \emph{real linear basis} of its Hermitian
part, where $d=43$ for $\mathcal A_{\rm kin}$ and $d=137$ for
$M_{11}\oplus M_4$. Suppose a specified process grammar is generated
by a family $\{\Phi^*_a\}_a$ of unital completely positive Heisenberg
maps on $M_N$ using finite compositions and convex mixtures. Then
\[
 \boxed{\begin{gathered}
 \text{Every generated process descends to }\mathcal A\\
 \iff\;\Phi^*_a(F_j)\in\mathcal A\\
 \text{for every generator }a\text{ and }j=1,\ldots,d.
 \end{gathered}}
\]
For every such process, the quotient on states intertwines its action
as in Theorem~\ref{p2b:thm:v293-operational-multiplicity-quotient}.
In particular, any channel with a Kraus representation whose operators
all belong to $\mathcal A$, and any unitary evolution generated by a
self-adjoint element of $\mathcal A$, automatically descend.
Testing merely a multiplicative generating subset of $\mathcal A$ is
\emph{not} sufficient for arbitrary completely positive maps.

Indeed, take $\mathcal A=M_2\otimes I_2\subset M_4$ and Pauli matrices
$X,Y,Z$. For any $0<\varepsilon<\tfrac14$ the map
\[
 \Phi^*(T)=\tfrac14\Tr(T)I_4
 +\varepsilon\Tr((Y\otimes I_2)T)(I_2\otimes Z)
\]
is unital and completely positive. It maps the Hermitian algebra
generators $X\otimes I_2$ and $Z\otimes I_2$ into $\mathcal A$,
but sends their product direction $Y\otimes I_2$ to
$4\varepsilon(I_2\otimes Z)\notin\mathcal A$.
This counterexample concerns a permitted mathematical CP map, not an
assertion that it is an occurring Born--Hopf process.
\end{proposition}
\begin{proof}
Linear preservation on the Hermitian basis is equivalent by
complexification to $\Phi_a^*(\mathcal A)\subseteq\mathcal A$.
The latter property is preserved by composition and convex mixture.
For Kraus operators $K_r\in\mathcal A$, each $K_r^*aK_r$
lies in $\mathcal A$; their sum does as well. Exponentiating a
self-adjoint $h\in\mathcal A$ remains in $\mathcal A$, giving the
unitary case.
Conversely every generator belongs to the generated process family,
so preservation by every word forces each generator to pass the test.
The equivalence with state-level quotient descent is the exact
Heisenberg-invariance criterion of
Theorem~\ref{p2b:thm:v293-operational-multiplicity-quotient}, applied
to the appropriate finite direct-sum algebra.

For the example, the unnormalized Choi matrix in the input/output
tensor ordering is
\[
 J(\Phi^*)=\tfrac14I_{16}
 +\varepsilon(Y\otimes I_2)^{\mathsf T}\otimes(I_2\otimes Z).
\]
The second factor is Hermitian with spectrum $\{+1,-1\}$, hence
$J(\Phi^*)\ge0$ for the stated range. The map is unital because
$\Tr(Y\otimes I_2)=0$; complete positivity follows from Choi
positivity. Pauli orthogonality gives zero on $X\otimes I_2$ and
$Z\otimes I_2$, while $\Tr[(Y\otimes I_2)^2]=4$ yields the displayed
nonblind image. Since $Y$ is a scalar multiple of the product $XZ$,
the final claim follows.
\end{proof}

\begin{theorem}[Canonical finite port-defect certificate and exact distinguishing witness]
\label{p2c:thm:v295-complete-port-twirl}
In the exact multiplicity-stable kinematic representation of
Theorem~\ref{p2c:thm:v294-native-multiplicity-algebra}, put
$H=\bigoplus_{\lambda}V_\lambda\otimes M_\lambda$,
$m_\lambda=\dim M_\lambda$, and
$\mathcal A=\bigoplus_\lambda\End(V_\lambda)\otimes I_{M_\lambda}$.
For $T\in\End(H)$ define the finite, representation-intrinsic map
\[
 \mathsf E_{\mathcal A}(T)
 :=\bigoplus_\lambda
 \frac{1}{m_\lambda}\Tr_{M_\lambda}
 (P_\lambda T P_\lambda)\otimes I_{M_\lambda}.
\]
It is the unique ordinary-trace-preserving conditional expectation onto
$\mathcal A$ and also the twirl over the compact multiplicity commutant
$\prod_\lambda U(M_\lambda)$ acting by conjugation.  It is a unital,
completely positive, idempotent $\mathcal A$-bimodule map and the
Hilbert--Schmidt orthogonal projection onto $\mathcal A$.
Consequently the nonnegative, basis-independent defect
\[
 \boxed{\Delta_{\mathcal A}(T)
 :=\|T-\mathsf E_{\mathcal A}(T)\|_{\rm HS}^2}
\]
vanishes if and only if $T\in\mathcal A$.  A specified collection of
additional frozen Hermitian test ports $\{T_a\}_{a\in J}$ is wholly
multiplicity blind if and only if every $\Delta_{\mathcal A}(T_a)=0$;
when $J$ is finite their sum is a single exact scalar certificate.
For a Hermitian $T$ with $\Delta_{\mathcal A}(T)>0$, set
$X=T-\mathsf E_{\mathcal A}(T)$ and $N=\dim H$. Then, for
$0<\varepsilon\leq (N\|X\|_{\rm op})^{-1}$,
\[
 \begin{gathered}
 \rho_\pm=I_H/N\pm\varepsilon X
 \quad\text{are density matrices with}\\
 Q(\rho_+)=Q(\rho_-),\\
 \boxed{\Tr[(\rho_+-\rho_-)T]
 =2\varepsilon\Delta_{\mathcal A}(T)>0,}
 \end{gathered}
\]
where $Q$ is the quotient of
Theorem~\ref{p2b:thm:v293-operational-multiplicity-quotient}.
If $\{F_j\}_{j=1}^{43}$ is any real Hermitian basis of $\mathcal A$,
then an arbitrary specified unital completely positive Heisenberg map
$\Phi^*$ descends to that quotient exactly when
\[
 \boxed{\sum_{j=1}^{43}
 \Delta_{\mathcal A}(\Phi^*(F_j))=0.}
\]
The analogous complete-coordinate statement uses the represented
$M_{11}\oplus M_4$ algebra, the two component multiplicity spaces of
Theorem~\ref{p2c:thm:v294-native-multiplicity-algebra}, and a
137-element Hermitian basis. In that setting the twirl acts on the
\emph{two coordinate components}, not separately on their six
irreducible labels. These certificates do not turn a mathematical
candidate port into an occurring GP2 test or process.
\end{theorem}
\begin{proof}
Conjugation by independent unitaries on the $M_\lambda$ fixes precisely
$\mathcal A$. Haar averaging kills inter-isotypic blocks by the independent
central phases, and on each surviving diagonal block its invariants are
$\End(V_\lambda)\otimes I_{M_\lambda}$ by Schur's lemma. Partial trace
identifies this average with the displayed formula. A unitary twirl is
unital and completely positive, preserves the ordinary trace, is
idempotent and $\mathcal A$-bimodular. Orthogonality follows from
$\Tr[(T-\mathsf E_{\mathcal A}T)^*A]=0$ for every $A\in\mathcal A$;
therefore any trace-preserving conditional expectation with this range
has the same Hilbert--Schmidt projection and is unique.
For Hermitian $T$, $X=X^*$ is trace zero and orthogonal to $\mathcal A$.
The stated bound makes both $I_H/N\pm\varepsilon X$ positive and trace
one. Their difference is in the kernel of the state quotient because
it has zero expectation against every $\mathcal A$ test, whereas their
$T$-expectations differ by
$2\varepsilon\Tr(XT)=2\varepsilon\Tr(X^2)$.
Finally, a completely positive Heisenberg generator preserves the whole
finite-dimensional $*$-algebra exactly when it preserves its real
Hermitian basis.  Nonnegativity of each defect makes vanishing of the
sum equivalent to every basis image lying in the algebra. Apply
Proposition~\ref{p2c:prop:v294-finite-process-descent}.
The two-component coordinate algebra obeys the same twirl and matrix-unit
argument, with its faithful-representation hypothesis retained.
\end{proof}


\begin{remark}[Typing boundary of coupling-graph closure]
\label{p2c:rem:coupling-graph-typing-boundary}
Theorem~\ref{p2c:thm:sector-coupling-graph-closure} is an operator-algebra
closure theorem.  It does not promote every generated bridge to an observable.
In particular, if a later central grading types a bridge as odd, connected
closure gives a graded field envelope; the parity-preserving observable
algebra is its even fixed point unless a separate parity-breaking event
interface is supplied.
\end{remark}

\begin{theorem}[Operational symmetry filtration and native state dimensions]
\label{p2c:thm:operational-symmetry-filtration}
For the generated intrinsic tuple put
$\mathcal A_X=\operatorname{alg}_{\mathbb C}^*(\mathcal A_{\rm kin},X^D)
=M_{11}\oplus M_4$.
The conjugation actions and normalized state-space affine dimensions are
\[
\begin{array}{c|c|c|c}
 \text{test algebra}&\text{conjugation kernel in }\Spin(4)
       &\text{effective action}&\dim_{\mathbb R}\text{ states}\\\hline
 \mathcal A_{\rm kin}&\mathbb Z_2\times\mathbb Z_2
       &SO(3)\times SO(3)&42\\
 \mathcal A_X&\mathbb Z_{2,\rm diag}&SO(4)&136\\
 M_{15}\ \text{if generated}&\{1\}&\Spin(4)&224
\end{array}
\]
Thus carrier faithfulness and faithfulness on a restricted test algebra
are different properties.  The affine restriction-map kernels from the
native coordinate state space to the operation state space, and from the
full matrix state space to the native one, have real dimensions $94$ and
$88$, respectively.  The latter directions are inter-sector coherences.
These are affine-kernel counts, realized by full-dimensional fibers over
interior states; positivity can shrink fibers over boundary states.
They do not select any extension or preparation.
\end{theorem}
\begin{proof}
The operation algebra is the direct sum of full irreducible blocks of
sizes $1,2,2,4,3,3$.  Central elements of $\Spin(4)$ act scalarly on
each block and are invisible there; conversely invisibility on both chiral
spinor blocks forces each $SU(2)$ factor to be central.  For
$\mathcal A_X=M_{11}\oplus M_4$, an invisible unitary must be scalar
on each of these two blocks.  The trivial sector fixes its geometric
scalar to one, both adjoints force the factor elements to be central,
and the nonzero scalar/vector edge forces their signs to agree.
The diagonal central subgroup is also sufficient.  On $M_{15}$ a scalar
action has scalar one because of the trivial sector; the two faithful
chiral spinors then force the identity element.  The state dimension of
$\bigoplus_jM_{n_j}$ is $\sum_j n_j^2-1$, giving
$43-1$, $137-1$, and $225-1$.  Restriction of normalized Hermitian
functionals is surjective, so subtraction gives the two affine-kernel
dimensions.  Around an interior density matrix, sufficiently small
Hermitian perturbations in these kernels remain positive; boundary
densities need not admit all these perturbations.
\end{proof}

\begin{corollary}[Native channels preserve the central population]
\label{p2c:cor:native-central-population-conservation}
Let $P_g,P_s$ be the central projections of $M_{11}\oplus M_4$.
Every trace-preserving channel with Kraus operators in this algebra
preserves $p=\operatorname{Tr}(\rho P_g)$, as does every Hamiltonian
evolution generated in the algebra.  Such dynamics cannot select a unique
global equilibrium population or force $p=11/15$ from arbitrary sources.
A population-changing claim requires its actual exchange operation.
\end{corollary}
\begin{proof}
For $\sum_\alpha K_\alpha^\dagger K_\alpha=I$ and
$[K_\alpha,P_g]=0$,
\[
 \operatorname{Tr}\!\left(P_g\sum_\alpha K_\alpha\rho K_\alpha^\dagger\right)
 =\operatorname{Tr}\!\left(\rho P_g
                \sum_\alpha K_\alpha^\dagger K_\alpha\right)
 =\operatorname{Tr}(\rho P_g).
\]
The unitary case follows by taking one Kraus operator.  Each invariant
finite block has a stationary state for a fixed channel, and their convex
combinations retain arbitrary $p$; for Hamiltonian evolution one may use
each block's normalized identity.  No unique global population follows.
\end{proof}

\paragraph{Alternative finite closure diagnostic.}
\label{p2c:rem:e9-coordinate-closure-diagnostic}
For the weaker former carrier signature, the least fixed-action closure of
\(\mathcal E_9\) is the exact four-case graph determined by the zero/nonzero
pattern of \((r_{1,0},r_{0,1})\).  Its projection to the coordinate atlas is
bijective and therefore selects no completion; an independently generated
requirement that the closure equal \(\mathcal E_9\) or \(\mathcal H_1\)
instead cuts out the corresponding fibre.  The complete statement and proof
are Proposition~\ref{p2c:prop:e9-coordinate-closure} in
Paper~II-F, Section~\ref{p2c:app:e9-diagnostic}.  This diagnostic neither reselects the
already derived physical host \(\mathcal H_p\cong\mathcal H_1\) nor turns
\(\mathcal E_9\) into a competing momentum carrier.

\subsection{Ambient coordinate completions and unique retractions}

\begin{definition}[Classical and ambient coordinate actions]
\label{p2:def:xp-operators}
On \(S^3(\Rmom)\subset\mathbb R^4\), let \(x_a\) be the real embedding
coordinates, so \(\sum_a x_a^2=\Rmom^2\).  Their scalar action is
\(M_{x_a}\) on \(L^2(S^3)\).  An \emph{ambient coordinate completion} is a
bounded self-adjoint \(\Spin(4)\)-vector
\[
\mathsf X=(\mathsf X_1,\ldots,\mathsf X_4)
\]
on \(\mathcal H_\infty\) whose diagonal scalar compression is \(M_x\).
For a fixed completion, the finite coordinates are
\[
\widehat X_a:=\Pi_m\mathsf X_a\Pi_m.
\]
No nonlinear identity such as
\(\widehat X_4=(\Rmom^2-\sum_{i=1}^3\widehat X_i^2)^{1/2}\) follows after
compression.
\end{definition}

\begin{remark}[Ambient multiplication does not reopen the seed frontier]
\label{p2c:rem:ambient-multiplication-seed-guard}
The multiplication operators \(M_{x_a}\) are smooth-parent/ambient operations
introduced for the II-C completion problem.  Their target lies in the smooth
function representation and is not a finite momentum-seed carrier sort.
Accordingly their use here does not retroactively enlarge the frozen GP2 seed
frontier of II-A and does not generate the symmetric \((1,1)\) sector as new
momentum-seed carrier content.
\end{remark}

\begin{proposition}[Existence and scalar-symbol nonuniqueness of ambient completions]
\label{p2:prop:ambient-completion-existence-nonuniqueness}
Let
\[
 \mathcal H_{\rm diag}
 :=
 \widehat{\bigoplus}_{j\in\frac12\mathbb Z_{\ge0}}
 V_j\otimes V_j
 \cong L^2(S^3),
 \qquad
 \mathcal H_\infty
 =
 \mathcal H_{\rm diag}\oplus\mathcal H_{\rm diag}^{\perp}.
\]
Then
\[
        \mathsf X_a^{(0)}:=M_{x_a}\oplus0
\]
is a bounded self-adjoint \(\Spin(4)\)-equivariant ambient coordinate
completion.

Let
\[
 \mathcal S_{\rm prot}
 =
 V_{\frac12,0}\oplus V_{0,\frac12}
 \subset\mathcal H_{\rm diag}^{\perp},
\]
let \(C_a\) be a Hermitian Euclidean Clifford vector on \(\mathcal S_{\rm prot}\), and
extend it by zero to \(\mathcal H_\infty\), writing the extension as
\(\widetilde C_a\).  For every \(\lambda\ge0\),
\[
        \mathsf X_a^{(\lambda)}
        :=
        \mathsf X_a^{(0)}+\lambda\widetilde C_a
\]
is another bounded self-adjoint equivariant ambient completion with the same
diagonal scalar compression \(M_{x_a}\).  Its protected compression satisfies
\[
 P_{\mathcal S_{\rm prot}}\mathsf X_a^{(\lambda)}P_{\mathcal S_{\rm prot}}
 =\lambda C_a,
 \qquad
 P_{\mathcal S_{\rm prot}}
 \left(\sum_{a=1}^4(\mathsf X_a^{(\lambda)})^2\right)
 P_{\mathcal S_{\rm prot}}
 =4\lambda^2 I_{\mathcal S_{\rm prot}}.
\]
Consequently the scalar symbol proves existence of an ambient completion but
determines neither the ambient action nor the protected radius.
\end{proposition}

\begin{proof}
Multiplication by the real bounded coordinate \(x_a\) is self-adjoint, has
norm at most \(\Rmom\), and transforms as the vector representation on the
Peter--Weyl diagonal.  Its zero extension to the invariant orthogonal
complement preserves those properties and has the required scalar
compression.  The Clifford vector is bounded, self-adjoint, and equivariant
on the invariant finite-dimensional block \(\mathcal S_{\rm prot}\); its zero extension
therefore has the same properties.  Since
\(\mathcal S_{\rm prot}\perp\mathcal H_{\rm diag}\), adding
\(\lambda\widetilde C_a\) does not change the diagonal scalar compression.
The final identity follows from
\(\sum_a C_a^2=4I_{\mathcal S_{\rm prot}}\).
\end{proof}

\begin{proposition}[Every compatible finite action has a bounded ambient extension]
\label{p2:prop:finite-action-ambient-extension}
Fix \(m\), and let
\[
        A=(A_1,\ldots,A_4)
\]
be a Hermitian \(\Spin(4)\)-equivariant vector on \(\mathcal H_m\).  Let
\(P_{\rm diag}\) project onto \(\mathcal H_{\rm diag}\), put
\(P_m^{\rm sc}:=P_{\rm diag}\Pi_m=\Pi_mP_{\rm diag}\), and suppose the
finite diagonal scalar block agrees with compressed multiplication:
\[
P_m^{\rm sc} A_aP_m^{\rm sc}
=
P_m^{\rm sc}\mathsf X_a^{(0)}P_m^{\rm sc}.
\]
Then \(A\) is the compression of a bounded self-adjoint equivariant ambient
completion.
\end{proposition}

\begin{proof}
On \(\mathcal H_m\) set
\[
        B_a:=A_a-\Pi_m\mathsf X_a^{(0)}\Pi_m
\]
and extend \(B_a\) by zero on \(\mathcal H_m^\perp\).  The finite-dimensional
operator \(B_a\) is bounded and self-adjoint.  Since \(\mathcal H_m\) and its
orthogonal complement are invariant, its zero extension is still an
equivariant vector.  The scalar-block hypothesis gives
\(P_{\rm diag}B_aP_{\rm diag}=0\).  Therefore
\[
        \mathsf X_a:=\mathsf X_a^{(0)}+B_a
\]
is a bounded self-adjoint equivariant ambient completion with scalar
compression \(M_{x_a}\), and
\[
        \Pi_m\mathsf X_a\Pi_m=A_a.
\]
The construction proves existence of an extension; it does not make that
extension unique.
\end{proof}

\begin{theorem}[Canonical baseline-relative zero-tail ambient section]
\label{p2c:thm:canonical-zero-tail-ambient-section}
Under the hypotheses of
Proposition~\ref{p2:prop:finite-action-ambient-extension}, let
\(\iota_m:\mathcal H_m\hookrightarrow\mathcal H_\infty\) be the canonical
inclusion, \(\iota_m^*=\Pi_m\), and put
\[
 B_a:=A_a-\Pi_m\mathsf X_a^{(0)}\Pi_m.
\]
Define
\[
 \boxed{
 s_m(A)_a:=\mathsf X_a^{(0)}
 +\iota_mB_a\iota_m^*.}
\]
Then \(s_m(A)\) is a bounded self-adjoint equivariant ambient completion,
has the declared scalar symbol, and is an affine right inverse of compression:
\[
 \boxed{\Pi_ms_m(A)_a\Pi_m=A_a,\qquad
 \pi_m\circ s_m=\operatorname{id}.}
\]
It is the unique lift whose correction
\(D_a=s_m(A)_a-\mathsf X_a^{(0)}\) has no domain or range outside the retained
band,
\[
 D_a=\Pi_mD_a\Pi_m.
\]
The section is canonical relative to the fixed pair
\((\Pi_m,\mathsf X^{(0)})\).  It is not a physical ambient-action selector
and does not assert commutativity, the exact sphere relation, locality, or a
single level-independent all-band reconstruction from one cutoff. A supplied
compatible family has a unique bounded ambient extension exactly when its
component norms are uniformly bounded; its compressions and baseline-relative
sections then converge strongly. Section~\ref{p2c:app:zero-tail-scope}
proves this criterion and records the conditional Hilbert--Schmidt minimum.
\end{theorem}

\begin{proof}
The finite-dimensional \(B_a\) is bounded and self-adjoint, and extension by
zero across the invariant orthogonal complement preserves equivariance.  Its
diagonal scalar block vanishes by compatibility, so adding it to
\(\mathsf X^{(0)}\) preserves the scalar symbol.  Compression gives
\[
 \Pi_ms_m(A)_a\Pi_m
 =\Pi_m\mathsf X_a^{(0)}\Pi_m+B_a=A_a.
\]
If another correction satisfies \(D_a=\Pi_mD_a\Pi_m\), the fixed compression
forces \(D_a=B_a\), proving uniqueness in the zero-tail class.  Affineness is
immediate.
\end{proof}

\begin{remark}[Two-stage coordinate-completion realization family]
\label{p2c:rem:two-stage-coordinate-realization-family}
Let \(\mathfrak M_{\rm coord}^{\rm amb}\) denote the compatible bounded
self-adjoint equivariant ambient completions with the declared scalar symbol.
Compression to the derived level-one carrier induces a surjective realization
map onto the finite atlas,
\[
 \boxed{
 \mathfrak M_{\rm coord}^{\rm amb}
 \xrightarrow{\ \pi_1\ }
 \mathfrak M_{\rm coord}(\mathcal H_1)/{\sim}
 \cong\mathbb R_{\ge0}^{3},
 \qquad
 \pi_1(\mathsf X)=(r_{1,0},r_{0,1},\varrho_{\rm sp}).}
\]
Proposition~\ref{p2:prop:finite-action-ambient-extension} proves that every
compatible finite point has at least one ambient lift.  The fibre of
\(\pi_1\) need not be a singleton: higher ambient Wigner--Eckart coefficients
and cycle data are not fixed by the finite compression.  Thus a
completion-dependent theorem must distinguish whether it needs only the
finite cone point or a full ambient lift.
Both \(\mathsf C_{\rm cone}\) and \(\mathsf C_D\) lie in the finite target,
but their ancestry statements differ.  Paper~I-B derives both geometric
classes and proves their inequivalence.  Corollaries~
\ref{p2c:cor:intrinsic-coordinate-class-corealization} and
\ref{p2c:cor:physical-fixed-radius-ancestry-exhaustion} physically inherit
only the intrinsic class from the occurring momentum root.  Surjectivity of
\(\pi_1\) supplies lifts of either mathematical class but selects no lift and
does not identify the two classes.
\end{remark}

\subparagraph{Ambient host factorization.}
The multiplicity-free symmetric-module classification, including
\(V_{\rm coord}\cong S_L\otimes S_R\subset\operatorname{Sym}^2(S_L\oplus S_R)\),
is proved in Paper~II-F, Section~\ref{p2c:app:deferred-host-infrastructure}.  It is host
infrastructure, not unrestricted physical GP2 closure.

\begin{theorem}[Ambient completion freedom]
\label{p2:thm:derived-ambient-coordinate-action}
Choose normalized intertwiners, including a fixed phase convention,
\[
T_{\mu\lambda}\in
\operatorname{Hom}_{\Spin(4)}
(V_{\rm coord}\otimes V_\lambda,V_\mu)
\]
for every allowed coordinate edge.  Because the ambient envelope is
multiplicity-free and, more specifically, the Clebsch--Gordan rule
\[
V_{\rm coord}\otimes(V_{j_L}\otimes V_{j_R})
=
\bigoplus_{\epsilon_L,\epsilon_R=\pm1}
V_{j_L+\epsilon_L/2}\otimes V_{j_R+\epsilon_R/2}
\]
(with negative labels omitted) contains every allowed target with
multiplicity one, any equivariant coordinate vector has, on the algebraic
direct sum, channel form
\[
\mathsf X
=
\sum_{\lambda,\mu}t_{\mu\lambda}T_{\mu\lambda}.
\]
These channel data define a bounded completion precisely when the algebraic
operator extends boundedly to the Hilbert direct sum.
Hermiticity relates the coefficients on reversed edges, but coefficients on
inequivalent unoriented edges remain independent.  Matching the scalar symbol
fixes only those combinations visible after diagonal scalar compression.
It does not fix the off-diagonal Wigner--Eckart coefficients.

Sector rephasings change individual edge phases.  Relative to the fixed
intertwiner convention, the complex product of edge coefficients around a
closed oriented channel cycle is invariant under those rephasings and may
vanish.  On the subgraph of nonzero coefficients, the phase of that product
is defined and gives the cycle sign or phase holonomy.  Consequently neither
equivariance, the symmetric-module factorization, nor the scalar symbol proves
a unique ambient non-diagonal action.
\end{theorem}

\begin{proof}
The displayed Clebsch--Gordan decomposition gives
\[
\dim\operatorname{Hom}_{\Spin(4)}
(V_{\rm coord}\otimes V_\lambda,V_\mu)\in\{0,1\}.
\]
The Wigner--Eckart theorem
\cite{Wigner1959,fulton1991representation} therefore gives one scalar
coefficient for each allowed edge, yielding the displayed sum.  Taking
adjoints reverses an edge and imposes the Hermiticity
relation, but does not relate unrelated edges.  A sector rephasing
\(V_\lambda\mapsto e^{i\theta_\lambda}V_\lambda\) sends
\(t_{\mu\lambda}\) to
\(e^{i(\theta_\mu-\theta_\lambda)}t_{\mu\lambda}\); the factors telescope
around a cycle.  Thus the complex cycle product is always invariant, while
its phase is meaningful only when the product is nonzero.  Scalar compression
sees only diagonal-to-diagonal paths, which proves the claimed limitation.
\end{proof}


\begin{corollary}[Protected phase freedom and normalization conventions]
\label{p2:cor:protected-phase-normalization-freedom}
Set
\[
\Gamma_{\rm sp}:=C_1C_2C_3C_4,
\qquad
D_a:=i\Gamma_{\rm sp}C_a,
\qquad
\{C_a,C_b\}=2\delta_{ab}I_4,
\]
in a fixed Euclidean chiral Clifford representation on
\(\mathcal H_{\rm sp}\).  Every Hermitian \(\Spin(4)\)-equivariant protected
vector has the form
\begin{equation}
\label{p2:eq:protected-two-parameter-vector}
\widehat X_a^{\rm sp}=\alpha C_a+\beta D_a,
\qquad
\alpha,\beta\in\mathbb R.
\end{equation}
Its phase-invariant scale is
\[
\varrho_{\rm sp}:=\sqrt{\alpha^2+\beta^2}.
\]
The relative chiral unitary
\begin{equation}
\label{p2:eq:relative-chiral-unitary}
U(\theta):=\exp\!\left(\frac{i\theta}{2}\Gamma_{\rm sp}\right)
\end{equation}
commutes with \(\Spin(4)\) and rotates the coefficient pair.  For
\(\varrho_{\rm sp}>0\), one can choose \(\theta\) so that
\[
U(\theta)\widehat X_a^{\rm sp}U(\theta)^\dagger=\varrho_{\rm sp} C_a.
\]
The same statement is trivial for \(\varrho_{\rm sp}=0\).  We call this the
\emph{relative chiral phase gauge} and write
\(\gamma:=\varrho_{\rm sp}\) only in that gauge.  It is not an \(O(4)\)
frame rotation.

The phase gauge removes only the direction of the pair
\((\alpha,\beta)\).  It does not determine its norm
\(\varrho_{\rm sp}\).  In particular, \(m=1\) alone does not imply any
normalization row of the following table.

The Paper-I-A protected-scale family and its ambient-cone member are therefore:
\[
\begin{array}{c|c|c|c}
\text{convention} & \varrho_{\rm sp} &
\text{phase-gauge }\gamma &
\sum_a(\widehat X_a^{\rm sp})^2\\
\hline
\text{Paper-I-A }\mathsf C_{\rm PI}(\rho_{\rm prot})
& \rho_{\rm prot}/2 & \rho_{\rm prot}/2 & \rho_{\rm prot}^2 I_4\\
\text{ambient cone }\mathsf C_{\rm cone}
& \Rmom/2 & \Rmom/2 & \Rmom^2 I_4
\end{array}
\]
The first row is the Paper-I-A family with independent
\(\rho_{\rm prot}>0\); \(m_p=1\) does not select a member.  The second row is
the scale-matched member
\(\mathsf C_{\rm cone}=\mathsf C_{\rm PI}(\Rmom)\) selected only after the
ambient common-cone-symbol ancestry is declared.  It is not the intrinsic
tangential class \(\mathsf C_D\).  A separate worked numerical calibration is recorded only in
Paper~II-F, Section~\ref{p2:app:X_spinor_block}; it is not used in the main argument.
\end{corollary}

\begin{proof}
The vector representation is
\((\tfrac12,\tfrac12)\), while the two protected chiral summands are
\((\tfrac12,0)\) and \((0,\tfrac12)\).  The Clebsch--Gordan rules give
\[
\dim_{\mathbb C}
\operatorname{Hom}_{\Spin(4)}
\!\left(
(\tfrac12,\tfrac12)\otimes(\tfrac12,0),(0,\tfrac12)
\right)=1
\]
and the analogous reverse intertwiner space also has complex dimension one;
there is no diagonal protected vector intertwiner.  Hermiticity identifies
the reverse coefficient with the adjoint coefficient.  One complex reduced
matrix element, hence two real parameters, remains.  The Hermitian vectors
\(C_a\) and \(D_a=i\Gamma_{\rm sp}C_a\) form a real basis for this space,
which proves \eqref{p2:eq:protected-two-parameter-vector}.

The Clifford relations imply
\[
\Gamma_{\rm sp}^\dagger=\Gamma_{\rm sp},
\qquad
\Gamma_{\rm sp}^2=I_4,
\qquad
\{\Gamma_{\rm sp},C_a\}=0.
\]
The chirality operator commutes with the even Clifford action generating
\(\Spin(4)\), so \(U(\theta)\) also commutes with \(\Spin(4)\).  Directly,
\begin{align*}
U(\theta)C_aU(\theta)^\dagger
&=\cos\theta\,C_a+\sin\theta\,D_a,\\
U(\theta)D_aU(\theta)^\dagger
&=-\sin\theta\,C_a+\cos\theta\,D_a.
\end{align*}
Thus
\[
(\alpha,\beta)\longmapsto
(\alpha\cos\theta-\beta\sin\theta,\,
 \alpha\sin\theta+\beta\cos\theta).
\]
For \(\varrho_{\rm sp}>0\), take
\(\cos\theta=\alpha/\varrho_{\rm sp}\) and
\(\sin\theta=-\beta/\varrho_{\rm sp}\).  This gives
\((\varrho_{\rm sp},0)\).
A mere \(O(4)\) change of the embedding frame acts on the vector index and
does not perform this continuous mixing of the Clifford vector and axial
vector.  Imposing a separate reflection-parity law could select one line in
the two-parameter space, but that would be an additional hypothesis.

Finally, the identities
\[
\{D_a,D_b\}=2\delta_{ab}I_4,
\qquad
\{C_a,D_b\}+\{D_a,C_b\}=0
\]
give
\(\sum_a(\widehat X_a^{\rm sp})^2=4\varrho_{\rm sp}^2I_4\).
The Paper-I-A protected condition
\(\sum_a(\widehat X_a^{\rm sp})^2=\rho_{\rm prot}^2I_4\) therefore fixes
\(\varrho_{\rm sp}=\rho_{\rm prot}/2\).  Substituting the optional
\(\rho_{\rm prot}=\Rmom\) gives the second, scale-matched row.
\end{proof}

\begin{remark}[Protected-scale owner firewall and generated normalization
graph]
\label{p2c:rem:protected-scale-owner-firewall}
The first row above is not three independent normalization assumptions.
Within the occurring Paper-I-A representative, the protected-radius relation
\[
 \sum_a(\widehat X_a^{\rm sp})^2
 =\rho_{\rm prot}^{\,2}I_4
\]
and the general identity
\(\sum_a(\widehat X_a^{\rm sp})^2
=4\varrho_{\rm sp}^{\,2}I_4\) generate the graph
\[
 \boxed{\varrho_{\rm sp}=\rho_{\rm prot}/2}
\]
on the positive owner line.  This is inherited same-owner normalization data,
not a selector of \(\rho_{\rm prot}\).

By contrast, Paper~I-B types \(\rho_{\rm prot}\) as an independent protected
finite-geometry owner and proves that the finite coordinate compression is not
the descendant of the scalar radius-\(\Rmom\) multiplication relation.
Therefore Paper~II-A's generated-owner normalization theorem cannot be invoked
to identify \(\rho_{\rm prot}\) with \(\Rmom\): ancestry and owner identity are
antecedents of that theorem, and numerical comparability is insufficient.
Absent an independently generated cross-owner calibration, the joint scale
atlas remains the relation-free product.  If the additional relation
\(\rho_{\rm prot}=\Rmom\) is supplied, its fibre product with the representative
graph gives \(\varrho_{\rm sp}=\Rmom/2\); it does not follow from
\(m_p=1\), equivariance, or compression.
\end{remark}

\begin{theorem}[Unique equivariant retraction relative to a fixed ambient action]
\label{p2:thm:finite-toeplitz-coordinate-uniqueness}
Let \(\iota_m:\mathcal H_m\hookrightarrow\mathcal H_\infty\) be the inclusion.
There is exactly one \(\Spin(4)\)-equivariant linear retraction
\[
R_m:\mathcal H_\infty^{\rm alg}\longrightarrow\mathcal H_m,
\qquad
R_m\iota_m=I_{\mathcal H_m},
\]
namely \(R_m=\Pi_m\) on the algebraic direct sum.  Its bounded Hilbert-space
extension is the orthogonal projection.  Therefore, for every fixed ambient
coordinate completion,
\[
\widehat X_a=R_m\mathsf X_a\iota_m
=\Pi_m\mathsf X_a\Pi_m
\]
is unique.  This proves uniqueness of compression, not uniqueness of
\(\mathsf X\).
\end{theorem}

\begin{proof}
On each retained irreducible \(V_\lambda\), the retraction identity forces
\(R_m|_{V_\lambda}=I\).  If \(V_\lambda\) is not retained, then
\[
\operatorname{Hom}_{\Spin(4)}(V_\lambda,\mathcal H_m)=0
\]
because the ambient decomposition is multiplicity-free and
\(\mathcal H_m\) contains no equivalent copy.  Hence \(R_m\) vanishes on every
omitted sector.  These two statements characterize \(\Pi_m\).
\end{proof}

\begin{proposition}[Exact compression and leakage identity]
\label{p2:prop:extended-coordinate-compression-identity}
Let \(Q_m=I-\Pi_m\) and
\(\mathcal R_\infty^2=\sum_a\mathsf X_a^2\).  For any fixed self-adjoint
ambient completion,
\begin{equation}
\label{p2:eq:extended-coordinate-compression-identity}
\sum_a\widehat X_a^2
=
\Pi_m\mathcal R_\infty^2\Pi_m
-
\sum_a\Pi_m\mathsf X_aQ_m\mathsf X_a\Pi_m.
\end{equation}
The subtracted operator is positive and satisfies
\begin{equation}
\label{p2:eq:extended-coordinate-leakage-norm}
\left\langle\psi,
\sum_a\Pi_m\mathsf X_aQ_m\mathsf X_a\Pi_m\psi
\right\rangle
=
\sum_a\|Q_m\mathsf X_a\Pi_m\psi\|^2
\end{equation}
for every \(\psi\in\mathcal H_\infty\).  If
\(\psi\in\mathcal H_m=\operatorname{Ran}\Pi_m\), the final \(\Pi_m\) may be
omitted.
\end{proposition}

\begin{proof}
Insert \(I=\Pi_m+Q_m\) between the two ambient factors and sum over \(a\).
Self-adjointness gives the norm formula and positivity.
\end{proof}

\subsection{Protected Clifford coordinates and normalization}

\begin{theorem}[Protected block with free normalization]
\label{p2:lem:minimal-spinor-block-closure}
\label{p2:lem:protected-vector-intertwiner-unique}
\label{p2:prop:equivariance-does-not-fix-protected-radius}
\label{p2:lem:protected-toeplitz-rme}
\label{p2:thm:position-operators}
\label{p2:thm:embedding-coordinate-operators}
\label{p2:prop:spinor-radius}
At \(m=1\), the compressed coordinate vector preserves
\[
\mathcal H_{\rm sp}
=
V_{\frac12,0}\oplus V_{0,\frac12}.
\]
Every Hermitian \(\Spin(4)\)-equivariant vector operator on this block has the
form
\begin{equation}
\label{p2:eq:protected-toeplitz-rme}
\widehat X_a^{\rm sp}
=\alpha C_a+\beta\,i\Gamma_{\rm sp}C_a,
\qquad
\alpha,\beta\in\mathbb R.
\end{equation}
Set \(\varrho_{\rm sp}=(\alpha^2+\beta^2)^{1/2}\).  Consequently
\begin{align}
\{\widehat X_a^{\rm sp},\widehat X_b^{\rm sp}\}
&=2\varrho_{\rm sp}^2\delta_{ab}I_4,\\
\sum_{a=1}^4(\widehat X_a^{\rm sp})^2
&=4\varrho_{\rm sp}^2I_4,
\label{p2:eq:protected-toeplitz-radius}\\
[\widehat X_a^{\rm sp},\widehat X_b^{\rm sp}]
&=\varrho_{\rm sp}^2[C_a,C_b].
\end{align}
The relative chiral unitary
\(U(\theta)=\exp(i\theta\Gamma_{\rm sp}/2)\) puts this family in the
phase gauge
\[
\widehat X_a^{\rm sp}=\gamma C_a,
\qquad
\gamma:=\varrho_{\rm sp}\ge0.
\]
Thus protected closure and the phase-invariant Clifford scale are derived,
while every numerical radius and coordinate--coordinate commutator scale
depends on the completion parameter \(\varrho_{\rm sp}\).
\end{theorem}

\begin{proof}
Proposition~\ref{p2:prop:h1-coordinate-coupling-graph} shows that the only
retained protected edge is the cross-chiral pair and that no retained edge
connects \(\mathcal H_{\rm sp}\) to \(\mathcal H_{\rm geo}\).  Its intertwiner
multiplicity is one in each direction.  Hermiticity relates the two
directions but leaves one complex coefficient, equivalently the real pair
\((\alpha,\beta)\) in
Corollary~\ref{p2:cor:protected-phase-normalization-freedom}.
Writing \(D_a=i\Gamma_{\rm sp}C_a\), direct Clifford algebra gives
\[
[D_a,D_b]=[C_a,C_b],
\qquad
[C_a,D_b]+[D_a,C_b]=0,
\]
as well as the anticommutator identities used in that corollary.  The three
displayed invariant formulas follow.  The same corollary proves the relative
chiral phase-gauge statement.
\end{proof}

\begin{definition}[Nondegenerate protected-coordinate stratum]
\label{p2c:def:nondegenerate-protected-coordinate-stratum}
Write
\[
 (\mathrm{ND}_{\rm sp}):\qquad \varrho_{\rm sp}>0.
\]
This is the nonzero stratum of the unresolved protected-coordinate completion
family.  It is not implied by the level-one carrier identity and does not fix
an absolute protected radius, a relative chiral representative, an ambient
Wigner--Eckart coefficient, or any Hurwitz-support datum.
\end{definition}

\begin{corollary}[Exact protected even algebra on the nondegenerate stratum]
\label{p2c:cor:nondegenerate-protected-clifford-orbit}
On \((\mathrm{ND}_{\rm sp})\), define the normalized protected coordinates
\[
 Y_a:=\varrho_{\rm sp}^{-1}\widehat X_a^{\rm sp}.
\]
Then
\[
 \boxed{\{Y_a,Y_b\}=2\delta_{ab}I_4.}
\]
For \(i=1,2,3\), put
\[
 e_i:=Y_iY_4.
\]
Writing \(D_a=i\Gamma_{\rm sp}C_a\), one has the exact represented identities
\[
 \boxed{e_i=C_iC_4,\qquad i=1,2,3.}
\]
Consequently
\[
 e_i^2=-I_4,
 \qquad
 e_ie_j+e_je_i=-2\delta_{ij}I_4,
\]
and the represented real even algebra is literally independent of both the
nonzero scale and the relative chiral phase:
\[
 \boxed{\operatorname{alg}_{\mathbb R}(I_4,e_1,e_2,e_3)
 =\operatorname{alg}_{\mathbb R}(I_4,C_1C_4,C_2C_4,C_3C_4)
 \cong \Cl(0,3).}
\]
The strata \(\varrho_{\rm sp}=0\) and \(\varrho_{\rm sp}>0\) remain distinct.
No conclusion is drawn here about the omitted completion data, an absolute
coordinate scale, the coordinate--coordinate commutator magnitude, or an
exceptional finite-support occurrence.
\end{corollary}

\begin{proof}
Theorem~\ref{p2:thm:position-operators} gives
\(\{\widehat X_a^{\rm sp},\widehat X_b^{\rm sp}\}
=2\varrho_{\rm sp}^2\delta_{ab}I_4\), so the normalized Euclidean Clifford
relations hold on \((\mathrm{ND}_{\rm sp})\).  Write
\[
 Y_a=aC_a+bD_a,
 \qquad
 a:=\frac{\alpha}{\varrho_{\rm sp}},\quad
 b:=\frac{\beta}{\varrho_{\rm sp}},\quad
 a^2+b^2=1.
\]
Because \(\Gamma_{\rm sp}C_a=-C_a\Gamma_{\rm sp}\) and
\(\Gamma_{\rm sp}^2=I\),
\[
 D_iD_4=C_iC_4,
 \qquad
 C_iD_4+D_iC_4=0.
\]
Hence
\[
 Y_iY_4
 =a^2C_iC_4+ab(C_iD_4+D_iC_4)+b^2D_iD_4
 =(a^2+b^2)C_iC_4
 =C_iC_4.
\]
The displayed Clifford relations and algebra equality follow immediately.
Thus normalization removes \(\varrho_{\rm sp}\), while passage to the even
products removes the relative chiral phase exactly rather than merely up to a
chosen conjugating representative.
\end{proof}

\begin{corollary}[Selector discharge for the normalized protected even algebra]
\label{p2c:cor:protected-clifford-selector-discharge}
For any two compatible finite coordinate completions \(\mathsf C,\mathsf C'\)
whose protected magnitudes satisfy
\(\varrho_{\rm sp}>0\) and \(\varrho'_{\rm sp}>0\), the normalized protected
even outputs coincide as represented real algebras:
\[
 \boxed{
 \mathcal A_{\rm even}^{\rm prot}(\mathsf C)
 =\mathcal A_{\rm even}^{\rm prot}(\mathsf C')
 =\operatorname{alg}_{\mathbb R}(I_4,C_1C_4,C_2C_4,C_3C_4)
 \cong\Cl(0,3).}
\]
Hence selecting \(\mathsf C_{\rm PI}\), a numerical value of the positive
protected magnitude, or a relative chiral phase is not required for this
normalized even-algebra output.  On the bare mathematical completion atlas,
the antecedent \((\mathrm{ND}_{\rm sp})\) remains a genuine nonvanishing domain
condition, and no raw coordinate coefficient, absolute radius, ambient lift,
or coordinate--coordinate commutator magnitude is fixed by this corollary.
\end{corollary}

\begin{corollary}[Full covariance of the protected coordinates]
\label{p2c:cor:protected-coordinate-covariance}
On a nonzero protected completion write
$\widehat X_a^{\rm sp}=\varrho_{\rm sp}\Gamma_a$ in its normalized
Clifford frame.  For every protected density operator, the real symmetric
coordinate covariance and its sharp trace bound are
\[
 V=\varrho_{\rm sp}^2I_4-\mu\mu^{\mathsf T},\qquad
 |\mu|\le\varrho_{\rm sp},\qquad
 \operatorname{Tr}V\ge3\varrho_{\rm sp}^2,
 \quad \mu_a=\operatorname{Tr}(\sigma\widehat X_a^{\rm sp}).
\]
Every mean vector in this ball occurs.  Its three transverse principal
variances are exactly $\varrho_{\rm sp}^2$; the fourth is
$\varrho_{\rm sp}^2-|\mu|^2$.
The intrinsic protected bound is $9\Rmom^2/16$.
For the separate Paper-I-A representative, the coefficient is
$\varrho_{\rm sp}=\rho_{\rm prot}/2$, and the bound is
$3\rho_{\rm prot}^2/4$.
\end{corollary}
\begin{proof}
Apply Paper~I-B, Proposition~\ref{p1:prop:protected-coordinate-covariance},
with its neutral Clifford coefficient $c=\varrho_{\rm sp}$.
The intrinsic value follows from
Theorem~\ref{p2c:thm:intrinsic-tangential-coordinate-class}.
\end{proof}

\begin{theorem}[Protected spectral projectors carry the midpoint connection]
\label{p2c:thm:protected-projector-midpoint}
On any positive protected Clifford stratum, let $n\in S^3(1)$ be a
coefficient direction and set $P_\pm(n)=(I\pm\Gamma(n))/2$.
In the chiral frame
\[
 \Gamma(n)=\begin{pmatrix}0&U(n)\\U(n)^\dagger&0\end{pmatrix},
 \qquad U(n)=n_4I+i\sum_{i=1}^3 n_i\sigma_i\in SU(2),
\]
both rank-two eigenbundles have global orthonormal frames
\[
 V_\pm(n)=\frac1{\sqrt2}\begin{pmatrix}\pm U(n)\\I\end{pmatrix}.
\]
Their projected connections are equivalent and in these frames coincide:
\[
 \boxed{A_+=A_-=\tfrac12\Theta},\qquad
 \boxed{F_+=F_-=-\tfrac14\Theta\wedge\Theta},\qquad
 \Theta=U^\dagger dU.
\]
This is the midpoint connection of Paper~I-C,
Theorem~\ref{p1b:thm:homogeneous}, pulled back along $n\mapsto U(n)$.
It has full restricted $SU(2)$ holonomy and, with I-B's orientation and
normalized Chern--Simons convention, class $1/2\pmod{\mathbb Z}$.
\end{theorem}
\begin{proof}
Direct multiplication gives $\Gamma(n)V_\pm=\pm V_\pm$ and
$V_\pm^\dagger V_\pm=I_2$.  Differentiation gives
$V_\pm^\dagger dV_\pm=U^\dagger dU/2$ for either sign.
The Maurer--Cartan identity $d\Theta=-\Theta\wedge\Theta$ gives the
curvature.  Its values span $\mathfrak{su}(2)$, so the holonomy theorem
applies.  The I-B transgression coordinate $3f^2-2f^3$ at $f=1/2$
is $1/2$; changing orientation changes its sign and leaves its class modulo
integers equal to $1/2$.  Fixed Clifford-basis changes induce unitary bundle
isomorphisms and do not change these conclusions.
\end{proof}
The bundles are trivial, while their projected connections are nonflat.
The two complementary signs alone produce no connection imbalance.
An actual adiabatic interpretation requires a supplied Hamiltonian with
these gapped eigenspaces and an actual parameter history.  A spacetime
pullback, gauge coupling, particle or physical helicity assignment retains
its downstream owner; the finite projector identity supplies none of them.

\begin{corollary}[Claim-relative discharge and persistence of
\(\mathrm{ND}_{\rm sp}\)]
\label{p2c:cor:joint-nondegeneracy-discharge}
There are two distinct joint-family situations.
\begin{enumerate}[label=(\roman*),leftmargin=2.2em]
\item Suppose a physical claim independently contains an occurring
protected-coordinate vector
\(\widehat X^{\rm sp}=(\widehat X_a^{\rm sp})_{a=1}^{4}\) whose frozen
claim-relative distinction is
\(\widehat X^{\rm sp}\neq0\), and suppose the licensed relation identifies
that vector with the protected block of the finite coordinate completion.
Then
\[
 \widehat X^{\rm sp}\neq0
 \quad\Longleftrightarrow\quad
 \varrho_{\rm sp}>0,
\]
so the joint zero-defect locus is the positive stratum
\[
 \boxed{\mathfrak Z_{\rm nz}
 =\{(r_{1,0},r_{0,1},\varrho_{\rm sp})
       \in\mathbb R_{\geq0}^{3}:\varrho_{\rm sp}>0\}.}
\]
On that claim, a second \((\mathrm{ND}_{\rm sp})\) assumption is redundant.
The same conclusion holds if the occurring completion is
\(\mathsf C_{\rm PI}\), because its independently owned
\(\rho_{\rm prot}>0\) and its representative relation give
\(\varrho_{\rm sp}=\rho_{\rm prot}/2>0\).
\item If no nonzero protected-coordinate occurrence and no occurring
representative with a generated positive normalization is present, then the
zero member \(\varrho_{\rm sp}=0\) preserves the finite coordinate typing,
the frozen scalar-symbol relation, and every licensed closure obligation that
does not explicitly require the protected edge.  Paper~II-A's zero-member
test therefore does not force nonvanishing.  In that situation
\((\mathrm{ND}_{\rm sp})\) remains an explicit theorem-domain predicate and
normalization by \(\varrho_{\rm sp}^{-1}\) is unavailable on the zero stratum.
\end{enumerate}
Neither branch selects \(r_{1,0}\), \(r_{0,1}\), a positive numerical value of
\(\varrho_{\rm sp}\), or an ambient lift.
\end{corollary}

\begin{proof}
Theorem~\ref{p2:thm:position-operators} gives
\[
 \sum_{a=1}^{4}(\widehat X_a^{\rm sp})^2
 =4\varrho_{\rm sp}^{\,2}I_4.
\]
Thus the protected vector is zero exactly when
\(\varrho_{\rm sp}=0\), proving the first equivalence and its fibre-product
locus.  The Paper-I-A representative implication follows from its displayed
normalization row.  In the absence of either occurring root, setting
\(\alpha=\beta=0\) leaves a Hermitian equivariant vector and does not disturb
the geometric scalar channel or any relation not requiring protected
nonvanishing.  The zero-member test then proves the second clause.
\end{proof}

\begin{proposition}[Invariant radius data for a fixed completion]
\label{p2:prop:sector-invariant-scalars}
For any fixed ambient completion, both
\[
\mathcal R_\infty^2=\sum_a\mathsf X_a^2
\quad\text{and}\quad
\widehat{\mathcal R}_m^2=\sum_a\widehat X_a^2
\]
are \(\Spin(4)\)-invariant and therefore scalar on every irreducible sector.
On the protected block,
$\widehat{\mathcal R}_1^2=4\varrho_{\rm sp}^2I_4$.
Scalar-symbol matching fixes the ambient scalar-sector value to $\Rmom^2$,
and the same value holds at every finite cutoff containing the
scalar-to-vector edge.
The remaining geometric-sector values depend on the geometric reduced
coefficients; on $\mathcal H_1$ they are the explicit entries of
Proposition~\ref{p2:prop:geo-coordinate-admissibility-checks}.
They are not determined by $\varrho_{\rm sp}$ alone.
\end{proposition}

\begin{proof}
Contracting two vector indices with the invariant Euclidean metric gives an
intertwiner.  Schur's lemma makes its restriction to each irreducible sector
scalar.  The protected value is
Theorem~\ref{p2:thm:position-operators}. For normalized constant $e_0$ and
first harmonics $u_a$, scalar matching gives $\widehat X_ae_0=(\Rmom/2)u_a$
whenever the cutoff contains that edge. Its only vector-type target is the
scalar first-harmonic sector, so $\sum_a\widehat X_a^2e_0=\Rmom^2e_0$.
The same argument applies to the ambient action. The other finite entries
are calculated in Proposition~\ref{p2:prop:geo-coordinate-admissibility-checks}.
\end{proof}

\begin{definition}[Completion-relative carrier scalar]
\label{p2:def:canonical-carrier-scale-scalar}
For a declared completion, define
\[
\widehat\Lambda_{\rm car}
:=
\Rmom^{-2}\sum_a\widehat X_a^2.
\]
The full spectrum is completion-relative, but the scalar-sector eigenvalue is
canonically $1$ whenever the cutoff contains the scalar-to-vector edge.
These operator eigenvalues do not by themselves define a position metric,
probability measure, or density convention.
\end{definition}

\begin{remark}[Status of radius and density data]
\label{p2:rem:radius-status}
\label{p2:cor:carrier-volume-density-factor}
An expression such as
\((\lambda_{\rm car}^{(s)})^{3/2}\) may be introduced as an additional
isotropic-density convention after a completion has been selected.  It is not
a derived position-space metric, probability measure, or Born-balance factor.
\end{remark}

\subsection{Finite commutators, leakage, and uncertainty bounds}

\begin{theorem}[Exact compact transvection--coordinate commutators]
\label{p2:thm:deformed-heisenberg}
\label{p2:prop:spinor-exact-commutators}
Let
\[
\widehat P_i=\frac{\hbar}{\Rmom}(L_i^+-L_i^-)
=\frac1{\Rmom}\widehat J_{i4}.
\]
On the ambient Hilbert space, \(\widehat P_i\) denotes the self-adjoint
generator of the corresponding one-parameter \(\Spin(4)\) action.  Classical
embedding multiplication preserves \(D(\widehat P_i)\), and on that domain
\[
[\widehat P_i,M_{x_j}]
=
-\frac{i\hbar}{\Rmom}\delta_{ij}M_{x_4},
\qquad
[\widehat P_i,M_{x_4}]
=
\frac{i\hbar}{\Rmom}M_{x_i}.
\]
Every bounded equivariant ambient coordinate completion satisfies
\(\mathsf X_aD(\widehat P_i)\subseteq D(\widehat P_i)\) and the same vector
commutator on \(D(\widehat P_i)\).  Since \(\widehat P_i\) preserves
\(\mathcal H_m\) and commutes
with \(\Pi_m\), finite compression gives exactly
\[
[\widehat P_i,\widehat X_j]
=
-\frac{i\hbar}{\Rmom}\delta_{ij}\widehat X_4,
\qquad
[\widehat P_i,\widehat X_4]
=
\frac{i\hbar}{\Rmom}\widehat X_i.
\]
These identities also hold on the protected block for every
\((\alpha,\beta)\); their coefficients are independent of the phase-invariant
scale \(\varrho_{\rm sp}\) and of the ambient normalization.
\end{theorem}

\begin{proof}
In the defining vector representation,
\[
[\widehat J_{\mu\nu},x_\rho]
=
-i\hbar(\delta_{\mu\rho}x_\nu-\delta_{\nu\rho}x_\mu).
\]
Taking \((\mu,\nu)=(i,4)\) proves the classical formulas.  Vector covariance
is exactly the same infinitesimal intertwining identity for any ambient
completion.  To make its domain content explicit, write
\[
U_i(t):=\exp(-it\widehat P_i/\hbar),\qquad
U_i(t)\mathsf X_aU_i(t)^\dagger
=R_i(t)_{ab}\mathsf X_b,
\]
where \(R_i(t)\) is the differentiable defining-vector matrix.  For
\(\psi\in D(\widehat P_i)\),
\[
U_i(t)\mathsf X_a\psi
=R_i(t)_{ab}\mathsf X_bU_i(t)\psi
\]
is differentiable at \(t=0\).  Hence
\(\mathsf X_a\psi\in D(\widehat P_i)\), and differentiation gives the
displayed ambient commutator there.  The same argument applies to the bounded
multipliers \(M_{x_a}\).  Finally,
\[
[\widehat P_i,\Pi_m\mathsf X_a\Pi_m]
=
\Pi_m[\widehat P_i,\mathsf X_a]\Pi_m
\]
because \(\Pi_m\) is equivariant.  No reduced matrix element enters.
\end{proof}

\begin{corollary}[Exact status of the finite coordinate-completion family]
\label{p2:cor:finite-coordinate-package-status}
The following are completion-independent: the projected channel graph, the
unique equivariant retraction \(R_m=\Pi_m\), protected closure, and the linear
transvection--coordinate commutators.  The following are
completion-dependent: off-diagonal ambient coefficients, invariant cycle
products and their sign or phase when nonzero,
\((\alpha,\beta)\), the invariant scale \(\varrho_{\rm sp}\), all protected radius
constants, coordinate--coordinate commutators, and their numerical Robertson
bounds.  The direction of \((\alpha,\beta)\) is removable by a relative
chiral unitary commuting with \(\Spin(4)\); it is not removed by an embedding
frame rotation.  The Paper-I-A value
\(\varrho_{\rm sp}=\rho_{\rm prot}/2\) is the specialization
\(\mathsf C=\mathsf C_{\rm PI}\); it is not a consequence of \(m=1\) alone, and
replacing \(\rho_{\rm prot}\) by \(\Rmom\) requires the optional scale match.
At \(m=1\) specifically, the projected graph is a forest, so no finite cycle
sign remains and the compatible finite class has the three nonnegative
magnitudes of
Corollary~\ref{p2:cor:h1-coordinate-classification}.  Ambient cycle data
remain possible on larger graphs.
\end{corollary}

\begin{proof}
The completion-independent graph, retraction, protected closure, and linear
commutator statements are respectively
Proposition~\ref{p2:prop:h1-coordinate-coupling-graph},
Theorem~\ref{p2:thm:finite-toeplitz-coordinate-uniqueness},
Theorem~\ref{p2:thm:position-operators}, and
Theorem~\ref{p2:thm:deformed-heisenberg}.  Ambient coefficient and cycle
freedom is Theorem~\ref{p2:thm:derived-ambient-coordinate-action}; the
protected parameters, phase gauge, and scale rows are
Corollary~\ref{p2:cor:protected-phase-normalization-freedom}.  Finally,
Corollary~\ref{p2:cor:h1-coordinate-classification} proves both the forest
classification at \(m=1\) and the existence of cycles on larger graphs.
\end{proof}

\begin{remark}[Nonlinear projection corrections]
\label{p2:rem:projection-corrections}
Equation~\eqref{p2:eq:extended-coordinate-compression-identity} controls
quadratic radius and coordinate-product leakage.  It does not modify the
linear commutators in
Theorem~\ref{p2:thm:deformed-heisenberg}.
\end{remark}
\subsubsection{Uncertainty bounds in the compact model}
\label{p2:sec:gup}

Compact-geometry uncertainty bounds follow from the classical deformed
commutator and its finite-level realizations.  At finite truncation, these bounds are localization estimates for the truncated model; they are not unqualified statements about the smooth parent outside the declared finite or fixed-band domain.  The Robertson lower
bound is controlled by state expectations of compact coordinate operators; it
is not a universal positive minimal-length bound.

\begin{theorem}[Robertson bound from the compact commutator]
\label{p2:thm:gup-general}
Let \(\psi\) be normalized and write
\(\omega_\psi(T)=\langle\psi,T\psi\rangle\).  For bounded self-adjoint
observables take any such vector.  For unbounded self-adjoint observables,
require \(\psi\) to lie in a common invariant dense domain
\(\mathcal D\subset D(AB)\cap D(BA)\) and require both variances to be finite.
The Robertson inequality \cite{Robertson1929} then gives
\[
\Delta_{\omega_\psi}(A)\,\Delta_{\omega_\psi}(B)
\ge
\frac12\,\bigl|\omega_\psi([A,B])\bigr|.
\]
For the classical compact embedding-coordinate--momentum pair on
\(L^2(S^3(\Rmom))\), one may take
\(\mathcal D=C^\infty(S^3)\), which is preserved by both coordinate
multiplication and the momentum generators.  Theorem~
\ref{p2:thm:deformed-heisenberg} then yields
\[
\Delta_{\omega_\psi}(M_{x_i})\,\Delta_{\omega_\psi}(\hat P_i)
\ge
\frac{\hbar}{2\Rmom}\,\bigl|\omega_\psi(M_{x_4})\bigr|.
\]
Moreover, on the protected spinor block $\mathcal H_{\mathrm{sp}}$, all operators are finite-dimensional and Proposition~\ref{p2:prop:spinor-exact-commutators} gives the exact finite-level analogue
\[
\Delta_{\omega_\psi}(\hat X_i^{\mathrm{sp}})\,
\Delta_{\omega_\psi}(\hat P_i)
\ge
\frac{\hbar}{2\Rmom}\,
\bigl|\omega_\psi(\hat X_4^{\mathrm{sp}})\bigr|.
\]
\end{theorem}

\begin{proof}
Use the Robertson inequality on the stated vector/domain:
\[
\Delta_{\omega_\psi}(A)\,\Delta_{\omega_\psi}(B)
\ge
\frac12\,\bigl|\omega_\psi([A,B])\bigr|.
\]
The content here is therefore only the compact-geometry commutator inserted into that inequality.  With
\(A=M_{x_i}\) and \(B=\hat P_i\), Theorem~\ref{p2:thm:deformed-heisenberg} gives
\[
[M_{x_i},\hat P_i]=\frac{i\hbar}{\Rmom}M_{x_4},
\]
which yields the displayed continuum compact bound.  On the protected spinor block,
Proposition~\ref{p2:prop:spinor-exact-commutators} gives
\[
[\hat X_i^{\mathrm{sp}},\hat P_i]=\frac{i\hbar}{\Rmom}\hat X_4^{\mathrm{sp}},
\]
and the same Robertson inequality gives the finite protected-block bound.
\end{proof}


\subparagraph{Spectral ownership.}
The heat, zeta, Dirac, and complete compact-spectrum statements are developed
in II-D.  They are not repeated or used to prove the coordinate results here.





\section{Intrinsic sectors, central witnesses and minimal bridge completion}
\label{p2c:sec:v352-center-bridge}
Work on the canonical one-copy module
$\mathcal H_1=G\oplus S$, $\dim G=11$, $\dim S=4$, with the
\emph{previously constituted} intrinsic
$\mathcal A_D=\operatorname{End}(G)\oplus\operatorname{End}(S)$.
The abstract-algebra theorem below does not independently license
new odd events.

\begin{proposition}[Central signs and an exact finite effect witness]
\label{p2c:prop:v352-center-witness}
Let $z_L=(-I,I)$, $z_R=(I,-I)$ and $z_D=(-I,-I)$ in
$\Spin(4)$.  Every intrinsic vector coordinate $X_a^D$ satisfies
\[
 \operatorname{Ad}_{U(z_L)}X_a^D
 =\operatorname{Ad}_{U(z_R)}X_a^D=-X_a^D,
 \qquad\operatorname{Ad}_{U(z_D)}X_a^D=X_a^D.
\]
For unit $e_0\in E_0\subset G$, $u\in E_1\subset G$,
put $\psi_\pm=(e_0\pm u)/\sqrt2$ and
$F_+=|\psi_+\rangle\langle\psi_+|\in\mathcal A_D$.
Then $U(z_L)\psi_+=\psi_-$ and
\[
 \langle\psi_+,F_+\psi_+\rangle=1,\qquad
 \langle U(z_L)\psi_+,F_+U(z_L)\psi_+\rangle=0.
\]
This is potential exact algebraic distinguishability, not an
assertion of a realized preparation, central operation or detector.
\end{proposition}
\begin{proof}
The six existing irreducible types have signs $(-1)^{2j_L}$
and $(-1)^{2j_R}$.  Every nonzero intrinsic coordinate edge is
of $(\tfrac12,\tfrac12)$ type, odd for $z_L,z_R$ and even for
$z_D$.  The trivial $E_0$ and vector $E_1$ have opposite
$z_L$ signs, so the two normalized superpositions are orthogonal.
The full $M_{11}$ block contains their rank-one effects.
\end{proof}

\begin{theorem}[A single nonzero geometric--spinorial bridge completes $M_{15}$]
\label{p2c:thm:v352-rankone-full-bridge}
If $0\ne B\in\operatorname{Hom}(S,G)$, then, with its off-diagonal
embedding in $\operatorname{End}(G\oplus S)$,
\[
 \boxed{C^*(\mathcal A_D,B,B^*)=M_{15}(\mathbb C)}.
\]
No lower bound on the nonzero singular values or on the rank of $B$
is needed.  With $\Gamma=I_G\oplus(-I_S)$, any such bridge obeys
$\Gamma B\Gamma=-B$; its inclusion reduces the remaining
conjugation kernel $\{1,z_D\}$ of the coordinate algebra to the
identity on the full algebra.
\end{theorem}
\begin{proof}
Choose a nonzero matrix entry of $B$.  Left and right multiplication
by the full matrix units already belonging to $M_{11}\oplus M_4$
produces all $11\times4$ off-diagonal matrix units.  Their
adjoints give the reverse units, so the generated algebra contains
all $15^2$ matrix units.  The parity relation is immediate from
$\Gamma=\operatorname{diag}(+I,-I)$.  The last kernel statement
uses the previously proved effective-symmetry filtration and the
presence of both chiral half-spinor types.
\end{proof}

\begin{proposition}[Scalar commutant obstruction for a constant source]
\label{p2c:prop:v352-static-source-commutant}
For a Hermitian operator $M$ on the \emph{one-copy} carrier,
\[
 [M,\mathcal A_D]=0\quad\Longrightarrow\quad
 \boxed{M=\mu_G I_{11}\oplus\mu_S I_4},
 \qquad \mu_G,\mu_S\in\mathbb R.
\]
Such a static fully coordinate-compatible source has at most two
distinct eigenvalues and singular-value magnitudes.  This is not a
statement about physical mass poles, the $48$-dimensional fermion
source module or a noncommuting and energy-dependent self-energy.
\end{proposition}
\begin{proof}
The commutant of each irreducibly represented full matrix algebra
consists of its scalar multiples; no cross-block map can commute
with both independent central projectors.
\end{proof}

\begin{remark}[Bridge licensing and generated-state distinction]
\label{p2c:rem:v352-bridge-licensing}
The closure above extends the earlier coupling-graph theorem by
showing that \emph{rank one} already suffices.  It neither inserts
$B$ into the frozen GP2 grammar nor upgrades a field-typed odd
operator to a measured event.  An odd-target constructor or an
independently occurring process/test must carry that obligation.
\end{remark}



\subsection{Amplified representations and their unmarked multiplicity gauge}
\begin{proposition}[Normalizer of a faithfully amplified intrinsic algebra]
\label{p2c:prop:v300-coordinate-normalizer}
For the represented intrinsic coordinate algebra of
Theorem~\ref{p2c:thm:v294-native-multiplicity-algebra},
\[
 H=(\mathbb C^{11}\otimes\mathbb C^{k_g})\oplus
 (\mathbb C^4\otimes\mathbb C^{k_s}),\qquad
 \mathcal A_D=(M_{11}\otimes I_{k_g})\oplus(M_4\otimes I_{k_s}),
\]
with $k_g,k_s\ge1$, its unitary normalizer is exactly
\[
 \boxed{\mathcal N_{U(H)}(\mathcal A_D)=
 \bigl(U(11)\otimes U(k_g)\bigr)\oplus
 \bigl(U(4)\otimes U(k_s)\bigr).}
\]
The right side denotes the images of tensor-product unitaries in
each summand, with redundant central phases identified.
Consequently, by Theorem~\ref{p4a:thm:v300-covariant-root-instrument},
strictly positive square-root readout updates followed by measurable
unitaries preserve \emph{all} intrinsic outcome-event histories
exactly when the unitary updates are of this blockwise product form
almost everywhere. Any genuine entangling visible/multiplicity
post-update outside the displayed normalizer produces a single-event
quotient-distinguishing witness.
\end{proposition}
\begin{proof}
A normalizing unitary induces an automorphism of $\mathcal A_D$.
Its two simple central matrix summands $M_{11}$ and $M_4$ are not
isomorphic, so the two central projections must be fixed. The unitary
therefore acts block diagonally. On a block
$M_n\otimes I_k$, every automorphism of $M_n$ is inner;
a corresponding $V_n\otimes I_k$ realizes it. The residual unitary
commutes with $M_n\otimes I_k$ and therefore lies in
$I_n\otimes U(k)$. The converse is immediate by conjugation.
Apply the exact post-update normalizer equivalence in
Theorem~\ref{p4a:thm:v300-covariant-root-instrument}.
\end{proof}

\begin{theorem}[Exact intrinsic algebra as the fixed points of the unmarked multiplicity gauge]
\label{p2c:thm:v301-intrinsic-gauge-fixedpoints}
Assume the faithfully represented intrinsic coordinate algebra of
Theorem~\ref{p2c:thm:v294-native-multiplicity-algebra}, on
\[
 H=(\mathbb C^{11}\otimes\mathbb C^{k_g})\oplus
   (\mathbb C^4\otimes\mathbb C^{k_s}),\qquad k_g,k_s\ge1.
\]
The \emph{implementation gauge} $K=U(k_g)\times U(k_s)$ acts by
$W_{u,v}=(I_{11}\otimes u)\oplus(I_4\otimes v)$. Its fixed-point
operator algebra is exactly the native intrinsic test algebra:
\[
 \boxed{\ \End(H)^{\operatorname{Ad}K}
   =(M_{11}\otimes I_{k_g})\oplus(M_4\otimes I_{k_s})
   =:\mathcal A_D.\ }
\]
The independent central phases of the two gauge factors remove
geometric--spinorial off-block operators; invariance under each full
unitary factor removes all within-block multiplicity directions.
This is an exact representation-theoretic fact, not the assertion
that $K$ is a new \emph{physical} transformation group or that every
physical process is automatically natural under it.
\end{theorem}
\begin{proof}
Let $X\in\End(H)$ commute with all $W_{u,v}$. Taking
$u=e^{it}I_{k_g}$ and $v=I_{k_s}$ for all $t$ forces both
geometric--spinorial off-diagonal corners of $X$ to vanish. Within
the geometric corner, the commutant of $I_{11}\otimes U(k_g)$ is
$M_{11}\otimes I_{k_g}$, by the full matrix-algebra commutant
identity (equivalently by elementary matrix units and Schur's
lemma). The same argument gives $M_4\otimes I_{k_s}$ for the
spinorial corner. Every operator in the displayed direct sum
commutes with the gauge, proving equality. The two phase factors
are independent even if $k_g=k_s$; this is why no off-block
intertwiner survives.
\end{proof}




\section{Generated reciprocal position, scalar bands and guarded observables}
\paragraph{Group-valued momentum interface.}
The standard group-multiplication, BCH, and fusion-rule consequences of the
already selected \(SU(2)\) momentum parent are summarized in
Paper~II-F, Section~\ref{p2c:app:group-valued-momentum}.  They are supplementary to the
finite-coordinate operator proof.

\begin{theorem}[Coordinate-generated reciprocal role and relation-free product]
\label{p2c:thm:coordinate-generated-reciprocal-role}
Let \(A=S^3_A(\Rmom)\) be the already constituted round parent and let
\[
 E_1^{\mathbb R}
 =\{f_u:A\to\mathbb R:\ f_u(a)=\langle u,a\rangle,\ u\in\mathbb R^4\}
\]
be its intrinsic first nonconstant real scalar eigenspace, whose
coordinate-free status is supplied by
Theorem~\ref{p2:thm:intrinsic-coordinate-sector}.  Define
\[
 \iota_1:A\longrightarrow(E_1^{\mathbb R})^*,
 \qquad
 \iota_1(a)(f):=f(a),
 \qquad
 B_A^{(1)}:=\iota_1(A).
\]
Then \(B_A^{(1)}\) is a compact separating evaluation spectrum and
\(\iota_1:A\to B_A^{(1)}\) is a homeomorphism.  For every constituted round
isometry \(T\), the contragredient map
\[
 T^\vee\bigl(\iota_1(a)\bigr):=\iota_1(T^{-1}a)
\]
is faithful and obeys the opposite composition law
\[
        (T_1T_2)^\vee=T_2^\vee T_1^\vee.
\]
Thus the evaluation/localization target carries an autonomous
opposite-variance role sort, and the frozen relation slot contains no
joint-support constraint on \(A\times B_A^{(1)}\).  Consequently
Theorem~\ref{p2:thm:gp2-generated-reciprocal-role-closure} constitutes
\[
        B_A^{(1)}\cong S^3_B,
        \qquad
        A\times B_A^{(1)}\cong S^3_A\times S^3_B
\]
as the reciprocal role and universal relation-free product.  The generated
exchange
\[
 \mathcal B_A(a,b)
 :=\bigl(\iota_1^{-1}(b),\iota_1(a)\bigr)
\]
is an involution.
\end{theorem}

\begin{proof}
If \(a\ne b\), choose \(u=a-b\).  Then
\[
 f_u(a)-f_u(b)=\langle a-b,a-b\rangle\ne0,
\]
so \(E_1^{\mathbb R}\) separates points and \(\iota_1\) is injective.  Its
image is compact, and a continuous bijection from compact \(A\) to the
Hausdorff space \(B_A^{(1)}\) is a homeomorphism.
The round metric may therefore be transported along \(\iota_1\) as
comparison geometry on \(B_A^{(1)}\).  This mathematical transport does not
fix a dimensionful local position--momentum normalization.

Every round isometry preserves \(E_1^{\mathbb R}\), since this space is an
intrinsic Laplace eigenspace.  With \(T^*f=f\circ T\),
\[
        (T_1T_2)^*=T_2^*T_1^*.
\]
This pullback representation is faithful: if \(T_1\ne T_2\), choose \(a\)
with \(T_1a\ne T_2a\) and then an \(f\in E_1^{\mathbb R}\) separating those
two points.  Thus \(T_1^*f\ne T_2^*f\).
On the evaluation spectrum the corresponding contragredient occurrence is
\[
 T^\vee=\iota_1\circ T^{-1}\circ\iota_1^{-1}.
\]
It is faithful, and
\[
 (T_1T_2)^\vee
 =\iota_1T_2^{-1}T_1^{-1}\iota_1^{-1}
 =T_2^\vee T_1^\vee,
\]
so the required operation occurrence has the proved opposite variance.

The source is the active round role, whereas
\(B_A^{(1)}\subset(E_1^{\mathbb R})^*\) is the tagged
evaluation/localization sort carrying the dual pullback action.  The
round-metric identification \(a\mapsto f_a\) gives the total comparison
\[
 \operatorname{ev}:A\times B_A^{(1)}\longrightarrow\mathbb R,
 \qquad
 \operatorname{ev}\bigl(a,\iota_1(z)\bigr)
 :=\iota_1(z)(f_a)=\langle a,z\rangle.
\]
It is defined for every pair.  The comparison isomorphism \(\iota_1\) is a
map between the two sorts, not a graph equation restricting their joint
support.  Hence the three antecedents of
Theorem~\ref{p2:thm:gp2-generated-reciprocal-role-closure} are proved, and
that theorem supplies the generated reciprocal role and Cartesian product.
Finally,
\[
 \mathcal B_A^2(a,b)
 =\mathcal B_A\bigl(\iota_1^{-1}(b),\iota_1(a)\bigr)
 =(a,b).
\]
\end{proof}

\begin{corollary}[The round \(S^3\) pair is one invertibly bridged co-realization]
\label{p2c:cor:s3-universal-corealization-bridge}
The two roles \(S^3_A\) and \(S^3_B\) belong to one rooted constructor
component.  In the terminology of
Theorem~\ref{p2a:thm:universal-corealization-bridge-normal-form}, its retained
edges include the evaluation embedding \(\iota_1\), the total pairing
\(\operatorname{ev}\), the contragredient operation
\(T\mapsto T^\vee\), and the Born exchange \(\mathcal B_A\).  The operation
edge is the evaluation-dual output of
Theorem~\ref{p2a:thm:evaluation-dual-corealization}, and the Born edge is its
invertible self-dual specialization from
Corollary~\ref{p2a:cor:born-invertible-corealization}.  Hence no independent
occurrence of \(S^3_B\), its pullback operations, their dynamics family, or
the exchange is required.

All these nodes have the occurring round parent and its intrinsic
\(E_1^{\mathbb R}\) evaluation spectrum as their ultimate rooted origin.
They do not have the same immediate constructor: role objectification,
evaluation dualization, relation-free product formation, and exchange remain
distinct typed steps.  The conclusion still supplies no dimensionful
position--momentum calibration, local phase-space interpretation, selected
Hamiltonian, or actual history.
\end{corollary}

\begin{proof}
The preceding theorem proves every antecedent and displays every named edge.
The cited Paper-II-A results retain the entire incidence component and its
dual operations from the same root.  Since \(\iota_1\) is an isomorphism and
the joint support is relation-free, the Born specialization applies.  None of
the excluded target sorts occurs in this constructor word.
\end{proof}

\begin{corollary}[Hereditary Born covariance on the generated round pair]
\label{p2c:cor:hereditary-born-covariance}
Let \(w\) be any later licensed constructor word whose complete frozen input
contains the generated round reciprocal pair
\(S_A^3\times S_B^3\) and whose consumed marks are fixed by the Born exchange
\(\mathcal B_A\) of
Theorem~\ref{p2c:thm:coordinate-generated-reciprocal-role}.  Then
Paper~II-A, Theorem~\ref{p2a:thm:generated-born-inheritance}, gives
\[
 \operatorname{ReqOut}(w,\mathcal B_A D)
 =\mathcal B_{\mathsf Q}\operatorname{ReqOut}(w,D)
\]
in the declared target sort.  If exchange instead moves a retained finite
level or another mark, only the typed reciprocal-branch transport of
Proposition~\ref{p2a:prop:typed-born-branch-transport} is available.

The natural scalar, Hodge, curl and spin-Dirac transport theorems below
identify concrete target actions for important descendants; they are not new
reciprocity premises.  None of these inheritances fixes a dimensionful
position--momentum calibration, chooses a Hamiltonian or state, actualizes a
history, or identifies compact momentum operators with reciprocal-position
differential operators.
\end{corollary}

\begin{proof}
The preceding theorem and
Corollary~\ref{p2c:cor:s3-universal-corealization-bridge} certify the generated
Born involution and its rooted origin.  The stated unmarkedness condition is
exactly the remaining hypothesis of Paper~II-A's generated Born-inheritance
theorem.  Its marked-branch alternative is the typed transport proposition.
\end{proof}


\begin{remark}[Physical reciprocal role versus mathematical cotangent certification]
\label{p2c:rem:reciprocal-role-logical-direction}
The reciprocal comparison role and the first-symbol cotangent pairing are
physical structural descendants of the Paper-II constructor chain. They do
not acquire physical status from Paper I-C. Independently, the pair-groupoid
tangent deformation of the same compact group $K$ has geometric zero fibre
$TK$. Its zero-fibre convolution C*-algebra is identified by fibrewise
Fourier transform with $C_0(T^*K)$, whose spectrum is $T^*K$.
This is the mathematical cotangent certification proved in Paper I-C and
compatible with the generated role. Paper III-A localizes the latter to the
canonical cotangent germ used by the dynamics chain. No finite-band selection,
Hamiltonian, state, history, or dimensionful calibration follows from the
Paper-I-B certification.
\end{remark}

\begin{remark}[Exact reciprocal-role boundary]
\label{p2c:rem:reciprocal-role-derived-boundary}
The theorem proves the legacy certification statuses
\((\mathrm{BR}_{\rm ind})\) and \((\mathrm{IND}_{xp})\); neither is an
additional II-C hypothesis.  It also shows that the metric-curry clause
\(\mathsf P_{\rm mc}\) is unnecessary for reciprocal-role existence,
opposite variance, relation-free composition, or Born exchange.  That clause
is only one optional G0 test-occurrence certificate by which the full family
of distance profiles \(\{d(z,\cdot):z\in A\}\) is realized as physical
readout tests.  It is not a foundational reciprocal-role premise, and G0 does
not require every apparatus to realize the full family.
The canonical first-symbol comparison is generated below and therefore is
not part of the open bridge ledger.  The symbol \((\mathrm{PAIR}_{xp})\) is
reserved hereafter for a genuinely stronger finite-band/local--global
comparison: such a claim must provide an explicit bridge map, its domain, and
an error topology; any dimensionful normalization additionally requires a
calibration.  A Hamiltonian interpretation and an actual physical history
remain later G0/dynamical occurrence certificates.  The scalar orientation
itself is inherited separately from the Paper-I-A Hopf/Reeb provenance chain.
\end{remark}

\begin{theorem}[Canonical first-symbol position--transvection pairing]
\label{p2c:thm:first-symbol-pairing}
Let \(a\in A\), put \(b=\iota_1(a)\), and transport the real first
harmonics by
\[
 E_{1,\mathbb R}^{x}
 :=\{\widetilde f_u=f_u\circ\iota_1^{-1}:u\in\mathbb R^4\}.
\]
The differential map
\[
 D_b:E_{1,\mathbb R}^{x}\longrightarrow T_b^*B_A^{(1)},
 \qquad D_b(\widetilde f)=d\widetilde f|_b,
\]
has kernel \(\mathbb R\widetilde f_a\) and induces a canonical isomorphism
\begin{equation}
\label{p2c:eq:first-symbol-radial-quotient}
 E_{1,\mathbb R}^{x}/\mathbb R\widetilde f_a
 \xrightarrow{\ \sim\ }T_b^*B_A^{(1)}.
\end{equation}
Let \(\mathfrak p_a\) be the transported symmetric-pair transvection
complement and let
\(\operatorname{ev}_a:\mathfrak p_a\xrightarrow{\sim}T_aA\) be the
canonical evaluation isomorphism of
Theorem~\ref{p2:thm:canonical-transvection-bridge}.  Then
\begin{equation}
\label{p2c:eq:first-symbol-transvection-pairing}
 \mathcal P_a([\widetilde f],X)
 :=d\widetilde f|_b\!\left(
 d\iota_{1,a}(\operatorname{ev}_aX)\right)
\end{equation}
is a nondegenerate isotropy-equivariant pairing between the two
three-dimensional standard modules.  Thus the first-symbol content formerly
hidden in \((\mathrm{PAIR}_{xp})\) needs no freely chosen frame or
intertwiner.  Because both roles are already generated occurring descendants
and the displayed pairing is their canonical typed comparison, the
Paper-II-A generated-role/co-realization contract makes this role-level
first-symbol comparison generated and co-realized.  The legacy tag
\((\mathrm{OBS}^{(1)}_{xp})\) denotes only a particular performed/read-out
instance of that already generated comparison.  No finite-band, product,
spectral, resolvent, or continuum comparison follows from this first-order
theorem.
\end{theorem}

\begin{proof}
For \(v\in T_aA=a^\perp\),
\(df_u|_a(v)=\langle u,v\rangle\).  Hence \(df_u|_a=0\) exactly for
\(u\in\mathbb Ra\), proving the kernel and rank-three quotient.  The
evaluation-spectrum diffeomorphism and the canonical transvection bridge
identify \(\mathfrak p_a\) with \(T_bB_A^{(1)}\).  Cotangent evaluation then
gives \(\mathcal P_a\).  Each map is isotropy-equivariant and invertible
after the radial quotient, so the pairing is nondegenerate.  The structural
co-realization statement is then inherited from the current Paper-II-A
generated-role contract; it does not require a second occurrence witness.  A
particular performance/readout remains a token-level G0 statement and is not
used in this derivation.
\end{proof}

\subsection{Scalar-band resolution and higher-harmonic growth}

\begin{definition}[Reciprocal scalar-resolution filtration]
\label{p2c:def:reciprocal-scalar-resolution}
Let \(E_k^x\) denote the degree-\(k\) scalar harmonic eigenspace for the round
comparison metric transported to the derived reciprocal role \(S^3_B\), and
put
\[
 \mathcal A_x^{\rm alg}
 :=\operatorname{alg}_{\mathbb C}^{*}(E_0^x\oplus E_1^x)
 \subset C(S^3_B,\mathbb C).
\]
For \(m\in\mathbb Z_{\ge0}\), define
\[
 \mathcal P_x(m)
 :=\operatorname{span}_{\mathbb C}
 \{f_1\cdots f_r:0\le r\le m,\ f_i\in E_1^x\},
\]
where the zero-fold product is the scalar unit.  For a finite-dimensional
algebraic position-test frontier \(F\subset\mathcal A_x^{\rm alg}\), define
\[
 m_x^{\rm res}(F)
 :=\min\{m:F\subseteq\mathcal P_x(m)\}.
\]
The algebra and its filtration are mathematical comparison/test
infrastructure.  This definition does not promote unrestricted pointwise
multiplication to a finite physical carrier constructor.
\end{definition}

\begin{theorem}[Exact reciprocal scalar bands]
\label{p2c:thm:exact-reciprocal-scalar-bands}
For every \(m\ge0\),
\[
 \boxed{\mathcal P_x(m)=\bigoplus_{k=0}^{m}E_k^x},
 \qquad
 \boxed{\dim_{\mathbb C}\mathcal P_x(m)
 =\frac{(m+1)(m+2)(2m+3)}6}.
\]
Moreover,
\[
 \mathcal P_x(m)\subset\mathcal P_x(m+1),
 \qquad
 \mathcal P_x(r)\mathcal P_x(s)\subseteq\mathcal P_x(r+s),
\]
and
\[
 \overline{\bigcup_{m\ge0}\mathcal P_x(m)}^{\|\cdot\|_\infty}
 =C(S^3_B,\mathbb C).
\]
In particular,
\[
 \mathcal P_x(1)=E_0^x\oplus E_1^x,
 \qquad
 \dim_{\mathbb C}\mathcal P_x(1)=5.
\]
This scalar band is not the full auxiliary \(\Spin(4)\) module
\(\mathcal H_1\), which has dimension \(15\).
\end{theorem}

\begin{proof}
Products of \(r\) elements of \(E_1^x\) span the restrictions to the round
sphere of homogeneous degree-\(r\) polynomials on \(\mathbb R^4\).  The
Fischer decomposition gives
\[
 \operatorname{Pol}_r(\mathbb R^4)
 =\bigoplus_{q=0}^{\lfloor r/2\rfloor}
 |x|^{2q}\operatorname{Harm}_{r-2q}(\mathbb R^4).
\]
On the round comparison sphere, the factor \(|x|^{2q}\) is constant.
Consequently degree-\(r\) coordinate products contain exactly the scalar
harmonic degrees \(r,r-2,\ldots\).  Taking every \(0\le r\le m\) yields every
degree \(0,\ldots,m\) and no larger degree, proving the orthogonal direct-sum
identity.

Since \(\dim E_k^x=(k+1)^2\),
\[
 \dim\mathcal P_x(m)
 =\sum_{k=0}^{m}(k+1)^2
 =\frac{(m+1)(m+2)(2m+3)}6.
\]
Nesting and the product filtration follow directly from polynomial degree.
The generated unital \(*\)-algebra is self-adjoint and separates points by
Theorem~\ref{p2c:thm:coordinate-generated-reciprocal-role}; complex
Stone--Weierstrass therefore proves the density statement.  Finally,
\(\dim E_0^x+\dim E_1^x=1+4=5\), whereas the imported momentum-carrier theorem
gives \(\dim\mathcal H_1=15\).
\end{proof}

\begin{theorem}[First-band rigidity and exact higher-harmonic detection]
\label{p2c:thm:v280-first-band-rigidity}
Let $B=S^3(R)$, $R>0$, carry its round metric, and denote its real
scalar harmonic spaces by $E_k^{\mathbb R}$.  For a continuous point map
$T:B\to B$, let $T^*f=f\circ T$.  Then
\[
 \boxed{T^*(E_0^{\mathbb R}\oplus E_1^{\mathbb R})
 \subseteq E_0^{\mathbb R}\oplus E_1^{\mathbb R}}
 \quad\Longleftrightarrow\quad
 \boxed{T\text{ is constant or }T\in O(4)|_B.}
 \tag{FB-RIG}
\]
The sharper condition $T^*E_1^{\mathbb R}\subseteq E_1^{\mathbb R}$
holds if and only if $T$ is a round orthogonal isometry.

Consequently, if $T$ is neither constant nor a round isometry, there
exist $f\in E_1^{\mathbb R}$ and $k\ge2$ with
\[
 \boxed{\operatorname{pr}_{E_k}(f\circ T)\ne0.}
 \tag{FB-LEAK}
\]
The projection here is an $L^2$ harmonic diagnostic, not a claim that
physical higher-band tests have already been constituted.  The theorem
classifies \emph{point maps}, not general stochastic, completely positive,
or operator-algebraic processes.
\end{theorem}
\begin{proof}
Identify $B$ with $\{x\in\mathbb R^4:|x|=R\}$.  Since $E_1^{\mathbb R}$
is spanned by the four coordinate restrictions, stability of
$E_0^{\mathbb R}\oplus E_1^{\mathbb R}$ is equivalent to
\[
 T(x)=c+Mx \qquad(|x|=R)
\]
for some $c\in\mathbb R^4$ and real $4\times4$ matrix $M$.  The sphere
condition $|c+Mx|^2=R^2$ for all $|x|=R$, compared at $x$ and $-x$,
gives $M^Tc=0$.  It follows that $x^TM^TMx$ is constant on the sphere,
so $M^TM=\lambda I$ for some $\lambda\ge0$.  If $\lambda=0$ then $M=0$
and $|c|=R$, making $T$ constant.  If $\lambda>0$, $M$ is invertible,
so $M^Tc=0$ forces $c=0$; the sphere equation forces $\lambda=1$.
Thus $M\in O(4)$.  Conversely, constant maps and orthogonal isometries
preserve $E_0^{\mathbb R}\oplus E_1^{\mathbb R}$.  Stability of $E_1$
alone forces $c=0$ and hence the isometry alternative.

If every degree-$k\ge2$ harmonic projection of every first-harmonic
pullback vanished, completeness of the round scalar harmonic expansion
would put each continuous coordinate pullback in $E_0\oplus E_1$ in
$L^2$, hence everywhere by continuity.  The first assertion then
would force $T$ to be constant or orthogonal, proving the contrapositive.
The harmonic decomposition is that of
Theorem~\ref{p2c:thm:exact-reciprocal-scalar-bands} and its complete
round-sphere $L^2$ spectrum.
\end{proof}


\begin{proposition}[Sharp degree-two leakage from a nonlinear point map]
\label{p2c:prop:v280-sharp-quadratic-band-leakage}
Let $B=S^3(R)\subset\mathbb R^4$ and fix an auxiliary unit-radius
vector $e\in\mathbb R^4$ with $|e|=R$.  Define
\[
 \boxed{Q_e(x):=\frac{2\langle x,e\rangle}{R^2}x-e,
 \qquad |x|=R.}
 \tag{Q2-WIT}
\]
Then $Q_e:B\to B$ is continuous, nonconstant and non-isometric, but
for every $w\in\mathbb R^4$ its first-harmonic pullback
$f_w\circ Q_e$, $f_w(y)=\langle w,y\rangle$, belongs to
$E_0^{\mathbb R}\oplus E_2^{\mathbb R}$.  Moreover if $w\perp e$
and $w\ne0$, the entire pullback is a nonzero element of
$E_2^{\mathbb R}$, with no degree-$k\ge3$ component.  Thus the
higher-band obstruction of
Theorem~\ref{p2c:thm:v280-first-band-rigidity} can begin
\emph{exactly in degree two}; no larger resolution is required for a
mathematical witness.  The auxiliary vector $e$ is a proof parameter,
not a constituted physical axis or point-map dynamics.
\end{proposition}
\begin{proof}
Writing $a=\langle x,e\rangle/R^2$, the defining equation gives
$|2ax-e|^2=4a^2R^2-4a\langle x,e\rangle+R^2=R^2$.
Also $Q_e(-x)=Q_e(x)$, hence this map cannot be an orthogonal
isometry, and $Q_e(e)=e$ while $Q_e(x)=-e$ whenever
$x\perp e$, so it is not constant.  For each $w$,
\[
 (f_w\circ Q_e)(x)
 =\frac{2}{R^2}\langle x,e\rangle\langle x,w\rangle
  -\langle e,w\rangle.
\]
The homogeneous quadratic factor is a sum of a trace part and a
traceless harmonic quadratic.  Thus its restriction is in
$E_0^{\mathbb R}\oplus E_2^{\mathbb R}$ and has no higher band.
If $w\perp e$, its Euclidean Laplacian is
$\Delta_{\mathbb R^4}[(x\cdot e)(x\cdot w)]=2e\cdot w=0$,
so the pullback is a nonzero degree-two spherical harmonic.
\end{proof}

\begin{corollary}[Evaluation-dual test closure and nonlinear operation firewall]
\label{p2c:cor:v280-first-band-dual-firewall}
For a continuous point map $T$ of the generated reciprocal round role,
choose the evaluation test sort $E=E_1^{\mathbb R}$.  The exact
mathematical existence criterion of
Theorem~\ref{p2a:thm:v280-evaluation-dual-stabilizer}, together with
Theorem~\ref{p2c:thm:v280-first-band-rigidity}, says that an
$E_1^{\mathbb R}$-valued evaluation dual exists exactly when $T$ is a
round isometry.  For a nonconstant non-isometric $T$, the pullback
necessarily has a nonzero degree-$k\ge2$ harmonic component and does
not close in $E_0\oplus E_1$.  An enlarged or approximate test target
must be declared if that output is to be consumed physically.

The presence of such a mathematical pullback never independently
licenses a constructor or performs an observation: physical co-realization
still requires a frozen typed origin certificate as in
Theorem~\ref{p2a:thm:evaluation-dual-corealization}.  The existing
isometry dual of Theorem~\ref{p2c:thm:coordinate-generated-reciprocal-role}
passes this exact test without assuming an additional finite-band pairing.
\end{corollary}
\begin{proof}
The evaluation image is exactly the first-harmonic function module.
The stabilizer criterion identifies dual existence with pullback stability;
first-band rigidity identifies that stabilizer with the round isometry
group.  The higher-band alternative and physical origin boundary are
respectively the second conclusion of that theorem and the last clause
of the Paper-II-A criterion.
\end{proof}

\begin{corollary}[Connected first-band-closed point dynamics is isometric]
\label{p2c:cor:v280-first-band-semigroup-rigidity}
Suppose $T_t:B\to B$ is a jointly continuous point-map semigroup,
$t\ge0$, $T_0=\mathrm{id}$, $T_{s+t}=T_s\circ T_t$, and each pullback
preserves $E_0^{\mathbb R}\oplus E_1^{\mathbb R}$.  Then each
$T_t\in SO(4)|_B$ and there exists one real skew-symmetric matrix $A$
with $T_t(x)=e^{tA}x$ for every $t\ge0$.  In particular, such a
point-evolution cannot be a nontrivial contracting or mixing map on the
round sphere without exciting higher scalar bands.  The mathematical
negative-time extension of $e^{tA}$ does not license a physical backward
process, nor does the result constrain non-point-map quantum channels.
\end{corollary}
\begin{proof}
By the preceding theorem every $T_t$ is either constant or an orthogonal
isometry.  These two families are disjoint and separated in the uniform
metric on $C(B,B)$: an orthogonal isometry is surjective and has uniform
distance $2R$ from every constant map of value in $B$.  Joint continuity
on compact $B$ gives continuous dependence on $t$ in this metric.
Starting from the identity, the connected parameter domain cannot enter
the constant family.  The determinant cannot change continuously from
$+1$, so $T_t\in SO(4)$.  The induced finite-dimensional orthogonal
semigroup is norm-continuous; Lemma~\ref{p2a:lem:repair-finite-generator-regularity}
gives $T_t=e^{tA}$ and orthogonality gives $A^T=-A$.
\end{proof}

\begin{theorem}[Generated all-band reciprocal covariance]
\label{p2c:thm:all-band-reciprocal-covariance}
Let \(T\) be a constituted round isometry of the momentum parent and let
\(T^\vee=\iota_1T^{-1}\iota_1^{-1}\) be its generated contragredient action on
\(S_B^3\).  Define
\[
 U_x(T)f:=f\circ T^\vee .
\]
Then
\[
 [U_x(T),\Delta_x]=0,
 \qquad
 U_x(T)E_k^x=E_k^x,
 \qquad
 U_x(T)\mathcal P_x(m)=\mathcal P_x(m)
\]
for every \(k,m\ge0\).  If \(X\) is an infinitesimal generated isometry and
\[
 D_X:=\left.\frac{d}{ds}\right|_{s=0}U_x(e^{sX}),
\]
then
\[
 [D_X,\Delta_x]=0,
 \qquad
 D_XE_k^x\subseteq E_k^x,
 \qquad
 D_X\mathcal P_x(m)\subseteq\mathcal P_x(m).
\]
Thus the reciprocal carrier already generates the same isometry/transvection
operation ancestry on every scalar harmonic band.  No finite-band
\((\mathrm{PAIR}_{xp})\) premise is required for this action.
\end{theorem}

\begin{proof}
The map \(T^\vee\) is a round isometry by the generated reciprocal-role
construction.  Corollary~\ref{p2c:cor:hereditary-born-covariance} supplies the
hereditary reciprocity rule for every unmarked descendant; the present
calculation identifies the concrete target action for the scalar harmonic
constructor.  The scalar Laplace--Beltrami operator is natural under round
isometries, so pullback by \(T^\vee\) commutes with \(\Delta_x\).  Hence each
Laplacian eigenspace \(E_k^x\), and therefore every direct sum
\(\mathcal P_x(m)=\bigoplus_{k\le m}E_k^x\), is invariant.  Differentiating
the commutation relation and the invariant subspaces at the identity gives the
infinitesimal statements.
\end{proof}


\begin{theorem}[Canonical reciprocal scalar-harmonic intertwiner]
\label{p2c:thm:canonical-reciprocal-harmonic-intertwiner}
Let $E_k^p$ denote the degree-$k$ scalar harmonic eigenspace on the
constituted round momentum parent $A=S_A^3$, and let $E_k^x$ be the reciprocal
position eigenspace of Definition~\ref{p2c:def:reciprocal-scalar-resolution}.
Define
\[
 J_{xp}:L^2(S_A^3)\longrightarrow L^2(S_B^3),
 \qquad
 J_{xp}f:=f\circ\iota_1^{-1}.
\]
With the round comparison metric and normalized round measure transported
along $\iota_1$, the map $J_{xp}$ is unitary and
\[
 \boxed{\Delta_xJ_{xp}=J_{xp}\Delta_p},
 \qquad
 \boxed{J_{xp}(E_k^p)=E_k^x}
 \quad(k\ge0).
\]
For a constituted momentum-parent round isometry $T$, define the scalar
momentum pullback $U_p(T)f:=f\circ T$.  The generated reciprocal action of
Theorem~\ref{p2c:thm:all-band-reciprocal-covariance} then obeys
\[
 U_x(T)J_{xp}=J_{xp}U_p(T^{-1}).
\]
Hence, for every constituted finite subgroup $\Gamma$ of the momentum-parent
round-isometry group, with group-average projectors
\[
 \Pi_\Gamma^p:=\frac1{|\Gamma|}\sum_{\gamma\in\Gamma}U_p(\gamma),
 \qquad
 \Pi_\Gamma^x:=\frac1{|\Gamma|}\sum_{\gamma\in\Gamma}U_x(\gamma),
\]
one has
\[
 \boxed{\Pi_\Gamma^xJ_{xp}=J_{xp}\Pi_\Gamma^p}.
\]
Consequently $J_{xp}$ restricts to a unitary isomorphism
\[
 E_k^p\cap\operatorname{Ran}\Pi_\Gamma^p
 \xrightarrow{\ \sim\ }
 E_k^x\cap\operatorname{Ran}\Pi_\Gamma^x
\]
for every $k$.  This is a scalar-harmonic reciprocal-role bridge.  It is not
a map from reciprocal-position observables to the compact momentum operators
on $\mathcal H_1$, does not fix a dimensionful position--momentum
calibration, and does not actualize a state or covariance.
\end{theorem}

\begin{proof}
Theorem~\ref{p2c:thm:coordinate-generated-reciprocal-role} supplies
$\iota_1$ and its transported round comparison metric and measure.
Theorem~\ref{p2c:thm:all-band-reciprocal-covariance} supplies
$T^\vee=\iota_1T^{-1}\iota_1^{-1}$. These are precisely the isometry and
inverse-action conventions of Paper~1B,
Theorem~\ref{p1b:thm:harmonic-transport}, which gives unitarity, all-band
Laplacian intertwining, and the finite-group projector identity.
Apply its invariant-subspace restriction to the constituted subgroup
$\Gamma$. The transported scalar-resolution family is the growing position tower defined
here, with the fixed momentum carrier retaining its separate role.
\end{proof}

\subsection{Transverse shells and transported differential operators}

\begin{definition}[Born-derived transverse shell windows]
\label{p2c:def:born-derived-transverse-shell-windows}
Transport the fixed Paper-I-A Hopf marking along the reciprocal comparison
isometry \(\iota_1:S_A^3\to S_B^3\), including its handedness when the
comparison reverses orientation, and denote the resulting metric Hopf base of
the reciprocal round role by \(S^2_{H,x}\).  For \(\mu\ge1\), define
\[
 \boxed{\mathcal T_{x,\mu}^{\perp}
 :=\bigoplus_{\ell=0}^{\mu-1}\mathscr Y_\ell(S^2_{H,x}).}
\]
This is a transverse scalar-symbol/test window on the already derived
reciprocal geometry.  It is not an independent position carrier, site set, or
microscopic lattice.
\end{definition}

\begin{theorem}[Born-transported transverse shell reconstruction]
\label{p2c:thm:born-transported-transverse-shell-reconstruction}
For every \(k\ge0\), the canonical reciprocal scalar-harmonic intertwiner and
Paper~I-A's Hopf-marked diagonal shell-symbol theorem induce an equivariant
isomorphism
\[
 \boxed{\mathfrak H_k^x:
 E_k^x\xrightarrow{\sim}\mathcal T_{x,k+1}^{\perp}.}
\]
Equivalently,
\[
 \boxed{\dim E_k^x=(k+1)^2
 =\dim\mathcal T_{x,k+1}^{\perp}.}
\]
For every \(m\ge1\), the exact scalar-band decomposition therefore gives the
frontier-quotient reconstruction
\[
 \boxed{
 \mathcal P_x(m)/\mathcal P_x(m-1)
 \xrightarrow{\sim}
 E_m^x
 \xrightarrow{\mathfrak H_m^x}
 \mathcal T_{x,m+1}^{\perp},}
\]
so the newly exposed degree-\(m\) reciprocal-position shell has exact
transverse rank
\[
 \boxed{N_{\partial x}(m):=(m+1)^2.}
\]
The theorem reconstructs one scalar spectral shell as a transverse symbol
window.  It does not identify the accumulated band \(\mathcal P_x(m)\) with a
single transverse window, does not count microscopic position sites or
physical states, and does not provide an all-band map into
\(\operatorname{End}(\mathcal H_1)\).
\end{theorem}

\begin{proof}
Theorem~\ref{p2c:thm:canonical-reciprocal-harmonic-intertwiner} gives the
unitary all-band Born transport
\[
 J_{xp}:E_k^p\xrightarrow{\sim}E_k^x.
\]
Transport the fixed Hopf marking with the same comparison geometry.  Applying
Paper~I-A, Theorem~\ref{p1:thm:hopf-marked-diagonal-shell-symbol}, on the
transported marked round sphere gives an equivariant shell-symbol
isomorphism from \(E_k^x\) onto
\(\bigoplus_{\ell=0}^{k}\mathscr Y_\ell(S^2_{H,x})\), which is
\(\mathcal T_{x,k+1}^{\perp}\) by
Definition~\ref{p2c:def:born-derived-transverse-shell-windows}.  Finally,
Theorem~\ref{p2c:thm:exact-reciprocal-scalar-bands} gives the orthogonal
splitting
\[
 \mathcal P_x(m)=\mathcal P_x(m-1)\oplus E_m^x,
\]
so the quotient is canonically represented by \(E_m^x\).  The rank formula
is \(\dim E_m^x=(m+1)^2\).
\end{proof}

\begin{remark}[Accumulated bulk rank versus frontier rank]
\label{p2c:rem:bulk-versus-frontier-rank}
The two exact growth laws are differently typed:
\[
 \dim\mathcal P_x(m)
 =\frac{(m+1)(m+2)(2m+3)}6=\Theta(m^3),
 \qquad
 N_{\partial x}(m)=(m+1)^2=\Theta(m^2).
\]
The first is cumulative scalar test capacity through degree \(m\); the second
is the rank of the newly exposed shell and its complete transverse symbol
window.  Neither quantity is, by itself, an entropy or a count of ontic
spatial cells.
\end{remark}

\begin{theorem}[Reciprocal transport of natural Hodge operators]
\label{p2c:thm:v244-reciprocal-natural-bundle-transport}
Let \(\iota_1:S_A^3\to S_B^3\) be the generated reciprocal round isometry
used in Theorem~\ref{p2c:thm:canonical-reciprocal-harmonic-intertwiner}, and
let \(\varepsilon_{\iota}\in\{+1,-1\}\) be its orientation sign.  For every
form degree \(r\), define
\[
 J_{xp}^{(r)}:=(\iota_1^{-1})^*:
 \Omega^r(S_A^3)\longrightarrow\Omega^r(S_B^3).
\]
Then the exterior derivative, codifferential, Hodge Laplacian and Hodge star
obey
\begin{align}
 d_xJ_{xp}^{(r)}&=J_{xp}^{(r+1)}d_p,\\
 \delta_xJ_{xp}^{(r)}&=J_{xp}^{(r-1)}\delta_p,\\
 &\boxed{\Delta_x^{(r)}J_{xp}^{(r)}=J_{xp}^{(r)}\Delta_p^{(r)}},\\
 *_xJ_{xp}^{(r)}&=\varepsilon_{\iota}
 J_{xp}^{(3-r)}*_p .
\end{align}
Consequently the coexact one-form subspace is transported unitarily and, for
\(C:=*d\),
\[
 \boxed{C_xJ_{xp}^{(1)}
 =\varepsilon_{\iota}J_{xp}^{(1)}C_p}.
\]
Thus an orientation-preserving reciprocal comparison preserves the two curl
branches separately, while an orientation-reversing comparison exchanges
their signs; in both cases \(C^2=\Delta_1^{\rm co}\) and the coexact
Laplacian spectrum are preserved.  This is natural-bundle transport on the
reciprocal comparison geometry.  It does not identify these differential
operators with compact momentum operators on \(\mathcal H_1\), actualize a
gauge field or state, or fix a dimensionful position--momentum calibration.
\end{theorem}

\begin{proof}
Pullback commutes with \(d\).  Since \(\iota_1\) is a round isometry, its
pullback is unitary on forms and intertwines the metric adjoint of \(d\),
which gives the codifferential identity.  Adding \(d\delta+\delta d\) gives
the Hodge-Laplacian identity.  The Hodge-star identity is the standard
naturality formula for an isometry, with the single orientation sign
\(\varepsilon_{\iota}\).  Combining it with the exterior-derivative identity
gives the curl relation.  Coexactness is preserved by the codifferential
intertwiner, and unitarity follows from transport of the round metric and
measure already used in the scalar theorem.
\end{proof}

\begin{corollary}[Reciprocal spin-Dirac transport on a compatible spin lift]
\label{p2c:cor:v244-reciprocal-spin-dirac-transport}
On an orientation-preserving branch, choose the compatible lift of
\(\iota_1\) to the unique round spin structures and denote the induced
unitary spinor transport by \(J_{xp}^{\rm spin}\).  Naturality of the spin
connection and Clifford multiplication gives
\[
 \boxed{D_xJ_{xp}^{\rm spin}=J_{xp}^{\rm spin}D_p}.
\]
Hence the reciprocal comparison transports the spin-Dirac eigenspaces without
introducing a fermion occurrence, a mass term, or a fermionic dynamics law.
\end{corollary}

\begin{corollary}[Harmonic $\mathrm{PAIR}_{xp}$ is generated]
\label{p2c:cor:harmonic-pairxp-generated}
The scalar-harmonic bridge
\[
 \boxed{\mathrm{PAIR}_{xp}^{\rm harm}:=J_{xp}}
\]
is generated by the already constituted reciprocal-role map.  The stronger
operator-valued $\mathrm{PAIR}_{xp}^{\rm op}$ from reciprocal-position
observables or harmonics to specified compact momentum operators, and any
continuum/local--global comparison, remain separately typed interfaces.
\end{corollary}

\begin{corollary}[Exact compressed derivation and guarded multiplication words]
\label{p2c:cor:guarded-transvection-multiplication}
Let $\Pi_m^x$ project onto the complete scalar band $\mathcal P_x(m)$,
let $D_X$ be an infinitesimal isometry derivation, and let $f$ be smooth.
Set $D_X^{(m)}=\Pi_m^xD_X\Pi_m^x$ and
$M_f^{(m)}=\Pi_m^xM_f\Pi_m^x$.
On $\mathcal P_x(m)$,
\[
 [D_X^{(m)},M_f^{(m)}]=M_{D_Xf}^{(m)}
\]
without a degree guard for this single commutator.
A multiplication word with symbol degrees $s_1,\ldots,s_r$ is unaffected by
truncation on $\mathcal P_x(K)$ whenever $m\ge K+\sum_js_j$.
The commutator assertion requires that the generator preserve the chosen
projection; it does not automatically apply to arbitrary quotient cutoffs.
\end{corollary}
\begin{proof}
Complete harmonic bands are isometry-invariant, so
$[D_X,\Pi_m^x]=0$ on the smooth domain. Therefore
\[
 [D_X^{(m)},M_f^{(m)}]
 =\Pi_m^x[D_X,M_f]\Pi_m^x
 =\Pi_m^xM_{D_Xf}\Pi_m^x.
\]
The result is a generated operation identity, not an identification of
position harmonics with compact momentum operators.
The word bound follows from degree addition and containment of every
intermediate product. In general compressed multiplication is not
multiplicative; this conclusion removes the guard only for the displayed
single derivation commutator.
\end{proof}

\begin{remark}[Narrowed meaning of \(\mathrm{PAIR}_{xp}\)]
\label{p2c:rem:narrowed-pairxp}
Theorems~\ref{p2c:thm:first-symbol-pairing},
\ref{p2c:thm:all-band-reciprocal-covariance}, and
\ref{p2c:thm:canonical-reciprocal-harmonic-intertwiner},
Theorem~\ref{p2c:thm:v244-reciprocal-natural-bundle-transport}, and
Corollary~\ref{p2c:cor:guarded-transvection-multiplication}, close the
first-symbol, shared-generator, scalar-harmonic, and natural Hodge/curl
reciprocal content.  Hereafter
the unresolved symbol \((\mathrm{PAIR}_{xp}^{\rm op})\) denotes only a stronger typed map that assigns
specified reciprocal-position observables or harmonics to specified compact
momentum operators, or a stronger continuum/local--global comparison.  Such a
claim must still state its source and target, covariance, domain, and error
topology, and any dimensionful normalization remains a calibration.
\end{remark}

\begin{corollary}[Derived mathematical tower and typed auxiliary family]
\label{p2c:cor:derived-mathematical-position-tower}
The family
\[
 \mathscr T_x^{\rm math}:=\{\mathcal P_x(m)\}_{m\ge1}
\]
is a theorem output, with canonical inclusions and product filtration.  In
particular the correctly typed auxiliary comparison family is
\[
 \mathsf K_{\rm res}(m,1)
 :=\bigl(\mathcal P_x(m),\mathcal H_p\bigr),
 \qquad \mathcal H_p\cong\mathcal H_1,\quad m_p=1.
\]
Accordingly \((\mathrm{TWR}_x)\) is retired as an existence premise.  This
does not assert that any level beyond the primitive seed occurs physically.
\end{corollary}

\begin{corollary}[Exact nonlinear guard band]
\label{p2c:cor:nonlinear-position-guard-band}
Let \(K,s_1,\ldots,s_r\ge0\) and
\(f_j\in\mathcal P_x(s_j)\).  Multiplication by the word
\(f_1\cdots f_r\) maps
\[
 \mathcal P_x(K)\longrightarrow
 \mathcal P_x\!\left(K+\sum_{j=1}^rs_j\right).
\]
Consequently every truncation level
\(m\ge K+\sum_js_j\) recovers that multiplication word exactly on the
input band \(\mathcal P_x(K)\).  At the terminal shell, where this buffer is
absent, no full-carrier product-norm claim is made.
\end{corollary}

\begin{proof}
Repeated application of
\(\mathcal P_x(r)\mathcal P_x(s)\subseteq\mathcal P_x(r+s)\) gives the
range inclusion.  The orthogonal truncation is the identity on the resulting
guard band, which proves exact double compression there.
\end{proof}

\begin{remark}[Compatibility with scalar fixed-band recovery]
\label{p2c:rem:scalar-band-recovery-export}
The space \(\mathcal P_x(K)\) is exactly the scalar band
\(\bigoplus_{k=0}^{K}E_k^x\) used by Paper~II-D.  Thus II-D's fixed-band
compression theorem can consume this filtration directly.  That downstream
recovery result is not a premise of the reciprocal-role or seed derivations
above.
\end{remark}

\subsection{Finite frontiers, guard bands and multiplicative obstructions}

\begin{corollary}[Derived reciprocal seed and persistent momentum descriptor]
\label{p2c:cor:derived-reciprocal-resolution-seed}
\label{p2c:cor:conditional-reciprocal-seed-branch}
The primitive coordinate-generated position-test frontier is
\[
        F_{x,0}:=E_0^x\oplus E_1^x,
\]
and its exact scalar resolution is
\[
        \boxed{m_x^{\rm res}(F_{x,0})=1}.
\]
Together with the derived momentum carrier, the correctly typed seed is
\[
 \boxed{
 \mathsf K_{\rm res}(1,1)
 :=\bigl(\mathcal P_x(1),\mathcal H_p\bigr),
 \qquad
 \mathcal H_p\cong\mathcal H_1,\quad m_p=1.}
\]
The first entry is a five-dimensional scalar reciprocal-resolution window;
the second is the fifteen-dimensional structural momentum-role carrier.  No
isomorphism \(\mathcal P_x(1)\cong\mathcal H_1\) is asserted.  The momentum
label \(m_p=1\) persists at every later cumulative constitutional stage at
which the frozen object \(\mathcal H_p\) is referenced.
\end{corollary}

\begin{proof}
The scalar unit belongs to \(E_0^x\), while the coordinate-generated
reciprocal construction retains the complete nonzero irreducible sector
\(E_1^x\).  Hence \(F_{x,0}\subseteq\mathcal P_x(1)\), but
\(F_{x,0}\not\subseteq\mathcal P_x(0)=E_0^x\), proving the minimum is one.
Paper~II-B supplies \(\mathcal H_p\cong\mathcal H_1\) and \(m_p=1\).
Paper~II-A's cumulative chronology preserves every frozen carrier distinction
and typing under \(P_n\hookrightarrow P_{n+1}\).  A later embedding,
projection, quotient, or transition involving a new carrier is a new
inter-carrier operation; it does not retype the representation descriptor of
the already frozen \(\mathcal H_p\).  Thus no separate
\((\mathrm{LVL}_{x,0}^{1})\) or \((\mathrm{FRZ}_p)\) premise is required.
\end{proof}

\begin{definition}[Finite reciprocal-position frontier history]
\label{p2c:def:finite-position-frontier-history}
The occurrence \((\mathrm{HIST}_x^{\rm fin})\) is an outcome-independent
physical history
\[
 \widetilde t\longmapsto F_x(\widetilde t),
 \qquad
 F_{x,0}\subseteq F_x(\widetilde t)\subset\mathcal A_x^{\rm alg},
\]
of finite-dimensional generated position-test frontiers.  Its resolution is
the derived statistic
\[
 m_x(\widetilde t):=m_x^{\rm res}(F_x(\widetilde t)).
\]
The history is \emph{retained} when
\(s\le t\Rightarrow F_x(s)\subseteq F_x(t)\), and is
\emph{complete-band} when
\(F_x(t)=\mathcal P_x(m_x(t))\).  A sparse frontier and the complete band
with the same value of \(m_x\) are not identified.
\end{definition}

\begin{proposition}[Least resolution of later generated position frontiers]
\label{p2c:prop:least-generated-position-resolution}
Let
\[
        F_{x,0}\subseteq F_n\subset\mathcal A_x^{\rm alg}
\]
be the complete finite-dimensional span of scalar reciprocal tests generated
and retained at constitutional stage \(n\).  Then
\[
 \boxed{
 m_x[n]
 :=m_x^{\rm res}(F_n)
 =\max\{k:\operatorname{pr}_{E_k^x}(F_n)\ne\{0\}\}}
\]
exists and is finite.  Paper~II-A's cumulative chronology gives
\(F_n\subseteq F_{n+1}\), and therefore
\[
        m_x[n]\le m_x[n+1].
\]
Writing \(m_n=m_x[n]\), an exact unfolding step occurs precisely when
\[
 \boxed{
 m_x[n+1]>m_n
 \iff
 \exists\,k>m_n:
 \operatorname{pr}_{E_k^x}(F_{n+1})\ne\{0\}.}
\]

For a nonempty finite comparison set
\(G\subset C(S^3_B,\mathbb C)\) not assumed algebraic, put
\[
 e_m(G):=\max_{f\in G}\inf_{g\in\mathcal P_x(m)}
             \|f-g\|_\infty.
\]
For every \(\varepsilon>0\), the controlled resolution
\[
 m_x(G;\varepsilon)
 :=\min\{m:e_m(G)\le\varepsilon\}
\]
also exists.  It is monotone under inclusion of finite comparison sets at
fixed \(\varepsilon\).
\end{proposition}

\begin{proof}
Choose a finite basis of \(F_n\).  Every basis element lies in the algebra
generated by \(E_0^x\oplus E_1^x\), hence has finite coordinate-polynomial
degree and lies in some \(\mathcal P_x(m)\).  The largest of these finitely
many degrees gives a containing band, so the defining subset of
\(\mathbb Z_{\ge0}\) is nonempty and has a least element.  The direct-sum
decomposition in
Theorem~\ref{p2c:thm:exact-reciprocal-scalar-bands} identifies that least
element with the highest nonzero harmonic projection.

Every retained position test in \(F_n\) belongs to the frozen presentation
\(P_n\).  Paper~II-A gives \(P_n\hookrightarrow P_{n+1}\) without deletion
or retyping, so \(F_n\subseteq F_{n+1}\).  Every band containing
\(F_{n+1}\) then contains \(F_n\), which proves monotonicity.  The same
highest-harmonic formula proves the displayed growth equivalence.

For the controlled version, Stone--Weierstrass density gives, for each
\(f\in G\), approximants from sufficiently large \(\mathcal P_x(m)\).
Because \(G\) is finite and the bands are nested, one common finite \(m\)
works for every \(f\in G\), so \(e_m(G)\downarrow0\).  The threshold set is
therefore nonempty and has a least integer.  Enlarging \(G\) can only increase
\(e_m(G)\), proving the final monotonicity statement.
\end{proof}


\begin{theorem}[Complete growth on an admitted one-step scalar rule]
\label{p2c:thm:v233-stagewise-band-growth}
Suppose a licensed stage rule with retained scalar outputs is exactly
\[
 F_0=\mathcal P_x(1),\qquad
 F_{n+1}=\operatorname{span}(F_n\cup E_1^xF_n).
\]
Then \(F_n=\mathcal P_x(n+1)\). Each step \(n\to n+1\) adds the
complete harmonic shell of dimension \((n+3)^2\).
Continued occurrence of this rule exhausts every finite band.
On such a history, completeness and strict resolution growth require
no independent assumptions about an integer trajectory.
\end{theorem}
\begin{proof}
Every monomial of degree at most \(m+1\) is either of degree at most
\(m\), or a coordinate times a monomial of degree at most \(m\).
Lemma~\ref{p1b:lem:v233-weighted-band-products} gives
\(\mathcal P_x(m+1)=\mathcal P_x(m)+
\operatorname{span}(E_1^x\mathcal P_x(m))\).
Induction proves the formula; \(\dim E_{n+2}=(n+3)^2\).
Each stage is finite, and continued iteration eventually includes
any fixed band.
\end{proof}

The rule includes output retention and continued iteration. It is
not unlimited multiplication closure within one stage. That closure
would contain \((q_0+iq_1)^r\) for every \(r\), nonzero harmonics
of distinct degrees, and would be infinite-dimensional.
Hopf pullback instead commutes with the round Laplacian and can act
within \(\mathcal P_x(1)\) indefinitely. It does not force this rule,
its repeated occurrence or the assignment of stages to physical
time. The fixed-frontier branch below remains admitted.


\begin{theorem}[Complete guard-band inclusion iff exact coordinate words]
\label{p2c:thm:v268-sharp-word-guard-equivalence}
Let $F\subset C^\infty(S^3_x(R_x),\mathbb C)$ be a finite-dimensional
reciprocal-position test subspace with
$\mathcal P_x(K)\subseteq F$, let $\Pi_F$ be the orthogonal
$L^2$-projection, and put
$\widehat X_a^F=\Pi_FM_{x_a}\Pi_F|_F$ for $a=1,\ldots,4$.
For any nonnegative integers $K,d$, the following conditions are
\emph{equivalent}:
\begin{enumerate}
\item $\mathcal P_x(K+d)\subseteq F$ (complete guard-band inclusion);
\item for every coordinate word $w=x_{a_s}\cdots x_{a_1}$ of
length $0\le s\le d$ and every $\psi\in\mathcal P_x(K)$,
\[
 w(\widehat X^F)\psi=M_{w}\psi
 \qquad\text{as vectors in }L^2(S^3_x).
\]
\end{enumerate}
No $\Spin(4)_x$-invariance of $F$ is required.  Under either
condition, the compression remainder is identically zero in operator
norm on the stated input band.  In addition, for any round Killing
field derivation $D_\xi$ and $K+d\ge1$,
\[
 \bigl[\Pi_FD_\xi\Pi_F,\widehat X_a^F\bigr]
 \big|_{\mathcal P_x(K+d-1)}
 =\bigl[D_\xi,M_{x_a}\bigr]
 \big|_{\mathcal P_x(K+d-1)}.
\]
Here $\mathcal P_x(K+d-1)$ is included in $F$; all operators act
on smooth, finite-band vectors in the displayed equality.
\end{theorem}
\begin{proof}
The coordinate polynomial filtration obeys
$M_{x_a}\mathcal P_x(r)\subseteq\mathcal P_x(r+1)$.
If the guard band is contained in $F$, every intermediate result of a
word with length $s\le d$ on $\mathcal P_x(K)$ lies in $F$.
Therefore no intermediate projection changes the word, proving (i)
$\Rightarrow$(ii) and the claimed zero operator norm.
Conversely, the left side of (ii) always has range in $F$.
The span of the right sides is \emph{exactly}
\[
 \operatorname{span}_{\substack{0\le s\le d\\a_1,\ldots,a_s}}
 M_{x_{a_s}}\cdots M_{x_{a_1}}\mathcal P_x(K)
 =\mathcal P_x(K+d).
\]
Indeed, each coordinate monomial of degree at most $K+d$ factors
into at most $d$ coordinate generators times a polynomial of degree
at most $K$; restricted polynomials span the full spherical-harmonic
filtration.  Hence (ii) implies (i).  For the commutator,
$D_\xi\mathcal P_x(r)\subseteq\mathcal P_x(r)$ for each $r$,
while $M_{x_a}\mathcal P_x(K+d-1)\subseteq\mathcal P_x(K+d)$.
Both terms in the commutator thus remain in $F$ throughout, so
all intervening compressions are identities.
\end{proof}

\begin{corollary}[Independent word-generation certificate and its occurrence boundary]
\label{p2c:cor:v268-primitive-word-guard}
For $N\ge1$, let the \emph{cumulative full coordinate-word span} of
$F_{x,0}=\mathcal P_x(1)$ be
\[
 \mathcal W_x^{(N)}(F_{x,0})
 :=\operatorname{span}\{M_{x_{a_s}}\cdots M_{x_{a_1}}f:
 0\le s\le N-1,\ f\in F_{x,0}\}.
\]
Then the exact algebraic identity is
\[
 \boxed{\mathcal W_x^{(N)}(F_{x,0})=\mathcal P_x(N).}
\]
Consequently an actual frontier $F_x(\widetilde t)$ that retains
\emph{all} these coordinate words through depth $N-1$ contains the
complete guard band $\mathcal P_x(N)$, independently of any claim
based only on its highest-resolution integer.
A sparse frontier with $m_x^{\rm res}\ge N$ need not contain the band.
On the separately occurring sharp spectral frontier of
Paper~III-F, Corollary~\ref{p3f:cor:sharp-frontier-discharges-generated-boundary},
the band inclusion follows from the spectral projector itself;
no word-generation occurrence is additionally required there.
\end{corollary}
\begin{proof}
Every generator word in the displayed set has polynomial degree at
most $N$, proving containment in $\mathcal P_x(N)$.  Conversely,
any coordinate monomial of degree $r\le N$ is a word of length at
most $N-1$ applied either to a primitive coordinate (if $r\ge1$)
or to $1$ (if $r=0$).  Such monomials span the restricted
polynomials of degree at most $N$, namely $\mathcal P_x(N)$.
No assumption about actual frontier histories enters this algebraic
identity.  Retention of the full word span gives the physical
inclusion, while the source definition of $m_x^{\rm res}$ only
records the largest nonzero harmonic projection, not completeness.
\end{proof}

\begin{theorem}[A single invariant scalar detects the entire reciprocal guard band]
\label{p2c:thm:v269-single-scalar-guard}
Fix $K,d\in\mathbb Z_{\ge0}$, let $P_K=\Pi_{\mathcal P_x(K)}$
be the orthogonal projector in $L^2(S^3_x(R_x))$, and let $F$ be
\emph{any} finite-dimensional position test space with projector $P_F$.
No initial band containment or $\Spin(4)_x$-invariance of $F$ is
assumed. Set $u_a=x_a/R_x$, $U_a=M_{u_a}$ ($1\le a\le4$), so that
$\sum_aU_a^2=I$, and introduce the finite-rank positive operator
\[
 B_{K,d}:=\sum_{s=0}^{d}\ \sum_{a_1,\ldots,a_s=1}^{4}
 U_{a_s}\cdots U_{a_1}P_K U_{a_1}\cdots U_{a_s},
 \qquad \text{the $s=0$ term being $P_K$}.
\]
The nonnegative \emph{single-scalar guard defect} is
\[
 \mathfrak L_{K,d}(F):=\operatorname{Tr}((I-P_F)B_{K,d})
 =\sum_{s=0}^{d}\sum_{a_1,\ldots,a_s}
 \big\|(I-P_F)U_{a_s}\cdots U_{a_1}P_K\big\|_{\rm HS}^{2}.
\]
Then, without any additional frontier-constructor hypothesis,
\[
 \boxed{\quad \mathfrak L_{K,d}(F)=0
 \quad\Longleftrightarrow\quad \mathcal P_x(K+d)\subseteq F.\quad}
\]
This scalar is frame-independent:
$\mathfrak L_{K,d}(gF)=\mathfrak L_{K,d}(F)$ for every
$g\in\Spin(4)_x$. Defining the continuous band deficit
\[
 \Delta_{K+d}(F):=\operatorname{Tr}
 \bigl((I-P_F)\Pi_{\mathcal P_x(K+d)}\bigr)\ge0,
\]
one has the \emph{quantitative, sharp-lower-constant} estimate
\[
 \boxed{\quad
 \frac{K+1}{2^d(K+d+1)}\,\Delta_{K+d}(F)
 \ \le\ \mathfrak L_{K,d}(F)
 \ \le\ (d+1)\,\Delta_{K+d}(F).\quad}
\]
The lower factor is the exact smallest eigenvalue of $B_{K,d}$ on
$\mathcal P_x(K+d)$, attained on the top shell $E_{K+d}^x$.
\end{theorem}
\begin{proof}
All operators $U_a$ are self-adjoint contractions.  With the unital
completely positive map $\mathcal T(A)=\sum_{a=1}^4U_aAU_a$,
$B_{K,d}=\sum_{s=0}^d\mathcal T^s(P_K)$; hence
$0\le\mathcal T^s(P_K)\le I$. Its range lies in
$\mathcal P_x(K+s)$. The trace equality is obtained by writing every
summand as $A_w A_w^*$, with $A_w=U_wP_K$, and cycling the finite-rank
trace. The union of all their ranges spans
\[
 \operatorname{span}_{0\le s\le d}\{U_w\mathcal P_x(K)\}
 =\mathcal P_x(K+d),
\]
by exact coordinate-polynomial factorization.  Because each summand
of the trace is nonnegative, the scalar vanishes exactly when all
these ranges are contained in $F$.  This proves the iff, including
that $P_K\subseteq F$ is a \emph{consequence} rather than a premise.

To quantify the defect, $B_{K,d}$ commutes with $\Spin(4)_x$ and is
supported on $\mathcal P_x(K+d)$. On normalized $S^3$, the addition
kernel of $\Pi_{E_j^x}$ is $(j+1)\mathrm U_j(t)$,
$t=(x\cdot y)/R_x^2$, where $\mathrm U_j$ is the Chebyshev polynomial
of the second kind. The kernel of $\mathcal T^s(P_K)$ is
\[
 t^s\sum_{k=0}^K(k+1)\mathrm U_k(t),\qquad
 t\mathrm U_k(t)=\tfrac12\bigl(\mathrm U_{k-1}(t)+
 \mathrm U_{k+1}(t)\bigr),\quad \mathrm U_{-1}=0.
\]
Write $B_{K,d}=\sum_{j=0}^{K+d}b_j\Pi_{E_j^x}$.
The $s=0$ term gives $b_j\ge1$ for $j\le K$. For
$K<j\le K+d$, the monotone up-step chain of length $j-K$
starting at $k=K$ contributes
\[
 b_j\ge\frac{K+1}{(j+1)2^{j-K}}
 \ge\frac{K+1}{(K+d+1)2^d}=:c_{K,d}.
\]
For $j=K+d$ only the length-$d$ all-up chain from $k=K$
contributes, so $b_{K+d}=c_{K,d}$ exactly. The unital-positivity
bound above gives $b_j\le d+1$. Therefore
\[
 c_{K,d}\Pi_{\mathcal P_x(K+d)}\ \le B_{K,d}\le
 (d+1)\Pi_{\mathcal P_x(K+d)}.
\]
Taking the positive trace against $I-P_F$ proves both inequalities,
even for frontiers mixing harmonic shells. The covariance of $P_K$
and the contracted sum over the four coordinate indices prove
$\Spin(4)_x$-invariance of $B_{K,d}$ and of the stated scalar under
$F\mapsto gF$.
\end{proof}

\begin{corollary}[Quantified norm error for a candidate incomplete frontier]
\label{p2c:cor:v269-compressed-word-stability}
If in addition $\mathcal P_x(K)\subseteq F$, let
$\widehat U_a^F=P_FU_aP_F|_F$ and let $w$ be any coordinate word of
length $s$ with $1\le s\le d$. Then
\[
 \left\|w(\widehat U^F)|_{\mathcal P_x(K)}
       -U_w|_{\mathcal P_x(K)}\right\|_{\rm op}
 \le s\,\sqrt{\mathfrak L_{K,d}(F)}.
\]
For unnormalized coordinates $x_a=R_x u_a$ the right side is
$sR_x^s\sqrt{\mathfrak L_{K,d}(F)}$. On a shell-complete frontier,
$\Delta_{K+d}(F)=\sum_{0\le j\le K+d,\,E_j^x\not\subset F}(j+1)^2$;
for a general non-invariant $F$ the deficit may be nonintegral.
The certificate classifies a \emph{given} $F$; it does not cause the
physical occurrence of the word-generated frontier, and measurement
of this scalar requires an independently licensed same-owner test.
\end{corollary}
\begin{proof}
Let $\psi\in\mathcal P_x(K)$ and compare the projected and
unprojected successive word prefixes. Initially their difference
vanishes because $P_K\subseteq F$. Each subsequent difference is
bounded by the preceding one plus
$\|(I-P_F)U_{a_r}\cdots U_{a_1}\psi\|$.
The latter is at most
$\|(I-P_F)U_{a_r}\cdots U_{a_1}P_K\|_{\rm HS}\|\psi\|$
and hence at most $\sqrt{\mathfrak L_{K,d}(F)}\|\psi\|$.
Summing over $r=1,\ldots,s$ gives the estimate. The shell formula
follows by taking traces against irreducible harmonic projectors.
\end{proof}

\begin{remark}[Exact stopping point for position unfolding]
\label{p2c:rem:position-unfolding-stopping-point}
Proposition~\ref{p2c:prop:least-generated-position-resolution} supplies a
resolution functor for content that a later theory actually generates.  It
does not generate that content or the G0 history witness mapping physical time
to constitutional stages.  If the later position dependency frontier is empty, the recursive
fixed-point property of
Corollary~\ref{p2:cor:recursive-gp-fixed-point} gives
\[
        F_n=F_{x,0},
        \qquad
        m_x[n]=1
\]
at every stage.  Thus the mathematical tower
\(\mathscr T_x^{\rm math}\) is already derived, and no
\((\mathrm{TWR}_x)\) premise remains.  The residual physical datum is the
frontier history \((\mathrm{HIST}_x^{\rm fin})\), not a freely chosen integer
function \((\mathrm{TRJ}_x)\).  A later dynamics theorem must generate a
nonzero \(E_k^x\) component above the current frontier and relate the actual
frontiers to physical time before a nonconstant \(m_x(\widetilde t)\) may be
claimed.  Retained inclusion then proves monotonicity; no separate
\((\mathrm{MON}_x)\) premise is needed on that branch.
\end{remark}

\subparagraph{Supplementary finite uncertainty and POVM outcomes.}
Finite Robertson bounds are sector-dependent.  On the protected block,
extremal commutator states saturate the bound, whereas coordinate eigenstates
exclude a universal state-independent positive coordinate-product minimum.
Normalized diagonal-\(SU(2)\)-covariant POVMs exist, but localization and
moment matching require separate interfaces.  The constructions and
obstructions are collected in
Paper~II-F, Section~\ref{p2c:app:extended-finite-coordinate-structure}.

\begin{remark}[Finite-exact and recovery export]
\label{p2c:rem:finite-exact-recovery-export}
The linear transvection--coordinate commutators, equivariant projections, and
finite leakage identities in this paper are finite-exact.  They provide the
operator input for the typed parent/finite recovery analysis of Paper~II-D.
No statement here identifies a momentum-side \(m_p\to\infty\) limit: the
structural momentum-role theorem imported from II-B is \(m_p=1\).  Any large-level
recovery therefore uses the reciprocal position factor and the explicit band
inclusion/compression maps of the common Paper-II contract.
\end{remark}


\begin{lemma}[Exact scalar-harmonic product span on the reciprocal round sphere]
\label{p2c:lem:v270-homogeneous-product-span}
For full degree-$k$ and degree-$s$ scalar harmonics on $S^3_x$,
\[
\boxed{\operatorname{span}(E_k^xE_s^x)
 =\bigoplus_{j=0}^{\min(k,s)} E_{k+s-2j}^x.}
\]
\end{lemma}
\begin{proof}
Under $S^3=SU(2)$, $E_k^x$ is the full coefficient space of the
spin-$k/2$ irreducible. Products of matrix coefficients span the
coefficient space of the tensor product. The Clebsch--Gordan rule
$V_{k/2}\otimes V_{s/2}=\bigoplus_{j=0}^{\min(k,s)}V_{(k+s-2j)/2}$
shows that every displayed harmonic component occurs and no other
component occurs.
\end{proof}

\begin{theorem}[Smallest invariant prefix guard for homogeneous observable words]
\label{p2c:thm:v270-homogeneous-prefix-guard}
Let $K,s_1,\ldots,s_r\ge0$, let $F\subset L^2(S^3_x)$ be a
finite-dimensional $\Spin(4)_x$-invariant subspace with $E_K^x\subset F$,
and define $R_0=\{K\}$,
\[
R_j=\bigcup_{k\in R_{j-1}}
 \{|k-s_j|,|k-s_j|+2,\ldots,k+s_j\},\qquad
F_{\min}=\bigoplus_{k\in\cup_{j=0}^{r}R_j}E_k^x.
\]
For every $u\in E_K^x$, every $f_j\in E_{s_j}^x$ and every prefix
$1\le j\le r$, all repeated $P_F$-compressions equal the corresponding
uncompressed multiplication prefixes \emph{if and only if}
$F_{\min}\subseteq F$.  In this universal homogeneous-prefix class
$F_{\min}$ is the smallest invariant exact-recovery frontier.
\end{theorem}
\begin{proof}
The span $W_0$ of possible inputs is $E_K^x$.  Inductively the span
$W_j$ of all length-$j$ products equals the span of
$E_{s_j}^xW_{j-1}$, and the previous lemma gives
$W_j=\bigoplus_{k\in R_j}E_k^x$. If every compressed prefix equals
its uncompressed output then each $W_j\subseteq F$, hence
$F_{\min}\subseteq F$. Conversely this inclusion makes every
intermediate compression the identity on its actual output; induction
proves exactness of every prefix.
\end{proof}
\begin{corollary}[Observable-specific sparse guard and differential compatibility]
\label{p2c:cor:v270-homogeneous-sparse-recovery}
One $E_2^x$ multiplier acting on an $E_2^x$ input needs only
$E_0^x\oplus E_2^x\oplus E_4^x$ (dimension $35$), rather than
$\mathcal P_x(4)$ (dimension $55$), for universal homogeneous
prefix exactness. Interspersed round Killing derivatives and
spectral multipliers $f(\Delta)$ preserve the same \emph{sufficient}
guard because they preserve each $E_k^x$. Such added operations may
weaken necessity when they annihilate sectors.  A physically occurring
frontier or actual multiplication word is \emph{not} inferred from
these purely mathematical comparisons.
\end{corollary}
\begin{proof}
The product-span lemma with $k=s=2$ gives degrees $0,2,4$.
Round Killing Lie derivatives commute with the Laplacian and
spectral multipliers act diagonally on harmonic degrees, so neither
creates a new band between multiplication steps.
\end{proof}

\begin{theorem}[Maximal multiplicative obstruction on the fixed finite carrier]
\label{p2c:thm:v270-maximal-covariant-multiplication-defect}
For any unital positive $\Spin(4)$-covariant linear map
$\Phi:C(S^3_x)\to\operatorname{End}(\mathcal H_1)$, define
\[
D(\Phi)=\sup_{\substack{f=f^*\in C(S^3_x)\\\|f\|_\infty\le1}}
 \|\Phi(f^2)-\Phi(f)^2\|.
\]
Then $\boxed{D(\Phi)=1}$.  Thus no such finite-carrier map can
implement the full global scalar function algebra multiplicatively.
This does not obstruct the exact scalar-parent multipliers or
observable-specific finite guard-band recovery.
\end{theorem}
\begin{proof}
By the Clebsch--Gordan decomposition of
$\operatorname{End}(\mathcal H_1)$, its diagonal scalar-harmonic
support is confined to $E_0^x\oplus E_1^x\oplus E_2^x$.
Equivariance and boundedness (automatic for positive maps) therefore
annihilate all higher-harmonic content. In ambient normalized
coordinates take $f_t(q)=\tanh(tq_1q_2q_3)$.
Reflection of one of the three coordinate factors omitted by each
monomial of total degree at most two shows that $f_t$ is orthogonal
to $E_0^x\oplus E_1^x\oplus E_2^x$; hence $\Phi(f_t)=0$.
Moreover $0\le f_t^2\le1$, and $f_t^2\to1$ almost everywhere, so
$a_t=\int f_t^2\,d\mu\to1$.
By $\Spin(4)$-covariance and transitivity,
$\int_G U_g\Phi(f_t^2)U_g^*\,dg=a_tI$;
consequently $\|\Phi(f_t^2)\|\ge a_t$, proving $D(\Phi)\ge1$.
A unital positive map out of a commutative $C^*$-algebra is completely
positive. Kadison--Schwarz gives
$0\le\Phi(f^2)-\Phi(f)^2\le\Phi(f^2)\le I$
for every real contraction $f$, so $D(\Phi)\le1$.
\end{proof}
\begin{proposition}[Explicit normalized cubic obstruction]
\label{p2c:prop:v270-cubic-defect-witness}
Write $z_1=q_1+iq_2$, $z_2=q_3+iq_4$ on $S^3$ and put
$h=(3\sqrt3/2)\operatorname{Re}(z_1^2z_2)$.
Then $h\in E_3$, $\|h\|_\infty=1$, $\Phi(h)=0$ for every
$\Phi$ of the previous theorem, and
\[
\boxed{\|\Phi(h^2)-\Phi(h)^2\|\ge 9/32.}
\]
\end{proposition}
\begin{proof}
The real part of the homogeneous holomorphic cubic is harmonic.
With $r=|z_1|^2$, $\max r\sqrt{1-r}=2/(3\sqrt3)$, so
$\|h\|_\infty=1$. On normalized $S^3$, $r$ is uniform on $[0,1]$
and the phase average of $\cos^2$ is $1/2$. Therefore
$\int h^2\,d\mu=(27/4)(1/2)\int_0^1r^2(1-r)\,dr=9/32$.
Since $\Phi(h)=0$, averaging $\Phi(h^2)$ over the group yields
$(9/32)I$, giving the lower bound.
\end{proof}
\begin{remark}[Scope of the homogeneous and global no-go results]
\label{p2c:rem:v270-local-global-scope}
The exact sparse guard is for \emph{homogeneous} input/symbol degrees;
requiring all degrees up to the same bounds restores the earlier
complete-band guard. The maximal scalar obstruction concerns global
covariant maps into a fixed finite $\mathcal H_1$, not the parent
$L^2(S^3)$, and supplies no occurrence claim for any measurement.
\end{remark}


\paragraph{Sharp-frontier history as a conditional downstream application.}
The physical-history boundary of the scalar frontier is discharged on the occurring
sharp-frontier branch by Paper~III-F,
Corollary~\ref{p3f:cor:sharp-frontier-generates-position-history}.
Its complete spectral frontier at the independently owned scale $\Lambda_F$
and retained scale history generate $F_x(t)$ and $m_x(t)$; no second
history token is needed. The remaining input on that branch is the physical
sharp-frontier word and its scale, with threshold reach where growth is
claimed. The spectral-action cutoff alone does not supply that word.
Local pairing is already derived; a stronger comparison must specify its
target, covariance, domain and error. The universal all-band target in the
fixed $\operatorname{End}(\mathcal H_1)$ is excluded, as stated above.
Paper~II-A, Section~\ref{p2a:sec:application-obligation-record}, distinguishes
such input and domain obligations from duplicate structural selectors.


\section{Born compatibility and operator-germ transport}
\label{paper2c:sec:organized-born}

\subsection{Operator-germ reciprocal pairing after common-symbol closure}
\label{p2c:sec:operator-germ-pairing-update}

Paper~I-A's round first-symbol theorem fixes the radial quotient of the finite
position-coordinate germ, while
Theorem~\ref{p2:thm:canonical-transvection-bridge} fixes the compact
transvection complement.  Their canonical operator relation appears after the
radial first-symbol quotient; it is not an identification of the global
four-vector representation with the compact Lie algebra.

Let $a\in A$, $b=\iota_1(a)\in B_A^{(1)}$, and let
\[
 \overline D_b:
 E_{1,\RR}^x/\RR\widetilde f_a
 \overset\sim\longrightarrow T_b^*B_A^{(1)}
\]
be the already derived first-symbol quotient.  With the round metric $h$ and
the symmetric-pair evaluation isomorphism
$\operatorname{ev}_a:\mathfrak p_a\to T_aA$, define
\begin{equation}
 \label{p2c:eq:theta-xp-local}
 \Theta_a
 :=\operatorname{ev}_a^{-1}\circ\sharp_{h,a}\circ
 (d\iota_{1,a})^*\circ\overline D_b.
\end{equation}

\begin{definition}[Moving infinitesimal transvection-operator bundle]
\label{p2c:def:moving-transvection-operator-bundle}
Let
\[
 U_1:\Spin(4)\longrightarrow\mathcal U(\mathcal H_1)
\]
be the retained unitary representation and let
\(dU_1:\mathfrak{spin}(4)\to\mathfrak u(\mathcal H_1)\) be its derived
representation.  In the Hermitian convention set
\[
 \mathscr L_1(Z):=i\,dU_1(Z),
 \qquad
 \widehat K(Z):=\frac{\hbar}{\Rmom}\,\mathscr L_1(Z).
\]
The sign and Lie-algebra basis normalization are those already fixed in
Definitions~\ref{p2:def:composition-versus-transvection} and
\ref{p2:def:compact-momentum-generators}; at the basepoint this means
\begin{equation}
 \label{p2c:eq:Khat-basepoint-normalization}
 \widehat K(\mathsf p_i)
 =\frac{\hbar}{\Rmom}(L_i^+-L_i^-)
 =\widehat P_i,
 \qquad i=1,2,3,
\end{equation}
for the corresponding normalized transvection basis
\(\{\mathsf p_i\}\subset\mathfrak p_e\).

For the transported complements \(\mathfrak p_a\), define
\[
 \mathscr P:=\bigsqcup_{a\in A}\{a\}\times\mathfrak p_a,
 \qquad
 \mathscr K_a:=\widehat K(\mathfrak p_a)
 \subset\operatorname{End}(\mathcal H_1)_{\rm sa},
 \qquad
 \mathscr K:=\bigsqcup_{a\in A}\{a\}\times\mathscr K_a.
\]
Thus \(\mathscr K\to A\) is the moving transvection-operator bundle.  No
new scale is introduced: its only dimensionful factor is the previously
derived momentum unit \(\hbar/\Rmom\).
\end{definition}

\begin{theorem}[Canonical local operator-valued position/transvection pairing]
\label{p2c:thm:operator-germ-xp-pairing}
The maps \eqref{p2c:eq:theta-xp-local} are isomorphisms
\[
 E_{1,\RR}^x/\RR\widetilde f_a\overset\sim\longrightarrow\mathfrak p_a
\]
which glue covariantly under the generated $\Spin(4)$ action.  The restrictions
\(\widehat K|_{\mathfrak p_a}:\mathfrak p_a\to\mathscr K_a\) are
isomorphisms, and
\[
 \boxed{
 \mathcal Q_a^{xp}
 :=\widehat K|_{\mathfrak p_a}\circ\Theta_a:
 E_{1,\RR}^x/\RR\widetilde f_a
 \overset\sim\longrightarrow\mathscr K_a}
\]
is a canonical first-order operator-germ pairing.  For every
\(g\in\Spin(4)\), it obeys
\begin{equation}
 \label{p2c:eq:operator-germ-covariance}
 \mathcal Q_{ga}^{xp}\bigl(g\cdot[\widetilde f]\bigr)
 =U_1(g)\mathcal Q_a^{xp}([\widetilde f])U_1(g)^*.
\end{equation}
Hence the local maps glue to an equivariant bundle isomorphism
from the radial first-symbol quotient bundle to \(\mathscr K\), and no
independent first-order \(\mathrm{PAIR}_{xp}\) map remains.
\end{theorem}

\begin{proof}
The kernel calculation in
Theorem~\ref{p2c:thm:first-symbol-pairing} makes \(\overline D_b\) an
isomorphism.  The pullback by the evaluation-spectrum diffeomorphism, the
round musical isomorphism, and
\(\operatorname{ev}_a^{-1}\) are isomorphisms and are equivariant under the
transported \(\Spin(4)\) action.  Therefore every \(\Theta_a\) is an
isomorphism and
\[
 \Theta_{ga}\bigl(g\cdot[\widetilde f]\bigr)
 =\operatorname{Ad}_g\Theta_a([\widetilde f]).
\]

Differentiating
\(U_1(g\exp(tZ)g^{-1})
=U_1(g)U_1(\exp(tZ))U_1(g)^*\) at \(t=0\) gives
\begin{equation}
 \label{p2c:eq:Khat-adjoint-covariance}
 \widehat K(\operatorname{Ad}_gZ)
 =U_1(g)\widehat K(Z)U_1(g)^*.
\end{equation}
The retained module contains both
\(V_{\frac12,0}\) and \(V_{0,\frac12}\), so the derived representation is
faithful on
\(\mathfrak{spin}(4)=\mathfrak{su}(2)_L\oplus\mathfrak{su}(2)_R\).
Consequently \(\widehat K|_{\mathfrak p_a}\) is injective and hence an
isomorphism onto the defined fibre \(\mathscr K_a\).  Composing the two
covariance identities proves
\eqref{p2c:eq:operator-germ-covariance} and the gluing statement.  Finally,
\eqref{p2c:eq:Khat-basepoint-normalization} recovers the already derived
operators \(\widehat P_i\), so the construction consumes neither a new
normalization nor an additional occurrence witness.
\end{proof}

\begin{corollary}[Canonical compact-transvection/cotangent-germ bridge]
\label{p2c:cor:compact-transvection-cotangent-germ}
For every $a\in A$ and $b=\iota_1(a)\in B_A^{(1)}$, the already derived
first-symbol quotient and moving operator-germ map compose to the canonical
isomorphism
\begin{equation}
 \label{p2c:eq:compact-transvection-cotangent-germ}
 \boxed{
 \mathfrak B_a
 :=\overline D_b\circ(\mathcal Q_a^{xp})^{-1}:
 \mathscr K_a\overset\sim\longrightarrow T_b^*B_A^{(1)}.}
\end{equation}
Equivalently, using Definition~\ref{p2c:def:moving-transvection-operator-bundle}
and Eq.~\eqref{p2c:eq:theta-xp-local},
\begin{equation}
 \label{p2c:eq:compact-transvection-cotangent-expanded}
 \boxed{
 \mathfrak B_a
 =\bigl((d\iota_{1,a})^*\bigr)^{-1}
 \circ\flat_{h,a}\circ\operatorname{ev}_a
 \circ\bigl(\widehat K|_{\mathfrak p_a}\bigr)^{-1}.}
\end{equation}
Thus, for any transported normalized transvection basis
$\{\mathsf p_i(a)\}$, if
$\widehat P_i(a):=\widehat K(\mathsf p_i(a))$ and
$e_b^i:=\mathfrak B_a(\widehat P_i(a))$, then
\[
 \mathfrak B_a\!\left(\sum_i c^i\widehat P_i(a)\right)
 =\sum_i c^i e_b^i.
\]
The statement is exact at the first-order generator-germ level.  It neither
identifies the global compact momentum carrier with the reciprocal cotangent
bundle nor supplies an injective all-resolution map
$\mathcal P_x(m)\to\End(\mathcal H_1)$.  Dimensionful reciprocal-position
calibration remains a Paper-III interface.
\end{corollary}

\begin{proof}
Theorem~\ref{p2c:thm:operator-germ-xp-pairing} makes
$\mathcal Q_a^{xp}$ an isomorphism onto $\mathscr K_a$, while
Theorem~\ref{p2c:thm:first-symbol-pairing} makes $\overline D_b$ an
isomorphism onto $T_b^*B_A^{(1)}$.  Their composition therefore proves
Eq.~\eqref{p2c:eq:compact-transvection-cotangent-germ}.  Inverting the factors
in the definition of $\Theta_a$ gives
Eq.~\eqref{p2c:eq:compact-transvection-cotangent-expanded}.  Covariance is
inherited from the two constituent bundle maps.  The final firewalls are
exactly Propositions~\ref{p2c:prop:no-fixed-global-xp-linear-map} and
\ref{p2c:prop:no-faithful-allband-xp-map}.
\end{proof}


\begin{remark}[The germ bridge does not globalize to a canonical symplectic structure]
\label{p2c:rem:global-symplectic-obstruction}
The local first-symbol comparison retains the global obstructions proved
in Paper~I-C.  Its native reciprocal two-form has
$d\Omega\ne0$ (Theorem~\ref{p1b:thm:transgression}), and the compact
$M_8$ admits no symplectic form (Proposition~\ref{p1b:prop:lcs}).
Thus the exact generator-germ isomorphism does not establish a global
symplectic or canonical cotangent identification.  Any closed canonical
form on a different target retains its actual target and comparison map.
\end{remark}

\begin{corollary}[The local first-symbol bridge generates the full native kinematic algebra]
\label{p2c:cor:v185-first-symbol-generates-akin}
For every base point $a$ in the reciprocal-position orbit,
\[
 \boxed{
 \operatorname{alg}_{\mathbb C}^{*}
 \bigl(\mathbf1,\operatorname{Ran}\mathcal Q_a^{xp}\bigr)
 =\mathcal A_{\rm kin}.}
\]
Thus the canonical local first-symbol/operator-germ bridge does not merely
supply three finite transvection operators: its generated finite test algebra
is exactly the already derived $43$-dimensional native kinematic algebra.  The
statement does not enlarge the algebra to $\operatorname{End}(\mathcal H_1)$
and does not identify the local position-symbol space with the full momentum
Hilbert carrier.
\end{corollary}
\begin{proof}
At one reference base point, the image of $\mathcal Q_a^{xp}$ is the
transvection span generated by the three $\widehat P_i$.  Proposition~\ref{p2:prop:momentum-commutation-relations}
gives
\[
 \widehat J_k
 =\frac{\Rmom^2}{i\hbar}\,[\widehat P_i,\widehat P_j]
 \qquad (\epsilon_{ijk}=1),
\]
so the unital $*$-algebra generated by the image also contains all three
rotations.  Definition~\ref{p2c:def:generated-kinematic-test-algebra} then
gives $\mathcal A_{\rm kin}$, while the reverse inclusion is immediate because
the image operators already belong to that algebra.  Equation~\ref{p2c:eq:operator-germ-covariance}
transports the equality to every base point by unitary conjugation.
\end{proof}

\begin{proposition}[No fixed global linear position-to-momentum intertwiner]
\label{p2c:prop:no-fixed-global-xp-linear-map}
There is no nonzero fixed $\Spin(4)$-equivariant linear map from the global
first-harmonic four-vector to the compact Killing algebra, because
\[
 \operatorname{Hom}_{\Spin(4)}\!\left((\tfrac12,\tfrac12),(1,0)\oplus(0,1)\right)=0.
\]
The moving bundle map above is therefore the correct covariant object.
\end{proposition}

\begin{proposition}[No faithful all-resolution map into the frozen momentum carrier]
\label{p2c:prop:no-faithful-allband-xp-map}
The reciprocal scalar band has
\[
 \dim_\CC\mathcal P_x(m)=\frac{(m+1)(m+2)(2m+3)}6,
\]
whereas $\dim_\CC\End(\mathcal H_1)=225$.  Already
$\dim\mathcal P_x(8)=285>225$.  Therefore no injective linear map
$\mathcal P_x(m)\to\End(\mathcal H_1)$ can exist for every $m$ while the
physical momentum carrier remains frozen at $m_p=1$.
If the map is also $\Spin(4)$-equivariant, the obstruction begins at degree
three. Indeed, the scalar harmonic multiplicities in $\End(\mathcal H_1)$
are $6,8,3$ in degrees $0,1,2$ and zero in every degree $k\ge3$.
Thus every such map annihilates $E_3^x\oplus\cdots\oplus E_m^x$ and has rank
at most $14$. The dimension-only degree-eight obstruction remains valid
without covariance.
\end{proposition}
\begin{proof}
The dimension bound is immediate. For the covariant refinement, write
$\mathcal H_1=V_{0,0}\oplus V_{1/2,1/2}\oplus V_{1,0}\oplus V_{0,1}
\oplus V_{1/2,0}\oplus V_{0,1/2}$.
Each summand has $j_L+j_R\le1$, so every type in
$\End(\mathcal H_1)$ has total highest weight at most two. A scalar harmonic
$E_k^x=V_{k/2,k/2}$ therefore cannot occur for $k\ge3$.
Clebsch--Gordan decomposition gives multiplicities $6,8,3$ for $k=0,1,2$.
Each source type occurs once, so the image dimension is at most $1+4+9=14$.
Corollary~\ref{p2c:cor:optimal-povm-harmonic-kernel} supplies a positive
covariant map attaining that rank.
\end{proof}

\begin{theorem}[Canonical native fourteen-dimensional quadratic operator-symbol copy]
\label{p2c:thm:v268-native-jordan-quadratic-copy}
On the already constituted physical carrier
$\mathcal H_p\cong\mathcal H_1$, let $R=\Rmom>0$, let $X_a^D$
($a=1,\ldots,4$) be the intrinsic $\Spin(4)$-vector coordinate tuple
of Theorem~\ref{p2c:thm:intrinsic-tangential-coordinate-class}, and put
\[
 Q_D=\sum_{c=1}^4(X_c^D)^2,\qquad
 T_{ab}^D=\frac12(X_a^DX_b^D+X_b^DX_a^D)
                     -\frac{\delta_{ab}}4 Q_D.
\]
Then the nine trace-free symmetric tensors $T_{ab}^D=T_{ba}^D$,
$\sum_aT_{aa}^D=0$, span one \emph{distinguished} faithful copy of
$E_2\cong V_{1,1}$ in $\End(\mathcal H_1)$.  In particular,
\[
 \boxed{\begin{gathered}
 \mathcal W_2^D:=\mathbb CI_{15}\oplus
      \operatorname{span}_{\mathbb C}\{X_a^D\}_{a=1}^4\\
      {}\oplus\operatorname{span}_{\mathbb C}\{T_{ab}^D\}_{a,b=1}^4\\
 \cong_{\Spin(4)} E_0\oplus E_1\oplus E_2=\mathcal P_2,\\
 \dim_{\mathbb C}\mathcal W_2^D=14.
 \end{gathered}}
\]
The identification is canonical \emph{relative to the native tuple}:
$1\mapsto I_{15}$, $q_a\mapsto X_a^D/R$, and
$q_aq_b-\delta_{ab}/4\mapsto36T_{ab}^D/R^2$ give a well-defined,
$\Spin(4)$-equivariant vector-space isomorphism from restricted unit-sphere
polynomials $\mathcal P_2$ to $\mathcal W_2^D$.
It is neither a multiplicative representation of the pointwise
scalar algebra nor an occurrence of a finite point-support role.
\end{theorem}
\begin{proof}
On the canonical geometric blocks
$\mathbb Ce_0\oplus V\oplus\Lambda^2V$,
$V=\operatorname{span}\{u_1,\ldots,u_4\}$,
the native operator has the fixed, nonzero edges
\[
 X_a^De_0=A u_a,\qquad
 X_a^Du_b=A\delta_{ab}e_0+B(e_a\wedge e_b),
 \qquad A=\frac R2,\quad B=\frac{\sqrt2 R}{3}.
\]
The adjoint edge is
$X_a^D(e_b\wedge e_c)=B(\delta_{ab}u_c-\delta_{ac}u_b)$.
Thus, with $P_V$ the canonical equivariant projection onto $V$,
\[
 \begin{split}
 P_V\frac12\{X_a^D,X_b^D\}P_V
 &= B^2\delta_{ab}I_V+
   (A^2-B^2)\frac{|u_a\rangle\langle u_b|+
                          |u_b\rangle\langle u_a|}{2},\\
 P_VQ_DP_V&=(A^2+3B^2)I_V=\frac{11R^2}{12}I_V,\\
 P_VT_{ab}^DP_V
 &=\frac{R^2}{36}\left(
    \frac{|u_a\rangle\langle u_b|+|u_b\rangle\langle u_a|}{2}
                   -\frac{\delta_{ab}}4I_V\right).
 \end{split}
\]
The last coefficient is nonzero since
$A^2-B^2=R^2(1/4-2/9)=R^2/36$.
Its displayed images span the entire nine-dimensional real symmetric
trace-free endomorphism space on $V$, and remain linearly independent
after complexification.  Tensor covariance makes their full images a
single copy of the irreducible $\Spin(4)$-type
$\operatorname{Sym}^2_0 V\cong(1,1)=E_2$.
The identity and the original coordinate tuple supply inequivalent
nonzero $E_0$ and $E_1$ images.  Their direct sum therefore has dimension
$1+4+9=14$.  The specified map respects the trace relation
$\sum_a(q_a^2-1/4)=\sum_aT_{aa}^D=0$ and is injective by the
three-type decomposition.

This does not identify the full diagonal-isotypic subspace of
$\End(\mathcal H_1)$ with $\mathcal P_2$: the multiplicities there
are $6,8,3$ for $E_0,E_1,E_2$, respectively.  It constructs one
native copy singled out by $X^D$, not an arbitrary section of those
multiplicity spaces.
\end{proof}

The surviving strong finite-band comparison problem is therefore a typed
refinement/quantization problem on the auxiliary reciprocal spectral tower,
not a new physical momentum-level trajectory.

\section{Intrinsic localization, effects and measurements}
\label{paper2c:sec:organized-localization}

\subsection{Canonical scalar and unsharp localization interfaces}
\label{p2c:sec:canonical-localization-interfaces}

Paper~I-A, Theorem~\ref{p1:thm:scalar-band-evaluation-povm}, supplies on every
finite scalar band \(\mathcal P_m^{\rm sc}\) the exact covariant evaluation
POVM
\[
 E_m(dx)=|K_{m,x}\rangle\langle K_{m,x}|\,d\mu_H(x),
\]
with identity resolution and scalar Toeplitz first moments.  Thus scalar-band
mathematical localization is derived.  It does not resolve the full retained
form/spin carrier and does not actualize an apparatus.

\begin{theorem}[Exact linear full-carrier POVM range for the intrinsic symbol]
\label{p2c:thm:unsharp-intrinsic-coordinate-povm}
Let \(X^D=(X_1^D,\ldots,X_4^D)\) be the intrinsic tangential coordinate tuple
on \(\mathcal H_1\), and let
\[
 d\mu_H:=\frac{d\operatorname{vol}}
 {\operatorname{Vol}(S^3(\Rmom))}.
\]
For
\[
 0\le t\le\frac12,
\]
define
\[
 F_t(x):=I+\frac{4t}{\Rmom^2}\sum_{a=1}^4x_aX_a^D,
 \qquad
 E_t(dx):=F_t(x)\,d\mu_H(x).
\]
Then \(E_t\) is a normalized \(\Spin(4)\)-covariant POVM on the full retained
carrier and
\[
 \boxed{\int E_t(dx)=I,\qquad
 \int x_aE_t(dx)=tX_a^D.}
\]
Moreover, the interval is exact for this linear density family:
\(F_t(x)\ge0\) for every \(x\in S^3(\Rmom)\) if and only if
\(0\le t\le\frac12\).  If signed \(t\) is allowed, the exact interval is
\(|t|\le\frac12\).
\end{theorem}

\begin{proof}
By \(\Spin(4)\)-covariance it is enough to take
\(x=\Rmom e_4\).  The explicit intrinsic edges give the orthogonal block
decomposition of \(X_4^D\):
\[
 \begin{aligned}
 e_0\longleftrightarrow u_4
 &\quad\text{with coefficient }\frac{\Rmom}{2},\\
 u_i\longleftrightarrow w_{4i}\quad(i=1,2,3)
 &\quad\text{with coefficient }\frac{\sqrt2\Rmom}{3}.
 \end{aligned}
\]
The three remaining geometric two-form directions are annihilated by
\(X_4^D\), while on the protected block
\(X_4^D=(\sqrt3\Rmom/4)C_4\).  Hence, with multiplicities displayed,
\[
 \operatorname{spec}(X_4^D)
 =\left\{
 \pm\frac{\Rmom}{2},\
 \pm\frac{\sqrt2\Rmom}{3}\ (\times3),\
 0\ (\times3),\
 \pm\frac{\sqrt3\Rmom}{4}\ (\times2)
 \right\}.
\]
Therefore
\[
 \left\|\sum_ax_aX_a^D\right\|
 =\Rmom\|X_4^D\|
 =\frac{\Rmom^2}{2}
\]
for every point on the sphere, and the spectrum is symmetric.  Thus the least
eigenvalue of \(F_t(x)\) is \(1-2|t|\).  This proves both sufficiency and
necessity of \(|t|\le\frac12\), including the stated nonnegative branch.

Normalized Haar symmetry gives
\[
 \int x_a\,d\mu_H=0,
 \qquad
 \int x_ax_b\,d\mu_H=\frac{\Rmom^2}{4}\delta_{ab},
\]
which proves normalization and the first-moment identity.  The contraction
\(\sum_ax_aX_a^D\) is equivariant, so the operator-valued measure is
covariant.
\end{proof}

\begin{theorem}[Optimal intrinsic classical first-moment radius]
\label{p2c:thm:optimal-intrinsic-povm-radius}
For the intrinsic coordinate quadruple $X^D$ on $\mathcal H_1$, the least
classical outcome radius of a normalized POVM with exact first moments
$X_a^D$ is $\sqrt3\,\Rmom$. A covariant POVM on that boundary sphere attains
the radius. Equivalently, a normalized POVM on the radius-$\Rmom$ ball with
moments $tX_a^D$ exists exactly for $|t|\le1/\sqrt3$.
A seven-outcome rule without a subgroup-support condition is proved in
Theorem~\ref{p2c:thm:v255-seven-outcome-attainment}.
An eight-outcome POVM on transported \(Q_8\) support attains this
radius, as does the original 24-outcome Hurwitz construction.
Among subgroup supports the attaining types are exactly the noncyclic
finite groups; \(Q_8\) has smallest cardinality, but there is no least
class under subgroup inclusion. The density and optimality proof are in
Paper~II-F, Section~\ref{p2c:app:optimal-classical-moment-radius};
Corollary~\ref{p2c:cor:v233-optimal-quaternion-povm} proves the
eight-outcome and subgroup statements.
The endpoint $1/2$ of the earlier theorem remains exact for its specified
linear density family.
\end{theorem}


\begin{theorem}[Sharp total coordinate variance on the full intrinsic carrier]
\label{p2c:thm:intrinsic-total-variance}
For every density operator $\sigma$ on the full fifteen-dimensional carrier,
the intrinsic tangential coordinates satisfy
\[
 \boxed{\sum_{a=1}^4\operatorname{Var}_\sigma(X_a^D)
       \ge\frac{1823}{4608}\Rmom^2.}
\]
The constant is sharp.  It is a full-carrier bound, rather than the
protected-block bound of
Corollary~\ref{p2c:cor:protected-coordinate-covariance}.
\end{theorem}
\begin{proof}
Put $Y_a=X_a^D/\Rmom$, $Q=\sum_aY_a^2$, and
$m_a=\operatorname{Tr}(\sigma Y_a)$.  Completion of a real square gives
\[
 \sum_a\operatorname{Var}_\sigma(Y_a)
 =\inf_{t\in\mathbb R^4}
   \operatorname{Tr}\bigl(\sigma[Q-2t\cdot Y+|t|^2I]\bigr).
\]
Taking the infimum over a product of the state space and $\mathbb R^4$,
and using round covariance, yields
\[
 \inf_\sigma\sum_a\operatorname{Var}_\sigma(Y_a)
 =\inf_{s\ge0}\lambda_{\min}(Q-2sY_4+s^2I).
\]
There is no minimax exchange.  The quadratic fingerprint and the exact
$Y_4$ blocks in the proof of
Theorem~\ref{p2c:thm:unsharp-intrinsic-coordinate-povm} give the complete joint table
\[
\begin{array}{c|c|c|c}
 \text{block}&\text{copies}&Q&Y_4\\\hline
 (e_0,u_4)&1&\operatorname{diag}(1,11/12)
    &\tfrac12\left(\begin{smallmatrix}0&1\\1&0\end{smallmatrix}\right)\\[3pt]
 (u_i,w_{4i})&3&\operatorname{diag}(11/12,4/9)
    &\tfrac{\sqrt2}{3}\left(\begin{smallmatrix}0&1\\1&0\end{smallmatrix}\right)\\[3pt]
 \text{uncoupled two-forms}&3&4/9&0\\
 \text{protected}&\text{dimension }4&(3/4)I_4&(\sqrt3/4)C_4
\end{array}
\]
in an edge-positive exterior gauge.  For a two-dimensional block with
diagonal entries $q_1,q_2$ and off-diagonal $Y_4$ coefficient $h$, its
lower eigenvalue is
\[
 s^2+a-\sqrt{b^2+4h^2s^2},\qquad
 a=(q_1+q_2)/2,\quad b=(q_1-q_2)/2.
\]
If $|b|\le2h^2$, minimizing over $s\ge0$ gives
$a-h^2-b^2/(4h^2)$.  Both coupled geometric blocks satisfy this condition.
The four branch minima, in table order, are
\[
 \frac{407}{576},\qquad\frac{1823}{4608},\qquad
 \frac49,\qquad\frac9{16}.
\]
Thus the second is the global minimum, including arbitrary mixed states
and arbitrary coherences.  It is attained, for $i\in\{1,2,3\}$, by
\[
 \psi=\frac{\sqrt{15}}8u_i+\frac78w_{4i}.
\]
The exterior action gives zero means in directions other than four and
$\langle Y_4\rangle=7\sqrt{30}/96$;
$\langle Q\rangle=1279/2304$ and
$|\langle Y\rangle|^2=245/1536$.  Their difference is $1823/4608$.
\end{proof}
Each individual coordinate still has sharp eigenstates.  This theorem
bounds their total operator variance; it does not select a preparation,
instrument, physical distance calibration or measured minimum length.

\begin{proposition}[Finite unit pairing and nonunique stronger reciprocal maps]
\label{p2c:prop:v233-unit-and-positive-maps}
The scalar position invariant is already \(\mathbb C1\) at every
finite complete resolution. Its reciprocal pullback sends \(1\) to
\(1\); into an already specified nonzero unital star algebra \(A\)
with unit-preserving action, the unique linear unital map is
\(\eta_A(z)=zI_A\).
Thus neither a large-resolution limit nor an arbitrary choice of
unit map is needed for these conclusions.

This does not determine a stronger observable-to-operator map.
There is an entire family
\[
 \Phi_c:C(S^3(1))\longrightarrow\operatorname{End}(\mathcal H_1),
 \qquad \tfrac13\le c\le1,
\]
of distinct unital completely positive \(\Spin(4)\)-covariant maps
with identical native first moments
\(\Phi_c(q_a)=X_a^D/(\sqrt3\Rmom)\).
Each has rank fourteen on every \(\mathcal P(m)\), \(m\ge2\), and
kernel \(\bigoplus_{k=3}^mE_k\) there.
They differ on quadratic functions. Positivity, covariance and these
first moments therefore do not select a unique extension.
\end{proposition}
\begin{proof}
The finite unit and its map are
Theorem~\ref{p1b:thm:v233-finite-scalar-unit} and
Proposition~\ref{p1b:prop:v233-unit-map}. The positive family,
that real slice's exact parameter interval, and its explicit quadratic
difference are proved in
Proposition~\ref{p2c:prop:v233-endpoint-map-family}.
\end{proof}

The canonical unit image is already in its target algebra. Neither
this map nor the positive family constitutes an additional singlet
sector, a physical measurement, or a mixed action contraction.
The operator-germ pairing remains the separately derived local
relation of Theorem~\ref{p2c:thm:operator-germ-xp-pairing}.


\begin{proposition}[Exact classical radius-\(\Rmom\) moment matching is impossible]
\label{p2c:prop:exact-intrinsic-classical-moment-no-go}
No normalized POVM with classical outcomes supported in the closed ball
\(B_{\Rmom}^4\), and hence none supported on \(S^3(\Rmom)\), has exact first
moments \(X_a^D\).  Equivalently, the former exact interface
\[
 \int x_aE(dx)=X_a^D
\]
is refuted on the current classical outcome domain.  The full Naimark
saturation proof is in
Paper~II-F, Section~\ref{p2c:app:exact-moment-saturation-obstruction}.
\end{proposition}

\begin{proof}
Paper~I-A gives
\(\sum_a(X_a^D)^2e_0=\Rmom^2e_0\), so the normalized constant saturates the
ball moment inequality.  Saturation for a classical POVM forces all first
moments to commute on that vector, whereas the explicit intrinsic edges give
\[
 [X_a^D,X_b^D]e_0
 =\frac{\sqrt2\Rmom^2}{3}w_{ab}\ne0
 \qquad(a\ne b).
\]
The two conclusions are incompatible.
\end{proof}

\begin{remark}[Mathematical localization versus one physical occurrence package]
The scalar evaluation, linear unsharp and radius-optimal full-carrier
families are derived mathematical POVMs.  Their existence, normalization, covariance, and displayed
first moments are closed.  A chosen apparatus/channel, prepared state,
sharpness or calibration \(t\), physical ownership of outcome/effect labels
as positions, performed instrument, and registered outcome are components of
one claim-relative G0 preparation/measurement occurrence package.  The tuple
is bookkeeping: its components are not generated from one another, and a
claim needs only the fields it actually consumes.  They are not separate
structural roots. For exact $X^D$ moments with classical outcomes, the
least radius is now fixed mathematically to $\sqrt3\Rmom$, and continuous,
seven-outcome, eight-outcome and 24-outcome constructions attain it. Physical use of those outcome
labels still needs the stated occurrence package; no instrument is selected.
\end{remark}


\subsubsection{Fixed momentum clock algebra and the degree-two position comparison}
\label{p2c:sec:v180-clock-comparison-band}

\begin{proposition}[Generated chiral operation algebras]
\label{p2c:prop:generated-chiral-operation-algebras}
Fix one handed chiral $SU(2)$ action containing the primitive Hopf circle.
The already derived full level-one momentum carrier and the complete position
bands restrict as
\[
 \mathcal H_p|_{SU(2)}=6V_0\oplus3V_{1/2}\oplus V_1,
 \qquad
 \mathcal P_x(m_x)|_{SU(2)}=
 \bigoplus_{k=0}^{m_x}(k+1)V_{k/2}.
\]
Consequently the algebras generated by the three corresponding chiral
generators are
\[
 \mathcal A_{{\rm clk},p}\cong\mathbb C\oplus M_2\oplus M_3,
 \qquad
 \mathcal A_{{\rm clk},x}(m_x)\cong
 \bigoplus_{k=0}^{m_x}M_{k+1}.
\]
For every $m_x\ge2$, block projection gives a unital $*$-epimorphism
\[
 \Phi_{m_x}:\mathcal A_{{\rm clk},x}(m_x)
 \twoheadrightarrow\mathcal A_{{\rm clk},p},\qquad
 \Phi_{m_x}(Q_i^x)=Q_i^p .
\]
At $m_x=2$ this is an algebra isomorphism.  It is not an isomorphism of the
represented Hilbert spaces or of their inherited trace states.
\end{proposition}
\begin{proof}
Apply Paper~1B, Theorem~\ref{p1b:prop:degree-two}, to the fixed
handed action on the full $m_p=1$ carrier and the reciprocal scalar band
constructed here. Its Casimir-block comparison gives the displayed
generator-preserving epimorphism. The inherited carrier multiplicities,
and therefore the represented traces, remain those of the two local
carriers.
\end{proof}

\begin{theorem}[Unique finite scalar comparison band]
\label{p2c:thm:position-to-momentum-chiral-quotient}
Define $m_{\rm cmp}$ to be the degree of a \emph{complete position scalar
band used only for the following comparison}.  Require simultaneously:
\begin{enumerate}
 \item a generator-preserving unital $*$-epimorphism from its chiral
 operation algebra onto that of the fixed physical momentum carrier
 $\mathcal H_p=\mathcal H_1$ with $m_p=1$;
 \item exact uniform $2I$ evaluation of $\int f^4d\sigma$ for every real
 $f\in\mathcal P_x(m_{\rm cmp})$.
\end{enumerate}
Then
\[
 \boxed{m_p=1\quad\hbox{and}\quad m_{\rm cmp}=2.}
\]
The first requirement forces $m_{\rm cmp}\ge2$, because the momentum clock
algebra contains the spin-one Casimir-$8$ block, absent from
$\mathcal P_x(1)$.  The second forces $m_{\rm cmp}\le2$: Paper~1B,
Theorems~\ref{p1b:thm:design} and~\ref{p1b:prop:degree-two}, give exact
$2I$ integration through degree eleven and its sharp comparison frontier.
For the
degree-three field $f=x_0^3$,
\[
 \frac1{120}\sum_{x\in2I}x_0^{12}=\frac{133}{4096}
 \ne\frac{33}{1024}
 =\int_{S^3}x_0^{12}\,d\sigma .
\]
Hence the two independent requirements meet uniquely at degree two.
\end{theorem}

\begin{remark}[The level-two comparison is not a second physical truncation]
\label{p2c:rem:v180-mp1-mcmp2-scope}
The primitive compact momentum truncation remains $m_p=1$.  The primitive
position seed may likewise begin with $m_x=1$, while the constituted position
frontier can later grow.  The symbol $m_{\rm cmp}=2$ names only the smallest
complete position band that both reproduces the full fixed-momentum chiral
clock algebra and is evaluated quartically exactly by the fixed $120$-point
$2I$ geometry.  It neither caps the physical $m_x$ frontier at two nor
equates the $14$-dimensional space $\mathcal P_x(2)$ with the
$15$-dimensional momentum carrier.
\end{remark}

\begin{corollary}[Quadratic specialization of the comparison band]
\label{p2c:cor:v191-quadratic-comparison-specialization}
The first nonlinear scalar descendant of the primitive position harmonic is
already
\[
 \operatorname{Sym}^2E_1^x\cong E_0^x\oplus E_2^x.
\]
Consequently the unique comparison band
$\mathcal P_x(2)=E_0^x\oplus E_1^x\oplus E_2^x$ is the first complete
position band containing both the primitive linear sector and its full
quadratic scalar descendant.  This is compatible with, but does not replace,
the two-sided uniqueness proof of
Theorem~\ref{p2c:thm:position-to-momentum-chiral-quotient}.
\end{corollary}

\begin{remark}[No mixed $x$--$p$ pairing from quadratic closure]
The preceding corollary concerns pointwise products inside the position
harmonic algebra.  It does not construct a mixed position--momentum
contraction, identify a position harmonic with a compact momentum operator,
or discharge an unspecified stronger $\mathrm{PAIR}_{xp}$ interface.
The unit-only specialization is derived by
Proposition~\ref{p2c:prop:v233-unit-and-positive-maps}; its stronger
positive extensions remain the explicitly nonunique family proved there.
\end{remark}


\section{Observable tests, effect frames and canonical readouts}
\label{paper2c:sec:organized-observables}

\subsection{Claim-relative effect frames and extension-independent readouts}
\label{p2c:sec:v2210-lit-effect-frames}
The localization families above already establish mathematical positivity,
covariance and the sharp classical outcome radius.  The question addressed
here is which source alternatives a supplied effect family separates.
This is a finite calculation on an actual claim space; it neither selects
an apparatus nor requires complete tomography for every G0 comparison.

\begin{theorem}[Exact visible-subspace frame certificate]
\label{p2c:thm:v2210-lit-effect-frame}
Let $E_1,\ldots,E_N$ be the effects of a supplied finite POVM on a finite
carrier, and let $V$ be a specified real linear subspace of Hermitian
source differences.  Give $V$ the Hilbert--Schmidt inner product and
write $\Pi_V$ for its orthogonal projection.  Set
\[
 \mathcal R_V(X)_j=\operatorname{Tr}(E_jX),\qquad
 \mathcal F_V(X)=\sum_{j=1}^N
       \operatorname{Tr}(E_jX)\,\Pi_VE_j.
\]
Then $\mathcal F_V=\mathcal R_V^*\mathcal R_V\ge0$ and
\[
 \ker\mathcal F_V=
 \{X\in V:\operatorname{Tr}(E_jX)=0\hbox{ for all }j\}.
\]
Thus the performed probabilities distinguish all source differences in
$V$ exactly when $\mathcal F_V$ is positive definite on $V$.
If its least eigenvalue is $\alpha_V>0$, then
\[
 \|\mathcal R_V(X)\|_2\ge\sqrt{\alpha_V}\,\|X\|_{\rm HS},
 \qquad X\in V,
\]
and a probability-difference error $\delta p$ changes the least-squares
reconstruction on $V$ by at most
$\|\delta p\|_2/\sqrt{\alpha_V}$.
\end{theorem}
\begin{proof}
For $X\in V$,
\[
 \langle X,\mathcal F_VX\rangle_{\rm HS}
 =\sum_j|\operatorname{Tr}(E_jX)|^2.
\]
This proves the factorization, positivity and kernel equality.  The
spectral theorem gives the lower bound.  When the kernel vanishes,
the smallest singular value of $\mathcal R_V$ is $\sqrt{\alpha_V}$;
its pseudoinverse has norm $1/\sqrt{\alpha_V}$, proving stability.
\end{proof}
This is the claim-subspace form of the covariant-effect frame method
\cite{LitGEODArPerSac2004Effects}.  The claim space and actual effects must
be specified independently of the resulting probabilities.  For a
continuous covariant POVM with an operator-valued density $E(x)$ and an
actual finite $L^2$ Gram integral, replace the sum by that integral;
the same proof applies.

\begin{corollary}[A licensed null frame direction gives an explicit ambiguity]
\label{p2c:cor:v2210-lit-state-ambiguity}
Suppose the claim permits every density operator on a $d$-dimensional
carrier.  If a nonzero traceless Hermitian $X$ lies in the kernel of the
effect frame, then, for
$0<\epsilon\le1/(d\|X\|)$,
\[
 \rho_\pm=I/d\pm\epsilon X
\]
are distinct admissible density operators with the same POVM
probabilities.  Restricted physical source families require only the
frame criterion on their actual difference space $V$; the displayed
full-state ambiguity is not an assertion that both states are prepared.
\end{corollary}
\begin{proof}
The operator bound gives $\rho_\pm\ge0$ and tracelessness gives unit
trace.  Each probability difference is
$2\epsilon\operatorname{Tr}(E_jX)=0$.  Distinctness follows from $X\ne0$.
\end{proof}

\begin{theorem}[Native central-sector refinement with exact moments]
\label{p2c:thm:v2211-central-sector-refinement}
On the derived carrier, let $P_g,P_s$ project onto
$\mathcal H_{\rm geo}$ and $\mathcal H_{\rm sp}$ of
Corollary~\ref{p2c:cor:present-coordinate-envelope-condition}.
Both projectors belong to $\mathcal A_{\rm kin}$: they are sums of its
generated irreducible-sector projectors.  They commute with every present
coordinate, and $P_g+P_s=I$.
Let $E_1,\ldots,E_N$ be a normalized POVM reducing these two sectors,
with classical outcomes $x_j\in\mathbb R^4$ and first moments
$X_a=\sum_jx_{j,a}E_j$.  Then
\[
 E_{j,t}:=P_tE_jP_t,\qquad t\in\{g,s\},
\]
with outcomes $(x_j,t)$ is a normalized POVM and
\[
 \boxed{\sum_{j,t}E_{j,t}=I,\qquad
        \sum_{j,t}x_{j,a}E_{j,t}=X_a,\qquad
        \sum_jE_{j,t}=P_t.}
\]
Consequently, for every density operator $\rho$,
\[
 \boxed{\sum_j\operatorname{Tr}(\rho E_{j,t})
       =\operatorname{Tr}(\rho P_t).}
\]
The tag therefore separates any two admitted source states with different
geometric/protected populations.  If the original POVM is covariant under
the constituted round action, the refinement is covariant with the tag fixed.
For finite supports, this retains the original transported-family covariance
rather than asserting invariance of a fixed finite support.

For every position-only function $f$ on the supplied outcomes one also has
\[
 \boxed{\sum_{j,t}f(x_j)E_{j,t}=\sum_jf(x_j)E_j.}
\]
Thus tagging adds a population readout without changing the original
position-symbol map, its quadratic aliasing, or its multiplication defects.
It does not assert complete source tomography or an apparatus occurrence.
\end{theorem}
\begin{proof}
The two projectors are generated sums of the polynomial projectors in
Theorem~\ref{p2c:thm:exact-level-one-kinematic-algebra}; no later grading
theorem is a premise.  The allowed coordinate graph reduces their ranges.
Each $P_tE_jP_t$ is positive.  Summing first over $j$ gives $P_t$, so
summing over the tags gives the identity.  Because $E_j$ reduces both
sectors, $\sum_tP_tE_jP_t=E_j$.  Multiplication by either $x_{j,a}$ or
$f(x_j)$ and summation proves the two moment/readout identities.
Taking the expectation of $\sum_jE_{j,t}=P_t$ proves the population
formula without any invariance condition on $\rho$.
The constituted round action commutes with each $P_t$, so conjugating
the refined effects transports the outcome and leaves the tag fixed.
These identities construct a mathematical measurement family; they neither
prepare a state nor assert that this family is performed.
\end{proof}

\begin{corollary}[Intrinsic endpoint populations and their exact frame]
\label{p2c:cor:v2211-tagged-intrinsic-populations}
For the endpoint density $K_c(q)$ of
Proposition~\ref{p2c:prop:v233-endpoint-map-family}, $1/3\le c\le1$,
use either the $N=8$ quaternion rule or the $N=24$ Hurwitz rule and set
\[
 E_j=\frac1N K_c(q_j),\qquad x_j=\sqrt3\Rmom q_j.
\]
The preceding theorem supplies $2N$ positive effects with the same exact
native moments $X_a^D$.  Put
\[
 \rho_p=\frac p{11}P_g+\frac{1-p}{4}P_s,\qquad0\le p\le1.
\]
The untagged rule gives $\operatorname{Tr}(\rho_pE_j)=1/N$ for every
$p$, while the refined rule gives exactly
\[
 \boxed{\operatorname{Tr}(\rho_pE_{j,g})=p/N,\qquad
        \operatorname{Tr}(\rho_pE_{j,s})=(1-p)/N.}
\]
On the one-dimensional claim space
\[
 V_{\rm pop}:=\mathbb R T,\qquad
 T:=P_g/11-P_s/4,
\]
the untagged effect frame is zero and the tagged effect frame of
Theorem~\ref{p2c:thm:v2210-lit-effect-frame} is
\[
 \boxed{\mathcal F_{V_{\rm pop}}
       =\frac{88}{15N}I_{V_{\rm pop}}.}
\]
Thus its exact positive frame constant is $11/15$ for the quaternion
rule and $11/45$ for the Hurwitz rule, in the stated unweighted Euclidean
outcome and Hilbert--Schmidt source norms.

There is also an exact classical comparison for the native center
$\mathcal Z=\mathbb CP_g\oplus\mathbb CP_s$:
\[
 \Phi_{\mathcal Z}(a_gP_g+a_sP_s)_{j,t}=a_t,\qquad
 \Psi_{\mathcal Z}(f)
       =\sum_{t\in\{g,s\}}\left(\frac1N\sum_j f_{j,t}\right)P_t.
\]
These maps are unital completely positive, $\Phi_{\mathcal Z}$ is
multiplicative, and
\[
 \boxed{\Psi_{\mathcal Z}\Phi_{\mathcal Z}=I_{\mathcal Z}.}
\]
The refined measurement map
$\mathsf M(f)=\sum_{j,t}f_{j,t}E_{j,t}$ obeys
$\mathsf M\Phi_{\mathcal Z}(A)=A$ for every $A\in\mathcal Z$.
Every consumed central word therefore has exact reconstruction and zero
multiplication defect in this comparison; no universal all-band
quantization is required.

This is population separation on the displayed claim space, not full
tomography: $2N\le48$ outcomes cannot separate the $136$-dimensional
traceless Hermitian source-difference space of $M_{11}\oplus M_4$.
An actual sector-tagged measurement still consumes its G0 occurrence
package, and a different source claim retains its own effect-frame test.
\end{corollary}
\begin{proof}
The density is block diagonal, so the refinement theorem applies.  Its
geometric trace is $11$: the scalar, vector and two-form traces are
$1,4,6$, since $\operatorname{Tr}(qq^*)=1$ and
$\operatorname{Tr}(2W_qW_q^*-I_6)=2\cdot3-6=0$.
The protected trace is $4$, since $\operatorname{Tr}C(q)=0$.
These identities give the untagged and tagged probabilities for every
$c$ and every node.  Exact normalization and moments are the finite
quadrature conclusions of
Corollaries~\ref{p2c:cor:v233-optimal-quaternion-povm}
and~\ref{p2c:cor:optimal-hurwitz-povm}.

Now $\|T\|_{\rm HS}^2=1/11+1/4=15/44$ and
\[
 \operatorname{Tr}(E_{j,g}T)=1/N,\qquad
 \operatorname{Tr}(E_{j,s}T)=-1/N.
\]
For the unit vector $T/\|T\|_{\rm HS}$, the sum of the $2N$ squared
frame coefficients is $(2/N)/(15/44)=88/(15N)$, proving the frame
formula.  The untagged coefficients vanish identically.

The two displayed classical maps preserve the unit and positivity;
complete positivity follows componentwise from scalar evaluation and
averaging on the two finite commutative algebras.  Substitution proves
the retraction, while pointwise products show multiplicativity of
$\Phi_{\mathcal Z}$.  The identity $\sum_jE_{j,t}=P_t$ proves the
refined measurement-map equality and its central-word version.  Finally,
normalized probabilities have at most $2N-1\le47$ independent real
differences, strictly below $11^2+4^2-1=136$.  This proves the stated
tomography boundary without asserting that every such source is prepared.
\end{proof}

\begin{corollary}[Maximal round-invariant sharp tag compatible with native moments]
\label{p2c:cor:maximal-native-sharp-tag}
Suppose a POVM with first moments $X_a^D$ jointly refines a sharp sector
tag, and its tag projections are round-invariant.  Those projections lie
in the center $\mathbb CP_g\oplus\mathbb CP_s$ of the native algebra.
The geometric/protected tag is therefore maximally fine under these
conditions; finer Casimir-sector tags cannot preserve all four moments.
\end{corollary}
\begin{proof}
Each positive joint effect under a tag projection is supported in that
projection, so all first moments commute with every tag projection.
Round invariance and multiplicity-free sector decomposition imply that
the tag projections commute with $\mathcal A_{\rm kin}$ as well.
They therefore commute with
$\operatorname{alg}_{\mathbb C}^*(\mathcal A_{\rm kin},X^D)=M_{11}\oplus M_4$
and belong to its center.  The preceding construction attains its two
minimal central projections.
\end{proof}

\begin{proposition}[First-moment agreement leaves a precise readout quotient]
\label{p2c:prop:v2210-lit-extension-quotient}
Let $\Phi_c$ be any supplied family of linear positive covariant maps
with the same native first moments, including the family of
Proposition~\ref{p2c:prop:v233-unit-and-positive-maps}.  On a specified
finite symbol space $S$, define
\[
 S_{\rm common}:=
 \bigcap_{c,c'}\ker(\Phi_c-\Phi_{c'}).
\]
Every operator readout of a symbol in $S_{\rm common}$ is
extension-independent.  For a specified family of physical source/effect
functionals $\ell_s$, the larger visible common space is exactly
\[
 S_{\rm vis}:=
 \bigcap_{s,c,c'}\ker\bigl(\ell_s\circ(\Phi_c-\Phi_{c'})\bigr).
\]
If the consumed functionals separate the occurring operator differences,
then $S_{\rm vis}=S_{\rm common}$.  Otherwise visible independence is
only the claim-relative quotient represented by $S_{\rm vis}$.
\end{proposition}
\begin{proof}
The first two statements are the necessary and sufficient zero-defect
conditions for all members to give the same operator or scalar.  Separation
forces any difference annihilated by every $\ell_s$ to be zero, yielding
the final equality.  The preceding effect-frame theorem gives a finite
rank test when the $\ell_s$ are actual effect evaluations.
\end{proof}
The family is not selected by this proposition.  A later claim is closed
on the common/visible subspace once its actual operator word and source
jets lie there.  Quadratic dependence outside that space remains an
explicit extension comparison, and the intrinsic coordinate fingerprint,
the sharp $\sqrt3\Rmom$ outcome radius and the frozen carrier are unchanged.


\section{Marked symmetry, central data and chiral limitations}
\label{paper2c:sec:organized-boundaries}

\subsection{Nonzero domain of the marked comparison and its connection}
\label{p2c:app:marked-comparison-domain}
This diagnostic spells out the domain of the exceptional-coordinate theorem imported in
the main comparison export; it introduces no new seed or physical branch.
For a unit Hopf axis $n$ and normalized even word $A_D(u,h)=r+v$,
quaternion matrix contraction gives
\[
 \|m\|^2=4r^2\bigl(|v|^2-(n\cdot v)^2\bigr),\qquad
 m\ne0\ \Longleftrightarrow\ r\ne0\ \hbox{and}\ n\times v\ne0.
\]
At $n=\mathbf i$, in the convention
$\rho_H(r+x\mathbf i+y\mathbf j+z\mathbf k)
=\left(\begin{smallmatrix}r+ix&y+iz\\-y+iz&r-ix\end{smallmatrix}\right)$,
with $\varepsilon=\left(\begin{smallmatrix}0&1\\-1&0\end{smallmatrix}\right)$,
the restriction of
$\operatorname{tr}(A)(\varepsilon A+(\varepsilon A)^T)/2$ to the positive
Hopf line has coefficient $2r(-y+iz)$, proving the norm formula; rotation
covariance proves the general expression.
Among the 24 Hurwitz units with $u=1,n=\mathbf i$, twenty are nonparallel
and sixteen give a nonzero quadratic: precisely the half-integer units,
all with squared norm $1/2$. The four nonparallel pure marks
$\pm\mathbf j,\pm\mathbf k$ have $r=0$ and zero quadratic. In particular
$h=\mathbf j$ is not in the simultaneous-conjugation orbit of $D_0$,
because conjugation preserves the scalar part.

For the complete tuple $D_0$, fixing its two independent marked axes fixes
the entire oriented frame. Its stabilizer is therefore $\{\pm1\}$ and
$\mathcal O=SU(2)/\{\pm1\}\cong SO(3)$. A temporary unit tensor frame
$e_g=(gs_0)\otimes gh_0g^{-1}$ obeys $e_{-g}=-e_g$, giving the associated
sign line $N_{\mathcal O}$. The quadratic has a constant nonzero coefficient
in these local unit frames; their transitions are constant signs. Declaring
them parallel defines a flat Hermitian connection with holonomy $-1$ around
the nontrivial fundamental loop. If another Hermitian connection differs
by $i\eta$, its action on the dual-square tensor differs by $-2i\eta m$,
forcing $\eta=0$. This parameter-orbit line is distinct from the
symmetric-space and chiral Hopf bundles used in the order-swap appendix.

\subsection{Scope of the canonical zero-tail ambient section}
\label{p2c:app:zero-tail-scope}
The zero-tail section is a positive mathematical right inverse relative to a
fixed cutoff and baseline.  The following details explain why it does not
become an absolute or physical selector.

\begin{proposition}[Conditional Hilbert--Schmidt minimum]
\label{paper2c:proposition:atlas-conditional-hilbert-schmidt-minimum-f448dc3}% stable atlas label; provenance v2_277

Restrict the lift fibre to componentwise Hilbert--Schmidt corrections
\(D=\mathsf X-\mathsf X^{(0)}\), with
\(\|D\|_{\mathrm{HS},4}^2=\sum_a\|D_a\|_{\mathrm{HS}}^2\).  Then the
zero-tail section is the unique minimum-norm correction.
\end{proposition}

\begin{proof}
With \(P=\Pi_m\), \(Q=I-P\), and \(PD_aP=B_a\), the four operator corners
are Hilbert--Schmidt orthogonal:
\[
 \|D_a\|_{\rm HS}^2
 =\|B_a\|_{\rm HS}^2+\|PD_aQ\|_{\rm HS}^2
 +\|QD_aP\|_{\rm HS}^2+\|QD_aQ\|_{\rm HS}^2.
\]
The unique minimum sets all three external corners to zero.  This is exactly
the correction used by \(s_m\).  The result is conditional because the
ambient signature declares no Hilbert--Schmidt cost as physical data.
\end{proof}

\begin{remark}[Operator norm, cutoff, and strong-relation boundaries]
Ordinary bounded-operator norm does not select the zero-tail lift in general.
If the forced finite correction \(B\ne0\), any sufficiently small nonzero
equivariant high-band block \(C\) gives
\(D^{(t)}=B\oplus tC\) with the same row norm as \(B\) for a nontrivial
interval of \(t\).  Moreover, for compatible actions \(A_m=\Pi_mA_n\Pi_m\),
one has \(\Pi_ms_n(A_n)\Pi_m=A_m\) but generally
\(s_n(A_n)\ne s_m(A_m)\) on \(\mathcal H_\infty\).  Compatibility already stabilizes every included matrix element. The next
proposition shows that a uniform component bound is exactly the remaining
condition for a bounded ambient extension of the supplied family. Finally,
compression alone does not
preserve undeclared nonlinear identities such as commutativity or the exact
sphere relation.  These are separate strong-coordinate obligations.
\end{remark}


\begin{proposition}[Bounded extension of a compatible family]
\label{p2c:prop:compatible-bounded-ambient-extension}
Let $\Pi_m\uparrow I$ strongly with dense algebraic union, and let
$A_{m,a}$ be Hermitian operators on the finite ranges, satisfying
$A_{m,a}=\Pi_mA_{n,a}\Pi_m$ for $n\ge m$.
There is a bounded self-adjoint ambient extension $A_a$ with these
compressions if and only if $C_a:=\sup_m\|A_{m,a}\|<\infty$ for every $a$.
It is unique, $\|A_a\|=C_a$, and $\Pi_mA_a\Pi_m\to A_a$ strongly.
Equivariance and scalar matching pass to the extension when imposed on
the compatible family. For a fixed bounded baseline $A^{(0)}$, its zero-tail
sections $A_a^{(0)}+\Pi_m(A_a-A_a^{(0)})\Pi_m$ also converge strongly to $A_a$.
\end{proposition}
\begin{proof}[Bounded extension of a compatible family]
For $u,v$ in the algebraic union define
$b_a(u,v)=\langle u,A_{n,a}v\rangle$ using any $n$ containing both vectors.
Compatibility makes the form well-defined and the uniform bound gives
$|b_a(u,v)|\le C_a\|u\|\|v\|$.
The represented bounded Hermitian operator has the required compressions.
Dense-union matrix elements prove uniqueness, and compression gives the
reverse norm inequality. Equivariance and scalar matching pass through the
same matrix elements. Finally
\[
 \|\Pi_mA_a\Pi_mv-A_av\|
 \le \|A_a\|\,\|\Pi_mv-v\|
       +\|(\Pi_m-I)A_av\|\longrightarrow0.
\]
The same estimate applied to $A_a-A_a^{(0)}$ proves convergence of the
zero-tail sections. Conversely every bounded extension uniformly bounds
its compressions.
\end{proof}

\begin{remark}[Necessity and scope of the bound]
Compatibility alone is insufficient. In the ambient representation choose
disjoint allowed vector edges $(n+1/2,n)\leftrightarrow(n,n+1/2)$,
$n=1,2,\ldots$, normalize the maximum component norm of each Hermitian edge
pair to one, and multiply the pair by $n$. Finite compressions of their
direct sum plus the scalar baseline are compatible and scalar-matched,
but their component norms cannot all be uniformly bounded.
Nor does strong convergence imply norm convergence: nonzero scalar
multiplication on the diffuse sphere is noncompact, so its finite-rank
compressions cannot converge in operator norm. The family itself remains
input data, and its bounded extension is not a physical-action selector.
\end{remark}

\subsection{Factorized position-area ownership}
\label{p2c:sec:v164-position-area-factorization}

\begin{theorem}[No monolithic reciprocal-position operator bridge]
\label{p2c:thm:v164-no-monolithic-pairxp}
The first-symbol pairing, operator-germ bundle, and scalar
reciprocal-position filtration constructed in this paper are the complete
GP1--GP3 local output.  There is no generated global linear intertwiner from
the reciprocal-position vector role to the momentum two-form role, since
\[
 \operatorname{Hom}_{\Spin(4)}
 \bigl((\tfrac12,\tfrac12),(1,0)\oplus(0,1)\bigr)=0.
\]
Consequently a universal global linear \(x\to p\) map is excluded, while the
local symbol/germ pairing remains derived.  The former monolithic
\(\mathrm{PAIR}_{xp}^{\rm op}\) premise must be factorized rather than
assumed.
\end{theorem}

\begin{corollary}[Generated position metric questions]
\label{p2c:cor:v164-generated-position-metric-questions}
Once a reciprocal-position chart and its same-action metric \(g_x\) occur,
the area of an already position-side surface is the intrinsic descendant
\[
 A_{g_x}(\Sigma)=\int_\Sigma dA_{g_x}.
\]
Its numerical unit presentation is a calibration, not a new pairing.  A
same-action gauge connection on the chart likewise defines
\[
 W_R(C)=\frac1{d_R}\operatorname{Tr}\operatorname{Hol}_{A_x}(C)
\]
without a global \(x\)-to-\(p\) operator isomorphism.  Large-loop state
limits and positive surface free energy remain downstream dynamics questions.
\end{corollary}


\subsection{Exact protected transvection--Clifford identification}
\label{p2c:sec:v187-transvection-clifford}
On the normalized protected block write
\[
 e_i=C_iC_4,
 \qquad e_i^2=-1,
 \qquad e_ie_j+e_je_i=-2\delta_{ij}.
\]

\begin{theorem}[Protected transvection--Clifford identity]
\label{p2c:thm:v187-transvection-clifford-identity}
On the protected nonzero coordinate stratum, with $e_i=C_iC_4$,
\[
 \boxed{\widehat P_i\big|_{\rm prot}
              =\frac{i\hbar}{2\Rmom}\,e_i.}
\]
\end{theorem}
\begin{proof}
In the compatible Pauli realization, with Pauli matrices $\sigma_i$,
\[
 C_i=\begin{pmatrix}0&-i\sigma_i\\i\sigma_i&0\end{pmatrix},
 \qquad C_4=\begin{pmatrix}0&I_2\\I_2&0\end{pmatrix}.
\]
Consequently
\[
 e_i=\operatorname{diag}(-i\sigma_i,i\sigma_i),
 \qquad
 (L_i^+-L_i^-)\big|_{\rm prot}
       =\frac12\operatorname{diag}(\sigma_i,-\sigma_i)
       =\frac i2 e_i.
\]
Substitute the already defined
$\widehat P_i=(\hbar/\Rmom)(L_i^+-L_i^-)$.
The normalized even generators are independent of the positive coordinate
magnitude and removable relative chiral phase, so the identity transports
to every representative of the protected class.
\end{proof}

\begin{corollary}[Common vector parent]
\label{p2c:cor:v187-common-vector-parent}
The three-dimensional vector component of the $M_8$ relative solder, the
compact transvection triple, and the protected $\Cl(0,3)$ vector generators
are the same generated vector geometry in successive presentations, with the
normalization fixed above.  Subsequent Clifford grade operations therefore
inherit parent provenance rather than requiring a new cross-sort
identification.
\end{corollary}


\subsection{Finite order-swap holonomy in the protected Clifford owner}
\label{p2c:sec:v188-clifford-order-swap}

\begin{theorem}[Central finite commutator and specified holonomy]
\label{p2c:thm:v188-clifford-central-commutator}
Let $e_i$ be the protected orthonormal generators with
$e_i^2=-1$ and $e_ie_j=-e_je_i$ for $i\ne j$. Then
\[
 e_ie_je_i^{-1}e_j^{-1}=-1,\qquad [e_i,e_j]_{\rm alg}=2e_ie_j.
\]
Under the protected transvection identity, these are the finite and
infinitesimal images of the same retained order-swap representation.
For the canonical symmetric-space connection on
$\Spin(4)\to\Spin(4)/\Spin(3)_{\rm diag}\cong S^3$,
the four right-appended exponential segments with velocities
$e_i,e_j,-e_i,-e_j$, each of parameter length $\pi/2$, form a horizontal
lift of a closed base loop with spin holonomy $-1$.
After either specified chiral projection, the corresponding Hopf connection
with fibre axis the complementary $e_k$ has the same holonomy for the
projected four-segment loop.
\end{theorem}
\begin{proof}
Since $e_i^{-1}=-e_i$, Clifford multiplication gives the displayed group
and algebra commutators. The connection on the symmetric-space bundle is
the stabilizer projection of the left Maurer--Cartan form. On a segment
$g_0\exp(tv)$ its value on the velocity is the stabilizer projection of $v$,
which vanishes for the four transvection velocities above.
The endpoint is
$\exp(\pi e_i/2)\exp(\pi e_j/2)\exp(-\pi e_i/2)\exp(-\pi e_j/2)=-1$,
in the diagonal stabilizer. Thus the base loop is closed and the spin
holonomy is $-1$.
In a chiral $SU(2)$ factor, the Hopf connection projects the Maurer--Cartan
form onto the chosen $e_k$ axis. The orthogonal $e_i,e_j$ velocities are
horizontal and the endpoint is again $-1$ in that fibre subgroup.
These statements specify the connections and loops; the common central
value alone does not identify their bundles with a different parameter-orbit
line bundle.
\end{proof}

\begin{corollary}[Curvature-bivector/normal-vector descent]
\label{p2c:cor:v188-curvature-bivector-normal}
Let $\omega=e_1e_2e_3$ be the retained orientation pseudoscalar.  The vector
part of the $M_8$ solder identifies the curvature two-plane module with
\[
 \Lambda^2\mathbb R^3\cong\Cl^{(2)}(0,3).
\]
For the oriented bivector $u_{ij}=e_ie_j$ and the positively complementary
index $k$,
\[
 \boxed{-u_{ij}\omega=e_k.}
\]
Thus the operation
\[
 \text{parent curvature plane}\longrightarrow
 \text{Clifford bivector}\longrightarrow
 \text{Hodge-dual protected vector}
\]
is canonical after orientation is fixed.  Selection of a particular plane is
still supplied only by the already generated orbit/marked member that owns it.
The oriented Clifford duality used here is Paper~1B,
Theorem~\ref{p1b:thm:protected}; the solder and marked plane retain their
local construction above.
\end{corollary}


\subsection{Reciprocal finite-subgroup invariant filtration}
\label{p2c:app:v191-finite-subgroup-invariant-filtration}

\begin{proposition}[Finite-subgroup projectors commute with every complete position band]
\label{p2c:prop:v191-finite-subgroup-band-commutation}
Let $\Gamma\subset SU(2)$ be a constituted finite subgroup acting on the
position harmonic tower through the already generated reciprocal isometries
$U_x(\gamma)$.  Define
\[
 \Pi_\Gamma^x:=\frac1{|\Gamma|}\sum_{\gamma\in\Gamma}U_x(\gamma),
 \qquad
 P_m^x:=\sum_{k=0}^mP_{E_k^x}.
\]
Because every reciprocal isometry commutes with the round scalar Laplacian,
\[
 \boxed{[\Pi_\Gamma^x,P_m^x]=0\qquad\text{for all }m.}
\]
Thus
\[
 \mathcal C_m^\Gamma
 :=\mathcal P_x(m)\cap\operatorname{Ran}\Pi_\Gamma^x
 =P_m^x\operatorname{Ran}\Pi_\Gamma^x
\]
is a canonical mathematical invariant filtration of the position tower.
The commutation is the complete-band clause of Paper~1B,
Theorem~\ref{p1b:thm:harmonic-transport}, applied to the generated
reciprocal isometry action.
\end{proposition}

\begin{corollary}[Binary-icosahedral rank jump]
\label{p2c:cor:v191-2i-position-rank-jump}
For $\Gamma=2I$, write
\[
 \frac{1-t^{60}}{(1-t^{12})(1-t^{20})(1-t^{30})}
 =\sum_{k\ge0}\mu_k^{2I}t^k.
\]
Then the invariant dimension in the round degree-$k$ scalar carrier is
$(k+1)\mu_k^{2I}$, and therefore
\[
 \dim\mathcal C_m^{2I}=\sum_{k=0}^m(k+1)\mu_k^{2I}.
\]
Since $\mu_k^{2I}=0$ for $0<k<12$ and $\mu_{12}^{2I}=1$,
\[
 \boxed{\dim(\mathcal C_m^{2I})_{\rm nonconst}=0\quad(m\le11),}
 \qquad
 \boxed{\dim(\mathcal C_{12}^{2I})_{\rm nonconst}=13.}
\]
Including the constant mode gives total dimension $14$ at $m=12$.
These dimensions are the one-sided $2I$ specialization of Paper~1B,
Theorem~\ref{p1b:thm:harmonic-transport}. The value $14$ here is the
invariant band at degree twelve, distinct from the complete degree-two
comparison band and from the fifteen-dimensional momentum carrier.
\end{corollary}

\begin{remark}[No quotient or state-gate promotion]
Physical position space is not thereby quotiented by $2I$.  The filtration is
an exact statement about the generated reciprocal action on position
harmonics.  Theorem~\ref{p2c:thm:canonical-reciprocal-harmonic-intertwiner}
structurally intertwines this filtration with the momentum finite-group
projector and preserves the harmonic degree on the bridge image.  It does not
actualize a primordial state/covariance, identify the background frontier
$m_x$ with that harmonic degree, identify the transported degree with an
observed sky multipole, or force an actual perturbation state to wait until
$m_x=12$.
\end{remark}



\subsection{Finite Born inversion on the rooted \texorpdfstring{$C_3$}{C3} transport}
\label{p2c:sec:v269-born-rooted-c3}

The generated reciprocal role carries the contragredient operation
$T^\vee=\iota_1T^{-1}\iota_1^{-1}$ and therefore the opposite composition
law.  Restricting this already generated inversion to the rooted finite
$C_3$ transport gives an exact finite Born-parity split.

\begin{theorem}[Finite Born inversion of the rooted trimer transport]
\label{p2c:thm:v269-born-finite-inversion}
Let $g$ generate the rooted $C_3$ transport and let
$X_G=\rho_G(g)$.  Under reciprocal/Born transport the group element is sent
to its inverse, hence
\[
 \boxed{
 \mathcal B_G(X_G)=\rho_G(g^{-1})=X_G^{-1}=X_G^\dagger.}
\]
Equivalently, adjoining the reciprocal involution to the rooted $C_3$ action
gives
\[
 C_3\rtimes_{\rm inv}C_2\cong D_3\cong S_3.
\]
\end{theorem}

\begin{proof}
The reciprocal-role construction gives the opposite-variance operation
$T\mapsto T^\vee$ by inversion.  Restriction to the constituted finite cyclic
subgroup preserves this rule.  Since $g^3=e$, the inverse of the represented
unitary generator is its adjoint.  The displayed semidirect product is the
standard presentation of the dihedral group of the triangle.
\end{proof}

\begin{corollary}[Finite Born parity of the primitive transport plane]
\label{p2c:cor:v269-born-parity-signline}
With
\[
 C_{1,G}=X_G+X_G^\dagger,
 \qquad
 C_{2,G}=-i(X_G-X_G^\dagger),
\]
one has
\[
 \boxed{\mathcal B_G(C_{1,G})=+C_{1,G},\qquad
 \mathcal B_G(C_{2,G})=-C_{2,G}.}
\]
Thus on the primitive nearest-neighbour transport space the even and odd
Born projections are precisely
\[
 \mathbb R C_{1,G}\oplus\mathbb R C_{2,G}.
\]
The $C_{2,G}$ line is the finite Born-relative/sign representation.
\end{corollary}

\begin{remark}[Scope of the identification]
\label{p2c:rem:v269-local-not-global}
This theorem identifies the induced action on the rooted finite transport
algebra.  It does not identify global Born factor exchange with global
manifold-orientation reversal, and it does not by itself assert that a
fermionic mass/source word consumes the Born-odd family line.  Such
factor-specific occurrence must be established on the physical source owner.
\end{remark}

\section{Exact intrinsic localization and finite position readout}
\subsection{Endpoint localization span and its measurement scope}
\label{p2c:sec:v226-family}

The current complete covariant endpoint family has outcomes $x=rq$ with
$r=\sqrt3R$, $q\in S^3(1)$, and the density
\begin{equation}\label{p2c:eq:v226-K}
 K_{a,\zeta}(q)=
 \begin{pmatrix}
 1&bq^*&0\\
 bq&B I+(a-B)qq^*&hW_q^*\\
 0&hW_q&I+(\Re\zeta)H_q+i(\Im\zeta)*H_q
 \end{pmatrix}\oplus(I+\gamma(q)),
\end{equation}
where
\[
 b=2/\sqrt3,\quad B=(4-a)/3,\quad h=4\sqrt6/9,
 \quad W_qv=q\wedge v,\quad H_q=2W_qW_q^*-I_6,
\]
and $*$ is the oriented four-dimensional Hodge star on bivectors. The
allowed parameters are
\[
 \frac43\le a\le\frac{20}{9},\quad
 |\zeta-u|\le1-u,\quad u=\frac{16}{9(4-a)}.
\]
This parameter domain has nonempty interior in three real coordinates.
The Hermitian native algebra has real dimension $11^2+4^2=137$.

\begin{theorem}[Exact span of all admissible endpoint effects]
\label{p2c:thm:v226-whole-family-rank}
For the continuous covariant family, with all admissible seeds and all
transported directions allowed, the real operator span is
\[
 \mathcal V=\mathbb{R} I\oplus E_1^{\rm fixed}\oplus
                 \bigl(E_2^{(V)}\oplus E_2^{(B)}\oplus E_2^{(*B)}\bigr),
 \qquad \boxed{\dim_\mathbb{R}\mathcal V=1+4+3\cdot9=32.}
\]
The symbols $E_2^{(\cdot)}$ denote independent operator copies of the
round $(1,1)$ representation, not extra fields. If the existing central
$g/s$ sector tags are retained, the total effect span instead has
\[
 \boxed{\dim_\mathbb{R}\mathcal V_{\rm tag}=2+4+4+3\cdot9=37.}
\]
Over an interior native state, the corresponding invisible Hermitian
state-difference spaces have dimensions $105$ and $100$, respectively.
\end{theorem}
\begin{proof}
The density is polynomial of degree at most two in $q$. Its Haar constant
is $I$ for every seed. The linear harmonic component is fixed by the
native first moments and has dimension four. The traceless quadratic part
of $qq^*$ is a copy of $E_2$ on $V$. On bivectors,
\[
 H_q=2\varepsilon_q\iota_q-I_6
\]
is a nonzero traceless quadratic harmonic with Haar mean zero, and so is
$i*H_q$. The two are independent over the reals: the first is real symmetric,
the second imaginary Hermitian with real skew coefficient. Each gives a
nonzero equivariant map from the irreducible nine-dimensional $E_2$;
therefore each map is injective. They are also independent of the $V$ copy
because their block supports differ. The parameter $a-B=(4a-4)/3$ varies
with $a$, while $\Re\zeta$ and $\Im\zeta$ vary independently in the
interior. Hence the family spans all three copies, and no other harmonic
component is possible. Distinct round types are independent; equivalently,
Haar projection separates degrees zero, one and two. This proves $32$.

Central refinement replaces the constant $I$ by independent $P_g,P_S$ and
separates the common linear coordinate component into geometric and spinor
copies. Both are nonzero four-vector maps and have disjoint supports.
The spinor density has no quadratic component, so the three quadratic
copies are unchanged. This proves $37$.

An admitted normalized native state has $136$ independent real differences.
An effect span of dimension $D$ containing $I$ separates at most $D-1$ of
them. The Hilbert--Schmidt annihilator has dimension $137-D$ and is
traceless because $I$ belongs to the span. Around the strictly positive
state $I/15$, every sufficiently small perturbation in that annihilator
is an admitted state difference. Thus these dimensions are attained locally.
\end{proof}

\begin{remark}[What has and has not been resolved]
Each fixed seed already has rank fourteen in the source manuscript. The
new $32/37$ ranks evaluate the \emph{union of the whole allowed family}.
They prove that changing the remaining seed parameters, transporting the
support, or adding the two existing tags cannot complete native state
tomography. Fixed finite quadratures can alias quadratic tests; the above
exact continuous spans are also attained by the Hurwitz 5-design for this
degree-two density, whereas the eight-node rule need not attain them.
\end{remark}
\subsection{Rigidity at the optimal position radius}
\label{p2c:sec:v226-outcomes}

\begin{theorem}[Protected equality for arbitrary optimal POVMs]
\label{p2c:thm:v226-equality}
Let $E_i$ be any finite POVM on $\mathcal H$, with outcomes
$x_i\in\overline B_{\mathbb{R}^4}(r)$, exact first moments $X_a^D$, and
$r=\sqrt3R$. No covariance or subgroup-support assumption is made.
Put $E_i^S=P_SE_iP_S$. If $E_i^S\ne0$, then
$|x_i|=r$, and for $n_i=x_i/r$,
\begin{equation}\label{p2c:eq:v226-equality}
 E_i^S=P_+(n_i)E_i^SP_+(n_i),\qquad
 P_+(n)=\tfrac12(I+\gamma(n)).
\end{equation}
The protected trace weights $w_i=\operatorname{Tr}_S(E_i^S)/4$ obey
\begin{equation}\label{p2c:eq:v226-design}
 \sum_iw_i=1,\qquad \sum_iw_in_i=0,\qquad
 \sum_iw_in_in_i^T=\tfrac14 I_4,
\end{equation}
with zero-weight outcomes omitted.
\end{theorem}
\begin{proof}
Set $c=\sqrt3R/4$, so $r=4c$. Exact moments and Clifford traces give
\[
 \sum_i\operatorname{Tr}_S((x_i\cdot\gamma)E_i^S)
 =\sum_a\operatorname{Tr}_S(\gamma_a c\gamma_a)=16c=4r.
\]
But $rI-x_i\cdot\gamma\succeq0$ and $\sum_i\operatorname{Tr}_S E_i^S=4$.
Thus every nonnegative summand
$\operatorname{Tr}_S((rI-x_i\cdot\gamma)E_i^S)$ is zero. A positive product has
zero trace only if the support of $E_i^S$ lies in the kernel of the first
factor. That kernel is empty when $|x_i|<r$, and otherwise is exactly
$\operatorname{Ran}P_+(n_i)$. This proves \eqref{p2c:eq:v226-equality}.

The Clifford identity gives
$P_+(n)\gamma_aP_+(n)=n_aP_+(n)$, hence
$\operatorname{Tr}_S(\gamma_aE_i^S)=n_{ia}\operatorname{Tr}_SE_i^S$.
Take the trace, then the $\gamma_a$ trace, of $\sum_iE_i^S=I_4$;
finally take the $\gamma_b$ trace of
$\sum_irn_{ia}E_i^S=c\gamma_a$. These are precisely
\eqref{p2c:eq:v226-design}.
\end{proof}

\begin{corollary}[Five outcomes cannot attain the intrinsic optimal radius]
\label{p2c:cor:v226-five}
Every finite POVM with exact intrinsic first moments and radius
$r=\sqrt3R$ has at least six outcomes. The lower bound is independent
of any upper-bound construction. The seven-outcome theorem below first gives the interval
$6\le N_{\min}\le7$; the later support-free sharp converse,
Theorem~\ref{p2c:thm:universal-geometric-rank6-radius},
subsequently excludes six outcomes at this endpoint.
\end{corollary}
\begin{proof}
The augmented vectors $\sqrt{w_i}(1,2n_i)^T\in\mathbb{R}^5$ have outer products
summing to $I_5$ by \eqref{p2c:eq:v226-design}. There must therefore be at least
five nonzero protected weights. If the whole POVM has five outcomes,
these five vectors are the columns of an orthogonal matrix. Their norms
and pairwise inner products force
\[
 w_i=\frac15,\qquad n_i\cdot n_j=-\frac14\quad(i\ne j).
\]
The support is a regular four-simplex. Its affine independence means that
normalization and first moments uniquely fix the \emph{full-carrier}
effects, not just their protected compressions:
\[
 E_i=\frac15\left(I+\frac4r n_i\cdot X^D\right).
\]
On the scalar/vector block, $n_i\cdot X^D$ has the eigenvalue $-R/2$.
Consequently each proposed effect has eigenvalue
\[
 \boxed{\frac15\left(1-\frac2{\sqrt3}\right)<0,}
\]
contradicting positivity. Five outcomes are impossible, as are fewer than
five. This proves the lower bound without an upper-bound premise.
\end{proof}

\subsection{Seven outcomes at the optimal intrinsic radius}
\label{p2c:sec:v255-seven-outcomes}

\begin{theorem}[Seven-outcome attainment of the intrinsic optimal radius]
\label{p2c:thm:v255-seven-outcome-attainment}
The intrinsic coordinate tuple $X_a^D$ on the retained carrier admits a
normalized seven-outcome POVM with exact first moments and all outcome radii
$\sqrt3\Rmom$. Consequently the arbitrary-support minimum obeys
\[
                 6\le N_{\min}\le7.
\]
No covariance, subgroup-support, or apparatus-occurrence assumption is
required. The later support-free sharp converse,
Theorem~\ref{p2c:thm:universal-geometric-rank6-radius},
resolves the endpoint count in favor of seven.
\end{theorem}
\begin{proof}
Write $\mathcal H=\mathcal H_G\oplus\mathcal H_S$, where
$\mathcal H_G=\mathbb C\oplus\mathbb C^4\oplus\Lambda^2\mathbb C^4$
and $\dim\mathcal H_S=4$. Use the intrinsic real sector-phase gauge and the
ordered orthonormal geometric basis
\[
 (e_0,u_1,u_2,u_3,u_4,w_{12},w_{13},w_{14},w_{23},w_{24},w_{34}),
 \qquad w_{ab}=e_a\wedge e_b\ (a<b).
\]
Put $Y_a=P_GX_a^DP_G$, $r=\sqrt3\Rmom$, and $Y(q)=\sum_aq_aY_a$.
The native tuple has scalar/vector coefficient $\Rmom/2$ and
vector/two-form coefficient $\sqrt2\Rmom/3$:
\[
 Y_ae_0=\tfrac{\Rmom}{2}u_a,\qquad
 Y_au_b=\tfrac{\Rmom}{2}\delta_{ab}e_0+
                   \tfrac{\sqrt2\Rmom}{3}e_a\wedge e_b,
\]
with the reverse edges fixed by self-adjointness. On $\mathcal H_S$ write
$X_a^D=(r/4)C_a$, with Hermitian Clifford matrices $C_a$.

Define positive diagonal resource operators on $\mathcal H_G$ by
\[
 A=\operatorname{diag}\left(
 \tfrac29,\tfrac38,\tfrac4{27},\tfrac15,\tfrac15,
 \tfrac12,\tfrac{10}{27},\tfrac{10}{27},0,0,0\right),
\]
let $B$ be its conjugate under the orthogonal interchange of directions $1,2$,
and put $D=I_G-A-B$. Explicitly,
\[
 D=\operatorname{diag}\left(
 \tfrac59,\tfrac{103}{216},\tfrac{103}{216},\tfrac35,\tfrac35,
 0,\tfrac{17}{27},\tfrac{17}{27},\tfrac{17}{27},\tfrac{17}{27},1\right).
\]
Choose seven unit directions
\[
 \mathcal Q=\{e_1,-e_1,e_2,-e_2,q_0,q_1,q_2\},\qquad
 q_j=\cos(2\pi j/3)e_3+\sin(2\pi j/3)e_4.
\]
For outcomes $x_q=rq$, define the geometric effects by
\[
 G_{\pm e_1}=\tfrac12(A\pm Y_1/r),\qquad
 G_{\pm e_2}=\tfrac12(B\pm Y_2/r),\qquad
 G_{q_j}=\tfrac13(D+2Y(q_j)/r).
\]
Define their protected parts by
\[
 S_{\pm e_1}=\tfrac18(I_S\pm C_1),\qquad
 S_{\pm e_2}=\tfrac18(I_S\pm C_2),\qquad
 S_{q_j}=\tfrac16(I_S+C(q_j)),
 \quad C(q)=\sum_aq_aC_a,
\]
and set $E_q=G_q\oplus S_q$.

We verify positivity exactly. The nontrivial $2\times2$ blocks of
$A\pm Y_1/r$ have diagonal pairs
\[
 (\tfrac29,\tfrac38),\qquad
 (\tfrac4{27},\tfrac12),\qquad
 (\tfrac15,\tfrac{10}{27})\quad\text{(twice)}.
\]
Their squared off-diagonal entries are respectively $1/12,2/27,2/27$,
so all four determinants vanish and all diagonal entries are nonnegative.
Unpaired entries are nonnegative too. Interchanging $1,2$ proves the same
for $B\pm Y_2/r$.
The operator $D$ is invariant under rotations in the $(3,4)$ plane, so
$D+2Y(q_j)/r$ is unitarily equivalent to $D+2Y_3/r$.
Its nontrivial $2\times2$ blocks have determinant margins
\[
 \tfrac59\tfrac35-\tfrac13=0,\qquad
 \tfrac{103}{216}\tfrac{17}{27}-\tfrac8{27}
       =\tfrac{23}{5832}>0\quad\text{(twice)},\qquad
 \tfrac35-\tfrac8{27}=\tfrac{41}{135}>0.
\]
All remaining diagonal entries are nonnegative. Thus every $G_q$ is
positive. Since $C(q)^2=I_S$ for unit $q$, every $S_q$ is positive too.

The elementary triangle identities
\[
 \sum_{j=0}^2q_j=0,\qquad
 \sum_{j=0}^2q_jq_j^T=\tfrac32
       (e_3e_3^T+e_4e_4^T)
\]
give $\sum_qG_q=A+B+D=I_G$ and
$\sum_qrq_aG_q=Y_a$ for $a=1,\ldots,4$.
The protected weights are $1/8$ on each of the four axial directions and
$1/6$ on each of the three triangle directions. They sum to one, have zero
weighted first vector moment, and have weighted second vector moment
$I_4/4$. Hence $\sum_qS_q=I_S$ and
$\sum_qrq_aS_q=(r/4)C_a$.
Combining the two summands proves normalization and every exact intrinsic
first moment on the full retained carrier. All seven protected effects are
nonzero, so the construction has seven actual mathematical outcomes.
The radius is optimal by Theorem~\ref{p2c:thm:optimal-intrinsic-povm-radius};
Corollary~\ref{p2c:cor:v226-five} supplies the lower bound six.
This construction provides a mathematical readout and does not supply an
occurring apparatus, instrument, preparation, or outcome token.
\end{proof}


\subsection{Universal chirality information and localization radius}
\label{p2c:sec:v226-radius-information}
Let $R=R_{\rm mom}$, $r_0=\sqrt3R$ and
$\chi=\gamma_1\gamma_2\gamma_3\gamma_4$ on the protected block.
Its sign follows the fixed orientation; both chiral sectors remain.
The operator $P_S\chi P_S$ is the coordinate word of
Theorem~\ref{p1:thm:v226-chirality}.

\begin{theorem}[Sharp chirality-radius information bound]
\label{p2c:thm:v226-chirality-radius}
Let a normalized finite POVM have assigned outcome coordinates
$x_i\in\overline B_{\mathbb R^4}(r)$ and exact moments $X_a^D$, with
$r\ge r_0$. For protected source states
$\rho_\pm=(I_4\pm\chi)/4$, its total-variation distinguishability satisfies
\[
 \boxed{\frac12\sum_i|p_i^+-p_i^-|
        \le\sqrt{1-\frac{r_0^2}{r^2}}.}
\]
For $\rho_s=(I_4+s\chi)/4$, $|s|<1$, the classical Fisher information
at $s=0$ satisfies
\[
 \boxed{\mathcal I_\chi(0)\le1-\frac{r_0^2}{r^2}.}
\]
Both bounds are sharp for the full intrinsic carrier when the complete
outcome label is retained. In particular every optimal-radius POVM is blind
to protected chirality, without a covariance or subgroup assumption.
It cannot be informationally complete on $\mathcal A_D$ for any number
of outcomes or additional labels. The assertions extend to continuous POVMs
using their protected trace measure.
\end{theorem}
\begin{proof}
Write $F_i=P_SE_iP_S$, $w_i=\operatorname{Tr}F_i/4$, and for $w_i>0$ set
$\tau_i=F_i/\operatorname{Tr}F_i$,
$u_{ia}=\operatorname{Tr}(\tau_i\gamma_a)$,
$v_i=\operatorname{Tr}(\tau_i\chi)$.
The five Hermitian matrices $\gamma_1,\ldots,\gamma_4,\chi$ anticommute and
square to the identity. Applying the operator-norm bound to their real
linear combinations gives $|u_i|^2+v_i^2\le1$.
The exact first moments give
\[
 \sum_iw_i=1,\qquad
 \sum_iw_i x_i\cdot u_i=r_0.
\]
Thus $\sum_iw_i|u_i|\ge r_0/r$, and Cauchy--Schwarz yields
\[
 \sum_iw_iv_i^2\le1-\sum_iw_i|u_i|^2
                      \le1-(r_0/r)^2.
\]
Since $p_i(s)=w_i(1+sv_i)$, total variation is
$\sum_iw_i|v_i|$ and $\mathcal I_\chi(0)=\sum_iw_iv_i^2$.
This proves both bounds. At $r=r_0$, $v_i=0$ for every effect, so the
nonzero native direction $P_S\chi P_S$ is invisible.

To attain the bounds, put $\lambda=r_0/r$,
$\beta=\sqrt{1-\lambda^2}$ and use the eight directions $q\in Q_8$ with
a binary label $t=\pm1$. If $K_g(q)$ is the existing positive geometric
endpoint density, use the sixteen effects
\[
 E_{q,t}=\frac1{16}\left(
 [I_g+\lambda(K_g(q)-I_g)]\oplus
 [I_4+\lambda\gamma(q)+t\beta\chi]\right),
 \qquad x_{q,t}=rq.
\]
The geometric part is a convex mixture. The protected part is positive
because $(\lambda\gamma(q)+t\beta\chi)^2=I_4$.
The $Q_8$ moments and the sum over $t$ prove normalization and the exact
native first moments. Protected conditional chirality is $v_{q,t}=t\beta$,
so both inequalities are equalities. Coarsening away $t$ can lose this
information. In the continuous case take the scalar probability measure
$\operatorname{Tr}_S E^S/4$ and its finite-matrix Radon--Nikodym density;
the same pointwise inequalities and integrals prove the assertions.
\end{proof}

\begin{theorem}[Unattained optimal radius for complete position tomography]
\label{p2c:thm:v226-ic-radius-infimum}
Among arbitrary mathematical POVMs with exact intrinsic first moments,
the infimum outcome radius allowing full native informational completeness
is $r_0=\sqrt3R$, and it is not attained. For every $r>r_0$ a finite POVM
with distinct position outcomes on $S^3(r)$ exists. Exactly $137$ outcomes
are sufficient and necessary on the full native state space, by
Theorem~\ref{p2c:thm:v258-minimal-complete-readout}.
A fixed noncovariant readout frame is allowed in this existence statement.
\end{theorem}
\begin{proof}
The lower bound is the existing optimal-radius theorem, and nonattainment
follows from Theorem~\ref{p2c:thm:v226-chirality-radius}.
For $r>r_0$ put $\lambda=r_0/r$ and let
$K^0(q)=I+\lambda(K_*(q)-I)$, with $K_*$ any existing positive endpoint
density of degree at most two. Then $K^0\succeq(1-\lambda)I$, is normalized,
and has first moments $X_a^D$ with outcome $rq$.
Choose an HS-orthonormal traceless native Hermitian basis $F_1,\ldots,F_{136}$
and 136 independent real scalar harmonics $f_k$ among degrees three to
seven. There are $16+25+36+49+64=190$ available modes.
Set $C=\sum_k\|f_k\|_\infty\|F_k\|$ and
\[
 K(q)=K^0(q)+\frac{1-\lambda}{2C}\sum_k f_k(q)F_k.
\]
This remains positive. Higher harmonic orthogonality preserves normalization
and first moments. If $\operatorname{Tr}(AK(q))$ vanishes for every $q$,
projection onto the chosen harmonics makes every $\operatorname{Tr}(AF_k)$
zero; the constant component then makes $\operatorname{Tr}A=0$.
Thus $A=0$ in the native algebra and the POVM is informationally complete.

All its density entries have degree at most seven. A positive spherical
cubature rule exact through degree fourteen preserves its normalization,
first moments and complete effect-frame Gram matrix. Such a rule exists
with at most $\dim\mathcal P_x(14)=15\cdot16\cdot31/6=1240$ nodes:
the Haar moment vector lies in the convex hull of the compact evaluation
vectors, and affine dependence permits elimination of weights until at
most 1240 remain. Merge coincident nodes. The preserved positive-definite
frame proves finite informational completeness at distinct position points.
This supplies a marked mathematical measurement, not an occurrence or an
unmarked covariant single-seed construction. The displayed cubature argument
gives the older $1240$-node upper bound; the independent construction of
Theorem~\ref{p2c:thm:v258-minimal-complete-readout} proves the exact count $137$.
\end{proof}

\subsection{The exact outcome count for complete intrinsic position readout}
\label{p2c:sec:v258-minimal-complete-readout}

\begin{theorem}[Minimal complete position POVMs at every strict radius slack]
\label{p2c:thm:v258-minimal-complete-readout}
Let $\mathcal A=M_{11}(\mathbb C)\oplus M_4(\mathbb C)$ be the native
coordinate algebra, let $X_a^D$ be its intrinsic coordinate tuple, and set
$r_0=\sqrt3\Rmom$. For every $r>r_0$ there is a normalized POVM
$\{E_j\}_{j=1}^{137}\subset\mathcal A$ such that
\[
 E_j\succ0,\qquad x_j\in S^3(r),\qquad x_j\ne x_k\ (j\ne k),
 \qquad \sum_jx_{j,a}E_j=X_a^D\quad(a=1,\ldots,4),
\]
and the $137$ effects form a real basis of $\mathcal A_{\rm sa}$.
Thus this position POVM is informationally complete on the full native
state space, and its outcome count is minimal. If $N_{\rm IC}(r)$ is the
least number of effects of a native informationally complete POVM with
these exact moments and outcomes in $\overline B^4(r)$, with infinity
denoting nonexistence, then
\begin{equation}
 \boxed{
 N_{\rm IC}(r)=
 \begin{cases}
  \infty,&0<r\le r_0,\\
  137,&r>r_0.
 \end{cases}}
 \label{p2c:eq:v258-exact-ic-count}
\end{equation}
The existence proof permits a fixed marked noncovariant readout frame.
It does not claim informational completeness on $M_{15}$ or supply an
apparatus, preparation, instrument, or registered outcome.
\end{theorem}
\begin{proof}
The real dimension of $\mathcal A_{\rm sa}$ is
$11^2+4^2=137$. A POVM with fewer effects has proper real operator span.
Because its span contains $I=\sum_jE_j$, a nonzero Hermitian element
$Z$ in its Hilbert--Schmidt annihilator is traceless. For sufficiently
small $\epsilon>0$, the distinct native states
$I/15\pm\epsilon Z$ are positive and give identical probabilities.
Consequently informational completeness requires at least $137$ effects.
The already proved radius bound excludes exact moments when $r<r_0$;
Theorem~\ref{p2c:thm:v226-chirality-radius} excludes full native
informational completeness at $r=r_0$. It remains to construct $137$
effects at an arbitrary $r>r_0$.

Write $d=137$ and $\lambda=r_0/r\in(0,1)$. Choose any positive intrinsic
endpoint density $K(q)$ admitted by
Corollary~\ref{p2c:cor:v233-optimal-quaternion-povm}, and enumerate its
eight quaternion nodes as $q_i\in\{\pm e_a:1\le a\le4\}$. The diluted
effects
\[
 W_i=\frac18\bigl[(1-\lambda)I+\lambda K(q_i)\bigr]
\]
are positive definite and satisfy
\[
 \sum_iW_i=I,\qquad
 \sum_iq_{i,a}W_i=\frac{X_a^D}{r},\qquad
 W_i\succeq\frac{1-\lambda}{8}I.
\]
The first two identities use the quaternion rule at radius $r_0$ and
$\sum_iq_i=0$; hence dilation of the classical radius has not changed
the requested operator moments.

Split $W_1$ into $18$ equal effects and every other $W_i$ into $17$ equal
effects. There are exactly $18+7\cdot17=d$ resulting effects, denoted
$F_j^0$, initially carrying repeated nodes $q_j^0$. They obey the same
five moment identities and the uniform strict positivity bound
\[
 F_j^0\succeq\mu I,\qquad \mu=\frac{1-\lambda}{144}>0.
\]
Let $T^0$ be the real $5\times d$ matrix with columns
$(1,(q_j^0)^T)^T$. Its row rank is five, because the eight distinct
centers are the four antipodal coordinate pairs. Define the fixed
operator-valued moment row by
\[
 M_0=I,\qquad M_a=X_a^D/r\quad(a=1,\ldots,4).
\]
These five Hermitian operators are linearly independent: on the protected
block their restrictions are $I_4$ and $(\lambda/4)C_a$; the Clifford
trace identities separate the scalar and all four vector coefficients.

We next make all $d$ positions distinct while retaining strict positivity
and every moment exactly. In the group over $q_i$, choose a unit vector
$v_i\perp q_i$ and use the real unit directions
\[
 q_{i,k}(u)=
 \frac{1-k^2u^2}{1+k^2u^2}\,q_i+
 \frac{2ku}{1+k^2u^2}\,v_i,
 \qquad 1\le k\le m_i,
\]
where $m_1=18$, $m_i=17$ for $i>1$, and $u>0$ will be small.
Within each group the stereographic parameter $ku$ is injective.
For $0<u\le1/1800$ these points are within distance $36u\le1/50$
of their respective centers, whose distinct pairwise distances are at
least $\sqrt2$. Thus all $d$ nodes are distinct. Let $T(u)$ have columns
$(1,q_j(u)^T)^T$. For sufficiently small $u$ it has row rank five, and
its real right inverse
\[
 L(u)=T(u)^T\bigl(T(u)T(u)^T\bigr)^{-1}
\]
is continuous at zero. Define the five operator defects and corrected
effects by
\[
 D_\alpha(u)=M_\alpha-\sum_jT_{\alpha j}(u)F_j^0,
 \qquad
 F_j(u)=F_j^0+\sum_{\alpha=0}^4L_{j\alpha}(u)D_\alpha(u).
\]
Since $T(u)L(u)=I_5$, one has exactly
\begin{equation}
 \sum_jT_{\alpha j}(u)F_j(u)=M_\alpha
 \qquad(\alpha=0,\ldots,4).
 \label{p2c:eq:v258-moment-correction}
\end{equation}
Every corrected effect is Hermitian and belongs to $\mathcal A$.
There is also an explicit positivity check. Here $D_0=0$ and
\[
 \|D_a(u)\|\le36u\quad(a=1,\ldots,4),\qquad
 \|F_j(u)-F_j^0\|
 \le36u\sum_{a=1}^4|L_{ja}(u)|.
\]
Indeed each $D_a$ is a real weighted sum of the positive effects $F_j^0$
with coefficients bounded in absolute value by $36u$; their sum is $I$.
Because $L(u)$ stays bounded near zero, choose a positive $u$ small
enough that $T(u)$ has row rank five and
\[
 36u\max_j\sum_{a=1}^4|L_{ja}(u)|\le\mu/2.
\]
Then $F_j(u)\succeq(\mu/2)I$ for every $j$. Such a $u$ can be found by
successive halving from $1/1800$. Taking the $v_i$ to be coordinate
vectors makes all node and right-inverse coefficients rational for
rational $u$. This completes the distinct-position, strictly positive
feasibility step; no generic positivity assertion is being used.

Fix these nodes and write $T=T(u)$. Choose a real basis of
$\mathcal A_{\rm sa}$ and use row coordinates for its operators. Let
$B_0\in\mathbb R^{d\times d}$ have the rows of the effects $F_j(u)$,
and let $M\in\mathbb R^{5\times d}$ have the rows of the five targets.
Equation~\eqref{p2c:eq:v258-moment-correction} is exactly
\[
 TB_0=M.
\]
Both $T$ and $M$ have row rank five. Extend their five rows separately
to real bases of $\mathbb R^d$, producing invertible matrices
$\widetilde T,\widetilde M\in\mathbb R^{d\times d}$. Then
\[
 B_*:=\widetilde T^{-1}\widetilde M
 \quad\hbox{is invertible and satisfies}\quad TB_*=M.
\]
Let $H_j$ be the Hermitian operators represented by its rows. They need
not be positive. Put
\[
 b=\max_j\|H_j-F_j(u)\|.
\]
If $b=0$, the current effects already form a basis. Otherwise put
$\eta=\min(1,\mu/(4b))>0$ and consider
\[
 E_j(t)=F_j(u)+t\bigl(H_j-F_j(u)\bigr),\qquad0\le t\le\eta.
\]
All these effects satisfy the five exact moment constraints, and
\[
 E_j(t)\succeq\frac\mu2I-tbI\succeq\frac\mu4I\succ0.
\]
Their coefficient determinant
$\det\bigl(B_0+t(B_*-B_0)\bigr)$ is a real polynomial of degree at most
$d$, and is nonzero at $t=1$. Thus it has at most $d$ distinct roots.
Among the $d+1=138$ distinct positive values
\[
 t_k=\frac{k\eta}{138},\qquad k=1,\ldots,138,
\]
at least one gives a nonzero determinant. Choose such a value. The
corresponding $137$ positive effects are linearly independent, sum to
$I$, and have exact moments $\sum_jq_{j,a}E_j=X_a^D/r$ at the already
fixed distinct directions. Assigning $x_j=rq_j$ proves the theorem on
$S^3(r)$, hence also in the radius-$r$ closed ball.
\end{proof}

\noindent\textit{Claim boundary.}
This sharpens the $1240$-node cubature construction of
Theorem~\ref{p2c:thm:v226-ic-radius-infimum} to the exact minimal number
$137$ at every positive radius slack. Its determinant selection chooses
a marked mathematical test frame within an already feasible moment
fiber; it is not a physical branch, source member, or apparatus selector.
The endpoint remains noninformationally complete, however many effects
are used. The separate six-versus-seven question concerns exact moments
at the endpoint without demanding complete state tomography.


\begin{corollary}[Complete native readout and coefficient export]
\label{p2c:cor:v226-complete-readout}
The 137-outcome measurement of Theorem~\ref{p1:thm:v226-tomography} and
the dual tests of Theorem~\ref{p1b:thm:v226-dual-coefficients} are available
on the same native algebra as the coordinates. A $137$-outcome complete measurement can retain every exact intrinsic
position moment at any specified radius $r>\sqrt3R$, by
Theorem~\ref{p2c:thm:v258-minimal-complete-readout}; complete tomography
at the endpoint remains impossible. Their mathematical
availability introduces no second momentum carrier or apparatus occurrence.
The full endpoint-family ranks 32 and 37 remain the sharper limitations of
that specified covariant family; the universal optimal-radius obstruction
above applies even outside it.
\end{corollary}
\begin{proof}
Apply the two provider theorems on the coordinate-generated algebra and
use the stated first-moment domains. The chirality obstruction is independent
of the seed-family rank calculation.
\end{proof}


\section{Native covariant readout spans and minimum-outcome converses}
\subsection{Covariant position, axis, and directed-tangent readout capacities}
\label{p2c:sec:covariant-direction-capacity}
Let $G=\operatorname{Spin}(4)$, $K=SU(2)_{\rm diag}$, $H=U(1)\subset K$
be the stabilizer of a directed unit tangent at the north pole, and
$N=N_K(H)$. On the intrinsic branch put
$\mathcal A_D=M_{11}(\mathbb C)\oplus M_4(\mathbb C)$.
For $J\in\{K,N,H\}$ set
\[
 \mathcal V_J=\operatorname{span}_{\mathbb R}
 \{U_gBU_g^*:g\in G,\;B\in(\mathcal A_D)_{\rm sa}^{J}\}.
\]
\begin{theorem}[Sharp homogeneous covariant native readout capacity]
\label{p2c:thm:covariant-stabilizer-capacity}
For $J=K,N,H$ one has, respectively,
\[
 \boxed{\dim_{\mathbb R}\mathcal V_K=65,\qquad
 \dim_{\mathbb R}\mathcal V_N=107,\qquad
 \dim_{\mathbb R}\mathcal V_H=137.}
\]
Every $G$-covariant POVM on $G/J$ with finitely many $G$-inert passive
labels has effects in $\mathcal V_J$. Thus position-only covariant
readouts have $72$ indistinguishable native Hermitian directions, while
position plus an unoriented tangent axis has $30$. No such invariant-span
obstruction remains for directed-tangent outcomes $G/H\simeq US^3$.
\end{theorem}
\begin{proof}
The isotypic components of $\mathcal A_{D,\rm sa}$, listed as
$(j_L,j_R):m$, are
\[
\begin{aligned}
 &(0,0):6,\quad (\tfrac12,\tfrac12):8,\quad (1,1):3,\\
 &(1,0):5,\quad(0,1):5,\quad(2,0):1,\quad(0,2):1,\\
 &(\tfrac32,\tfrac12):2,\quad(\tfrac12,\tfrac32):2.
\end{aligned}
\]
The decomposition follows by Clebsch--Gordan on the four geometric
modules $V_{0,0},V_{1/2,1/2},V_{1,0},V_{0,1}$ and the two spinor modules
$V_{1/2,0},V_{0,1/2}$, taking endomorphisms of the two blocks separately.
Restriction to $K$ is $V_{j_L,j_R}|_K=\bigoplus_{J=|j_L-j_R|}^{j_L+j_R}V_J$.
The $K$-fixed dimension is one exactly for $j_L=j_R$.
Every integer $K$-spin has one $H$-fixed weight-zero vector; under the
normalizer Weyl involution that vector has eigenvalue $(-1)^J$, so the
$N$-fixed multiplicity is the number of even $J$ in the displayed range.
The orbit span of all $J$-fixed seeds consists exactly of the irreducible
isotypic components having a nonzero $J$-fixed vector. Summing their
real dimensions gives $65,107,137$. For a covariant POVM, its invariant
trace measure is Haar; absolute continuity of its matrix elements and
covariance yield a $J$-fixed positive seed on each passive-labelled orbit.
The resulting effect span lies in $\mathcal V_J$. For a nonzero
traceless Hermitian $Z\in\mathcal V_J^\perp$, the states
$I_{15}/15\pm\eta Z$ are positive for
$0<\eta<1/(15\|Z\|)$ and indistinguishable by all those effects.
\end{proof}

\begin{theorem}[Exact moment-preserving saturation and minimal passive labels]
\label{p2c:thm:covariant-direction-complete-povm}
Put $r_0=\sqrt3R_{\rm mom}$. At every $r>r_0$ there exist normalized
$G$-covariant native POVMs on
$(G/J)\times\{1,\ldots,L\}$, with position outcome $x=rq$,
exact first moments $\int x_a\,dE=X_a^D$, and effect span
$\mathcal V_J$, for precisely the sharp minimal passive label numbers
\[
 \boxed{L_{K,\min}=8,\qquad L_{N,\min}=8,\qquad L_{H,\min}=6.}
\]
In particular the $J=H$ construction is fully informationally complete
on the $137$-dimensional native Hermitian operator space while retaining
full $\operatorname{Spin}(4)$ covariance. At $r\le r_0$ such complete
exact-moment POVMs do not exist, by the established optimal-radius and
chirality-radius theorems of Paper II-C.
\end{theorem}
\begin{proof}
Write $h_{\lambda,J}=\dim(V_\lambda^J)$ and $m_\lambda$ for the above
isotypic multiplicity. The multiplicity rank of $L$ labelled $J$-fixed
seeds is bounded by $Lh_{\lambda,J}$. The maximal lower bounds
$\lceil m_\lambda/h_{\lambda,J}\rceil$ over visible types are $8,8,6$.
For sufficiency, let $K_*(q)\succeq0$ be the pre-existing normalized
covariant optimal-radius intrinsic POVM density of Paper II-C. Set
$\lambda=r_0/r$ and
$K_r(q)=(1-\lambda)I+\lambda K_*(q)$, so
$K_r\succeq(1-\lambda)I>0$ and
$\int rq_aK_r\,d\mu=X_a^D$.
For $J=H$ choose five Hermitian $H$-fixed seeds $B_1,\ldots,B_5$
such that their group orbits together with $I$ span $\mathcal A_{D,\rm sa}$.
They exist because every nontrivial type has
$m_\lambda\le5h_{\lambda,H}$ and the five seeds can span the complement
of $I$ in the six-dimensional invariant component. For $J=K,N$, take
seven $J$-fixed seeds whose orbits together with $K_r$ span $\mathcal V_J$;
all types admit seven seed multiplicity directions except the eightfold
first-harmonic type, on which the nonzero first moments already supply
one independent direction. Choose the seven seed coefficients in its
complement. These choices can be made independently in distinct
isotypic components.
With $L=6$ or $8$, let $B_L=-\sum_{j<L}B_j$ and define
\[
 D_j(gJ)=L^{-1}K_r(gK)+\epsilon U_gB_jU_g^*,\qquad j=1,\ldots,L.
\]
Each density is well-defined and positive when
$0<\epsilon<(1-\lambda)/(L\max_j\|B_j\|)$.
Summing over labels gives $K_r(gK)$ pointwise, preserving normalization
and exact first moments after Haar integration. Differences
$D_j-D_L=\epsilon U_g(B_j+\sum_{i<L}B_i)U_g^*$ span the individual
seed orbits because $I_{L-1}+\mathbf1\mathbf1^T$ is invertible.
Thus the effects span exactly $\mathcal V_J$.
\end{proof}

\begin{lemma}[No globally equivariant directed-tangent selector]
\label{p2c:lem:no-equivariant-tangent-section}
The natural bundle $G/H\to G/K$ admits no $G$-equivariant section.
Among nontrivial connected homogeneous $K=SU(2)$ auxiliary fibres,
$K/H\cong S^2$ has the least possible positive dimension, two.
\end{lemma}
\begin{proof}
A $G$-equivariant section would select a $K$-fixed point in the
transitive $K$-space $K/H$, which has none. Connected $K$ acts trivially
on every discrete finite orbit, and a one-dimensional nontrivial
homogeneous orbit would require a two-dimensional Lie subalgebra of
$\mathfrak{su}(2)$, which does not exist. The maximal torus stabilizer
has codimension two, proving the minimum.
\end{proof}

\begin{remark}[Scope firewall]
These are mathematical existence and span results on the intrinsic
branch, not physical detector selection. The directed-tangent subgroup
$U(1)\subset SU(2)_{\rm diag}$ differs from Paper II-E's left-Hopf
$U(1)_H\subset SU(2)_L$. Universal GP3 rigidity and the Paper II-F
negative boundaries remain unchanged.
\end{remark}


\begin{corollary}[Radius-conditioned instability of native state reconstruction]
\label{p2c:cor:v270-radius-reconstruction-instability}
Let $r>r_0$, and let $\mathcal M$ be any normalized exact-first-moment
native POVM whose assigned position outcomes have length at most $r$,
including POVMs with additional directional and passive labels.  Suppose
$\mathcal R$ exactly reconstructs the native protected states from the
resulting probability laws and is Lipschitz from classical total variation
to quantum trace distance.  Then
\[
 \boxed{\operatorname{Lip}(\mathcal R)\ge
 \left(1-\frac{r_0^2}{r^2}\right)^{-1/2}.}
\]
The bound diverges as $r\downarrow r_0$, notwithstanding existence of
informationally complete mathematical measurements for every $r>r_0$.
\end{corollary}
\begin{proof}
Take the orthogonal protected states $\rho_\pm=(I_4\pm\chi)/4$,
extended by zero on the geometric carrier.  Their quantum trace distance
is one.  The sharp chirality-radius theorem
\ref{p2c:thm:v226-chirality-radius} bounds the total-variation distance
of their measurement laws by $\sqrt{1-r_0^2/r^2}$, irrespective of the
extra outcome labels.  Any exact left inverse must increase that pairwise
distance by at least the displayed ratio.
\end{proof}

\begin{proposition}[Finite exact outcomes cannot carry full continuous round covariance]
\label{p2c:prop:v270-finite-outcome-covariance-no-go}
No finite discrete outcome space bearing a continuous action of connected
$G=\Spin(4)$ can carry a normalized $G$-covariant POVM with nonzero
four-vector first moments and a $G$-equivariant nonzero assigned position
vector at each outcome.  This does not prohibit finite, marked,
noncovariant $137$-effect native frames.
\end{proposition}
\begin{proof}
A continuous action of a connected group on a finite discrete set is
trivial.  The equivariance of each position label then makes it a fixed
vector of the standard $G$-action on $\mathbb R^4$, whose only fixed
vector is zero.  Hence every first moment vanishes, contradicting the
nonzero intrinsic tuple $X^D$.  The cited finite construction does not
claim continuous $G$-covariance of its outcome labels.
\end{proof}

\begin{remark}[Exact mathematical scope of the direction-position extension]
\label{p2c:rem:v270-direction-physical-firewall}
The capacities $65/107/137$, minimum passive-orbit channel counts
$8/8/6$, and strict-radius six-channel directed-tangent construction are
existence and obstruction results on $\mathcal A_D$.  Six channels are six
\emph{continuous homogeneous orbit families}, not six atomic effects,
and do not contradict the $137$-effect finite-outcome theorem
\ref{p2c:thm:v258-minimal-complete-readout}.  Nor do they select a
specific POVM, CP instrument, detector, state, physical preparation, or
outcome.  The tangent stabilizer $U(1)_{\rm diag}$ is distinct from the
retained source-Hopf subgroup $U(1)_{\rm H}\subset SU(2)_L$.
A physical realization requires a separately generated, same-action
apparatus/effect incidence and G0-typed occurrence.  The existing
Paper II-F countermodels and the universal GP3 whole-candidate
comparison obligation remain unchanged.
\end{remark}

\subsection{Converse-derived outcome thresholds and observable readout}
\label{p2c:sec:v266-converse-readout}

Throughout this subsection $R=\Rmom>0$, $r_0=\sqrt3R$,
$\mathcal H=\mathcal H_G\oplus\mathcal H_S$ with dimensions $11+4$,
$X_a^D=Y_a\oplus(r_0/4)\gamma_a$, and
$\widetilde\chi=0_G\oplus\chi$, $\chi=\gamma_1\gamma_2\gamma_3\gamma_4$.
The radius $r$ below is a bound on \emph{assigned classical outcome labels};
none of these finite constructions licenses a physical apparatus or a
calibration of $r$ as a spacetime distance.

\begin{theorem}[Sharp five-outcome radius and inverse quadratic fingerprint]
\label{p2c:thm:v266-five-threshold}
Let $X=(X_1,\ldots,X_4)$ be a Hermitian four-tuple with
$I,X_1,\ldots,X_4$ real-linearly independent, separately traceless
$X_a$, and direction-independent symmetric spectrum of $X(n)=\sum_an_aX_a$
with $\|X(n)\|=\kappa>0$ for every unit $n\in\mathbb R^4$.
An exact-first-moment POVM with at most five outcomes $x_i\in\overline
B^4(r)$ exists if and only if
\[
\boxed{r\ge4\kappa.}
\]
Five outcomes are necessary. At $r=4\kappa$ their positions necessarily
form a regular four-simplex, up to an orthogonal rotation and relabelling,
and the effects are uniquely
\[
E_i=\frac15\left(I+\kappa^{-1}X(n_i)\right),\qquad
n_i\cdot n_j=\frac{5\delta_{ij}-1}{4}.
\]
For Paper~II-C's scalar-matched compatible coordinate forest, if $Q_+$,
$Q_-$ denote the two three-dimensional irreducible-sector eigenvalues
of $Q=\sum_aX_a^2$, and $Q_S$ the protected-sector eigenvalue, then
\[
\kappa=\max\left\{\frac R2,\frac12\sqrt{Q_++Q_-},
\frac12\sqrt{Q_S}\right\},\qquad
\boxed{r_5^\star=2\max\{R,\sqrt{Q_++Q_-},\sqrt{Q_S}\}}.
\]
In particular $r_5^\star(X^D)=2R$; the alternative Paper~I-A cone
class also has $r_5^\star(X^{\rm cone})=2R$.
\end{theorem}
\begin{proof}
Independence of $I,X_1,\ldots,X_4$ forces at least five effects. For
exactly five, the $5\times5$ matrix with columns $(1,x_i/r)$ is invertible;
otherwise the five moment operators would be dependent. Its inverse yields
unique barycentric affine functions
$\ell_i(z)=w_i+g_i\cdot z/r$, $\ell_i(x_j)=\delta_{ij}$, and hence
$E_i=w_iI+r^{-1}g_i\cdot X$. Since $\operatorname{Tr}X_a=0$,
$w_i=\operatorname{Tr}E_i/\dim\mathcal H>0$ and $\sum_iw_i=1$.
The symmetric directional spectrum gives
$\lambda_{\min}(E_i)=w_i-\kappa|g_i|/r$.
At its own outcome,
$1=w_i+g_i\cdot(x_i/r)\le w_i+|g_i|$, so positivity implies
$r\ge\kappa(1-w_i)/w_i$. Some $w_i\le1/5$, proving $r\ge4\kappa$.
Conversely five regular-simplex unit vectors obey
$\sum_i n_i=0$ and $\sum_in_in_i^T=5I_4/4$.
Then $E_i=(I+(4/r)X(n_i))/5$ are positive for $r\ge4\kappa$,
normalize, and have the prescribed first moments. At equality, the preceding
bound for \emph{each} $i$ forces $w_i\ge1/5$ and thus $w_i=1/5$;
saturation gives $|x_i|=r$, $g_i=(4/5)n_i$, and the stated simplex
inner products. This proof does not presuppose that the effects lie in
the native block-diagonal algebra.

For the compatible forest, a directional coordinate acts by one
scalar--longitudinal-vector edge $R/2$, three transverse-vector--bivector
edges of common magnitude
$h=\sqrt{(r_+^2+r_-^2)/2}$, three geometric zero modes, and a
protected Clifford edge $\rho$. The respective spectra are
$\{\pm R/2,\pm h\ (\mathrm{each\ threefold}),0\ (\mathrm{threefold}),
\pm\rho\ (\mathrm{each\ twice})\}$.
The sector fingerprint gives $Q_+=2r_+^2$, $Q_-=2r_-^2$ and
$Q_S=4\rho^2$. Substitution proves the invariant radius formula.
At the intrinsic point $r_+=r_-=\sqrt2R/3$, $\rho=\sqrt3R/4$;
at the cone point all three equal $R/2$. Both have $\kappa=R/2$.
\end{proof}

\begin{theorem}[Six intrinsic outcomes on orthogonal trines]
\label{p2c:thm:v266-six-trines}
There is a native exact-first-moment six-outcome POVM at each
\[
\boxed{r\ge r_6:=\frac{\sqrt{34}}3R.}
\]
Its six unit directions are the two orthogonal planar trines
$n_j=(\cos\theta_j,\sin\theta_j,0,0)$ and
$m_j=(0,0,\cos\theta_j,\sin\theta_j)$,
$\theta_j=2\pi j/3$, with labels $x_v=rv$. In the geometric ordered
basis $(e_0,u_1,u_2,u_3,u_4,w_{12},w_{13},w_{14},w_{23},w_{24},w_{34})$,
set
\[
A=\operatorname{diag}\left(\frac12,\frac9{17},\frac9{17},
\frac8{17},\frac8{17},1,\frac12,\frac12,\frac12,\frac12,0\right),
\quad B=I_G-A.
\]
For $v=n_j$ take $A_v=A$, and for $v=m_j$ take $A_v=B$. Then
\[
\boxed{E_v=\frac13\left(A_v+\frac2rY(v)\right)\oplus
\frac16\left(I_S+\frac{r_0}{r}\gamma(v)\right).}
\]
For this \emph{fixed six-point support}, even unrestricted geometric
positive effect assignments cannot have $r<r_6$. The global, support-free
geometric effect-rank-six obstruction is proved in
Theorem~\ref{p2c:thm:universal-geometric-rank6-radius} below.
\end{theorem}
\begin{proof}
The two planar trines separately sum to zero and obey
$\sum_{j}n_jn_j^T=3P_{12}/2$, $\sum_{j}m_jm_j^T=3P_{34}/2$.
They imply $\sum_vE_v=I$ and
$\sum_vrv_aE_v=X_a^D$. The Clifford part is positive because
$r\ge r_6>r_0$. For the geometric part, natural planar rotations
reduce positivity to $A+2Y_1/r$ and $B+2Y_3/r$.
At $r=r_6$, the nontrivial $2\times2$ principal blocks of the former
have diagonal pairs $(1/2,9/17)$, $(9/17,1)$, and twice
$(8/17,1/2)$. Their squared off-diagonal entries are respectively
$9/34$, $4/17$, $4/17$, so their determinants are
$0$, $5/17$, $0$, $0$. Remaining diagonal entries are nonnegative.
The latter matrix is obtained by exchanging the coordinate planes.
Reducing the off-diagonal coefficients by increasing $r$ preserves
positivity.

For sharpness relative to this support, average any geometric solution
over the finite two-trine symmetry group (including the interchange of
planes). The averaged effects are still positive and have identical
prescribed moments; their three-effect resource sum in the first plane
is scalar on the six distinct form-sector representations. Plane
interchange and $A+B=I_G$ force its scalar and mixed-bivector entries
to $1/2$, and vector-plane entries to $a$ and $b=1-a$.
The trine moment equations then fix the first-plane unit-node effect
to $(A+2Y_1/r)/3$. Its scalar--vector and
transverse-vector--mixed-bivector principal minors require
$a\ge2R^2/r^2$ and $b\ge16R^2/(9r^2)$.
Therefore $1=a+b\ge34R^2/(9r^2)$. The displayed construction
attains equality. The averaging is a mathematical $O(4)$ action on
forms, not a claimed physical spinor control.
\end{proof}

\begin{lemma}[Positive two-amplitude Radon--exterior kernel]
\label{p2c:lem:universal-rank6-exterior-kernel}
Let $R=\Rmom$, let $V=\mathbb R^4$, and represent the geometric carrier
as $\mathcal H_G=\mathbb C\oplus\mathbb C^4\oplus\Lambda^2\mathbb C^4$
with native coordinate operators
\[
 Y(u)=\frac R2L_u+\frac{\sqrt2R}{3}(J_u+J_u^*),\qquad
 J_u(v)=u\wedge v,
\]
where $L_u$ exchanges scalar and vector components. For any $s,t\ge0$
and distinct unit $n,p\in S^3$, define
\[
 d=1-n\cdot p,\quad v=n-p,\quad k=n\wedge p,\qquad
 K_{s,t}(n,p)=\frac{1}{d}\big|(-sd,v,tk)\big\rangle
                      \big\langle(-sd,v,tk)\big|\succeq0.
\]
When $n=p$, assign the same positive matrix to both orientations,
\[
 K_{s,t}(n,n):=2\big|(0,v,-t\,n\wedge v)\big\rangle
                   \big\langle(0,v,-t\,n\wedge v)\big|,
\]
for any fixed unit $v\perp n$. Then for arbitrary finitely supported
probability measures $\mu,\nu$ on $S^3$ with the same barycenter,
\begin{equation}\label{p2c:eq:rank6-exterior-balance}
 \sum_{i,j}\mu_i\nu_j K_{s,t}(n_i,p_j)
 =\sum_{i,j}\mu_i\nu_j K_{s,t}(p_j,n_i).
\end{equation}
Moreover, for all $n,p$, including coincident points with the above
convention,
\begin{align}
 \operatorname{Tr}K_{s,t}(n,p)
 &=2+s^2+t^2+(t^2-s^2)n\cdot p,\label{p2c:eq:rank6-exterior-trace}\\
 \operatorname{Tr}[Y(u)K_{s,t}(n,p)]
 &=R\left[-\left(s+\frac{2\sqrt2}{3}t\right)u\cdot n
 +\left(s-\frac{2\sqrt2}{3}t\right)u\cdot p\right].
 \label{p2c:eq:rank6-exterior-moment}
\end{align}
\end{lemma}
\begin{proof}
For $n\ne p$, the definition for $K_{s,t}(p,n)$ is the rank-one
projector of $(sd,v,tk)/\sqrt d$, since $p-n=-v$ and
$p\wedge n=-k$. Subtraction cancels every block except the scalar--vector
and scalar--bivector cross blocks and equals
\[
 K_{s,t}(n,p)-K_{s,t}(p,n)
 =-2s\left(\lvert e_0\rangle\langle(0,v,tk)\rvert
           +\lvert(0,v,tk)\rangle\langle e_0\rvert\right).
\]
Averaging with $\mu_i\nu_j$ makes the vector part vanish by equality
of barycenters and makes the bivector part vanish because
$(\sum\mu_in_i)\wedge(\sum\nu_jp_j)=m\wedge m=0$.
For $n=p$ the paired matrices are equal by definition.
Using $\|n-p\|^2=2(1-n\cdot p)$ and
$\|n\wedge p\|^2=1-(n\cdot p)^2$ gives the trace identity.
The scalar--vector coupling contributes $-sR\,u\cdot(n-p)$;
the vector--bivector coupling contributes
$-(2\sqrt2Rt/3)u\cdot(n+p)$. Summing proves the moment identity.
For coincident points the chosen tangent $v$ has
$\|v\|=\|n\wedge v\|=1$ and the scalar component is zero;
therefore the trace is $2+2t^2$ and the moment is
$-4\sqrt2Rt\,u\cdot n/3$, exactly the limits of the two identities.
\end{proof}

\begin{theorem}[Sharp universal geometric effect-rank-six radius]
\label{p2c:thm:universal-geometric-rank6-radius}
Let $G_i\succeq0$ be an arbitrary finite positive-operator-valued measure
on the eleven-dimensional geometric carrier, with positions
$x_i\in\overline B^4(r)$, normalized and with exact first moments
\[
 \sum_iG_i=I_G,\qquad \sum_i x_{ia}G_i=Y_a\quad(a=1,\ldots,4).
\]
If $\dim_{\mathbb R}\operatorname{span}\{G_i\}\le6$, then
\[
 \boxed{r\ge r_6:=\frac{\sqrt{34}}3R.}
\]
The bound is attained at $r=r_6$ by the geometric compressions of
Theorem~\ref{p2c:thm:v266-six-trines}. In particular, no upper bound
on the number of labels, no distinctness hypothesis, no finite-covariance
condition, and no protected-Clifford support condition is needed.
\end{theorem}
\begin{proof}
Every interior point $x_i/r$ in the unit ball is the barycenter of
an antipodal pair of unit vectors. Splitting $G_i$ proportionally between
those labels preserves positivity, the operator span, normalization and
moments; hence assume $x_i=rn_i$ with $\|n_i\|=1$.
The independent operators $I_G,Y_1,\dots,Y_4$ belong to the real effect
span, so its dimension $g$ is at least five. If $g=5$, then
$G_i=\alpha_iI_G+Y(\beta_i)$. Because the minimum eigenvalue of
$Y(u)$ is $-R\|u\|/2$, positivity gives
$\alpha_i\ge(R/2)\|\beta_i\|$. The moment relations imply
$\sum_i n_i\cdot\beta_i=4/r$, whence
\[
 1=\sum_i\alpha_i\ge\frac R2\sum_i\|\beta_i\|
 \ge\frac{2R}{r},
\]
so $r\ge2R>r_6$.

For $g=6$ choose any extra real Hermitian generator $H$ and expand
$G_i=\alpha_iI_G+Y(\beta_i)+c_iH$.
The exact zeroth and first moments imply
\[
\begin{aligned}
 &\sum_i\alpha_i=1,\quad \sum_i\alpha_in_i=0,\quad \sum_i\beta_i=0,\\
 &\sum_i n_i\beta_i^{T}=I_4/r,\quad \sum_ic_i=0,\quad \sum_ic_in_i=0.
\end{aligned}
\]
The positive and negative coefficient sets
$P=\{i:c_i>0\}$, $M=\{j:c_j<0\}$ are nonempty. Put
\[
 C=\sum_{i\in P}c_i=-\sum_{j\in M}c_j>0,\quad
 \mu_i=c_i/C\ (i\in P),\quad \nu_j=-c_j/C\ (j\in M).
\]
Their barycenters agree, with common value $m$.
For arbitrary $s,t\ge0$ define positive dual matrices
\[
 D_i=\begin{cases}
 C^{-1}\sum_{j\in M}\nu_j K_{s,t}(n_i,n_j),&c_i\ge0,\\
 C^{-1}\sum_{j\in P}\mu_j K_{s,t}(n_i,n_j),&c_i<0.
 \end{cases}
\]
The first branch also handles zero coefficients. The kernel balance
identity~\eqref{p2c:eq:rank6-exterior-balance} implies
$\sum_ic_iD_i=0$ *as an operator*.
By the two kernel trace identities, for every $i$,
\begin{align*}
 \operatorname{Tr}D_i&=C^{-1}\left[2+s^2+t^2+(t^2-s^2)n_i\cdot m\right],\\
 \operatorname{Tr}[Y(u)D_i]&=\frac RC\left[
 -\left(s+\frac{2\sqrt2}{3}t\right)u\cdot n_i
 +\left(s-\frac{2\sqrt2}{3}t\right)u\cdot m\right].
\end{align*}
Since $\operatorname{Tr}(G_iD_i)\ge0$ for all $i$, the moment relations
cancel all terms involving $H$, $m$, and the unknown coefficients except
the fixed coordinate trace, giving the exact dual bound
\[
 0\le\sum_i\operatorname{Tr}(G_iD_i)
 =\frac1C\left[2+s^2+t^2
 -\frac{4R}{r}\left(s+\frac{2\sqrt2}{3}t\right)\right].
\]
The maximum of
$4(s+(2\sqrt2/3)t)/(2+s^2+t^2)$ over $s,t\ge0$ is
$\sqrt{34}/3$, attained at
$s=3\sqrt2/\sqrt{17}$, $t=4/\sqrt{17}$.
Consequently $r\ge\sqrt{34}R/3$. The six orthogonal-trines effects
attain this radius, and their geometric span is six-dimensional because
the rank-five case requires $r\ge2R$. This proves both necessity
and sufficiency. 
\end{proof}


\begin{corollary}[Globally sharp native intrinsic outcome counts]
\label{p2c:cor:v266-outcome-count}
For normalized exact intrinsic first moments at positions in
$\overline B^4(r)$, with unrestricted positive effects, the minimum
number of outcomes is
\[
\boxed{N_{\min}(r)=\begin{cases}
\infty,&0<r<\sqrt3R,\\
7,&\sqrt3R\le r<\sqrt{34}R/3,\\
6,&\sqrt{34}R/3\le r<2R,\\
5,&r\ge2R.
\end{cases}}
\]
The cone tuple instead has global exact-moment minimum radius $2R$,
attained with five outcomes. Its five-outcome threshold agrees with the
intrinsic tuple despite their unequal quadratic fingerprints.
\end{corollary}
\begin{proof}
The universal intrinsic protected trace witness from
Theorem~\ref{p2c:thm:optimal-intrinsic-povm-radius} gives $r\ge r_0$.
Theorem~\ref{p2c:thm:v266-five-threshold} excludes at most five
outcomes below $2R$ and attains five at or above it.
Theorem~\ref{p2c:thm:v266-six-trines} attains six at or above $r_6$.
Theorem~\ref{p2c:thm:universal-geometric-rank6-radius} excludes
all native POVMs with at most six nonzero effects for $r<r_6$:
their geometric compressions have real effect rank at most six.
The seven-outcome rule in
Theorem~\ref{p2c:thm:v255-seven-outcome-attainment} extends to all
$r\ge r_0$: if $w_i$ are its protected trace weights and
$\lambda=r_0/r$, then
$E_i(r)=\lambda E_i(r_0)+(1-\lambda)w_iI$ with labels $rn_i$
remains positive and retains every moment, since
$\sum_iw_in_i=0$. The cone protected radius witness gives
$r\ge4(R/2)=2R$, achieved by the simplex construction.
\end{proof}



% BEGIN PROPOSED GEOMETRIC POVM EXACTIFICATION (Paper II-C manuscript successor v2.177)
% These are exact mathematical statements; no physical instrument occurrence is asserted.
\begin{theorem}[Range-rank-preserving exact finite reduction of a general POVM]
\label{p2c:thm:range-rank-povm-exactification}
Let $X\subset\mathbb R^4$ be compact and let $G$ be any countably
additive positive-operator-valued Borel measure on a finite-dimensional
complex Hilbert space $\mathcal K$, with $G(X)=I_{\mathcal K}$.
Put
\[
 E=\operatorname{span}_{\mathbb R}\{G(B):B\subset X\text{ Borel}\},
 \qquad g=\dim_{\mathbb R} E.
\]
Then there are $k\le5g$ labels $x_1,\ldots,x_k\in X$ and positive
operators $G_j\in E$ such that
\[
 \sum_{j=1}^kG_j=I_{\mathcal K},\qquad
 \sum_{j=1}^k x_{ja}G_j=\int_X x_a\,dG(x)
 \quad(a=1,\ldots,4).
\]
Consequently the exact zeroth and four first moments of every
continuous-outcome native finite-matrix POVM have a finite positive
presentation with no increase in real effect-span dimension.  This
assertion concerns those five operator moments only; it neither
preserves arbitrary higher moments nor selects a physical detector,
state-update instrument or finite event record.
\end{theorem}
\begin{proof}
Write $d=\dim_{\mathbb C}\mathcal K$ and
$P(B)=d^{-1}\operatorname{Tr}G(B)$.  This is a probability measure.
Since $G(B)\succeq0$, each scalar matrix coefficient of $G$ is
absolutely continuous with respect to $P$.  Finite-dimensional
Radon--Nikodym theory therefore gives
\[
 G(B)=d\int_B F(x)\,dP(x),\quad
 F(x)\succeq0,\quad \operatorname{Tr}F(x)=1
 \quad(P\text{-almost everywhere}).
\]
Furthermore $F(x)\in E$ almost everywhere: apply each of a finite
basis of real linear functionals annihilating $E$ to this identity
for all Borel $B$.  Define the bounded measurable vector
\[
 Z(x)=(F(x),x_1F(x),\ldots,x_4F(x))\in E^5.
\]
Its closed essential range $K_Z$ is compact.  Indeed $X$ is compact
and all positive trace-one matrices form a compact set, while
$(x,F)\mapsto(F,x_1F,\ldots,x_4F)$ is continuous.  Every $Z\in K_Z$
has the form $(F,x_1F,\ldots,x_4F)$ for some $x\in X$ and positive
trace-one $F\in E$.  Also $K_Z$ lies in the affine hyperplane
$\operatorname{Tr}Z_0=1$ of the real vector space $E^5$, hence has
affine dimension at most $5g-1$.
The barycenter $\int Z\,dP$ lies in the convex hull of $K_Z$.
Carath\'eodory's theorem in that affine hull supplies $k\le5g$
points $Z_j=(F_j,x_{j1}F_j,\ldots,x_{j4}F_j)$ and nonnegative
weights $w_j$ summing to one with
$\int Z\,dP=\sum_j w_jZ_j$.  The effects $G_j=dw_jF_j$
are positive, belong to $E$, and satisfy every displayed identity.
\end{proof}

\begin{corollary}[Unconditional continuous-label intrinsic effect-rank thresholds]
\label{p2c:cor:continuous-povm-effect-rank-hierarchy}
On the full native intrinsic coordinate carrier, allow arbitrary
countably additive positive-operator-valued Borel measures supported in
$\overline B^4(r)$ whose four first moments equal $X_a^D$.  If $g$ is
the real dimension of their operator-value span, the sharp minimum is
\[
 \boxed{g_{\min}^{\mathrm{Bor}}(r)=
 \begin{cases}
 \infty, & 0<r<\sqrt3R,\\
 7, & \sqrt3R\le r<\sqrt{34}R/3,\\
 6, & \sqrt{34}R/3\le r<2R,\\
 5, & r\ge2R.
 \end{cases}}
\]
Here $R=\Rmom$ and $\infty$ means no exact POVM exists on that ball.
In particular, even infinitely many labels cannot reduce the native
real effect-span dimension below seven at the optimal radius.
\end{corollary}
\begin{proof}
The preceding theorem converts any proposed Borel POVM of real effect
rank $g$ into a finite POVM of rank at most $g$ with exactly the same
four coordinate moments.  The protected trace witness
(Theorem~\ref{p2c:thm:optimal-intrinsic-povm-radius}) excludes
$r<\sqrt3R$.  The finite five-effect-span threshold
(Theorem~\ref{p2c:thm:v266-five-threshold}), applied to real effect
spans, excludes $g\le5$ for $r<2R$; alternatively the $g=5$ argument
in Theorem~\ref{p2c:thm:universal-geometric-rank6-radius} gives this
on geometric compression.  The universal geometric
rank-six theorem excludes $g\le6$ for $r<\sqrt{34}R/3$ after
compressing to the eleven-dimensional geometric summand.
The finite native seven-, six- and five-outcome constructions in
Corollary~\ref{p2c:cor:v266-outcome-count} give the corresponding
upper bounds on $g$.  Since all four $X_a^D$ and the identity are
linearly independent, $g\ge5$ whenever an exact POVM exists.
\end{proof}

\begin{corollary}[Strict qualitative noise floor for rank-six measurements]
\label{p2c:cor:rank6-universal-noise-floor}
For every fixed $0<r<\sqrt{34}R/3$ there is a number
$\delta(r,R)>0$ such that every normalized native Borel POVM
supported in $\overline B^4(r)$ with real effect rank at most six
obeys
\[
 \max_{1\le a\le4}
 \left\|\int x_a\,dG(x)-X_a^D\right\|_{\mathrm{op}}
 \ge\delta(r,R).
\]
The bound is qualitative: no explicit value of $\delta$ is claimed.
For $r\ge\sqrt{34}R/3$ the best such residual is zero.
\end{corollary}
\begin{proof}
Theorem~\ref{p2c:thm:range-rank-povm-exactification}
reduces each admissible POVM to at most thirty positive effects,
without changing the residual or increasing its effect rank.
Pad with zero effects to obtain thirty labels.  All tuples of
positive operators $0\preceq G_j\preceq I$, with $\sum_jG_j=I$,
and labels $x_j\in\overline B^4(r)$ form a compact set.
The real-rank-at-most-six condition is closed, since every
seven-by-seven minor of their vectorized real Hermitian coefficient
matrix vanishes.  The maximum operator-norm moment residual is
continuous, so it attains a minimum on this compact set.  A zero
minimum would produce an exact rank-six native POVM at
$r<\sqrt{34}R/3$, contradicting the protected and geometric
lower bounds in
Corollary~\ref{p2c:cor:continuous-povm-effect-rank-hierarchy}.
Hence the minimum is strictly positive.
For $r\ge\sqrt{34}R/3$, the finite native six-outcome construction
is an exact admissible example, so the infimum is zero.
\end{proof}
% END PROPOSED GEOMETRIC POVM EXACTIFICATION


\begin{theorem}[Rigidity of every hypothetical six-outcome optimal intrinsic-position POVM]
\label{p2c:thm:vGC18-six-endpoint-rigidity}
Let $r_0=\sqrt3R_{\rm mom}$ and consider the intrinsic exact first-moment
POVM problem for $X_a^D=Y_a\oplus(r_0/4)\gamma_a$ on
$\mathcal H_G\oplus\mathcal H_S$, of dimensions $11+4$.
If a POVM with exactly six nonzero effects and outcome norms at most $r_0$
exists, then: (i) all six protected compressions are nonzero, and each
outcome equals $r_0n_i$ with $|n_i|=1$; (ii) its protected trace weights
satisfy $0<w_i<1/5$; (iii) up to orthogonal transport and relabelling there
is a real unit $t\in\mathbb R^6$ with $t_i\ne0$,
$w_i=(1-t_i^2)/5$, $\sum_i t_i\sqrt{w_i}=0$, and
\[
 n_i\cdot n_j=-\tfrac14\left(1+\frac{t_it_j}{\sqrt{w_iw_j}}\right),\quad i\ne j.
\]
The signs of $t$ partition the labels as $2+4$ or $3+3$. The protected
compressions are uniquely $E_i^S=2w_iP_+(n_i)$, while after native
central pinching every geometric compression is
\[
 E_i^G=w_i\left(I_G+\frac4{r_0}Y(n_i)\right)+t_i\sqrt{w_i}\,H
\]
for one common, necessarily indefinite Hermitian $11\times11$ matrix $H$.
Thus six-outcome existence is equivalent to positivity of these six
matrices for some such $(t,H)$. Equal protected weights $w_i=1/6$
are impossible at $r_0$. Every opposite-sign pair satisfies
\[\boxed{a_i^2+a_j^2-8a_ia_j+2\le0},\quad a_i=\sqrt{w_i}/|t_i|.\]
\end{theorem}
\begin{proof}
The existing saturated Clifford support theorem
\ref{p2c:thm:v226-equality} gives a spherical weighted two-design for
all protected-nonzero outcomes. Five protected-nonzero effects would
form a regular four-simplex. If the sixth geometric compression is
$H_6\succeq0$ and its label is $x_6$, simplex inversion forces
\[
 E_i^G=\frac15\left(I_G+\frac4{r_0}Y(n_i)\right)
 -\frac15\left(1+\frac4{r_0}n_i\cdot x_6\right)H_6,\quad 1\le i\le5.
\]
The five displayed subtraction coefficients sum to one, so at least one
is positive. The corresponding baseline matrix has a strictly negative
eigenvalue $(1-2/\sqrt3)/5$ from the native scalar--vector direction,
contradicting $E_i^G\succeq0$. Thus all six protected compressions
are nonzero.

Let $F$ be the $5\times6$ real frame with $i$th column
$\sqrt{w_i}(1,2n_i)^T$. The design gives $FF^T=I_5$; if $t$ is its
unit kernel vector, then $F^TF=I_6-tt^T$. The diagonal and off-diagonal
entries yield the stated weights and inner products. Every geometric
solution differs from the canonical moment solution
$w_i(I_G+4Y(n_i)/r_0)$ by $t_i\sqrt{w_i}H$, because the moment map has
one-dimensional kernel. If any $t_i=0$, that effect would equal its
nonpositive canonical term, so $t_i\ne0$ and $0<w_i<1/5$.

Analogously, every protected compression differs from
$2w_iP_+(n_i)$ by $t_i\sqrt{w_i}K$ with one Hermitian $K$.
The support saturation forces the range of $K$ into all six
$\operatorname{Ran}P_+(n_i)$, whose common intersection is zero:
for distinct $n,m$, $\gamma(n-m)^2=|n-m|^2I$.
Hence $K=0$. Pinching to the native central blocks preserves the
positive effects, their normalization, and their first moments.

The relation $\sum_i t_i\sqrt{w_i}=0$ requires both signs. Within any
one sign class, the Gram formula gives $n_i\cdot n_j<-1/4$.
Five such vectors would have negative squared norm after summing them.
Thus the only sign cardinalities are $2+4$ and $3+3$.
Opposite signs of the nonzero correction coefficients and the negative
eigenspaces of every canonical geometric term force the common $H$
to have both positive and negative expectations. Conversely any
admissible $(t,H)$ satisfying the displayed semidefinite inequalities
produces a valid centrally pinched POVM: complete the orthonormal
pair $\sqrt w,t$ in $\mathbb R^6$ by four orthonormal rows and divide
the remaining frame columns by $2\sqrt{w_i}$.

A further exact necessary condition is available without selecting the
common Hermitian correction. For any opposite signs $t_i>0>t_j$, put
$a_i=\sqrt{w_i}/|t_i|$ and $a_j=\sqrt{w_j}/|t_j|$.
Dividing the two geometric positive effects by their positive
correction coefficients and adding cancels $H$, giving
$a_i(I_G+4Y(n_i)/r_0)+a_j(I_G+4Y(n_j)/r_0)\succeq0$.
The direction-independent geometric norm $\|Y(n)\|=R_{\rm mom}|n|/2$
and the displayed Gram identity therefore imply
\[
\boxed{a_i^2+a_j^2-8a_ia_j+2\le0}\quad(t_it_j<0).
\]
This is necessary but not sufficient for the common-matrix semidefinite
problem.

Finally $w_i=1/6$ forces $t_i=\pm1/\sqrt6$ with three signs of each kind.
The Gram formula then yields two orthogonal planar regular trines,
whose fixed-support sharp radius is $\sqrt{34}R_{\rm mom}/3>r_0$
by Theorem~\ref{p2c:thm:v266-six-trines}. This is impossible at $r_0$.
\end{proof}
\begin{remark}[Conditional six-outcome rigidity and its resolved feasibility]
The preceding theorem gives necessary conditions for a hypothetical
six-outcome endpoint solution; these are retained as structural
identities. The support-free rank obstruction in
Theorem~\ref{p2c:thm:universal-geometric-rank6-radius} now excludes
all such effects at $r_0<r_6$, including asymmetric configurations.
No occurrence of a physical detector, instrument, or outcome follows.
\end{remark}

\begin{theorem}[Sharp antipodal--tetrahedral six-outcome radius]
\label{p2c:thm:gc19-antipodal-tetra}
Let $R=\Rmom>0$, and let $X^D=Y\oplus (\sqrt3 R/4)\gamma$
be the native intrinsic four-position tuple on
$\mathcal H_G\oplus\mathcal H_S$ of Paper~II-C. Fix six unit directions
$\mathcal Q=\{e_1,-e_1,q_1,q_2,q_3,q_4\}\subset\mathbb R^4$,
where the $q_j$ form a regular tetrahedron in $e_1^\perp$:
$\sum_jq_j=0$, $q_i\cdot q_j=-1/3$ for $i\ne j$.
Allow arbitrary positive effects $E_q$ on the full carrier, with
$\sum_qE_q=I$ and $\sum_q(rq_a)E_q=X_a^D$ for $a=1,\ldots,4$.
Then such a POVM exists precisely when
\[
\boxed{r\ge r_{\mathrm{tet}}
 =\frac{\sqrt{85+\sqrt{2617}}}{6}\,R.}
\]
In particular this fixed support does not attain the intrinsic optimal
radius $\sqrt3R$ and is strictly less economical in radius than the
previously established two-orthogonal-trines six-outcome rule:
$r_{\mathrm{tet}}>\sqrt{34}R/3$.
\end{theorem}
\begin{proof}
By central pinching it suffices to analyze the positive geometric effects
$G_q=P_GE_qP_G$ independently. Put $s=(r/R)^2$. The geometric tuple is
$O(4)$-equivariant in the native form basis
$\mathbb R\oplus\mathbb R^4\oplus\Lambda^2\mathbb R^4$.
Average any feasible geometric POVM over reflection $e_1\mapsto-e_1$
and the full tetrahedral symmetry group $S_4$ on $e_1^\perp$.
This averaging preserves positivity, normalization and exact moments.
The total resource $A=G_{e_1}+G_{-e_1}$ is invariant under this group.
The five mutually inequivalent real representation blocks are
\[
\langle e_0\rangle\oplus\langle u_1\rangle
\oplus V_3\oplus(e_1\wedge V_3)\oplus\Lambda^2V_3
\quad(1+1+3+3+3=11),
\]
so the averaged resource is
\[
A=\operatorname{diag}(a_0,a_1,a_VI_3,a_MI_3,a_WI_3),\quad
0\le a_\bullet\le1\quad(\bullet=0,1,V,M,W).
\]
Put $B=I_G-A$. The moment equations uniquely determine
\[
G_{\pm e_1}=\tfrac12(A\pm Y(e_1)/r),\qquad
G_{q_j}=\tfrac14(B+3Y(q_j)/r).
\]
The native scalar--vector and vector--two-form coefficients are $R/2$
and $\sqrt2R/3$, respectively. Each indicated operator is positive
if and only if its $2\times2$ coupled blocks are positive, equivalently
\begin{align*}
a_0a_1&\ge\frac1{4s}, &a_Va_M&\ge\frac2{9s},\\
(1-a_0)(1-a_V)&\ge\frac9{4s},
 &(1-a_1)(1-a_M)&\ge\frac2s,\\
(1-a_V)(1-a_W)&\ge\frac2s.
\end{align*}
For necessity, the first and third inequalities give
$a_0\ge1/(4s)$ and $1-a_0\ge9/(4s)$, so $s\ge5/2$.
At equality the required $a_1=1$ contradicts the fourth inequality;
thus $s>5/2$. Set $v=a_V$. To make the fourth inequality easiest to
satisfy, the first three bounds force
\[
a_0\le1-\frac9{4s(1-v)},\quad
 a_1\ge\frac1{4sa_0},\quad a_M\ge\frac2{9sv}.
\]
Their combination with the fourth inequality requires
\[
\left(1-\frac1{4s\left(1-\frac9{4s(1-v)}\right)}\right)
\left(1-\frac2{9sv}\right)-\frac2s\ge0.
\]
The positive denominator is $9sv\,[4s(1-v)-9]$ and its numerator is
\[
N_s(v)=(-36s^2+81s)v^2+(36s^2-154s+160)v+20-8s.
\]
The coefficient of $v^2$ is negative for $s>5/2$, and
\[
\operatorname{disc}_v N_s
 =4(2s-5)(9s-20)(18s^2-85s+64).
\]
For $5/2<s<(85+\sqrt{2617})/36$ this discriminant is
negative, hence $N_s(v)<0$ for every real $v$, a contradiction.

Conversely set $s_*=(85+\sqrt{2617})/36$ and choose the exact resource
\[
\boxed{
\begin{aligned}
a_0&=\frac{\sqrt{2617}-21}{128},&
a_1&=\frac{\sqrt{2617}-13}{136},\\
a_V&=\frac{\sqrt{2617}-21}{136},&
a_M&=\frac{\sqrt{2617}-13}{144},&a_W&=0.
\end{aligned}}
\]
All entries lie in $[0,1]$. The first four determinant inequalities
hold with equality and the fifth holds strictly. The displayed $G_q$
are therefore positive, normalized and have exact geometric first moments
at $r=\sqrt{s_*}R$. For the protected block use weights
$w_{\pm e_1}=1/8$ and $w_{q_j}=3/16$, which obey
$\sum_qw_q=1$, $\sum_qw_qq=0$,
$\sum_qw_qqq^T=I_4/4$. Put
\[
S_q=w_q\bigl(I_S+(\sqrt3R/r)\gamma(q)\bigr).
\]
These effects are positive because $r\ge\sqrt3R$, and they normalize
and reproduce $(\sqrt3R/4)\gamma_a$. Thus $E_q=G_q\oplus S_q$
attain all target moments. For $r>\sqrt{s_*}R$ the same resource $A$
remains feasible because the off-diagonal couplings decrease with $1/r$;
the same protected formula remains positive. This proves necessity and
attainment. Finally $s_*-34/9=(\sqrt{2617}-51)/36>0$.
\end{proof}

\begin{corollary}[Elimination of all two-level endpoint sign classes]
\label{p2c:cor:gc19-two-level-no-go}
At $r_0=\sqrt3R$ no six-outcome exact-moment POVM can have a Naimark
complement $t$ whose strictly positive entries share one magnitude and
whose strictly negative entries share another magnitude. Here $t$ is
as in the endpoint reduction of Paper00-GraphClosureDiscovery v1.8,
$w_i=(1-t_i^2)/5$ and
$\sum_it_i\sqrt{w_i}=0$.
\end{corollary}
\begin{proof}
The existing endpoint reduction allows only sign partitions $3+3$
and $2+4$ (up to global sign). In the $3+3$ case, the norm and
centering equations force both magnitudes to be $1/\sqrt6$;
the corresponding support consists of two orthogonal trines, already
excluded at $r_0$ by Theorem~\ref{p2c:thm:v266-six-trines}.
In the $2+4$ case, write $t=(p,p,-q,-q,-q,-q)$, $p,q>0$.
The two equations
$2p^2+4q^2=1$ and $p\sqrt{1-p^2}=2q\sqrt{1-q^2}$
give $p^2=3/8$, $q^2=1/16$.
The Naimark Gram formula
$n_i\cdot n_j=-(1+t_it_j/\sqrt{w_iw_j})/4$ for $i\ne j$
then forces an antipodal pair orthogonal to a regular tetrahedron.
The preceding theorem excludes this support at $r_0$.
\end{proof}

\noindent\emph{Claim boundary.} This is an unconditional mathematical
fixed-support threshold and a no-go for two-level endpoint complements;
this fixed-support result alone would not settle unrestricted asymmetric
six-versus-seven feasibility.  The independent support-free converse,
Theorem~\ref{p2c:thm:universal-geometric-rank6-radius}, does settle it:
at $r_0=\sqrt3R<\sqrt{34}R/3$ at least seven outcomes are necessary.
It does not assert a physically occurring measurement or apparatus.

\begin{theorem}[Antipodal-free six-outcome endpoint]
\label{p2c:thm:gc20-antipodal-free-six}
Let $X^D_a=Y_a\oplus(\sqrt3R/4)\gamma_a$ be the native intrinsic
position tuple of Paper~II-C, $R=\Rmom>0$.
There is no normalized six-outcome positive-operator-valued measure,
with exact first moments $\sum_i x_{i,a}E_i=X^D_a$ and labels
$\lvert x_i\rvert\le\sqrt3R$, whose six protected-nonzero labels
include any antipodal pair. The conclusion permits arbitrary positive
geometric effects, arbitrary noncentral off-diagonal effect entries,
unequal protected weights and arbitrary orientation of the antipodal pair.
Moreover, every protected weighted spherical two-design of six points
in $S^3$ with an antipodal pair is, up to rotation and relabelling,
necessarily a latitude tetrahedral design
\[
 n_+=e_1,\quad n_-=-e_1,\quad
 n_j=z e_1+\sqrt{1-z^2}\,q_j\quad(j=1,\dots,4),
 \qquad -\tfrac14<z<\tfrac14,
\]
where $q_i\cdot q_j=-1/3$ for $i\ne j$ and the weights are uniquely
\[
 w_+=\frac{1-4z}{8(1-z)},\qquad
 w_-=\frac{1+4z}{8(1+z)},\qquad
 w_j=\frac3{16(1-z^2)}.
\]
At $z=0$ this includes the symmetric antipodal--tetrahedral support;
the entire unequal-weight tilted family is now excluded at the endpoint.
\end{theorem}
\begin{proof}
The protected equality theorem
\ref{p2c:thm:v226-equality} states that all labels with positive protected
compression lie on the sphere and their $w_i>0$, $n_i=x_i/(\sqrt3R)$
obey $\sum w_i=1$, $\sum w_i n_i=0$ and
$\sum w_i n_i n_i^T=I_4/4$. If one effect has zero protected
compression, the proposed six-effect construction reduces to five
protected nonzero labels, whose geometric extension is excluded by the
separately supplied five-plus-one endpoint argument; the latter follows
by regular-simplex inversion and the negative eigenvalue of
$(I_G+4Y(n)/(\sqrt3R))/5$. It therefore suffices to handle six
strictly positive protected weights. Central pinching permits
$E_i=G_i\oplus S_i$, with $G_i\succeq0$ and the geometric moments
unchanged.

Set $R=1$ in the remaining algebra. The $5\times6$ real matrix with
columns $\sqrt{w_i}(1,2n_i)^T$ has orthonormal rows. Its column Gram
matrix is $I_6-tt^T$ for a unit one-dimensional Naimark complement
$t\in\mathbb R^6$. Choose $n_+=e_1$ and $n_-=-e_1$. The $+,-$
Gram entry gives $t_+t_-=3\sqrt{w_+w_-}>0$.
For any remaining label $n_j$, writing $z_j=e_1\cdot n_j$, its two
Gram entries give
\[
 -t_+t_j=\sqrt{w_+w_j}(1+4z_j),\qquad
 -t_-t_j=\sqrt{w_-w_j}(1-4z_j).
\]
No $t_j$ can vanish, since then $z_j$ would equal both $-1/4$
and $+1/4$. Since $t_+t_->0$, these identities imply that
$(1+4z_j)/(1-4z_j)$ is independent of $j$. Thus $z_j=z$ for all
four residual labels. Put $s=\sqrt{1-z^2}$ and
$n_j=z e_1+s q_j$ with $q_j\in S^2$. Their transverse moments are
$\sum_j w_jq_j=0$, $\sum_jw_j q_j q_j^T=I_3/(4s^2)$.
The $3\times4$ frame with columns $2s\sqrt{w_j}\,q_j$ has
orthonormal rows and normalized null vector
$(\sqrt{w_j})_j/\sqrt{W}$, $W=\sum_jw_j$. Its diagonal column
Gram equations yield $4s^2w_j=1-w_j/W$, hence all four $w_j$
are equal. The same Gram relation gives
$w_j=3/(16s^2)$ and $q_i\cdot q_j=-1/3$ for $i\ne j$.
The axial zeroth, first and second moments now give the displayed
$w_\pm$ and $-1/4<z<1/4$.

It remains to prove that no positive geometric effects extend this
protected design at $r=\sqrt3$ (not merely that a particular ansatz
fails). Write $A=G_+$, $B=G_-$ and $C_j=G_j$, and let $S=A+B$.
The exact moment equations invert uniquely to
\[
 \begin{aligned}
 2A&=(1+z)S-zI_G+Y(e_1)/r,\\
 2B&=(1-z)S+zI_G-Y(e_1)/r,\\
 4C_j&=I_G-S+\frac3{rs}Y(q_j).
 \end{aligned}
\]
Since $A,B,C_j\succeq0$ and sum to $I_G$, one has
$0\preceq S\preceq I_G$.
Average over the tetrahedral rotation group on the transverse
three-space, simultaneously permuting the $C_j$. This does not assume
a physical covariance law: it replaces a hypothetical positive
solution by another positive solution of the same finite moment
identities. The scalar and longitudinal vector diagonal entries of $S$
are $a,b$, and the diagonal entries on each of the transverse-vector
and mixed-bivector irreducible three-spaces are $v,m$, with
$0\le a,b,v,m\le1$. Possible invariant off-diagonal blocks need
not vanish.

In the tetra effect $4C_j$, the native $Y(q_j)$ scalar--transverse
vector coefficient is $1/2$; the longitudinal vector--mixed
bivector coefficient is $\sqrt2/3$. Positivity of these two
principal $2\times2$ blocks forces
\[
 (1-a)(1-v)\ge\frac9{4r^2s^2},\qquad
 (1-b)(1-m)\ge\frac2{r^2s^2}.
\]
At $r^2=3$, using $v,m\ge0$, one obtains
\[
 a\le a_*:=\frac{1-4z^2}{4(1-z^2)},\qquad
 b\le b_*:=\frac{1-3z^2}{3(1-z^2)}.
\]
Write $h=\langle e_0,S u_1\rangle\in\mathbb C$ and
$c=1/(2r)$. The two pair effects have scalar--longitudinal-vector
principal blocks whose positivity gives, with
$U=(1+z)a-z$, $U'=(1+z)b-z$,
$V=(1-z)a+z$, $V'=(1-z)b+z$,
\[
 |(1+z)h+c|\le\sqrt{UU'},\qquad
 |(1-z)h-c|\le\sqrt{VV'}.
\]
Eliminating $h$ by the triangle inequality and applying weighted
Cauchy--Schwarz yields
\begin{align*}
 \frac1r&\le(1-z)\sqrt{UU'}+(1+z)\sqrt{VV'}\\
 &\le 2\sqrt{[(1-z^2)a+z^2][(1-z^2)b+z^2]}\\
 &\le 2\sqrt{\tfrac14\cdot\tfrac13}=\frac1{\sqrt3}=\frac1r.
\end{align*}
All inequalities must be equalities. Equality in the last bound
forces $a=a_*$ and $b=b_*$, which satisfy
$b_*-a_*=1/[12(1-z^2)]>0$.
Equality in weighted Cauchy--Schwarz requires $UV'=VU'$; however
$UV'-VU'=2z(a-b)$. Thus $z=0$.
At $z=0$, equality forces $a=1/4$, $b=1/3$. The two tetra
minor inequalities then force $v=m=0$.
As $S\succeq0$, each corresponding transverse-vector and
mixed-bivector row/column of $S$ vanishes. But
$\langle u_2,Y(e_1)w_{12}\rangle=\sqrt2/3\ne0$, so the pair
$2A=S+Y(e_1)/r$ contains a $2\times2$ principal block with
zero diagonal and nonzero off-diagonal, contradicting $A\succeq0$.
The contradiction is independent of the choice of reference $e_1$
and of every noninvariant effect: averaging was a necessary operation
on any hypothetical feasible solution.
\end{proof}

\begin{corollary}[Exclusion of a fourfold Naimark sign-class degeneracy]
\label{p2c:cor:gc20-fourfold-no-go}
In the six-column endpoint Parseval parameterization with
$F^TF=I_6-tt^T$, no feasible six-effect endpoint rule has four
Naimark entries identical (in particular, four same-sign components
of equal magnitude). Equivalently, no surviving candidate belongs
to the $2+4$ sign partition with four equal magnitudes in its
four-member class, even if the two opposite-sign magnitudes differ.
\end{corollary}
\begin{proof}
Four equal $t_j$ give four equal weights
$w_j=(1-t_j^2)/5$ and constant mutual directional inner product
$n_i\cdot n_j=-(1+t_i t_j/\sqrt{w_iw_j})/4$.
Their three differences span a three-dimensional subspace of
$\mathbb R^4$. Each of the two remaining directions has constant
inner product with the four equal-parameter directions and so is
orthogonal to all three differences. Therefore both lie on the
remaining one-dimensional axis. They cannot coincide: merging equal
outcome labels would produce a five-effect exact endpoint rule,
excluded by Corollary~\ref{p2c:cor:v226-five}. Hence they are
antipodal, excluded by the theorem above.
\end{proof}

\begin{corollary}[Uniform angular gap for hypothetical six-outcome rules]
\label{p2c:cor:gc20-angular-gap}
If six-outcome exact-position POVMs exist at $r_0=\sqrt3R$, their
outcome directions necessarily admit a universal uniform pairwise
angular separation from both coincidence and antipodality: there
exists $\varepsilon>0$ such that
$-1+\varepsilon\le n_i\cdot n_j\le1-\varepsilon$ for every
pair $i\ne j$ in every such rule.
This existential constant is not numerically determined by the proof.
\end{corollary}
\begin{proof}
The space of six effects $0\preceq E_i\preceq I$ and six labels in
the closed sphere $S^3(r_0)$ is compact in finite dimensions.
Normalization and exact operator first-moment identities define a
closed subset. The protected equality theorem and the previous
five-plus-one exclusion guarantee that every feasible limit retains
positive protected weight at every label. A coincident pair can be
merged into five effects, impossible by the five-outcome lower bound;
an antipodal pair is impossible by the theorem above. The continuous
minimum pairwise distances from these two closed degeneracies therefore
have strictly positive minima on the compact feasible set. If this
set is empty, the conditional statement is vacuously true.
\end{proof}


\begin{theorem}[Sharp chirality information with seven or six labels]
\label{p2c:thm:v266-sharp-chirality-support}
Let $\beta(r)=\sqrt{1-r_0^2/r^2}$. The sharp bounds in
Theorem~\ref{p2c:thm:v226-chirality-radius},
$D_{\mathrm{TV}}\le\beta$ and $\mathcal I_\chi(0)\le\beta^2$,
are attained by a \emph{seven-outcome} exact-position POVM for every
$r\ge r_0$ and a \emph{six-outcome} exact-position POVM for every
$r\ge r_6$. No five-outcome exact-position POVM can distinguish
$\rho_+=(I_S+\chi)/4$ from $\rho_-=(I_S-\chi)/4$, at any radius.
\end{theorem}
\begin{proof}
Take the geometric endpoint effects $G_i^{(0)}$ from
Theorem~\ref{p2c:thm:v255-seven-outcome-attainment}, with four axial
weights $w_i=1/8$ and three trine weights $w_i=1/6$. Give the four
axial nodes signs $s_i=+1$ and the three trine nodes $s_i=-1$;
$\sum_iw_is_i=\sum_iw_is_in_i=0$. For $\lambda=r_0/r$ put
\[
E_i(r)=\bigl[\lambda G_i^{(0)}+(1-\lambda)w_iI_G\bigr]
\oplus\bigl[w_i(I_S+\lambda\gamma(n_i)+s_i\beta\chi)\bigr].
\]
Geometric positivity is convex; protected positivity follows from
$\{\chi,\gamma(n_i)\}=0$ and $\lambda^2+\beta^2=1$.
The weighted two-design proves normalization and all four first moments.
For $\rho_s=(I_S+s\chi)/4$, the output probabilities are
$p_i(s)=w_i(1+s_is\beta)$; therefore
$D_{\mathrm{TV}}=\beta$ and $\mathcal I_\chi(0)=\beta^2$.
For six outcomes use the geometric effects of
Theorem~\ref{p2c:thm:v266-six-trines} and replace their protected parts by
\[
S_v=\frac16\bigl[I_S+\lambda\gamma(v)+t_v\beta\chi\bigr],
\qquad t_v=+1\text{ on the first trine},\quad -1\text{ on the second}.
\]
The same Clifford positivity and two-trine moments apply; now
$p_v(s)=(1+t_vs\beta)/6$. The five-outcome no-go follows from
Theorem~\ref{p2c:thm:v266-minimal-protected-readout} below.
The six-label sharpness claim in the $r\ge r_6$ range is a minimum
\emph{among POVMs conveying positive sharp chirality information};
at $r=r_0$ the sharp upper bound itself vanishes.
\end{proof}

\begin{theorem}[Minimal direct native protected readout]
\label{p2c:thm:v266-minimal-protected-readout}
A five-outcome exact-position POVM, even when its effects are initially
allowed in all of $M_{15}(\mathbb C)$, has precisely the real test span
$\operatorname{span}_{\mathbb R}\{I,X_1^D,\ldots,X_4^D\}$.
Consequently it reads neither $P_S$ nor $\widetilde\chi$ directly.
For every $r\ge2R$ there are exact-position native POVMs with:
(i) six outcomes directly reconstructing \emph{either} $P_S$ or
$\widetilde\chi$ together with the four coordinate moments;
(ii) seven outcomes directly and independently reconstructing
\emph{both}. These counts are minimal for the respective operator
systems. At such radii a five-outcome readout is intrinsically blind
to protected chirality regardless of the existence of coordinate words
representing $\chi$ and $P_S$.
\end{theorem}
\begin{proof}
For five effects the invertible affine moment matrix in the proof of
Theorem~\ref{p2c:thm:v266-five-threshold} forces each effect to lie
in the stated five-dimensional span. The five moment operators
also lie in the effect span, giving equality. Their restrictions on
$\mathcal H_S$ are scalar and Clifford-vector terms. Orthogonality of
$\chi$ to $I_S$ and the $\gamma_a$ excludes $\widetilde\chi$;
separate block traces exclude $P_S$.

Let $v_i$ be the two orthogonal trines of
Theorem~\ref{p2c:thm:v266-six-trines}, with $t_i=+1$ on the first
and $-1$ on the second, and let $\lambda=r_0/r$, $\beta=\sqrt{1-\lambda^2}$.
Then at $r\ge2R$ the following are normalized exact-position POVMs:
\[
E_i^{(\chi)}=\frac16\left[I+\frac4rX^D(v_i)+t_i\beta\widetilde\chi\right],
\quad
E_i^{(P)}=\frac16\left[I+\frac4rX^D(v_i)+t_i\delta P_S\right],
\quad 0<\delta\le1-\lambda.
\]
Their contrasts satisfy
$\sum_it_iE_i^{(\chi)}=\beta\widetilde\chi$ and
$\sum_it_iE_i^{(P)}=\delta P_S$.
On the geometric block all effects are positive because
$\|Y(v_i)\|=R/2$; on the protected block, respectively,
$(\lambda\gamma(v_i)+t_i\beta\chi)^2=I_S$ and
$\lambda+\delta\le1$. The two-trine first and second moments
prove exact normalization and coordinates.

For joint access use the seven endpoint nodes and weights $w_i$
from Theorem~\ref{p2c:thm:v255-seven-outcome-attainment}, ordered as
$(+e_1,-e_1,+e_2,-e_2,q_0,q_1,q_2)$, and let
\[
a=(1,1,-1,-1,0,0,0),\qquad b=(1,1,1,1,-1,-1,-1),
\quad\delta=\eta=\frac{1-\lambda}{4}>0.
\]
The weighted contrasts have zero scalar and vector moments, are
mutually orthogonal in the weighted scalar pairing, and satisfy
$\sum_iw_ia_i^2=1/2$, $\sum_iw_ib_i^2=1$.
The seven positive effects
\[
\boxed{E_i=w_i\left[I+\frac4rX^D(n_i)+
\delta a_iP_S+\eta b_i\widetilde\chi\right]}
\]
have exact first moments and
$\sum_ia_iE_i=(\delta/2)P_S$,
$\sum_ib_iE_i=\eta\widetilde\chi$.
Geometric positivity follows from $r\ge2R$. On the protected block
the lowest eigenvalue is at least
$w_i(1-\delta-\lambda-\eta)=\frac12w_i(1-\lambda)>0$.
The seven target Hermitian operators are independent, so no POVM with
fewer than seven effects can span all of them; the analogous six-target
independence proves six-effect necessity for either individual extension.
This result establishes finite operator accessibility, not an actual
source state or selected detector.
\end{proof}

\begin{proposition}[Two-shot quadratic fingerprint and algebra/span distinction]
\label{p2c:prop:v266-repeat-fingerprint}
At $r\ge2R$, let $n_1,\ldots,n_5$ be any regular four-simplex and set
$E_i=(I+(4/r)X^D(n_i))/5$. They obey
\[
\boxed{\sum_{i=1}^5E_i^2=\frac15I+\frac4{5r^2}Q,\qquad
Q=\sum_{a=1}^4(X_a^D)^2.}
\]
For the specific repeated \emph{L\"uders} instrument
$K_i=\sqrt{E_i}$, the probability that the second result equals the
first is $1/5+4\langle Q\rangle/(5r^2)$. At $r=2R$ this is $7/20$
for the intrinsic normalized protected maximally mixed state, versus
$2/5$ for the cone-class protected state with the same nominal radius.
The one-shot simplex probabilities for these two benchmark states are
both $1/5$. For any \emph{native} exact-position POVM, whatever its
number of effects, $C^*(E_i)=C^*(X_a^D)=M_{11}\oplus M_4$; this
algebraic equality does not make a five-effect one-shot experiment
informationally complete.
\end{proposition}
\begin{proof}
Use $\sum_in_i=0$ and $\sum_in_in_i^T=(5/4)I_4$ to expand the squares.
For the specified L\"uders update the two-step return event has effect
$\sum_i\sqrt{E_i}E_i\sqrt{E_i}=\sum_iE_i^2$.
The protected quadratic eigenvalue is $3R^2/4$ for the intrinsic tuple,
$R^2$ for the cone tuple. At $r=2R$ these give the stated probabilities.
The comparison assumes a fixed known preparation and the specified
repeated instrument; unknown preparations need not be distinguished.
For a native exact-moment POVM, each $E_i\in C^*(X_a^D)$ by hypothesis,
while $X_a^D=\sum_ix_{ia}E_i$, hence both generated $C^*$-algebras coincide.
In contrast, one-setting statistics depend only on the \emph{real linear}
span of effects, of dimension five for the simplex, not on their products.
\end{proof}

\begin{remark}[Remaining finite and physical obligations]
\label{p2c:rem:v266-readout-firewall}
The support-free sharp converse in
Theorem~\ref{p2c:thm:universal-geometric-rank6-radius}
has resolved this interval: the unrestricted minimum is seven;
see Corollary~\ref{p2c:cor:v266-outcome-count}. The earlier
six-effect protected direct-readout ambiguity below $2R$ is excluded
by the sharp dimension-and-positivity converse of
Theorem~\ref{p2c:thm:v353-sharp-protected-readout} (integrated below);
more than six outcomes may still admit other tests.
The quadratic contrast is an identity for a \emph{specified} sequential
L\"uders instrument and not an instrument-independent two-event law.
Nothing here selects an action-generated apparatus, state preparation,
charged source, detector record, or physical distance calibration;
Paper~II-A's consumer/occurrence firewall remains in force.
\end{remark}


\subsection{Central offsets, state weights and complete readout}
\label{p2c:sec:v227-central-readout}
For the intrinsic algebra $M_{11}(\mathbb C)\oplus M_4(\mathbb C)$,
Theorem~\ref{p1b:thm:v227-central-readouts} identifies precisely the two
block-scalar Hamiltonian offsets invisible to its even Heisenberg dynamics
at fixed density/support. Fixed sectorwise canonical states have the same
invariance. Neither result determines their central probabilities, and
neither identifies the protected chirality with a central observable.

\begin{corollary}[Symmetry of a supplied canonical readout state]
\label{p2c:cor:v227-canonical-covariance}
If an actual readout Hamiltonian is equivariant, its normalized full-ground
or Gibbs density is invariant. At fixed occupied-block weights the same
holds blockwise, without a gap between the scalar ground energies of distinct
blocks. An arbitrary pure orientation in a degenerate ground space remains
a different state and need not be invariant.
\end{corollary}
\begin{proof}
Use Theorem~\ref{p1b:thm:v227-state-symmetry} and
Corollary~\ref{p1b:cor:v227-block-covariance}. The latter consumes only
within-block spectral projectors and fixed probabilities.
\end{proof}
The 137-coefficient reconstruction and the optimal-radius chirality blindness
of the preceding closure keep their original measurement domains. Symmetry
of the density does not supply an informationally complete endpoint POVM
or select a primordial population.




\section{Sharp endpoint readout bounds}
\label{p2c:sec:v355-sharp-readout-resources}
The theorems below use only the fixed intrinsic finite coordinate tuple,
its native algebra $\mathcal A_D=M_{11}(\mathbb C)\oplus M_4(\mathbb C)$,
and the already proved finite exact-position existence and support statements.
They are mathematical capability and no-go statements, not assertions that a
particular apparatus, preparation, measurement update or physical source is
selected by G0 or GP2. Outcome labels may include passive tags at equal
position, as in the existing exact-position readout convention.
The 16-dimensional protected Hermitian test space is not a new carrier.


\begin{theorem}[Sharp minimum-outcome protected direct readout]
\label{p2c:thm:v353-sharp-protected-readout}
Write $R=\Rmom$, $\mathcal H=\mathcal H_G\oplus\mathcal H_S$
with $\dim\mathcal H_G=11$ and $\dim\mathcal H_S=4$.
Let $T_1,\ldots,T_k$ be any linearly independent Hermitian operators
supported on $\mathcal H_S$, where $1\le k\le 16$.
A normalized exact-first-moment POVM on outcome positions
$\|x_i\|\le r$ whose real effect span contains
$I$, all four $X^D_a$, and every $T_j$ needs at least $5+k$
effects. With exactly $5+k$ effects, such a POVM exists if and only if
\[\boxed{r\ge 2R}.\]
Distinct effects may carry the same assigned position and distinguishable
passive labels. At $r=2R$, a $21$-effect POVM directly reconstructing the
full Hermitian protected $M_4(\mathbb C)$ exists.
\end{theorem}
\begin{proof}
Restriction to $\mathcal H_G$ proves that the $5+k$ target operators
are linearly independent. For precisely $5+k$ effects, their real span
therefore equals the target span and each geometric restriction has the form
$G_i=\alpha_i I_G+Y(\beta_i)$, where
$Y(b)=\sum_a b_aY_a$. By the symmetric directional spectrum
$\lambda_{\min}(Y(b))=-R\|b\|/2$, positivity implies
$\alpha_i\ge R\|\beta_i\|/2$. Normalization gives
$\sum_i\alpha_i=1$. Exact first moments and independence imply
$\sum_i x_{ia}\beta_{ib}=\delta_{ab}$. Taking the trace in the four
coordinate indices yields
\[
4=\sum_i x_i\cdot\beta_i\le r\sum_i\|\beta_i\|
\le (2r/R)\sum_i\alpha_i=2r/R,
\]
whence $r\ge 2R$.
For attainment, take the regular four-simplex directions $n_0,\ldots,n_4$
with $\sum n_j=0$, $\sum n_jn_j^\top=5I_4/4$, and
$F_j=[I+(4/r)X^D(n_j)]/5$. These are an exact-position POVM for
$r\ge 2R$. On the protected block
$F_j\succeq c_rI_S$ with
$c_r=[1-\sqrt3R/r]/5>0$. Let $w=(k+1)^{-1}$ and choose
$\epsilon>0$ small enough that
$\epsilon\max\{\|\sum_{j=1}^kT_j\|,\max_j\|T_j\|\}<wc_r$.
Replace $F_0$ by $E_0=wF_0-\epsilon\sum_jT_j$ and
$E_j=wF_0+\epsilon T_j$ for $1\le j\le k$, all with position
$rn_0$ and distinct passive labels. Retain $F_1,\ldots,F_4$.
All effects stay positive, their sum and all exact first moments are unchanged,
and $T_j=\epsilon^{-1}(E_j-w\sum_{\ell=0}^kE_\ell)$
lies in their real span. This gives exactly $5+k$ effects at every
$r\ge 2R$, including equality.
\end{proof}
\begin{corollary}[Outcome-count penalty below the simplex radius]
\label{p2c:cor:v353-readout-penalty}
For the above $k$ protected directions and $r<2R$, at least $6+k$
effects are required for any direct exact-position POVM, if one exists.
The count $5+k$ is attainable for every $r\ge2R$. This does not settle
the unrestricted below-$2R$ minimum for larger outcome counts.
\end{corollary}


\subsection{Endpoint quotient tomography and complete chiral erasure}
\label{p2c:sec:v354-endpoint-quotient}
Let $\chi=\gamma_1\gamma_2\gamma_3\gamma_4$, $\widetilde\chi=0_G\oplus\chi$,
$r_0=\sqrt3R_{\rm mom}$, and $P_+(n)=(I_4+\gamma(n))/2$.

\begin{theorem}[Unique protected endpoint support obstruction]
\label{p2c:thm:v354-unique-protected-obstruction}
For every finite positive weighted design $(w_i,n_i)$ with
$\sum_iw_i=1$, $\sum_iw_in_i=0$, and
$\sum_iw_in_in_i^T=I_4/4$,
\[
\operatorname{span}_{\mathbb R}\{P_+(n_i)HP_+(n_i):H=H^\dagger\}
=\{T=T^\dagger:\operatorname{Tr}(T\chi)=0\}.
\]
This real support span has dimension $15$.
\end{theorem}
\begin{proof}
By Clifford anticommutation, $P_+(n_i)\chi P_+(n_i)=0$.
If $Z=Z^\dagger$ annihilates every support subspace, then
$P_+(n_i)ZP_+(n_i)=0$. Sum over the design to get
$Z+\frac14\sum_a\gamma_aZ\gamma_a=0$.
Taking the Hilbert--Schmidt inner product with $Z$ gives
$\frac18\sum_a\|Z+\gamma_aZ\gamma_a\|_2^2=0$.
Thus $\{Z,\gamma_a\}=0$ for every $a$. The product $\chi Z$
commutes with all four generators and is scalar by irreducibility
of complex $\mathrm{Cl}_4$. Thus $Z$ is a real multiple of $\chi$.
\end{proof}

\begin{corollary}[Sharp 22-effect endpoint protected quotient frame]
\label{p2c:cor:v354-protected-22}
The existing seven-outcome endpoint POVM, with full rank protected
compressions $E_i^S=2w_iP_+(n_i)$ on their supports, has a positive
position-preserving tagged refinement with $7+15=22$ effects whose
real effect span contains every protected Hermitian operator orthogonal
to $\chi$. The minimum number of effects for this target is exactly $22$.
\end{corollary}
\begin{proof}
Choose fifteen independent locally supported protected Hermitian
operators from the preceding theorem. Split each parent effect into
$k_i+1$ same-position children whose sum is the parent and whose
protected contrasts recover its $k_i$ assigned operators, with
sufficiently small coefficients. This preserves positivity because
$E_i^S$ is strictly positive on its rank-two support.
The total effect count is $7+\sum_i k_i=22$.
For the lower bound, let $\mathcal W$ denote the real span of any
candidate native endpoint POVM whose effect span contains
$\mathcal T=\{0_G\oplus T:T=T^*,\ \operatorname{Tr}(T\chi)=0\}$.
Under endpoint Clifford support saturation, all protected effects
are orthogonal to $\chi$, so the kernel of geometric compression
$\mathcal W\to\operatorname{span}_{\mathbb R}\{G_i\}$ equals
$\mathcal T$ (both inclusion directions follow from the required
readout span). By the universal rank-six radius theorem
\ref{p2c:thm:universal-geometric-rank6-radius}, at
$r_0=\sqrt3R<\sqrt{34}R/3$ the geometric real effect rank is
at least seven. Rank-nullity gives
$N\ge\dim\mathcal W\ge7+\dim\mathcal T=22$.
The positive 22-effect tagged refinement attains the bound.
This is an operator-span result, not a physical occurrence assertion.
\end{proof}

\begin{theorem}[Complete instrument erasure of the endpoint chiral family]
\label{p2c:thm:v354-chiral-instrument-erasure}
For any first exact-position POVM at $r=r_0$ and any CP instrument
implementing it, including arbitrary output systems,
\[
\mathcal I_i\bigl((I_4+s\chi)/4\bigr)
=\mathcal I_i\bigl((I_4+t\chi)/4\bigr)
\qquad(s,t\in[-1,1])
\]
for every first outcome $i$. Hence all later adaptive/postselected
protocols have identical histories on the entire initial chiral family.
\end{theorem}
\begin{proof}
If $K_{i\alpha}$ is a Kraus operator, its protected-input restriction
$T_{i\alpha}=K_{i\alpha}P_S$ satisfies
$T_{i\alpha}^\dagger T_{i\alpha}\le E_i^S$, and the existing
endpoint equality gives $\operatorname{supp}E_i^S\subseteq P_+(n_i)$.
Thus $T_{i\alpha}=T_{i\alpha}P_+(n_i)$ and
$T_{i\alpha}\chi T_{i\alpha}^\dagger=0$.
Every Kraus contribution is independent of $s$; so is their sum.
\end{proof}

\begin{theorem}[Exact native endpoint rank and finite attainment]
\label{p2c:thm:v354-exact-native-rank}
Let $D_{\max}(r)$ denote the largest real effect-span dimension
among finite native POVMs with exact first moments $X_a^D$ and
position labels in the closed ball of radius $r$. Then
\[
\boxed{
D_{\max}(r)=
\begin{cases}
\text{infeasible}, & r<\sqrt3R_{\rm mom},\\
136, & r=\sqrt3R_{\rm mom},\\
137, & r>\sqrt3R_{\rm mom}.
\end{cases}}
\]
At the endpoint a 136-outcome POVM spans the entire real hyperplane
$\{Z\in(\mathcal A_D)_{\rm sa}:\operatorname{Tr}(Z\widetilde\chi)=0\}$.
Its exact minimal endpoint outcome count is $\boxed{136}$.
\end{theorem}
\begin{proof}
At $r=r_0$ every protected effect is orthogonal to $\chi$;
therefore rank is at most $136$.
For attainment, use the existing seven geometric endpoint effects
$G_i$ in the ordered basis
$(e_0,u_1,u_2,u_3,u_4,w_{12},w_{13},w_{14},w_{23},w_{24},w_{34})$.
The compressed maps $Z\mapsto(G_iZG_i)_i$ have real symmetric and
imaginary-skew ranks $65$ and $54$, respectively. The exact
finite-field certificate in
the exact v2\_354 certificate under \texttt{Proofs/ReadoutQuotient/}
reduces their algebraic coefficients modulo $97$ with
$\sqrt2\mapsto14$, $\sqrt3\mapsto10$; the two explicit null elements
$|w_{12}\rangle\langle w_{34}|\pm|w_{34}\rangle\langle w_{12}|$
(with $i$ multiplying the skew choice) give the matching upper bound.
Thus the original supported geometric space has dimension $119$.

Rotate the $(e_1,e_3)$ plane by $c=3/5,s=4/5$, with geometric lift
$U_G=1\oplus O\oplus\Lambda^2O$. It sends
$w_{12}\mapsto cw_{12}-sw_{23}$ and
$w_{34}\mapsto-sw_{14}+cw_{34}$. Hence the original and rotated
2-dimensional nullspaces intersect trivially; the combined fourteen
geometric supports span the full real $121$-dimensional Hermitian space.
The protected supports span $\chi^\perp$ of dimension $15$ by the
preceding theorem. The disjoint half-and-half union of the original
and rotated seven-outcome native POVMs is normalized and first-moment
exact. Its supported native real span has dimension $121+15=136$.
The fourteen parent effects span exactly nine real dimensions:
upper bound $7+7-5=9$ from their common $I,X_a^D$,
and exact lower bound nine from their geometric compressions
(using the same modular certificate).

For any finite POVM of $m$ effects spanning $b$ dimensions, if its
locally supported Hermitian directions span a $d$-dimensional space,
select $d-b$ directions extending the effects to a basis and split
each parent into $k_i+1$ sufficiently small positive same-position
contrasts. The resulting POVM has $m+d-b$ effects and spans all $d$
directions. Here $m+d-b=14+136-9=141$, which proves $D_{\max}(r_0)=136$
but does not yet optimize the outcome count.

For sharp outcome-count attainment, replace those fourteen parents with
the following \emph{eight independent} parents. Identify the geometric
carrier with $\mathbb C\oplus V\oplus W$, $V=\mathbb C^4$,
$W=\Lambda^2 V$, and put
$L_n=|e_0\rangle\langle n|+|n\rangle\langle e_0|$,
$J_n(v)=n\wedge v$, $B_n=J_n+J_n^*$,
$P_n=|n\rangle\langle n|$, and $Q_n=J_nJ_n^*$.
For $n\in\mathcal N=\{\pm e_1,\pm e_2,\pm e_3,\pm e_4\}$, set
\[
T(n)=\begin{pmatrix}
1 & a\langle n| & 0\\
a|n\rangle & \frac56 I_V+\frac23P_n & bJ_n^*\\
0 & bJ_n & \frac12 I_W+Q_n
\end{pmatrix},\qquad
a=\frac2{\sqrt3},\quad b=\frac{4\sqrt2}{3\sqrt3},
\]
and assign $x_n=r_0n$ and effects
\[
E_n=G_n\oplus S_n,\qquad
G_n=\tfrac18T(n),\qquad
S_n=\tfrac18(I_4+\gamma(n))=\tfrac14P_+(n).
\]
The native geometric coordinate is
$Y(n)=\frac R2L_n+\frac{\sqrt2R}{3}B_n$. Thus the odd terms
of the $\pm n$ pairs give exactly $\sum_nx_{na}E_n=X_a^D$.
For normalization, $\sum_{a=1}^4P_{e_a}=I_V$ and
$\sum_{a=1}^4Q_{e_a}=2I_W$, while all odd terms cancel,
so $\sum_nE_n=I_{15}$.

In the decomposition $V=\mathbb Cn\oplus n^\perp$ and
$W=(n\wedge n^\perp)\oplus\Lambda^2n^\perp$,
the geometric seed is the direct sum of
$\left(\begin{smallmatrix}1&a\\a&3/2\end{smallmatrix}\right)$,
three copies of
$\left(\begin{smallmatrix}5/6&b\\b&3/2\end{smallmatrix}\right)$,
and three copies of $(1/2)$. Their two nontrivial determinants are
$1/6$ and $7/108$, respectively, so $G_n\succ0$ on all eleven
geometric dimensions; $S_n$ is strictly positive on its rank-two
protected support.

These eight parents are real-linearly independent. Indeed, in any
vanishing linear combination the scalar--vector geometric entries
force equal coefficients for $+e_a$ and $-e_a$. The scalar entry then
forces their common coefficients $c_a$ to sum to zero; the vector
block becomes $(2/3)\operatorname{diag}(c_1,c_2,c_3,c_4)$ and
forces all $c_a=0$.
As every $G_n$ has full geometric rank, the locally supported real
geometric space is all of $(M_{11})_{\rm sa}$ of dimension $121$.
The eight equally weighted protected supports form a spherical
two-design, so Theorem~\ref{p2c:thm:v354-unique-protected-obstruction}
gives the complete fifteen-dimensional protected $\chi^\perp$ span.
Thus the locally supported native target has dimension $d=136$.
The preceding positive supported-refinement construction applies with
$m=b=8$, producing exactly $m+d-b=136$ positive same-position tagged
effects spanning the endpoint hyperplane. Conversely, spanning a real
136-dimensional space requires at least 136 effects. Therefore
$\boxed{N_{\min}^{\rm endpoint\ quotient}=136}$.
The choice of eight coordinate axes is geometrically compatible with
the quaternion-unit vertices of Paper I-A, but is not a physically
selected detector support or an identification of position with
momentum-cell vertices.

The earlier optimal-radius bound rules out smaller $r$, and
Paper II-C's existing $137$-outcome theorem gives full rank for
all $r>r_0$.
\end{proof}



\begin{remark}[Endpoint-quotient and constructor frontier]
\label{p2c:rem:v355-endpoint-physical-firewall}
The endpoint effect-span quotient is a \emph{real operator-vector-space}
quotient, not a quotient $C^*$-algebra. The values $136$ and $137$ are
algebraic distinguishability dimensions for the already constituted
coordinate realization. They are not measured event counts or detector
predictions, and the support-preserving splittings use freely tagged
mathematical outcomes. No GP2-licensed physical process, no new central-odd
$11/4$ bridge, and no universal GP3 Hilbert-multiplicity-one theorem is
inferred. The exact mathematical minimum of endpoint quotient-tomographic outcomes
is $136$; the unrestricted native exact-position count at
$\sqrt3 R\le r<\sqrt{34}R/3$ is exactly seven, and the complete
protected endpoint quotient requires exactly 22 effects. These three
sharp optimization results do not certify physical measurement
occurrence or apparatus selection.
\end{remark}


\section{Covariant finite instruments and amplified intrinsic readout}
\subsection{Amplified exact-coordinate effects and event reconstruction}
\begin{corollary}[Exact 137-effect one-step test on replicated intrinsic coordinates]
\label{p2c:cor:v298-amplified-coordinate-ic-one-step}
Assume the complete intrinsic coordinate algebra is represented faithfully
as $M_{11}\oplus M_4$ and an implementation has the componentwise
multiplicities characterized by
Theorem~\ref{p2c:thm:v294-native-multiplicity-algebra}:
\[
H=(\mathbb C^{11}\otimes\mathbb C^{k_g})\oplus
  (\mathbb C^4\otimes\mathbb C^{k_s}),\qquad k_g,k_s\ge1.
\]
For every $r>\sqrt3\Rmom$, the mathematical $137$-effect intrinsic
position POVM of
Theorem~\ref{p2c:thm:v258-minimal-complete-readout}
amplifies by $E_j\mapsto (E_{j,g}\otimes I_{k_g})\oplus
(E_{j,s}\otimes I_{k_s})$ to a positive normalized POVM on $H$.
It preserves the amplified exact first coordinate moments and spans the
entire represented coordinate-algebra Hermitian space, still of real
dimension $137$.

If these $137$ effects and $b$ CP outcome maps are actually registered,
all their finite adaptive outcome probabilities are multiplicity-blind
if and only if the $137b$ Heisenberg pullbacks of the effects lie in the
represented coordinate algebra. Failure produces a single-outcome,
terminal-effect distinguishing experiment by
Theorem~\ref{p4a:thm:v298-ic-one-step-outcome-descent}.
This is not a proof that the marked position POVM or its instruments occur,
nor a proof of physical endpoint multiplicity one.
\end{corollary}
\begin{proof}
The componentwise faithful algebra representation maps every positive
effect to a positive effect and preserves $I$, sums, and products.
Applying it to the original $137$ positive effects preserves their
normalization and all four first-moment equalities. Faithfulness makes
their $137$-dimensional real span exactly the represented algebra's
Hermitian part. Apply the one-step theorem to that effect frame.
The additional registration and GP3 claims are not consequences of
faithful representation.
\end{proof}


\begin{theorem}[Six-direction ceiling for any finite exactly Spin(4)-covariant outcome space]
\label{p2c:thm:v299-finite-covariant-six-ceiling}
Let $G=\operatorname{Spin}(4)$ act on the six inequivalent generated
momentum types, each with any positive multiplicity on which the
$G$-action is trivial. Take as represented test algebra either the
multiplicity-blind six-block kinematic algebra
$\mathcal A_{\rm kin}=\bigoplus_{\lambda\in T_p}
 \operatorname{End}(V_\lambda)\otimes I_{M_\lambda}$ or the
faithfully implemented intrinsic coordinate algebra
$\mathcal A_D=(M_{11}\otimes I_{k_g})\oplus(M_4\otimes I_{k_s})$.
Suppose a normalized finite-outcome POVM
$\{E_\omega:\omega\in\Omega\}\subset\mathcal A$ is exactly
$G$-covariant for a group action $\pi:G\to\operatorname{Sym}(\Omega)$:
\[
 U_gE_\omega U_g^*=E_{\pi(g)\omega}.
\]
Then every $E_\omega$ is $G$-invariant and
\[
\boxed{\dim_{\mathbb R}\operatorname{span}_{\mathbb R}
 \{E_\omega\}\leq\dim_{\mathbb R}\mathcal A_{\rm sa}^{G}=6.}
\]
The ceiling is sharp: the six orthogonal irreducible-type projections
form a normalized, invariant mathematical six-effect PVM spanning that
fixed algebra. Consequently no finite exactly $G$-covariant POVM
with outcome labels permuted by $G$ is informationally complete for
$\mathcal A_{\rm kin}$ (Hermitian dimension $43$) or $\mathcal A_D$
(Hermitian dimension $137$). The conclusion does not assume that
$\pi$ is continuous. It neither forbids a finite marked
noncovariant readout nor a continuously labelled covariant POVM or
an equivariant family of finite readouts indexed by an external frame.
\end{theorem}
\begin{proof}
Put $m=|\operatorname{Sym}(\Omega)|$. Every element $g$ of the
connected compact Lie group $G$ lies in a maximal torus and has an
$m$-th root $h$ in that torus. Therefore
$\pi(g)=\pi(h)^m=e$; the finite permutation action is trivial even
without a continuity hypothesis. Covariance makes every effect
commute with $U(G)$. On each of the six pairwise inequivalent
irreducible $G$-types, Schur's lemma makes an element of either
specified multiplicity-blind algebra a scalar; in the intrinsic
algebra its four geometric and two spinorial types are still
separate irreducible $G$-summands. Hence the fixed Hermitian algebra
has exactly the six central type projectors as a real basis.
These projectors are positive, sum to the identity, and attain the
six-dimensional span. Since $43,137>6$, neither such finite POVM is
informationally complete. This strengthens the first-moment
obstruction of Proposition~\ref{p2c:prop:v270-finite-outcome-covariance-no-go}
but does not constrain the separate continuous homogeneous-outcome
constructions of Theorem~\ref{p2c:thm:covariant-direction-complete-povm}.
\end{proof}

\begin{corollary}[Finite event basis inside the continuous directed-tangent readout]
\label{p2c:cor:v299-directed-tangent-137-event-basis}
At strict position-label radius $r>\sqrt3\Rmom$, choose the
mathematical six-passive-channel $G$-covariant informationally
complete directed-tangent POVM $E$ on
$\Omega=(G/U(1)_{\rm diag})\times\{1,\ldots,6\}$ from
Theorem~\ref{p2c:thm:covariant-direction-complete-povm}.
There exist $137$ Borel events $B_1,\ldots,B_{137}\subset\Omega$ whose
effects $E(B_j)$ form a real basis of $\mathcal A_{D,\rm sa}$.
One may choose the events as sufficiently small open neighborhoods
inside individual continuous-outcome channels. The events need not
be disjoint and are \emph{not} asserted to constitute a finite-outcome
POVM. For any specified finite family of $b$ registered outcome maps,
these event effects supply the $137b$ finite membership tests of
Theorem~\ref{p4a:thm:v299-continuous-ic-event-test}, without
breaking exact covariance of the original continuous POVM.
Physical registration of this continuous readout, its radius/owner
identification, and the outcome maps remain separate hypotheses.
\end{corollary}
\begin{proof}
The six-channel construction has continuous, positive effect
densities $D_a(gU(1)_{\rm diag})$ with full-support invariant measure,
and their real span is all of the $137$-dimensional native Hermitian
algebra. Each event effect is an integral of these densities and
lies in their span. Conversely, for a point $x$ in channel $a$,
averages $E(B)/\mu(B)$ over shrinking positive-measure open
neighborhoods of $x$ converge in operator norm to $D_a(x)$.
The span of all event effects is finite-dimensional and closed, so
it contains every density and is the entire algebra. Select $137$
linearly independent event effects; by continuity they may be
chosen from sufficiently small neighborhoods of a set of density
values that forms a basis. Apply the source-independent one-step
event criterion of Theorem~\ref{p4a:thm:v299-continuous-ic-event-test}.
Nothing in this construction selects a physically occurring detector.
\end{proof}

\begin{corollary}[Covariant mathematical instrument on the directed-tangent POVM]
\label{p2c:cor:v300-directed-tangent-luders}
Fix the $r>\sqrt3\Rmom$ mathematical six-channel directed-tangent
POVM of Theorem~\ref{p2c:thm:covariant-direction-complete-povm}.
On either the one-copy carrier or a faithful multiplicity-amplified
coordinate implementation of
Theorem~\ref{p2c:thm:v294-native-multiplicity-algebra}, its continuous
positive effect densities $D_a(q)$ are elements of the represented
intrinsic algebra $\mathcal A_D$ and obey $D_a(q)\succeq\eta I$
for some $\eta>0$ on the compact six-channel outcome space.
The event maps
\[
 \boxed{\ \mathcal I_B^{\rm root}(\rho)=\int_B
 D(\omega)^{1/2}\rho D(\omega)^{1/2}\,d\mu(\omega)\ }
\]
therefore define an exactly $\Spin(4)$-covariant \emph{mathematical}
CP instrument with the original position first moments, and every
one of its outcome maps preserves $\mathcal A_D$.
All $137$ basis-event pullback defects of
Corollary~\ref{p2c:cor:v299-directed-tangent-137-event-basis}
vanish for every event map of this constructed instrument; its
registered-grammar compatibility is consequently exact \emph{if}
that update rule is later licensed. This does not imply that an
apparatus realizes this instrument, that the multiplicity quotient is
a generated GP3 comparison, or that scalar-source/action ports are
preserved.
\end{corollary}
\begin{proof}
The proof of Theorem~\ref{p2c:thm:covariant-direction-complete-povm}
gives continuous densities with a strict positive lower bound at the
stated radius. The represented coordinate algebra is closed under
positive functional calculus, so its density square roots remain
inside that algebra. Apply the instrument and covariance theorem
\ref{p4a:thm:v300-covariant-root-instrument}, which proves
normalization, exact effect equality, covariance, and the quotient
invariance of each conditional map. Because the POVM itself is
unchanged, its four position first moments are unchanged.
Every event basis effect belongs to $\mathcal A_D$, hence all its
Heisenberg pullbacks stay there, giving zero exact defects.
\end{proof}


\section{Petersen--Hopf finite comparisons and their exact blind spaces}
\label{p2c:sec:v318-axis-frame-readout}
Assume the conditional occurring and embedded $\mathsf{Supp}_{2I}$
role of Paper~II-A.  Choose representatives of its ten unoriented
threefold Hopf axes $u_{ij}\in\mathbb R^3$, where $1\le i<j\le5$,
and put
\[
 Q_{ij}=u_{ij}u_{ij}^{\mathsf T}-\tfrac13 I_3,
 \qquad \mathcal Q=\operatorname{Sym}^2_0(\mathbb R^3).
\]
The projectors $Q_{ij}$ are independent of the choices of signs.

\begin{theorem}[Tight finite quadrupole comparison frame]
\label{p2c:thm:v318-quadrupole-tomography}
The ten axis quadrupoles form a tight, redundant frame of the real
five-dimensional comparison space $\mathcal Q$:
\[
 \sum_{i<j}Q_{ij}=0,\quad
 \sum_{i<j}\operatorname{Tr}(XQ_{ij})Q_{ij}=\frac43X,
 \quad X\in\mathcal Q.
\]
Thus the ten exact scalar comparisons $\operatorname{Tr}(XQ_{ij})$
reconstruct an arbitrary $X\in\mathcal Q$ by the canonical formula
\[
 \boxed{X=\frac34\sum_{i<j}
   \operatorname{Tr}(XQ_{ij})Q_{ij}.}
\]
No orientation, five-core label, physical test occurrence, scalar
calibration or density matrix is selected by the identity.
\end{theorem}
\begin{proof}
Paper~I-B's exact ten-axis quadrupole Gram matrix equals
$\frac43P_{\rm quad}$, where $P_{\rm quad}$ is the rank-five
orthogonal projector in the ten-axis coefficient space.  Therefore
the analysis/synthesis composition on $\mathcal Q$ is
$\frac43I_{\mathcal Q}$.  The vanishing sum is the five-core
star-incidence relation (or the $2I$-invariant traceless-tensor
average).  Both identities are independent of axis sign.
\end{proof}

\begin{theorem}[Finite real-state POVM and the complex obstruction]
\label{p2c:thm:v318-real-petersen-povm}
On the auxiliary \emph{real} three-dimensional Hopf-vector space,
define the ten effects
\[
 F_{ij}=\frac3{10}u_{ij}u_{ij}^{\mathsf T}\succeq0,
 \qquad\sum_{i<j}F_{ij}=I_3.
\]
For every real symmetric, positive, trace-one operator $R$, with
$p_{ij}=\operatorname{Tr}(RF_{ij})$, these probabilities determine $R$
exactly:
\[
 \boxed{R=\frac13I_3+\frac52\sum_{i<j}
                  (p_{ij}-\tfrac1{10})Q_{ij}.}
\]
The ten outcomes are $A_5$-covariant as a finite orbit, not fully
covariant under the connected $SU(2)$ with a nontrivial continuous
action on ten discrete labels.  They are informationally complete on
the six-real-dimensional symmetric-operator space but do not determine
a general complex Hermitian $3\times3$ density operator: the three
imaginary antisymmetric Hermitian directions have identically zero
expectation against every $F_{ij}$.
These auxiliary effects are not identified with the intrinsic
$\mathcal A_D=M_{11}(\mathbb C)\oplus M_4(\mathbb C)$ coordinate
readout or with a physically occurring apparatus.
\end{theorem}
\begin{proof}
The ten-axis tight-vector identity is
$\sum_{i<j}u_{ij}u_{ij}^{\mathsf T}=\frac{10}{3}I_3$.
Writing $R=I_3/3+X$ with $X\in\mathcal Q$ gives
$p_{ij}-1/10=(3/10)\operatorname{Tr}(XQ_{ij})$.
The preceding theorem therefore yields the reconstruction formula.
A real symmetric effect has zero trace against $iA$ for every real
antisymmetric $A$, proving the complex no-go.  A connected group
cannot act nontrivially and continuously on a finite discrete
outcome set, proving the continuous-covariance limitation.
\end{proof}

\begin{theorem}[Sharp five-axis comparison-set classification]
\label{p2c:thm:v318-five-probe-criterion}
Choose five of the ten unordered core pairs, $S\subset\binom{[5]}2$
with $|S|=5$.  These five axis quadrupoles span $\mathcal Q$ if and
only if the graph on the five \emph{core labels} whose edge set is
the five omitted pairs $T=\binom{[5]}2\setminus S$ is connected and
contains an odd cycle.  Exactly $162$ of the $252$ five-subsets satisfy
this condition: twelve have an omitted $5$-cycle, and $150$ have an
omitted triangle with two attached tree edges.  The remaining $90$
have rank four.  In particular five arbitrary axis comparisons are
not always a basis.  This is a statement about five linear comparison
\emph{settings}, not a normalized five-outcome POVM.
\end{theorem}
\begin{proof}
Write $D:\mathbb R^5\to\mathbb R^{10}$ for the unsigned
core-to-axis incidence map $(Df)_{ij}=f_i+f_j$.
The exceptional-geometry sequence in Paper~I-B identifies
$\ker(\sum a_{ij}Q_{ij})=\operatorname{im}D$.
A linear dependence among the five $Q_{ij}$ indexed by $S$
is equivalent to a nonzero $f$ obeying $f_i+f_j=0$ on every
omitted edge in $T$.  Such $f$ exists exactly when the omitted
edge graph has a bipartite connected component.  With five vertices
and five edges, having no bipartite component is equivalent to
being connected and odd-unicyclic.  The sole odd cycle has length
three or five.  There are $(5-1)!/2=12$ five-cycles.  For each of
the $\binom53=10$ triangles, the two remaining labelled vertices
can be attached to the triangle in $3^2+2\cdot3=15$ rooted-forest
patterns, giving $150$ further bases.  A direct exact incidence-rank
check establishes rank four for each of the other $90$ subsets.
\end{proof}

\begin{proposition}[Reciprocal product-frame limitation]
\label{p2c:prop:v318-real-reciprocal-blind-nine}
If the two distinct quaternionic factors of $M_8$ carry matched,
separately constituted $2I$ packets and the \emph{auxiliary real}
three-dimensional Hopf-vector model is taken in both factors,
then the one-hundred product effects $F_{ij}^{(x)}\otimes
F_{kl}^{(p)}$ span a $36$-dimensional real operator subspace,
\[
 \operatorname{Sym}^2(\mathbb R^3_x)\otimes
 \operatorname{Sym}^2(\mathbb R^3_p).
\]
It is a proper subspace of the $45$-dimensional real symmetric
operator space on $\mathbb R^3_x\otimes\mathbb R^3_p$; the
nine-dimensional $\Lambda^2\mathbb R^3_x\otimes
\Lambda^2\mathbb R^3_p$ sector is invisible.
This $36+9$ comparison identity is distinct from the $225=120+105$
exchange-parity decomposition of the \emph{fifteen-edge} tensor
carrier in Paper~I-C.  Neither constructs the native physical
coordinate instrument on $\mathcal H_1$.
\end{proposition}
\begin{proof}
The ten local effects span the six-dimensional real symmetric
operator subspace by the preceding theorem.  Their tensor products
span the $6\cdot6=36$ dimensional tensor product.  The canonical
splitting
$\operatorname{Sym}^2(V\otimes W)
\cong(\operatorname{Sym}^2V\otimes\operatorname{Sym}^2W)
\oplus(\Lambda^2V\otimes\Lambda^2W)$
for real vector spaces $V,W$ yields the complementary dimension
$3\cdot3=9$.
\end{proof}


\section{Auxiliary complex effect frames and robust finite reconstruction}


\subsection{Directed Hopf-axis completion of the auxiliary complex readout}
\label{p2c:sec:cc12-complex-axis-completion}
\begin{theorem}[Twenty-outcome covariant complex informational completion]
\label{p2c:thm:cc12-twenty-complex-povm}
For a supplied marked $2I$ embedding, take the twenty positive rank-one
matrices $H_{ij}$ of
Theorem~\ref{p1:thm:cc12-directed-coherent-gram}. Then
\[
 E_{ij}=\frac3{20}H_{ij}\succeq0,\qquad
 \boxed{\sum_{i\ne j}E_{ij}=I_3.}
\]
The twenty effects form an $A_5$-covariant \emph{auxiliary} POVM on
$\mathbb C^3$, informationally complete on the nine-real-dimensional
space $\operatorname{Herm}_3(\mathbb C)$, not merely its real symmetric
subspace. For $\rho\succeq0$, $\operatorname{Tr}\rho=1$, let
$p_{ij}=\operatorname{Tr}(\rho E_{ij})$. Write uniquely
\[
 \rho=\tfrac13 I_3+Q+iA(v),\quad Q=Q^{\mathsf T},\quad
 \operatorname{tr}Q=0,\qquad
 Q_{ij}=n_{ij}n_{ij}^{\mathsf T}-\tfrac13 I_3\ (i<j).
\]
Then the exact inversion is
\[
 \boxed{\begin{aligned}
 v&=\sum_{i<j}(p_{ij}-p_{ji})\,n_{ij},\\
 Q&=-5\sum_{i<j}(p_{ij}+p_{ji}-\tfrac1{10})Q_{ij},\\
 \rho&=\tfrac13I_3+Q+iA(v).
 \end{aligned}}
\]
For every normalized input the five outgoing and five incoming
marginals are separately uniform:
\[
 \boxed{\sum_{j\ne i}p_{ij}
       =\sum_{j\ne i}p_{ji}=\tfrac15\quad(1\le i\le5).}
\]
\end{theorem}
\begin{proof}
Each core's oriented tetrahedron has zero first and isotropic second
moments. Therefore $\sum_{j\ne i}H_{ij}=\tfrac43I$, giving both the
POVM normalization and the outgoing margin; the incoming statement
follows from $H_{ji}=\overline{H_{ij}}$. For
$X=aI+Q+iA(v)$ one computes
\[
 \operatorname{Tr}(XH_{ij})=
 a-\tfrac12n_{ij}^{\mathsf T}Qn_{ij}+v\cdot n_{ij}.
\]
Thus $p_{ij}-p_{ji}=\tfrac3{10}(v\cdot n_{ij})$ and
$p_{ij}+p_{ji}=\tfrac1{10}-\tfrac3{20}
 n_{ij}^{\mathsf T}Qn_{ij}$.
The ten-axis vector frame $\sum_{i<j}n_{ij}n_{ij}^{\mathsf T}
=\tfrac{10}{3}I$ and the pre-existing quadrupole frame
$\sum_{i<j}(n_{ij}^{\mathsf T}Qn_{ij})Q_{ij}=\tfrac43Q$
then yield the inverse. This also proves the nine-dimensional span.
\end{proof}

\begin{proposition}[All eleven homogeneous null comparisons and the exact positive image]
\label{p2c:prop:cc12-directed-probability-constraints}
Among the twenty real numbers $(p_{ij})_{i\ne j}$ produced by an
\emph{arbitrary} Hermitian matrix, there are precisely eleven
independent homogeneous linear relations. Eight express uniformity
of outgoing and incoming marginals relative to their common total:
\[
 \sum_{j\ne i}p_{ij}=\frac15\sum_{k\ne l}p_{kl},\qquad
 \sum_{j\ne i}p_{ji}=\frac15\sum_{k\ne l}p_{kl}.
\]
Three more remove the $\mathbf3'$ (Galois-conjugate) signed hand:
\[
 \boxed{\sum_{i<j}(p_{ij}-p_{ji})\,n_{ij}^{\sigma}=0\in\mathbb R^3,}
\]
where $n^\sigma$ is the fixed $\sqrt5\mapsto-\sqrt5$
conjugate of the exact ten-axis coordinate array.
Adding $\sum p_{ij}=1$ cuts out an eight-dimensional affine
subspace of $\mathbb R^{20}$. Within that subspace the image of
\emph{positive} trace-one density operators consists exactly of
those $p$ whose reconstructed $(Q,v)$ satisfy
\[
 \boxed{\begin{aligned}
 &\operatorname{tr}Q^2+2\|v\|^2\le\frac23,\\
 &\frac1{27}-\frac16\operatorname{tr}Q^2-\frac13\|v\|^2
     +\frac13\operatorname{tr}Q^3-v^{\mathsf T}Qv\ge0.
 \end{aligned}}
\]
These two inequalities characterize matrix positivity on this
trace-one Hermitian slice. In particular nonnegative outcome values
alone are not asserted sufficient.
\end{proposition}
\begin{proof}
In $\mathbb R^{20}$ the kernel of the coherent synthesis is
$\mathbf3'\oplus2\mathbf4$, of dimension eleven.
The two uniform-marginal families yield two independent copies of
$\mathbf4$; the signed conjugate-axis equation is the $\mathbf3'$
annihilator because $UU^{\sigma\,\mathsf T}=0$. This accounts for
all eleven independent constraints. For trace-one $\rho$,
\[
 \frac12(1-\operatorname{tr}\rho^2)
 =\frac13-\frac12\operatorname{tr}Q^2-\|v\|^2,
\quad
\det\rho=\frac1{27}-\frac16\operatorname{tr}Q^2
 -\frac13\|v\|^2+\frac13\operatorname{tr}Q^3-v^TQv.
\]
The eigenvalues of a Hermitian $3\times3$ matrix are real; with sum
one they are all nonnegative iff their second elementary symmetric
polynomial and product are nonnegative. These are exactly the two
inequalities.
\end{proof}

\begin{corollary}[One marked tetrahedral core gives a fourteen-effect completion]
\label{p2c:cor:cc12-fourteen-core-povm}
Fix one parent core $i$ and retain the ten unoriented real-axis
projections $R_a=u_au_a^{\mathsf T}$. The fourteen effects
\[
 \boxed{F_a=\tfrac3{20}R_a\quad(a=1,\ldots,10),\qquad
 T_j=\tfrac38H_{ij}\quad(j\ne i)}
\]
sum to $I_3$ and span $\operatorname{Herm}_3(\mathbb C)$.
Within the class consisting of the fixed real ten-axis orbit plus
positive-weight circular-polarisation effects of the form $H_n$,
\emph{at least four} additional effects are necessary for
informational completeness and a normalized POVM: cancellation of
the total imaginary part forces $\sum_{r=1}^m t_rn_r=0$, and if
$m\le3$ such a positive dependence leaves the $n_r$ in a space
of dimension at most two. The tetrahedral four attain this bound.
This optimality is relative to this fixed family, not to arbitrary
complex-qutrit POVMs.
\end{corollary}
\begin{proof}
The ten real effects sum to $\tfrac12I_3$ by the inherited
$10/3$-frame identity. The four local circular effects likewise sum
to $\tfrac12I_3$. Real comparisons span $\mathbf1\oplus\mathbf5$;
subtracting their real parts from the four circular effects leaves
$iA(n_{ij})/2$, whose $n_{ij}$ span $\mathbb R^3$.
The stated positive-dependence argument proves the restricted
minimality.
\end{proof}

\begin{proposition}[Two-factor completion of the nine real blind channels]
\label{p2c:prop:cc12-reciprocal-blind-nine-completion}
On the separately constituted auxiliary $\mathbb C^3_x$ and
$\mathbb C^3_p$ comparison factors, the $20\times20=400$ product
effects span all $81$ real dimensions of
$\operatorname{Herm}(\mathbb C^3_x\otimes\mathbb C^3_p)$.
Restricting to real symmetric two-factor density operators, the
previous ten-axis product span of dimension $36$ has the exact
nine-dimensional complement
$(i\mathfrak{so}(3)_x)\otimes(i\mathfrak{so}(3)_p)$.
It is resolved by the directed odd comparisons. Write
$B(n)=iA(n)$ and for real symmetric $R$ let
$K_{ab}=\operatorname{Tr}[R\,B(u_a)\otimes B(u_b)]$.
Then its orthogonal component in that nine-space is
\[
 \boxed{R_{oo}=\frac9{400}\sum_{a,b=1}^{10}
  K_{ab}\,B(u_a)\otimes B(u_b).}
\]
Here the ten-axis tight-frame factor is
$\sum_{a=1}^{10}|B(u_a)\rangle\!\langle B(u_a)|
 =\tfrac{20}{3}P_{\rm odd}$ in the real Hilbert--Schmidt
metric, so the two-factor inverse coefficient is $(3/20)^2=9/400$.
The data $K_{ab}$ can be formed as real linear combinations
of the $400$ product-effect expectations, since each $B(u_a)$ belongs
to their one-factor real span. This does not actualize any
reciprocal-position apparatus on $M_8$.
\end{proposition}
\begin{proof}
Local informational completeness implies full product completeness.
The real-symmetric tensor identity
$\operatorname{Sym}^2(V\otimes W)=
(\operatorname{Sym}^2V\otimes\operatorname{Sym}^2W)
\oplus(\Lambda^2V\otimes\Lambda^2W)$
gives $36+9$. The ten-axis vector frame and
$\operatorname{Tr}(B(u)B(v))=2u\cdot v$ yield
$\sum_a\langle B(u_a),X\rangle B(u_a)=\tfrac{20}{3}X$
for $X\in i\mathfrak{so}(3)$, hence the tensor-product frame
constant $(20/3)^2$ and its inverse $9/400$.
\end{proof}

\begin{remark}[Mathematical rather than physical measurement]
\label{p2c:rem:cc12-readout-occurrence-firewall}
All effects are mathematical operators on the auxiliary
$\mathbb C^3$ Hopf-vector comparison module. The 20 outcomes are
$A_5$-covariant after choosing the packet embedding, not a
continuous finite-outcome $SU(2)$-covariant native apparatus.
The 14-effect completion additionally marks one tetrahedral core.
Nothing here asserts these outcomes occur on $\mathcal H_1$ or on
the $M_8$ reciprocal tensor product.
\end{remark}

\subsection{Exact condition numbers and robust reconstruction of the twenty-effect comparison}
\label{p2c:sec:cc16-directed-tomography-stability}
\begin{theorem}[Sharp noise amplification of the coherent twenty-effect nonet]
\label{p2c:thm:cc16-directed-tomography-stability}
Let $H_{ij}$ and $F_{ij}=\frac3{20}H_{ij}$ be the twenty
rank-one matrices and effects of the proposed coherent-core results
\ref{p1:thm:cc12-directed-coherent-gram} and
\ref{p2c:thm:cc12-twenty-complex-povm}.
Equip $\operatorname{Herm}_3(\mathbb C)$ with the real
Hilbert--Schmidt norm and $\mathbb R^{20}$ with its Euclidean norm.
The normalized linear probability map
$\mathcal T(X)=(\operatorname{Tr}(XF_{ij}))_{i\ne j}$
has exact positive squared singular values
\[
 \boxed{\operatorname{spec}(\mathcal T^*\mathcal T)
 =\{(3/20)^{(1)},\,(3/40)^{(3)},\,(3/200)^{(5)}\}.}
\]
Its full condition number is $\sqrt{10}$, and its condition
number restricted to the eight-dimensional traceless slice is
$\sqrt5$. If $\tilde p=\mathcal T(\rho)+\eta$ for any Hermitian
trace-one $\rho$ and any real measurement error $\eta$, then
$\tilde\rho=\mathcal T^+\tilde p$ obeys
\[
 \boxed{\|\tilde\rho-\rho\|_{HS}\le
 \sqrt{200/3}\,\|\eta\|_2.}
\]
The constant $\sqrt{200/3}$ is optimal for the full linear inverse,
and is attained along a quintet direction. When $\sum_i\tilde p_i=1$,
$\tilde\rho$ has trace one. Its nearest-Hilbert--Schmidt
positive trace-one projection $\rho_{\rm phys}$ satisfies the same
bound whenever the true $\rho$ is a density matrix.
\end{theorem}
\begin{proof}
The unweighted synthesis $S:e_{ij}\mapsto H_{ij}$ of the earlier
core-pair theorem has nonzero squared singular values
$20/3,10/3,2/3$ with multiplicities $1,3,5$. Since
$\mathcal T=(3/20)S^*$ in the real Hilbert--Schmidt pairing,
its squared singular values are $(3/20)^2$ times those three
numbers, giving the displayed spectrum. A Moore--Penrose inverse
has norm equal to the reciprocal of the smallest positive singular
value, $\sqrt{200/3}$. Orthogonality of singlet, triplet and
quintet gives the two condition numbers. The quintet eigenspace
attains the operator norm. The constant-vector output contains the
entire trace mode, so unit-sum observations return unit-trace matrices.
Metric projection onto the closed convex density-matrix set is
nonexpansive for Hilbert--Schmidt norm. Since the true $\rho$ belongs
to that set, the final inequality follows.
\end{proof}
\begin{corollary}[Noise-aware physical-state incompatibility tests]
\label{p2c:cor:cc16-tomography-incompatibility}
If observations $\tilde p$ have a certified Euclidean error bound
$\|\eta\|_2\le\varepsilon$, then any compatible density state must
satisfy both
\[
 \|(I-\mathcal T\mathcal T^+)\tilde p\|_2\le\varepsilon,
 \qquad \lambda_{\min}(\mathcal T^+\tilde p)
 \ge-\sqrt{200/3}\,\varepsilon.
\]
A violation of either inequality excludes the auxiliary
positive-qutrit model at that error tolerance; passing the
inequalities alone does not establish compatibility or physical
occurrence. If $N$ independent outcomes of this exact POVM are sampled,
then, for $0<\alpha<1$, the empirical frequency vector satisfies
\[
 \boxed{\|\hat\rho-\rho\|_{HS}
 \le\sqrt{\frac{2000}{3N}\log\!\left(\frac{40}{\alpha}\right)}}
 \quad\text{with probability at least }1-\alpha,
\]
where $\hat\rho$ is the trace-one, positivity-projected least-squares
reconstruction. The bound is conservative and assumes the exact
POVM, independent trials and no systematic calibration error.
\end{corollary}
\begin{proof}
Distance to the linear image is a necessary compatibility condition.
For any candidate with probability error at most $\varepsilon$, the
preceding theorem and Weyl's eigenvalue perturbation inequality give
the second necessary condition. For the sampling statement, each
coordinate empirical frequency obeys the two-sided Hoeffding bound
$\Pr(|\hat p_i-p_i|>u)\le2e^{-2Nu^2}$.
The union bound over 20 coordinates gives
$\|\hat p-p\|_2\le\sqrt{20\log(40/\alpha)/(2N)}$ with probability
$\ge1-\alpha$. Multiply by $\sqrt{200/3}$ and apply the
nonexpansive positivity projection.
\end{proof}


\section{F40 cycle frames and compatible comparison packets}
\subsection{Complete cycle frame and its blind spaces}
\label{p2:sec:v345-f40-observable-frame}
The conditional forty-vertex quaternionic cover of Paper~I-B
has a twenty-one-dimensional real cycle space $\mathcal Z$.
Its three readouts are actual finite functions of the edge data:
right quotient $q^{-1}q'$, left quotient $q'q^{-1}$, and the
ordered transition of the uniquely omitted tetrahedral core.
By Theorem~\ref{p1:thm:v345-f40-triple-frame}, the single,
paired and triple restrictions have ranks $11$, $16$ and $21$.

\begin{proposition}[Complete geometric readout is not a physical instrument]
\label{p2:prop:v345-f40-readout-scope}
The joint comparison $\mathcal R$ satisfies
\[
 (4-\sqrt{10})\|z\|^2\leq\|\mathcal Rz\|^2
 \leq(6+\sqrt5)\|z\|^2\quad(z\in\mathcal Z).
\]
Hence $\mathcal R$ has a bounded left inverse, while no pair
of these three native readout types is injective on $\mathcal Z$.
The symmetric quotient $T_R+T_L$ cannot distinguish either
Hodge-mirror component, whereas $T_R-T_L$ has nonzero response
in both.  None of these linear chain readouts alone establishes
positive effects, an implemented detector, a probability law,
or a selection of an occurring $2I$ support.
\end{proposition}
\begin{proof}
Apply the exact frame spectrum and left/right Gram identities of
Paper~I-B.  The upper bound is the largest root $6+\sqrt5$;
the lower bound is the smallest root $4-\sqrt{10}$.
\end{proof}

\subsection{Compatible packets and instrument eligibility}
\label{p2:sec:v349-f40-compatible-packets}
\begin{proposition}[Exact readout eligibility test without a detector postulate]
\label{p2:prop:v349-f40-readout-compatibility}
Suppose a proposed comparison packet reports twelve right-shell,
twelve left-shell and twenty directed core-arrow numbers, all
intended to originate from \emph{one} F40 cycle.  Those 44
numbers represent a unique common cycle precisely when the
23 independent homogeneous equations of
Theorem~\ref{p1c:thm:v349-f40-23-equations} hold.  Equivalently,
for the native $44\times21$ frame $\mathcal R$, its Gram
$F=\mathcal R^*\mathcal R$ and the exact polynomial inverse
in Proposition~\ref{p1c:prop:v349-f40-rational-frame-inverse},
\[
 P_{44}d=d,\qquad P_{44}=\mathcal R F^{-1}\mathcal R^*,
 \qquad z=F^{-1}\mathcal R^*d.
\]
If only a single twisted-25 endomorphism comparison channel
is retained, it cannot faithfully encode the full cycle space
while preserving the diagonal $A_5$ action: its rank is at most
sixteen.  The Petersen five-core quintet provides the minimal
additive mathematical repair, with no new constitutional
postulate.
\end{proposition}
\begin{proof}
Apply the exact rank-nullity criterion for the full native
frame, and the $A_5$ multiplicity obstruction and repair from
Paper~I-B.  The inverse is positive-definite on the cycle
space by the certified frame spectrum.
\end{proof}
\begin{remark}[Operational limitation]
\label{p2:rem:v349-f40-detector-guard}
A compatible comparison packet is \emph{not} thereby an
implemented POVM, an admissible preparation, a detector or a
physical source.  A packet violating these equations cannot
be modeled as the triple readout of one F40 cycle, but it might
describe other degrees of freedom or an independently specified
physical measurement.  GP3 remains a fundamental constitutional
principle, not a missing multiplicity-one theorem.
\end{remark}


% BEGIN FRAMEWORK v3.067: COMPLETE SCALAR COORDINATE COMPARISON
\section{Finite scalar-coordinate algebra and information boundaries}
\label{p2c:sec:v3067-scalar-toeplitz-occurrence}
% Proposed mathematical insertion for Paper II-C, NOT YET ADOPTED.
% Source: Paper00-SharpFrontierOccurrence-v1_2.md, section 1--5.
% The physically test-typed corollary is conditional; the theorem itself
% is exact on the mathematically constituted full scalar band.

\begin{theorem}[Full finite scalar algebra from compressed round coordinates]
\label{p2c:thm:finite-scalar-coordinate-full-algebra}
Fix $K\ge0$, the unit-round reciprocal scalar harmonic bank
$F_K=\mathcal P_x(K)=\bigoplus_{j=0}^K E_j^x$, and the projector
$P_K$ onto it. For normalized coordinate functions
$u_a=x_a/R_x$ ($a=1,\ldots,4$), define the self-adjoint finite
coordinate compressions
\[
 T_a:=P_K M_{u_a}P_K|_{F_K}.
\]
Then
\[
 \boxed{C^*(I,T_1,T_2,T_3,T_4)=\operatorname{End}_{\mathbb C}(F_K).}
\]
For every $K\ge0$ their quadratic boundary defect is
\[
 \boxed{I_{F_K}-\sum_{a=1}^4T_a^2
 =\frac{K+2}{2(K+1)}\,\Pi_{E_K^x}.}
\]
For $K\ge1$, the degree-$K$ to degree-$(K-1)$ coupling identity is
\[
 \sum_{a=1}^4
 \bigl(\Pi_{E_K^x}T_a\Pi_{E_{K-1}^x}\bigr)
 \bigl(\Pi_{E_{K-1}^x}T_a\Pi_{E_K^x}\bigr)
 =\frac{K}{2(K+1)}\,\Pi_{E_K^x}.
\]
Consequently all scalar harmonic projectors and the finite-band
round Laplacian belong to this same finite coordinate algebra. This
is a statement on $F_K$; it does not constitute $F_K$ as a physical
instantaneous test carrier or replace the momentum level-one algebra.
\end{theorem}
\begin{proof}
A homogeneous solid harmonic $h_j$ of degree $j$ in four real
coordinates has
\[
 u_a h_j=(u_ah_j)_{j+1}+\frac{1}{2(j+1)}\partial_a h_j,
\]
where the two components have harmonic degrees $j+1$ and $j-1$.
The identity
$\sum_a\|\partial_a h_j\|_2^2=2j(j+1)\|h_j\|_2^2$
shows that the total squared lowering coefficient is
$j/[2(j+1)]$ and the total squared raising coefficient is
$(j+2)/[2(j+1)]$. Since $\sum_aM_{u_a}^2=I$, the only
coordinate leakage from $F_K$ goes from degree $K$ to degree $K+1$.
Taking the adjoint of this leakage map gives the displayed boundary
identity. The same lowering identity gives the displayed nonzero
adjacent-shell coupling.

Proceed by induction on $K$. For $K=0$, $F_0$ has dimension one.
For $K\ge1$, the boundary identity puts both $\Pi_{E_K^x}$ and
$\Pi_{F_{K-1}}$ into the algebra generated by the $T_a$. In the
$F_{K-1}$ corner, the compressed operators are those at level $K-1$;
by induction the corner contains $\operatorname{End}(F_{K-1})$.
The nonzero adjacent-shell coupling has ranges spanning $E_K^x$.
Multiplying these coupling operators by all lower-band rank-one
operators produces every top-to-lower matrix unit; adjoints and
products then produce all top-shell matrix units. The full matrix
algebra $\operatorname{End}(F_K)$ is generated.
\end{proof}

\begin{corollary}[Coordinate-only harmonic-shell recovery and intrinsic resolution]
\label{p2c:cor:finite-coordinate-intrinsic-resolution}
For $K\ge1$, let $X_a=aT_a$ on the complete round band $F_K$, where
$a>0$ is the homogeneous round radius. Then
\[
 \operatorname{Spec}\Bigl(\sum_aX_a^2\Bigr)
 =\left\{
 a^2\text{ with multiplicity }\dim F_{K-1},\quad
 a^2\frac{K}{2(K+1)}\text{ with multiplicity }(K+1)^2
 \right\}.
\]
Writing $r$ for the ratio of the smaller to the larger spectral value,
\[
 \boxed{K=\frac{2r}{1-2r},\qquad
 a=\sqrt{\lambda_{\max}(\sum_a X_a^2)}.}
\]
In normalized coordinates, the boundary defect has integer trace
\[
 \boxed{\operatorname{Tr}\bigl(I-\sum_aT_a^2\bigr)
 =\frac{(K+1)(K+2)}2.}
\]
All harmonic shell projectors are recovered recursively by applying
the boundary identity in each lower-band corner; consequently the
finite scalar kinetic operator
$\Delta_{F_K}=a^{-2}\sum_{j=0}^K j(j+2)\Pi_{E_j^x}$
and its finite functional calculus lie in $C^*(T_1,\ldots,T_4)$.
These are intrinsic mathematical and generated-role readout identities,
not an independent physical occurrence of a sharp test frontier.
\end{corollary}

\begin{theorem}[Irreducibility is equivalent to a complete invariant band]
\label{p2c:thm:finite-coordinate-complete-band-criterion}
Let $F\subset L^2(S_x^3)$ be any finite-dimensional
$\operatorname{Spin}(4)_x$-invariant scalar subspace containing
$E_0^x\oplus E_1^x$. Put $T_a^F=P_FM_{u_a}P_F|_F$. Then
\[
 \boxed{\{T_1^F,T_2^F,T_3^F,T_4^F\}'=\mathbb C I_F
 \quad\Longleftrightarrow\quad F=\mathcal P_x(K)
 \text{ for some finite }K\ge1.}
\]
\end{theorem}
\begin{proof}
Invariance and the multiplicity-one decomposition of scalar harmonics
imply $F=\bigoplus_{j\in J}E_j^x$ for some finite set $J$ containing
0 and 1. Coordinate multiplication has only degree $j\mapsto j\pm1$
components. If $J$ has a gap, the sum of harmonics below that gap and
the sum above it are both nonzero reducing subspaces of all $T_a^F$;
their block projection is a nonscalar member of the commutant. If $J$
has no gap, $J=\{0,\ldots,K\}$, and the full-matrix theorem above
makes the commutant scalar.
\end{proof}

\begin{remark}[Precisely typed occurrence boundary]
\label{p2c:rem:finite-coordinate-algebra-occurrence-boundary}
On an independently certified scalar TEST-sort finite-band occurrence
whose compressed coordinate family is a frozen licensed test word,
Paper~II-A's generated-image co-realization and Paper~IV-A's finite
physical test-algebra closure may retain the displayed full matrix
algebra, shell projectors and resolution readouts **without further
independent test-role roots**. No conclusion is drawn that such an
initial finite test role, its physical-time history, its particular
measurement tokens, or Paper~III-F's independently scaled sharp
frontier $\mathsf{SF}_{\Lambda_F}$ occurs. Finite algebra closure
alone also does not select the level $K$, since the theorem holds for
every finite $K$.
\end{remark}

% Research-only insertion candidate for Paper II-C / Paper IV-A.
% Source baseline: PapersMinimal v3_059, Paper II-C v2_173, Paper I-A v5_180.
% No physical occurrence, instrument, sharp frontier or same-action readout is added.
\begin{theorem}[Boundary-supported scalar Toeplitz commutators]
\label{p2c:thm:research-v13-scalar-boundary-commutator}
On the unit round scalar band $F_K=\bigoplus_{k=0}^KE_k$,
$K\ge1$, let $P_K$ be the orthogonal projection,
$T_a=P_KM_{u_a}P_K$, and $P_{[K]}=\Pi_{E_K}$.
With $L_{ab}=u_a\partial_b-u_b\partial_a$ one has
\[
 [T_a,T_b]=\frac1{2(K+1)}P_{[K]}L_{ab}P_{[K]},
 \quad
 -\sum_{a<b}[T_a,T_b]^2
  =\frac{K(K+2)}{4(K+1)^2}P_{[K]}.
\]
Furthermore, setting $D_K=I-\sum_aT_a^2$, one has
\[
 D_K=\frac{K+2}{2(K+1)}P_{[K]},
 \quad
 -\sum_{a<b}[T_a,T_b]^2=\frac{K}{2(K+1)}D_K,
 \quad
 \|[T_1,T_2]\|=\frac{K}{2(K+1)}.
\]
All these are mathematical scalar-band operator identities. They
reconstruct the finite $\mathfrak{spin}(4)$ generators and their
Casimir from the compressed coordinates on a supplied band, without
actualizing a physical test or action Hamiltonian.
\end{theorem}
\begin{proof}
For $h_K$ homogeneous harmonic of degree $K$,
$R_bh_K=u_bh_K-|u|^2\partial_bh_K/[2(K+1)]$
is its degree-$(K+1)$ raising part. Multiplication by $u_a$
lowers a harmonic $g_{K+1}$ via $\partial_ag_{K+1}/[2(K+2)]$.
Insert $I=P_K+(I-P_K)$ in the commuting ambient coordinate
products. The resulting commutator is supported only on $E_K$ and
there equals
\[
 -\frac{\partial_aR_b-\partial_bR_a}{2(K+2)}
 =\frac{u_a\partial_b-u_b\partial_a}{2(K+1)}.
\]
The standard round angular Casimir equals $K(K+2)$ on $E_K$.
The squared-coordinate defect follows by the harmonic raising norm,
as in the established scalar Toeplitz leakage theorem. The final
norm statement follows from the weight-$K$ vector
$(u_1+iu_2)^K$ and the bound $|m|\le K$ for weights in $E_K$.
\end{proof}

\begin{proposition}[No multiplicative coordinate inheritance across scalar resolutions]
\label{p2c:prop:research-v13-toeplitz-UCP-nesting}
For $K\ge1$ no unital $*$-homomorphism
$\iota:\End(F_K)\to\End(F_{K+1})$ sends all four $T_a^{(K)}$
to $T_a^{(K+1)}$. The canonical inclusion
$V_K:F_K\hookrightarrow F_{K+1}$ instead gives a unital
completely positive map
\[
 \Phi_{K+1\to K}(A)=V_K^*AV_K,
 \qquad \Phi_{K+1\to K}(T_a^{(K+1)})=T_a^{(K)},
\]
with exact quadratic multiplicativity defect
\[
 \sum_a\left(\Phi((T_a^{(K+1)})^2)-\Phi(T_a^{(K+1)})^2\right)
 =\frac{K+2}{2(K+1)}P_{[K]}.
\]
For the canonical scalar evaluation POVM of Paper I-A,
$Q_K(f)=P_KM_fP_K$ satisfies
$\Phi_{K+1\to K}\circ Q_{K+1}=Q_K$ for every bounded
scalar symbol $f$.
\end{proposition}
\begin{proof}
A unital $*$-homomorphism preserving coordinates would preserve the
polynomial $D_K=I-\sum_a(T_a^{(K)})^2$. It would therefore send the
projection $P_{[K]}$ to
$c_{K+1}P_{[K+1]}/c_K$, where
$c_j=(j+2)/[2(j+1)]$. This is nonzero and not a projection,
a contradiction. The compression identity follows from insertion
of $P_{K+1}=P_K+P_{[K+1]}$ in the coordinate products and the
boundary leakage identity. The compatibility with $Q_K$ follows
from $P_KP_{K+1}=P_K$.
\end{proof}

\begin{theorem}[Exact information rank of the canonical scalar evaluation POVM]
\label{p2c:thm:research-v13-evaluation-effect-rank}
Let $D_K=\dim_{\mathbb C}F_K=(K+1)(K+2)(2K+3)/6$, and let
$E_K(du)=|\kappa_{K,u}\rangle\langle\kappa_{K,u}|\,d\mu(u)$
be the canonical normalized rank-one scalar-band evaluation POVM.
Then the real linear span of its effect densities in
$\operatorname{Herm}(F_K)$ has dimension exactly $D_{2K}$, with
Hilbert--Schmidt blind dimension exactly $D_K^2-D_{2K}$.
For the specified rank-one L\"uders instrument, repeated uses at
any number of stages retain exactly this same original-state effect
span; they do not provide complete state tomography for $K\ge1$.
\end{theorem}
\begin{proof}
In a real orthonormal harmonic basis, the evaluation map sends an
operator $A$ to $\sum_{ij}A_{ij}\phi_i(u)\phi_j(u)$, a real scalar
polynomial of degree at most $2K$. All monomials up to degree $2K$
are products of two monomials of degree at most $K$, so this map has
image the full real polynomial space $\mathcal P(2K)$ and rank
$D_{2K}$. Hilbert--Schmidt duality gives the effect-span and blind
counts. Writing $|\psi_u\rangle=|\kappa_{K,u}\rangle/\sqrt{D_K}$,
the effect density is $e_u=D_K|\psi_u\rangle\langle\psi_u|$.
The L\"uders two-step probability density is
\[
 p_\rho(u,v)=p_\rho(u)D_K|\langle\psi_v,\psi_u\rangle|^2,
\]
and induction gives the same Markov factorization for any number
of steps. Thus every repeated-history original-state effect is a
scalar multiple of $e_u$ at the first outcome.
\end{proof}

% This insertion is a scope refinement of mathematical/operational
% distinction, not a proof of Paper III-F's sharp-frontier occurrence.

% END FRAMEWORK v3.067: COMPLETE SCALAR COORDINATE COMPARISON

\section{Conclusion and mathematical export boundary}
\label{p2c:sec:final-conclusion}
The generated reciprocal position role has a five-dimensional level-one
scalar-resolution seed, distinct from the fixed fifteen-dimensional
momentum carrier. Its intrinsic finite coordinate and kinematic
constructions yield the exact operator algebras and compact
transvection/coordinate relations established above. In particular, the
selected intrinsic algebra $\mathcal A_D=M_{11}(\mathbb C)\oplus
M_4(\mathbb C)$ is not the full $M_{15}(\mathbb C)$ envelope. The
coupling-graph and rank-one bridge results identify which additional
operator data would algebraically close that gap; they do not generate
such a bridge as a physical process.

The finite reciprocal-harmonic filtration, first-symbol pairing,
transported shell windows and guarded multiplication rules are exact
on their named domains. The intrinsic tangential class $\mathsf C_D$
and the radial cone class $\mathsf C_{\rm cone}$ remain distinct, and
ambient completions require their stated independent data.
The exact localization and covariant readout theorems provide
separate attainment, span and no-go results. The seven-outcome
exact-position construction, $22$-effect protected endpoint comparison,
$136$-dimensional endpoint real quotient, and $137$-dimensional
strict-radius-slack native Hermitian span refer to different targets,
and none of these numerical statements can be substituted for another.

Conditional exceptional finite frames, auxiliary complex POVMs,
F40 packet constraints, and stability estimates supply mathematical
comparisons, not an independent source or detector. In particular,
the marked directed-tangent readout does not supply its own globally
chosen orientation, and finite exact covariance alone does not produce
a continuous physical instrument. A physical preparation, measurement
update, time history or calibration still requires the corresponding
claim-relative geometric occurrence and action data. The conditional
sharp-frontier history supplied downstream does not retrospectively
turn the first-order position-role derivation into an unconditional
physical trajectory. The paper exports precisely these finite
mathematical capabilities and their stated boundaries.

\printbibliography


% BEGIN FRAMEWORK v3.068 GEOMETRY BRIDGE SOURCE INTEGRATION
\section{Fixed-Hurwitz support and finite endpoint moment sharpness}
\noindent\emph{Scope.} The statements below retain their declared finite geometric, source-chart, and action-relative hypotheses. They do not assert physical particle decay or unconditional source occupation.
\begin{theorem}[Sharp endpoint outcome minimum on every Hurwitz 24-cell]
\label{p2c:thm:v368-hurwitz-endpoint-eight}
For the native intrinsic tuple
$X^D=Y\oplus (\sqrt3\Rmom/4)\gamma$ and the classical endpoint
radius $r_0=\sqrt3\Rmom$, require every outcome direction to belong
to a fixed isometric Hurwitz set
\[
V_H=\{\pm e_a:1\le a\le4\}\ \cup\
\{\tfrac12(s_1,s_2,s_3,s_4):s_a=\pm1\}.
\]
Every finite POVM with exact intrinsic first moments then has at least
eight nonzero effects, and the bound is attained by the source-derived
$Q_8$ endpoint rule. Hence $N_{\min}^{V_H}(r_0)=8$. The same holds
on each of the 25 bi-translated Hurwitz chamber-centre blocks of
Theorem~\ref{p1b:thm:cc17-facet-double-coset}.
\end{theorem}
\begin{proof}
The optimal-radius protected equality of
Theorem~\ref{p2c:thm:v226-equality} yields a positive weighted
spherical two-design: $\sum_iw_i=1$, $\sum_iw_iq_i=0$ and
$\sum_iw_iq_iq_i^T=I_4/4$. Combine equal directions.
Let $p$ be the weight on the eight axial $\pm e_a$ points and
$h=1-p$ the weight on sixteen sign vertices $s/2$.
If $0<p<1$, each coordinate axis carries axial weight $p/4>0$,
so at least four axial points occur. The half-coordinate weighted
moment identity $\sum_s w_s ss^T=hI_4$ has rank four, so at least
four half-coordinate points occur. Thus at least eight.
If $p=1$, first and second moments force weight $1/8$ at
both $+e_a$ and $-e_a$ for every $a$, hence eight.
If $p=0$, write $\pi(s)$ for the probability on the sign cube.
The zero first and off-diagonal second moments annihilate every
Walsh character of degree one or two. Fourier inversion gives
\[
16\pi(s)=1+\epsilon(s)(b+\mathbf a\cdot s),\qquad
\epsilon(s)=s_1s_2s_3s_4,\quad b\in[-1,1].
\]
As $\epsilon(-s)=\epsilon(s)$, every antipodal sign pair has mass
$\pi(s)+\pi(-s)=[1+\epsilon(s)b]/8$. If $-1<b<1$, each of the
eight sign antipodal pairs has positive mass, so at least eight
points occur. If $b=\pm1$, parity is deterministic and the cubic
coefficients satisfy
$a_j=\mathbb E[\epsilon s_j]=\pm\mathbb E[s_j]=0$.
The distribution is then uniform on all eight signs of that parity.
Thus every possible positive weighted two-design uses at least eight
Hurwitz vertices.
For $q\in Q_8=\{\pm e_a\}$ the source kernel
$E_q=K_c(q)/8$ is positive, normalized and has exact moments for
$1/3\le c\le1$ by
Corollary~\ref{p2c:cor:v233-optimal-quaternion-povm}.
Transport by the existing $\Spin(4)$ action proves the same sharp
statement on every isometric 24-cell block.
\end{proof}

\begin{corollary}[Three octet readouts and fixed finite-group symmetry minima]
\label{p2c:cor:v368-octet-and-finite-symmetry}
For $H=2T$ and $Q=Q_8\triangleleft H$, choose
$\omega=(1+i+j+k)/2$. Then
$H=Q\dot\cup\omega Q\dot\cup\omega^2Q$, and every octet
supports the positive eight-outcome POVM $E_q=K_c(q)/8$.
The $24$-outcome Hurwitz POVM $K_c(q)/24$ is their equal mixture.
Under an explicitly required finite left-$H$ covariance on classical
endpoint outcomes, the exact minimum is $24$; under an explicitly
required finite left-$2I$ covariance it is $120$.
The same attaining rules are also covariant under the corresponding
finite left-and-right action.
\end{corollary}
\begin{proof}
The three cosets are left translates of the spherical
three-design $Q_8$, so polynomial degree-two positivity kernels
have exact degree-three cubature for normalization and first moments.
The finite left-action $g:q\mapsto gq$ is free on $S^3$;
any finite invariant nonempty support is a union of full orbits
of at least $|H|=24$ or $|2I|=120$ elements. The positive kernels
$K_c(q)/|H|$ on $H$ and $K_c(q)/120$ on $2I$ attain those
bounds because the two supports are spherical five- and eleven-designs,
respectively. The subgroup support is stable under both its left and
right actions, proving the last claim.
\end{proof}

\begin{corollary}[Exact finite endpoint operator-valued scalar moment ports]
\label{p2c:cor:v368-finite-moment-ports}
Let $K_c(q)$ be the already generated quadratic native endpoint kernel,
$1/3\le c\le1$, and let $\nu$ be the positive, round cubature
on either the Hurwitz 24-set (degree five), the binary-icosahedral
120-set (degree eleven), or the 720-point vertex/facet support
of Theorem~\ref{p1b:thm:cc17-positive-vertex-facet-19}
(degree nineteen, weights $1/504$ and $2/1575$).
For every scalar restricted polynomial $f$ of degree at most
$t-2$, where $t=5,11,19$ is the corresponding design strength,
\[
\sum_{q\in\mathrm{supp}\nu}\nu_q f(q)K_c(q)
=\int_{S^3}f(q)K_c(q)\,d\sigma(q).
\]
On the respective complete scalar test bands $\mathcal P_{t-2}$,
the finite maps have the same rank $14$ and the same harmonic kernel
$\bigoplus_{n=3}^{t-2}E_n$ as the continuous map.
\end{corollary}
\begin{proof}
Every entry of $f(q)K_c(q)$ has coordinate degree at most
$\deg f+2\le t$, and the positive cubature is scalar-exact through
degree $t$. Apply it entrywise. The rank and kernel follow from
the continuous native harmonic-kernel theorem
\ref{p2c:cor:optimal-povm-harmonic-kernel}.
\end{proof}

\begin{remark}[No imported physical occurrence]
The sharp support-specific minimum $8$ must be distinguished from the
already established unrestricted endpoint minimum $7$ of
Corollary~\ref{p2c:cor:v266-outcome-count}. Finite-group covariance is a
specified mathematical premise, not an apparatus-selection theorem.
None of these results generates an instrument, physical event token,
charged mass/source occurrence, decay width, sky transfer, or new
momentum level. The 720 test-degree-17 statement is operator-valued
integration exactness, not scalar evaluation isometry on $\mathcal P_{17}$.
\end{remark}
% END FRAMEWORK v3.068 GEOMETRY BRIDGE SOURCE INTEGRATION

\end{document}
